
\begin{theorem}[Printed Theorem 1.1]\label{mcrt-auto-theorem-1-1}\dcref{1.1}\source{matsumura-commutative-ring-theory:page-0017}\uses{}
\begin{verbatim}
If I is a proper ideal then there exists at least one maximal
ideal containing   I.

   An ideal P of A for which A/P is an integral domain       is called a prime
ideal. In other words, P is prime if it satisfies
   (i) P # A and (ii) x,y$P+xy$P         for x,y~A
A field is an integral domain, so that a maximal ideal is prime.
  If I and J are ideals and P a prime ideal, then


Indeed, taking x~l and YEJ with x,y$P, we have XYEIJ but xy#P.
   A subset S of A is multiplicative if it satisfies
   (i) x,y~S+xy~S,      and (ii) 1~s;
(here condition (ii) is not crucial: given a subset S satisfying (i), there will
usually not be any essential change on replacing S by Su (1)). If I is
an ideal disjoint from S, then exactly as in the proof of Theorem 1 we see
that the set of ideals containing I and disjoint from S has a maximal
element. If P is an ideal which is maximal among ideals disjoint from S
then P is prime. For if x$P, y#P, then since P + xA and P + yA both
meet S, the product (P + xA) (P + yA) also meets S. However,
         (P$xA)(P+yA)cP+xyA,
so that we must have xy$P. We have thus obtained the following theorem.
\end{verbatim}
\end{theorem}
\begin{proof}
\notready
\end{proof}


\begin{theorem}[Printed Theorem 1.2]\label{mcrt-auto-theorem-1-2}\dcref{1.2}\source{matsumura-commutative-ring-theory:page-0017}\uses{}
\begin{verbatim}
Let S be a multiplicative      set and I an ideal disjoint from S;
then there exists a prime ideal containing I and disjoint from S.
   If I is an ideal of A then the set of elements of A, some power of which
belongs to I, is an ideal of A (for X?EZ and y?~I=>(x + J$?+~-~EI           and
Ideals                                                                  3
91
(ax)?EZ). This set is called the radical of I, and is sometimes written JZ:
          JZ = (aEAla?eZ     for some n > O}.
If P is a prime ideal containing Z then X?EZ c P implies that XEP, and
hence ,,/I c P; conversely, if x#JZ then S, = {1,x,x2,. . .} is a multi-
plicative set disjoint from I, and by the previous theorem there exists a
prime ideal containing Z and not containing x. Thus, the radical of Z is
the intersection of all prime ideals containing I:
         JI=      ,nlp.
                    2
In particular if we take Z = (0) then J(O) is the set of all nilpotent elements
of A, and is called the nilradical of A; we will write nil(A) for this. Then
nil(A) is intersection of all the prime ideals of A. When nil(A) = 0 we say
that A is reduced. For any ring A we write Ared for A/nil(A);A,,, is of
course reduced.
    The intersection of all maximal ideals of a ring A( # 0) is called the
Jacobson radical, or simply the radical of A, and written rad(A). If
xerad(A) then for any aEA, 1 + ax is an element of A not contained in
any maximal ideal, and is therefore a unit of A by Theorem 1. Conversely
if XEA has the property that 1 + Ax consists entirely of units of A then
xErad(A) (prove this!).
    A ring having just one maximal ideal is called a local ring, and a
(non-zero) ring having only finitely many maximal ideals a semilocalring.
We often express the fact that A is a local ring with maximal ideal m by
saying that (A, m) is a local ring; if this happens then the field k = A/m is
called the residue field of A. We will say that (A, m, k) is a local ring to
mean that A is a local ring, m = rad(A) and k = A/m. If (A, m) is a local
ring then the elements of A not contained in m are units; conversely a
(non-zero) ring A whose nan-units form an ideal is a local ring.
    In general the product II? of two ideals I, I? is contained in Z n I?, but
does not necessarily coincide with it. However, if Z + I? = (1) (in which case
we say that Z and I? are coprime), then II? = 1 nZ?; indeed, then
Znl? = (ZnZ?)(Z + I?) c ZZ? c ZnZ?. Moreover, if Z and I?, as well as Z
and I? are coprime, then I and I?I? are coprime:
          (1) = (I + Z?)(Z + Z?) c 1 + Z?Z? C (1).
By induction we obtain the following theorem.
\end{verbatim}
\end{theorem}
\begin{proof}
\notready
\end{proof}


\begin{theorem}[Printed Theorem 1.3]\label{mcrt-auto-theorem-1-3}\dcref{1.3}\source{matsumura-commutative-ring-theory:page-0018}\uses{}
\begin{verbatim}
If I,, I,, . . . ,Z, are ideals which arc coprime in pairs then
          z,z,. . . Z,=IlnZ,n~~~nZ,.
In particular if A is a semilocal ring and m,, . . . m, are all of its maximal
ideals then
4         Commutative rings and modules

   Furthermore, if I + I? = (1) then A/II? N A/I x A/I?. To see this it is
enough to prove that the natural injective map from A/If? = A/I r?l I? to
A/I x A/I? is surjective; taking etzl, e/El? such that e + e?= 1, we
have ae? + a?e = a (mod I) ae?+ a?e E a? (mod 1?) for any a, u?EA,
giving the surjectivity. By induction we get the following theorem.
\end{verbatim}
\end{theorem}
\begin{proof}
\notready
\end{proof}


\begin{theorem}[Printed Theorem 1.4]\label{mcrt-auto-theorem-1-4}\dcref{1.4}\source{matsumura-commutative-ring-theory:page-0019}\uses{}
\begin{verbatim}
If I,,. . , I, are ideals which are coprime in pairs then
         AJI,.   .I, z AJI, x ... x AJI,,.
\end{verbatim}
\end{theorem}
\begin{proof}
\notready
\end{proof}


\begin{theorem}[Printed Theorem 2.1]\label{mcrt-auto-theorem-2-1}\dcref{2.1}\source{matsumura-commutative-ring-theory:page-0022}\uses{}
\begin{verbatim}
Suppose that M is an A-module generated by n elements,
 and that cpEHom,(M, M); let I be an ideal of A such that q(M) c ZM.
 Then there is a relation of the form
    (**) (p~+alcp?-l+~~~+u,-,~+u,=O,
 with Ui~Z? for 1 d i 6 n (where both sides are considered as endomorph-
 isms of M).
\end{verbatim}
\end{theorem}
\begin{proof}
\begin{verbatim}
Let M = Ao, + ... + AU,; by the assumption cp(M) c IM there exist
 Uij~l such that cp(oi) = cJ= i uijwj. This can be rewritten

          jil   ((~6, - 6,) uj = 0 (for 1 d i < n),

where aij is the Kronecker symbol. The coefficients of this system of linear
equations can be viewed as a square matrix ((~6~~- aij) of elements of A?[q],
the commutative subring of the endomorphism ring E of M generated by
the image A? of A together with cp; let b, denote its (i,j)th cofactor, and
8         Commutative   rings and modules

d its determinant. By multiplying the above equation through by b, and
summing over i, we get do, = 0 for 1 d k d n. Hence d.M = 0, so that d = 0
as an element of E. Expanding the determinant d gives a relation of the
form (**). W
\end{verbatim}
\end{proof}


\begin{theorem}[Printed Theorem 2.2]\label{mcrt-auto-theorem-2-2}\dcref{2.2}\source{matsumura-commutative-ring-theory:page-0023}\uses{}
\begin{verbatim}
NAK). Let M be a finite A-module and I an ideal of A. If
M = IM    then there exists aEA such that aM = 0 and a = 1 mod I. If in
addition I c rad (A) then M = 0.
\end{verbatim}
\end{theorem}
\begin{proof}
\begin{verbatim}
Setting cp= 1, in the previous theorem gives the relation a =
l+a,+    .. . + a, = 0 as endomorphisms   of M, that is aM = 0, and
a = 1 modI. If I c rad(A) then a is a unit of A, so that on multiplying
both sides of aM = 0 by a-l we get M L 0. n
\end{verbatim}
\end{proof}


\begin{theorem}[Printed Theorem 2.3]\label{mcrt-auto-theorem-2-3}\dcref{2.3}\source{matsumura-commutative-ring-theory:page-0023}\uses{}
\begin{verbatim}
Let (A, m, k) be a local ring and M a finite A-module; set
R = M/mM. Now M is a finite-dimensional      vector space over k, and we
?2       Modules

write n for its dimension. Then:
   (i) If we take a basis {tii,..., U,) for M over k, and choose an
inverse image UiE M of each y then {ul ,. . . ,u,} is a minimal basis
of M;
   (ii) conversely every minimal basis of M is obtained in this way, and
so has n elements.
   (iii) If {q,. . .,u,} and {ui ,. . . ,u,} are both minimal bases of M, and
ui = 1 UijUj with Uij~A then det (aij) is a unit of A, SO that (aij) is an
invertible matrix.
\end{verbatim}
\end{theorem}
\begin{proof}
\begin{verbatim}
(i) M = xAui + mM, and M is finitely generated (hence also
M/xAui),       so that by the above corollary M = xAui. If {ul,. .,u,} is
not minimal, so that, for example, {u2,. . . , un} already generates M
then (&,..., tin} generates fi, which is a contradiction.                 Hence
{ul, _. . , u,> is a minimal basis.
   (ii) If {ui, . . . , u,} is a minimal basis of M and we set Ui for the image
of ui in M, then u,,. . . , U, generate A, and are linearly independent
over k; indeed, otherwise some proper subset of {tii , . . . ,U,,,} would
be a basis of M, and then by (i) a proper subset of {ui,. . ,u,} would
generate M, which is a contradiction.
                                                                  - -
   (iii) Write Zij for the image in k of aij, so that Ui = Caijuj holds in
M. Since (aij) is the matrix transforming one basis of the vector space
I$ into another, its determinant is non-zero. Since det (aij) modm =
det (aij) # 0 it follows that det (aij) is a unit of A. By Cramer?s formula
the inverse matrix of (aij) exists as a matrix with entries in A. n
   We give another interesting application of NAK, the proof of which is
due to Vasconcelos [2].
\end{verbatim}
\end{proof}


\begin{theorem}[Printed Theorem 2.4]\label{mcrt-auto-theorem-2-4}\dcref{2.4}\source{matsumura-commutative-ring-theory:page-0024}\uses{}
\begin{verbatim}
Let A be a ring and M a finite A-module. If f:M -M             is
an A-linear map and f is surjective then f is also injective, and is thus
an automorphism     of M.
\end{verbatim}
\end{theorem}
\begin{proof}
\begin{verbatim}
Since f commutes with scalar multiplication       by elements of A, we
can view M as an A[X]-module        by setting X.m = f(m) for meM. Then by
assumption XM = M, so that by NAK there exists YEA[X] such that
(1 +XY)M    =O. Now        for uEKer(f)        we have 0 = (1 + XY)(u) =
u + Yf(u) = u, so that f is injective.     n
\end{verbatim}
\end{proof}


\begin{theorem}[Printed Theorem 2.5]\label{mcrt-auto-theorem-2-5}\dcref{2.5}\source{matsumura-commutative-ring-theory:page-0024}\uses{}
\begin{verbatim}
Let (A,m) be a local ring; then a projective module over A
is free (for the definition of projective module, see Appendix B, p. 277).
\end{verbatim}
\end{theorem}
\begin{proof}
\begin{verbatim}
This is easy when M is finite: choose a minimal basis ol,. . . , w, of
M and define a surjective map cp:F -+ M from the free module F =
Ae, O...@Ae, to M by cp(xaiei) = c a,~,,.. ?f1 we set K = Ker(cp) then, from
10       Commutative      rings and modules

the minimal     basis property,
         Capi     = 0 s- aiEm      for all i.
Thus K c mF. Because M is projective, there exists $: M -F     such that
F = $(M)@    K, and it follows that K = mK. On the other hand, K is a
quotient of F, therefore finite over A, so that K = 0 by NAK and F N M.
  The result was proved by Kaplansky [2] without the assumption that
M is finite. He proves first of all the following lemma, which holds for
any ring (possibly non-commutative).

Lemma 1. Let R be any ring, and F an R-module which is a direct sum of
countably generated submodules; if M is an arbitrary direct summand ofF
then M is also a direct sum of countably generated submodules.
Proof of Lemma 1. Suppose that F = M @ N, and that F = oloA E,, where
each E, is countably generated. By translinite induction, we construct a
well-ordered family {F,} of submodules of F with the following properties:
      (i) ifa<Bthen        F,cF,,
     (ii) F = ua F,,
   (iii) if c1is a limiting ordinal then F, = UBcaFB,
   (iv) F,, ,/F, is countably generated,
     (v) F,= M,@N,,          where M,= MnF,,   N,= NnF,,
   (vi) each F, is a direct sum of E, taken over a suitable subset of A.
We now construct such a family {FE). Firstly, set F, = (0). For an
ordinal cc,assume that F, has been defined for all ordinals p < CI. If a is
a limiting ordinal, set F, = Us<= F,. If tl is of the form a=,&+ 1, let Q1
be any one of the E, not contained in F, (if F, = F then the construction
stops at Fp). Take a set xrl, x12,. . . of generators of Qi, and decompose
xi1 into its M- and N-components; now let Q, be the direct sum of the
finitely many E, which are necessary to write each of these two components
in the decomposition F = GE,, and let xzl, xz2,. . . be generators of Qz.
Next decompose xi2 into its M- and N-components, let Q3 be the direct
sum of the finitely many E, needed to write these components, and let
x3l,     x32,*..  be generators of Q3. Then carry out the same procedure with
xzl, getting xbl, xa2,. . . , then do the same for x13. Carrying out the same
procedure for each of the xij in the order x~l,x~2,xz1,x~3,x22,x31,...
we get ?countably many elements xij. We let F, be the submodule of F
generated by F, and the xii, and this satisfies$all our requirements. This
gives the family {Fb} .
      Now M = u M,, with each M, a direct summand of F, and M,, 1 =) M,,
so that M, is also a direct summand of M,+l. Moreover,
         F,+ JFn = W,+ ,lMJOW,+            JNa),
?2       Modules                                                                11

and hence M,, ,/M, is countably generated. Thus we can write
        M a+l = M,@M&+,,         with Mi,, countably generated.
When a is a limit ordinal, since M, = up <aM,, we set Mh = 0. Then finally
we can write
        M = @ Mi with Mb (at most) countably generated. n
   Of course aolfree module satisfies the assumption of Lemma 1, so that,
in particular, we see that any projective module is a direct sum of countably
generated projective modules. Thus in the proof of Theorem 2.5 we can
assume that M is countably generated.
Lemma 2. Let M be a projective module over a local ring A, and XEM.
Then there exists a direct summand of M containing x which is a free
module.
Proof of Lemma 2. We write M as a direct summand of a free module
F = M @ N. Choose a basis B = {ui}ipl of F such that the given element
x has the minimum         possible number of non-zero coordinates when
expressed in this basis. Then if x = ulal + ... + ~,a,, with 0 # Ui~A, we have
          ai4 1 Aaj for i = 1, 2,. . . , n;
             J#i
indeed, if, say, a, = 1; - i biai then x = C;- i (ui + u,b,)q, which contradicts
the choice of B. Now set ui = yi + zi with y,gM and z,EN; then
          X=&Ui=piyi.
If we write yi = I;= 1cijuj + ti, with ti linear combinations of elements of
B other than ui, . . . , u,, we get relations a, = CJ= I ajcji, and, hence, in view
of what we have seen above, we must have
          1 - ci+m       and cijEm for i #j.
It follows that the matrix (cij) has an inverse (this can be seen from the
fact that the determinant is s 1 mod m, or by elimination). Thus replacing
Ul,..., unbyy,,...,     y, in B, we still have a basis of F. Hence, F, = CYiA is
a direct summand of F, and hence also of M, and satisfies all the
requirements of Lemma 2. n
   To prove the theorem, let M be a countably generated projective module,
M=co,A+o,~+....              By Lemma 2, there exists a free module F, such
that o,EF,, and M = F, GM,, where M, is a projective module. Let w;
be the M ,-component of o2 in the decomposition M = F I @M 1, and take
a free module F, such that o;EF~ and M, = F, GM,, where M, is a
Projective module. Let oj be the M,-component                 of wj in M = F, @
F, 0 M,; proceeding in the same way, we get
         M=F,@F,@...,
so that M is a free module.       n
12       Commutative    rings and modules

   We say that an A-module M # 0 is a simple module if it has no
submodules other than 0 and M itself For any 0 # UEM, we then have
M = Ao. Now Ao N A/ann(o), but in order for this to be simple, arm(o)
must be a maximal ideal of A. Hence, any simple A-module is isomorphic
to A/m with m a maximal ideal, and conversely an A-module of this form
is simple. If M is an A-module, a chain

of submodules of M is called a composition series of M if every Mi/Mi+ 1 is
simple; r is called the length of the composition series. If a composition series
of M exists, its length is an invariant of M independent of the choice of
composition series. More precisely, if M has a composition series of length
r, and if MI N, I>... =3N, is a strictly descending chain of sub-
modules, then we have s < r. This invariance corresponds to part of the
basic JordanHolder        theorem in group theory, but it can easily be proved
on its own by induction, and the reader might like to do this as an exercise.
The length of a composition series of M is called the length of M, and written
2(M); if M does not have a composition series we set l(M) = co. A necessary
and sufficient condition for the existence of a composition series of M is that
the submodules of M should satisfy both the ascending and descending
chain conditions (for which see 93). In general, if N c M is a submodule,
we have
         l(M) = l(N) + l(M/N).
IfO+M,     -M2     -..?   - M, -+ 0 is an exact sequence of A-modules and
each Mi has finite length then
         i$l (- 1)?NMJ= 0.
   If m is a maximal ideal of A and is finitely generated over A then
1(A/m?) < co. In fact,
          1(A/m?) = 1(A/m) + l(m/m2) + ... + l(mv-l/mv);
now each m?/m i+l is a finite-dimensional        vector space over the field
k = A/m, and since its A-submodules are the same thing as its vector
subspaces, ,(mi/mifl)     is equal to the dimension of m?/m?+ l as k-vector
space. (This shows that A/m? is an Artinian ring, see $3.)
   Considering /(A/m?) for all v, we get a function of v which is intimately
related to the ring structure of A, and which also plays a role in the resolution
of singularities in algebraic and complex analytic geometry; this is studied in
Chapter 5.
   We say that an A-module M is of finite presentation if there exists an
exact sequence of the form
         A*-Aq+M+O.
?2         Modules                                                                  13

This means that M can be generated by 4 elements ol,. . . , wq in such a way
that the module R = {(a,,. .., a,)~A~lCa~o~ = 0) of linear relations
holding between the oi can be generated by p elements.
\end{verbatim}
\end{proof}


\begin{theorem}[Printed Theorem 2.6]\label{mcrt-auto-theorem-2-6}\dcref{2.6}\source{matsumura-commutative-ring-theory:page-0028}\uses{}
\begin{verbatim}
Let A be a ring, and suppose that M is an A-module               of finite
presentation. If
           O+K-N-M-0
is an exact sequence and N is finitely generated then so is K.
\end{verbatim}
\end{theorem}
\begin{proof}
\begin{verbatim}
By assumption     there exists an exact sequence of the form
L, 2  L, L M -+ 0, where L, and L2 are free modules of finite rank.
From this we get the following commutative diagram (see Appendix B):




   If we write N = At, + ... + A<,,, then there exist ui~L, such that
cp(li)=f(ui).  Set 5: = ti -a(~,); then (p(&)=O, so that we can write
I$ = @(vi) with ~],EK. Let us now prove that
       K=j3(L2)+4,+..++Aq,.
For any ~EK, set $(r]) = xaiti;        then
        $(S - Cailli) = Cai(ti - 5;) = cc(Cuiui)3
and since 0 = (~a(1 aiui) = f(C a,~,.), we can write cuiui = g(u) with UEL,.
Now
        IcIPf?) = crS(u)= 4x Vi) = $411- 1 uiUi),
SOthat q = /I(u) + 1 uiqi, and this proves our assertion.           w
\end{verbatim}
\end{proof}


\begin{theorem}[Printed Theorem 3.1]\label{mcrt-auto-theorem-3-1}\dcref{3.1}\source{matsumura-commutative-ring-theory:page-0030}\uses{}
\begin{verbatim}
Let A be a ring and M an A-module.
  (i) Let M? c M be a submodule and 9: M + M/M? the natural map. If
N, and N, are submodules of M such that N, c N,, N, n M? = N, n M?
and cp(N,) = cp(N,) then N, = Nz.
   (ii) Let 0 -+ M? -M      + M? + 0 be an exact sequence of A-modules;
if M? and M? are both Noetherian (or both Artinian), then so is M.
   (iii) Let M be a finite A-module; then if A is Noetherian (or Artinian),
so is M.
\end{verbatim}
\end{theorem}
\begin{proof}
\begin{verbatim}
(i) is easy, and we leave it to the reader.
   (ii) is obtained by applying (i) to an ascending (respectively descending)
chain of submodules of M.
   (iii) If M is generated by n elements then it is a quotient of the free module
A?, so that it is enough to show that A? is Noetherian (respectively Artinian).
However, this is clear from (ii) by induction on n. w
   For a module M, it is equivalent to say that M has both the a.c.c. and
the d.c.c., or that M has finite length. Indeed, if I(M) < 00 then I(M i) < 1(M,)
for any two distinct submodules MI c M2 c M, so that the two chain
conditions are clear. Conversely, if M has the d.c.c. then we let M, be a
minimal non-zero submodule of M, let M, be a minimal element among all
submodules of M strictly containing M, , and proceed in the same way to
obtain an ascending chain 0 = MO c M 1 c M, c ...; if M also has the a.c.c.
then this chain must stop by arriving at M, so that M has a composition
series,
   Every submodule of the Z-module Z is of the form nZ, so that Z is
Noetherian, but not Artinian. Let p be a prime, and write W for the Z-
module of rational numbers whose denominator is a power of p; then
the Z-module W/Z is not Noetherian, but it is Artinian, since every proper
submodule of W/Z is either 0 or is generated by p-? for n = 1,2,. . . . This
shows that the a.c.c. and d.c.c. for modules are independent conditions,
but this is not the case for rings, as shown by the following result.
16         Commutative        rings and modules
\end{verbatim}
\end{proof}


\begin{theorem}[Printed Theorem 3.2]\label{mcrt-auto-theorem-3-2}\dcref{3.2}\source{matsumura-commutative-ring-theory:page-0031}\uses{}
\begin{verbatim}
Y. Akizuki). An Artinian ring is Noetherian.
\end{verbatim}
\end{theorem}
\begin{proof}
\begin{verbatim}
Let A be an Artinian ring. It is sufficient to prove that A has finite
length as an A-module. First of all, A has only finitely many maximal ideals.
Indeed,ifp,,p,,...     is an infinite set of distinct maximal ideals then it is easy
 to see that p1 3 plpz =) p1p2p3?. is an infinite descending chain of ideals,
 which contradicts the assumption. Thus, we let pl, pz,. . . , pr be all the
 maximal ideals of A and set I = p1p2.. . p, = rad (A). The descending chain
 III2    2 ?.. stops after finitely many steps, so that there is an s such that
 I? = Z?+i. If we set (0:I?) = J then
           (J:Z) = ((0:P):I) = (O:r+l) = J;
 let?s prove that J = A. By contradiction, suppose that J # A; then there exists
 an ideal J? which is minimal among all ideals strictly bigger than J. For any
 XEJ? - J we have J? = Ax + J. Now I = rad (A) and J # J?, so that by NAK
J? # Ix + J, and hence by minimality of J? we have Ix + J = J, and this gives
 Ix c J. Thus XE(J:I) = J, which is a contradiction. Therefore J = A, so that
 I? = 0. Now consider the chain of ideals



Let M and Mpi be any two consecutive terms in this chain; then MIMpi is a
vector space over the field A/p,, and since it is Artinian, it must be tinite-
dimensional. Hence, l(M/Mp,) < co, and therefore the sum 1(A) of these
terms is also finite. a
\end{verbatim}
\end{proof}


\begin{theorem}[Printed Theorem 3.3]\label{mcrt-auto-theorem-3-3}\dcref{3.3}\source{matsumura-commutative-ring-theory:page-0031}\uses{}
\begin{verbatim}
If A is Noetherian then so are A[X] and A[Xj.
\end{verbatim}
\end{theorem}
\begin{proof}
\begin{verbatim}
The statement for A[X] is the well-known Hilbert basis theorem
(see, for example Lang, Algebra, or [AM], p. 81), and we omit the proof.
We now briefly run through the proof for A[Xj.              Set B = A[XJ, and
let I be an ideal of B; we will prove that I is finitely generated. Write
Z(r) for the ideal of A formed by the leading coefficients a, of f = a,X? +
a *+1 xr+1 + . . ? as f runs through In X?B; then we have
          I(0) c 1(l) C 1(2) C . . . .
Since A is Noetherian, there is an s such that I(s) = I(s + 1) = ?.; moreover,
each Z(i) is finitely generated. For each i with 0 < i < s we take finitely many
elements aivEA generating Z(i), and choose giVEI nXiB having Ui, as the
coefficient of Xi. These giV now generate 1. Indeed, for f~1 we can take
a linear combination        go of the g,,? with coefficients in A such that
?3       Chain conditions                                                      17

f -geEZn       XB, then take a linear combination   gi of the glV with
coefficients in A such that f - go - g,EZnX?B,    and proceeding in the
same way we get
           f-go-g1     --...-g,ElnX?+?B.
Now Z(s + 1) = Z(s), so we can take a linear combination       gs+ i of the Xg,,
with coefficients in A such that
        f-g90-g1--~~-gs+lEznX?+2B.
We now proceed in the same way to get gs+ *,. . . For i < s, each gi
is a linear combination        of the giy with coefficients in A, and, for
i > s, a combination      of the elements XiPsgsV. For each i > s we write
gi = ~vaivXi-?g,,,     and then for each v we set h,, = ~im,saivXi~S; h, is an
element of B, and
         f=gO+...+gs-l+ChYgSY.             n .

   A ring A[b,,. . .,b,] which is finitely generated as a ring over a
Noetherian ring A is a quotient of a polynomial ring A[X,, . . . ,X,1, and
so by the Hilbert basis theorem is again Noetherian. We now give some
other criteria for a ring to be Noetherian.
\end{verbatim}
\end{proof}


\begin{theorem}[Printed Theorem 3.4]\label{mcrt-auto-theorem-3-4}\dcref{3.4}\source{matsumura-commutative-ring-theory:page-0032}\uses{}
\begin{verbatim}
I. S. Cohen). If all the prime ideals of a ring A are finitely
generated then A is Noetherian.
\end{verbatim}
\end{theorem}
\begin{proof}
\begin{verbatim}
Write r for the set of ideals of A which are not finitely generated.
If r # fa then by Zorn?s lemma r contains a maximal element I. Then I is
not a prime ideal, so that there are elements x, YEA with x$Z, ~$1 but
xy~l. Now Z+Ay is bigger than I, and hence is finitely generated, so
that we can choose ui,. . .,u,EZ such that
           Z+Ay=(u,      ,..., u,,y).
Moreover,      Z:y= {aEAlayeZ}        contains x, and is thus bigger than
I, so it has a finite system of generators             (ur,. .,u,}. Finally, it is
easy to check that Z = (ui ,..., y,u,y,..      .,v,y); hence, Z$r, which is a
contradiction.    Therefore r = 121. n
\end{verbatim}
\end{proof}


\begin{theorem}[Printed Theorem 3.5]\label{mcrt-auto-theorem-3-5}\dcref{3.5}\source{matsumura-commutative-ring-theory:page-0032}\uses{}
\begin{verbatim}
Let A be a ring and M an A-module. Then if M is a
Noetherian module, A/ann(M) is a Noetherian ring.
\end{verbatim}
\end{theorem}
\begin{proof}
\begin{verbatim}
If we set A = Alann (M) and view M as an A-module, then the
submodules of M as an A-module or A-module coincide, so that M is
also Noetherian as an A-module. We can thus replace A by 2, and then
arm(M) = (0). Now letting M = Ao, + ... + Aw,, we can embed A in M?
by means of the map a++(au r, . . . , aw,). By Theorem 1, M? is a Noetherian
module, so that its submodule A is also Noetherian. (This theorem can
be expressed by saying that a ring having a faithful Noetherian module
is Noetherian.)  n
18       Commutative    rings and modules
\end{verbatim}
\end{proof}


\begin{theorem}[Printed Theorem 3.6]\label{mcrt-auto-theorem-3-6}\dcref{3.6}\source{matsumura-commutative-ring-theory:page-0033}\uses{}
\begin{verbatim}
E. Formanek Cl]). Let A be a ring, and B an A-module
which is finitely generated and faithful over A. Assume that the set of
submodules of B of the form 1B with I an ideal of A satisfies the a.c.c.;
then A is Noetherian.
\end{verbatim}
\end{theorem}
\begin{proof}
\begin{verbatim}
It will be enough to show that B is a Noetherian A-module. By
contradiction, suppose that it is not; then the set
             IB I is an ideal of A and B/IB is                               /
           i    1non-Noetherian   as A-module   i
contains (0) and so is non-empty, so that by assumption it contains a
maximal element. Let ZB be one such maximal element; then replacing B
by BJIB and A by A/ann(B/IB)         we see that we can assume that B is a
non-Noetherian      A-module, but for any non-zero ideal I of A the quotient
BjZB is Noetherian.
   Next we set
          I = (NIN is a submodule of B and B/N is a faithful A-module}.
If B = Ab, + *a. + Ab, then for a submodule N of B,
         NElYoVaaA-0,          {ab, ,..., ab,} $N.
From this, one sees at once that Zorn?s lemma applies to I; hence there
exists a maximal element N, of I. If B/N,, is Noetherian then A is a
Noetherian     ring, and thus B is Noetherian,         which contradicts our
hypothesis. It follows that on replacing B by B/N, we arrive at a module
B with the following properties:
   (1) B is non-Noetherian    as an A-module;
   (2) for any ideal I # (0) of A, B/IB is Noetherian;
   (3) for any submodule N # (0) of B, B/N is not faithful as an A-module.
   Now let N be any non-zero submodule of B. By (3) there is an element
aeA with a # 0 such that a(B/N) = 0, that is such that aB c N. By (2)
B/aB is a Noetherian module, so that N/aB is finitely generated; but since
B is finitely generated so is aB, and hence N itself is finitely generated.
Thus, B is a Noetherian module, which contradicts (1). n
   As a corollary of this theorem we get the following result.
\end{verbatim}
\end{proof}


\begin{theorem}[Printed Theorem 3.7]\label{mcrt-auto-theorem-3-7}\dcref{3.7}\source{matsumura-commutative-ring-theory:page-0033}\uses{}
\begin{verbatim}
(i) (Eakin-Nagata    theorem). Let B be a Noetherian ring, and Aa subring
of B such that B is finite over A; then A is also a Noetherian ring.
    (ii) Let B be a non-commutative    ring whose right ideals have the a.c.c.,
and let A be a commutative subring of B. If B is finitely generated as a
left A-module then A is a Noetherian ring.
    (iii) Let B be a non-commutative     ring whose two-sided ideals have the
a.c.c., and let A be a subring contained in the centre of B; if B is finitely
generated as an A-module then A is a Noetherian ring.
93        Chain conditions                                                          19
\end{verbatim}
\end{theorem}
\begin{proof}
\begin{verbatim}
B has a unit, so is faithful as an A-module.              Hence it is enough to
apply the previous theorem.      n
\end{verbatim}
\end{proof}


\begin{theorem}[Printed Theorem 4.1]\label{mcrt-auto-theorem-4-1}\dcref{4.1}\source{matsumura-commutative-ring-theory:page-0037}\uses{}
\begin{verbatim}
(i) All the ideals of A, are of the form IA,, with I an ideal of A.
   (ii) Every prime ideal of A, is of the form pA, with p a prime ideal of
A disjoint from S, and conversely, pA, is prime in A, for every such p;
exactly the same holds for primary ideals.
\end{verbatim}
\end{theorem}
\begin{proof}
\begin{verbatim}
(i) If J is an ideal of A,, set I = .Zn A. If x = a/sEJ then
x.f(s) = f(a)~J, so that ael, and then x = (l/s)*f(a)~ZA~. The converse
inclusion IA, c J is obvious, so that J = IA,.
   (ii) If P is a prime ideal of A, and we set p = Pn A, then p is a prime
ideal of A, and from the above proof P = pA,. Moreover, since P does
not contain units of A,, we have p nS = a. Conversely, if p is a prime
ideal of A disjoint from S then
        ab
        -.-EPA,     with   s,teS=>rab~p      for some    r~s,
        s t
and since r$p we must have aEp or bEp, so that a/s or bItEpA,. One
also sees easily that l$pAs, so that pAs is a prime ideal of A,.
   For primary ideals the argument is exactly the same: if p is a primary
ideal of A disjoint from S and if rabep with reS, then since no power
of r is in p we have abEp. From this we get either u/s~pA, or (b/t)?EpAs
for some n. n
\end{verbatim}
\end{proof}


\begin{theorem}[Printed Theorem 4.2]\label{mcrt-auto-theorem-4-2}\dcref{4.2}\source{matsumura-commutative-ring-theory:page-0038}\uses{}
\begin{verbatim}
Localisation commutes with passing to quotients by ideals.
More precisely, let A be a ring, S c A a multiplicative set, I an ideal of A
and S the image of S in A/I; then
          AdzAs = (4%
\end{verbatim}
\end{theorem}
\begin{proof}
\begin{verbatim}
Both sides have the universal property for ring homomorphisms
g:A -+ C such that
      (1) every element of S maps to a unit of C,
and (2) every element of Z maps to 0;
the isomorphism follows by the uniqueness of the solution to a universal
mapping problem. In concrete terms the isomorphism is given by
          afs mod IA,++@,        where ti=u+Z,    S=s+Z.    n

   In particular, if p is a prime ideal of A then
          A&A, 21(A/P),,.
The left-hand side is the re.sidue field of the local ring A,, whereas the
right-hand side is the field of fractions of the integral domain A/p. This
field is written K(P) and called the residue field of p.
\end{verbatim}
\end{proof}


\begin{theorem}[Printed Theorem 4.3]\label{mcrt-auto-theorem-4-3}\dcref{4.3}\source{matsumura-commutative-ring-theory:page-0038}\uses{}
\begin{verbatim}
Let A be a ring, S c A a multiplicative             set, and
f:A -A,     the canonical map. If B is a ring, with ring homomorphisms
g:A --+B and h:B -A,      satisfying
     (1) f= 4,
and (2) for every b~Z3 there exists s~S such that g(s).kg(A),      then A,
can also be regarded as a ring of fractions of B. More precisely,
         A, = B,(,, = B,, where T= {teBlh(t)r is~- a unit of A,).
\end{verbatim}
\end{theorem}
\begin{proof}
\begin{verbatim}
We can factorise h as B -+B,        -+ A,; write CI:B=-       A, for
24       Prime ideals

the second of these maps. Now g(S) c T, so that the composite A -+ B
-+ B, factorises as A --+ A, -B,;   write fi:As -B,   for the second of
these maps. Then

so that E/J = 1, the identity map of A,. Moreover by assumption, for kB
there exist aEA and SES such that bg(s) = g(a). Hence, fl(a/s) = g(a)/g(s) =
b/l. In particular for tcT, if we take UEA, such that t/l =b(u) then
u = @p(u)= a(t/l) = h(t), so that u is a unit of A,. Hence, b/t = fl(a/s)b(u-?),
and /3 is surjective. Thus CI and fi are mutually inverse, giving an
isomorphism A, N B,. The fact that A, 2: Bqcs, can be proved similarly.
(Alternatively, this follows since T is the saturation of the multiplicative
set g(S). The reader should check this fo,r himself.) w

Corollary 1. If p is a prime ideal of A, S = A - p and B satisfies the
conditions of the theorem, then setting P = pA, n B we have A, = B,.
Proof. Under these circumstances the T in the theorem is exactly B - P.
Corollary 2. Let S c A be a multiplicative    set not containing any zero-
divisors; then A can be viewed as a subring of A,, and for any intermediate
ring A c B c A,, the ring A, is a ring of fractions of B.
Corollary 3. If S and T are two multiplicative subsets of A with S c T,
then writing T? for the image of T in A,, we have (AS)T, = A,.
Corollary 4. If S c A is a multiplicative  set and P is a prime ideal of A
disjoint from S then (AS)PA,7
                            = A,. In particular if P c Q are prime ideals of
A, then
         (A&Ap = A,.
\end{verbatim}
\end{proof}


\begin{theorem}[Printed Theorem 4.4]\label{mcrt-auto-theorem-4-4}\dcref{4.4}\source{matsumura-commutative-ring-theory:page-0041}\uses{}
\begin{verbatim}
M, E M QA As.
\end{verbatim}
\end{theorem}
\begin{proof}
\begin{verbatim}
The map M x As -M,             defined by (m, a/s)wam/s      is A-bilinear,
so that there exists a linear map cz:M @As ----) M, such that a(m @ a/s) =
am/s. Conversely we can define 8: MS -+ M @ As by /?(m/s) = m @ (l/s);
indeed, if m/s = m?Js? then ts?m = tsm? for some t ES, and so
         m @ (l/s) = m @ (ts?ltss?) = ts?m @ (l/tss?) = tsm? @3(l/tss?)
                   = m? Q ( 1is?).
Now it is easy to check that CI and /3 are mutually inverse As-linear
maps. Hence, MS and M @AAs are isomorphic as As-modules.
\end{verbatim}
\end{proof}


\begin{theorem}[Printed Theorem 4.5]\label{mcrt-auto-theorem-4-5}\dcref{4.5}\source{matsumura-commutative-ring-theory:page-0041}\uses{}
\begin{verbatim}
M-M,        is an exact (covariant) functor from the category
of A-modules to the category of As-modules. That is, for a morphism of
A-modules f:M - N there is a morphism fs: MS --+ N, of As-modules
such that
         (id), = id (where id is the identity map of M or Ms),
         (Sf)s = Ysfsi
 and such that an exact sequence 0 +M?-M-M?+0                   goes into an
exact sequence 0 + Mk - MS --+ Mg --t 0.
\end{verbatim}
\end{theorem}
\begin{proof}
\begin{verbatim}
To prove the exactness of 0 + M; -MS       on the last line, view M?
 as a submodule of M; then for XEM? and SES,
            x/s=0 in M,otx     =0 for some tES
                          ox/s =0 in Ms,
as required. The remaining assertions follow from the properties of the
 tensor product (see Appendix A) and from the previous theorem. (Of
course they can easily be proved directly.)    n
    It follows from this that localisation commutes with @ and with Tor,
and we will treat all this together in the section on flatness in 47.
    Let A be a ring, M an A-module and p~Spec A. There are at least
two interpretations of what it should mean that some property of A or M
holds ?locally at p?. Namely, this could mean that A, (or MJ has the
property, or it could mean that A, (or M,,) has the property for all q in
some neighbourhood U of p in Spec A. The first of these is more commonly
used, but there are cases when the two interpretations coincide. In any
case, we now prove a number of theorems which assert that a local property
implies a global one.
?4      Localisation   and Spec of a ring                                27
\end{verbatim}
\end{proof}


\begin{theorem}[Printed Theorem 4.6]\label{mcrt-auto-theorem-4-6}\dcref{4.6}\source{matsumura-commutative-ring-theory:page-0042}\uses{}
\begin{verbatim}
Let A be a ring, M an A-module and XEM. If x = 0 in M,
      for every maximal ideal p of A, then x = 0.
\end{verbatim}
\end{theorem}
\begin{proof}
\begin{verbatim}
x = 0 in M,osx = 0 for some .SEA - poann (x) + p.
    . However, if 1 $ann(x) then by Theorem 1.1, there must exist a maximal
      ideal containing ann (x). Therefore 1Eann (x), that is x = 0. H
\end{verbatim}
\end{proof}


\begin{theorem}[Printed Theorem 4.8]\label{mcrt-auto-theorem-4-8}\dcref{4.8}\source{matsumura-commutative-ring-theory:page-0042}\uses{}
\begin{verbatim}
Let A be a ring and M a finite A-module. If M OAIc(m) = 0
      for every maximal ideal m then M = 0.
\end{verbatim}
\end{theorem}
\begin{proof}
\begin{verbatim}
I      = A,/mA,,   so that M @~c(m) = M,ImM,,,,   and by NAK
      (Theorem 2.2), M 0 JC(~) = OoM,,, = 0. Thus the assertion follows from
\end{verbatim}
\end{proof}


\begin{theorem}[Printed Theorem 4.9]\label{mcrt-auto-theorem-4-9}\dcref{4.9}\source{matsumura-commutative-ring-theory:page-0042}\uses{}
\begin{verbatim}
Let f: A --t B be a homomorphism     of rings, and M a finite
      B-modules; if M OAic(p) = 0 for every p&pec A, then M = 0.
\end{verbatim}
\end{theorem}
\begin{proof}
\begin{verbatim}
If M # 0 then by Theorem 6 there is a maximal ideal P of B such
      that M, # 0, so that by NAK, M,/PM,     # 0. If we now set p = Pn A then,
      since pM, c PM,, we have M,/pM,       # 0. Set T = B -P and S = A - p;
      then the localisation Ms = M, of M as an A-module and the localisation
      Mf,s, of M as a B-module coincide (both of them are {m/slm # M, SES} ).
      We have f(S) c T, so that
              MP = MT = O-f/&       = (Mph,
      and hence
                MP/PM,J = (M,IW,h      = Wf 0, ~P))T;
      it follows that M OA+) # 0.       w
28         Prime ideals
\end{verbatim}
\end{proof}


\begin{theorem}[Printed Theorem 4.10]\label{mcrt-auto-theorem-4-10}\dcref{4.10}\source{matsumura-commutative-ring-theory:page-0043}\uses{}
\begin{verbatim}
Let A be a ring and M a finite .4-module.
  (i) For any non-negative integer r set
           U, = {pESpec AIM, can be generated over A, by r elements};
then U, is an open subset of Spec A.
   (ii) If M is a module of finite presentation then the set
           U, = {p ESpec A IM, is a free AD-module}
is open in Spec A.
\end{verbatim}
\end{theorem}
\begin{proof}
\begin{verbatim}
(i) Suppose that A, = M,o, +. . + A,w,. Each oi is of the form
Oi = mi/si with SiEA - p and mgM, but since si is a unit of A, we can
replace Wi by mi, and so assume that oi is (the image in M, of) an element
of M. Define a linear map cp:A? --+ M from the direct sum of r copies of
A to M by (a,, . . , a,)++caiwi,  and write C for the cokernel of cp.Localising
the exact sequence A?-+ M -+ C+O at a prime ideal q, we get an
exact sequence
           A;-M,-C,+O,
and when q = p we get C, = 0. C is a quotient of M, so is finitely generated,
so that the support Supp(C) is a closed set, and hence there is an open
neighbourhood I/ of p such that C, = 0 for qE V. This means that V c U,.
(In short, if Oi,. . . , O,EM generate M, at p then they generate M, for all q
in a neighbourhood of p.)
   (ii) Suppose that M, is a free A,-module, and let wi, . . . , w, be a basis.
As in (i) there is no loss of generality in assuming that oi~M. Moreover,
if we choose a suitable D(a) as a neighbourhood of p in Spec A, ml,. . . , O,
generate M, for every qeD(a). Thus, replacing A by A, and M by M,
we can assume that the elements w 1,. . . , 0, satisfy M, = C A ,a, for every
prime ideal q of A. Then by Theorem 6,
          M/~Ao,=O,           thatisM=Aw,+...+Aw,.
(We think of replacing A by A, as shrinking Spec A down to the
neighbourhood o(a) of p.) Now, defining cp:A? -M         as above, and letting
K be its kernel, we get the exact sequence
          ,0-+K--+A?---+M+0;
moreover, K, = 0. By Theorem 2.6, K is finitely generated, so that applying
?4         Localisation    and Spec     qf a ring                                           29

(i) with   r = 0, we have that      K, = 0 for every q in a neighbourhood                V of
p; this gives (A,)? N M,,      SOthat I/ c U,.        w
\end{verbatim}
\end{proof}


\begin{theorem}[Printed Theorem 5.1]\label{mcrt-auto-theorem-5-1}\dcref{5.1}\source{matsumura-commutative-ring-theory:page-0046}\uses{}
\begin{verbatim}
Let k be a field, L an algebraic extension of k and
al,...,  a,&      then
  (i) k[a,, . . . ,~l =kh,...,q,).
  (ii) Write q:k[X,,      . . . ,X,] --+ k(a,, . . , a,) for the homomorphism
over k which maps Xi to cci; then Ker cp is the maximal ideal generated
by n elements of the form fl(X,), f2(X1, X,),. . . ,f,,(X,, . . . ,X,), where
each .fi can be taken to be manic in Xi.
\end{verbatim}
\end{theorem}
\begin{proof}
\begin{verbatim}
Let fi(X) be the minimal polynomial of a1 over k; then (fl(X,)) is a
32         Prime ideals

maximal ideal of k[X,], so that k[aJ z k[X,]/(fl(X,)) is a field, and hence
k[~~] = @a,). Now let (p2(X) be the minimal polynomial of a2 over k(cr,);
then since k(cr,) = k[~r], the coefficients of pDz can be expressed as
polynomials in a1, and there is a polynomial f,~k[X,,  X,], manic in X,,
such that (p2(X2) = .fi(ai, X,). Thus
           kC~,>4 = 4dC~l                                                   =kt~l>Q ?v k(a,)Cxzl/(fi(crl,x,)).
Proceeding         in     the                          same way, for        l<i<n                          there     is   an
fi(X,,...,Xi)EkCXl,...,                                X,], manic in Xi, such that
           4a 1). . . ) ai] = k(a,). . . ) ai)
                            -k(cx,,...,cci-,)[xil/(.l;(crl,...,ai-,,Xi)).


Now if P(X)Ek[X,,.            . ,X,] is in the kernel of q, we have q(P) =
P(a 1,..., sr,)=O, so that P(a,,. . .,a,-r,X,,)       is divisible by f,,(c~~,. ,
CI,,-~ ,X,); dividing       P(X,, . , X,) as a polynomial           in X, by the
manic polynomial f,,(X, , . . , X,) and letting R,(X I , . . . , X,) be the remain-
der, we can write P = QJ,, + R,, with Rn(al,. . . ,cI,_ l,Xn) = 0. Similarly,
dividing R,(X, , . . . , X,) as a polynomial in X, - 1 by f,, - I(X 1, . . , X, - 1) and
letting R, - ,(X1,. . . , X,) be the remainder, we get
           R,=Q,-If,-I                            +k1,
with

           L,(a       1,...,a,-2,Xn-1,Xn)=0.
Proceeding in the same way we get P = 1 Qifi + R, with R(X,, . . . , X,) = 0;
that is R = 0 and P = 1 Qifi, so that Kercp=(f,,f,   ,..., f,,). n
   The following theorem can be regarded as a converse of Theorem 1, (i).
\end{verbatim}
\end{proof}


\begin{theorem}[Printed Theorem 5.3]\label{mcrt-auto-theorem-5-3}\dcref{5.3}\source{matsumura-commutative-ring-theory:page-0048}\uses{}
\begin{verbatim}
Let k be a field, and let m be any maximal ideal of the
   polynomial ring k[X, , . . , X,]; then the residue class field k[X,, . . . , X,1/m
   is algebraic over k. Hence m can be generated by n elements, and in
   particular if k is algebraically closed then m is of the form m =
   (Xl-al,...,       X, - a,) for aiEk.
\end{verbatim}
\end{theorem}
\begin{proof}
\begin{verbatim}
Set k[X,, . . . , X,1/m = K, and write. cli for the image of Xi in K; then
   K = k[a, , . . . , a,]. By the previous theorem, since K is a field it is algebraic
   over k, and then by Theorem 1, (ii), m is generated by n elements. If k is
   algebraically closed then k = K, so that each Xi is congruent modulo m to
   some aiEk; then (Xi -al,...,X,-a,J               cm.     On the other hand
   (X1-%...,       X, - a,) is obviously a maximal ideal, so that equality must
   hold. a
     Let k be a field and k its algebraic closure. Suppose that CDc
   kCx i,. . . ,X,1 is a subset. An n-tuple CI= (a,, . . . , a,) of elements a,Ekis an
   algebraic zero of @ if it satisfies f(a) = 0 for every f(X)&.
\end{verbatim}
\end{proof}


\begin{theorem}[Printed Theorem 5.4]\label{mcrt-auto-theorem-5-4}\dcref{5.4}\source{matsumura-commutative-ring-theory:page-0048}\uses{}
\begin{verbatim}
The Hilbert Nullstellensatz).
     (i) If @ is a subset of k[X l,. . . ,X,1 which does not have any algebraic
 ~ zeros then the ideal generated by @ contains 1.
         (ii) Given a subset CDof k[X,, . . . ,X,] and an element fek[X,, . . .,X,1,
      suppose that f vanishes at every algebraic zero of Q. Then some power off
  : belongs to the ideal generated by a?, that is there exist v > 0,
      BiEk[X,,     , , X,] and hi& such that f? = xgihi.
     ?Proof. (i) Let Z be the ideal generated by @; if 141 then there exists a
      maximal ideal m containing I. By the previous theorem, k[X, , . . . , X,1/m
 ?i: is algebraic over k, so that it has a k-linear isomorphic embedding 8 into
 ? k. If we set &Xi mod m) = ai then for all g(X)Em we have 0 = @g(X)) =
tj,.
34            Prime ideals

gb,     9 *. ., a,,), and
                    therefore c(= (~1~).. . , a,) is an algebraic zero of m, and
hence also of @. This contradicts the hypothesis; hence, 1~1.
   (ii) Inside k[X, , . . . , X,, Y] we consider the set 0 u (1 - Yj(X)}; then this
set has no algebraic zeros, so that by (i) it generates the ideal (1). Therefore
there exists a relation of the form
         1 = 1 Pi@, Yh(W + Q(X, VU - Yf(X)),
with h,(X)&. This is an identity in Xi,. . . , X, and Y, so that it still holds if
we substitute Y = f(X)-?. Hence we have
         l = 1 pi(x, f - ?)hi(X),
SO that multiplying    through by a suitable power of f and clearing
denominators      gives f? = c g,(X)h,(X),      with g,gk[X, , , . , X,]      and
high.        n
\end{verbatim}
\end{theorem}
\begin{proof}
\notready
\end{proof}


\begin{theorem}[Printed Theorem 5.5]\label{mcrt-auto-theorem-5-5}\dcref{5.5}\source{matsumura-commutative-ring-theory:page-0049}\uses{}
\begin{verbatim}
Let k be a field, A a ring which is finitely generated over k, and
I a proper ideal of A; then the radical of I is the intersection of all maximal
ideals containing I, that is ,/I = nlcmm.
\end{verbatim}
\end{theorem}
\begin{proof}
\begin{verbatim}
Let A = k[a,, . . . , a,], so that A is a quotient of k[X,,.               .,X,1.
Considering the inverse image of I in k[X] reduces to the case A = k[X],
and the assertion follows from Theorem 4, (ii). n
     Compared to the result ,/I = nIcP P proved in $1, the conclusion of
Theorem 5 is much stronger. It is equivalent to the condition on a ring
that every prime ideal P should be expressible as an intersection of maximal
ideals. Rings for which this holds are called Hilbert rings or Jacobson
rings, and they have been studied independently by 0. Goldmann [l] and
W. Krull [7]. See also Kaplansky [K] and Bourbaki [B-5].
\end{verbatim}
\end{proof}


\begin{theorem}[Printed Theorem 5.6]\label{mcrt-auto-theorem-5-6}\dcref{5.6}\source{matsumura-commutative-ring-theory:page-0049}\uses{}
\begin{verbatim}
Let k be a field and A an integral domain which is finitely
generated over k; then
             dim A = tr.deg, A.
\end{verbatim}
\end{theorem}
\begin{proof}
\begin{verbatim}
Let A = k[X 1,. . . , X,1/P, and set r = tr.deg, A. To prove that
r > dim A it is enough to show that if P and Q are prime ideals of
k[X] = k[X,, . . . , X,] with Q 3 P and Q #P then
             tr. deg, k [Xl/Q < tr. deg, k[X]/P.
The k-algebra homomorphism           k [Xl/P -+ k [Xl/Q is onto, so that tr. deg,
k(X)/Q d tr. deg,k [Xl/P is obvious. Suppose that equality holds. Let
k[X]/P = k[a,, . . . , a,] and k[X]/Q = k[pl,. . ,&,]; we can assume that
/Ii,. . . , p,. is a transcendence basis for k(j3) over k. Then a,, . . . , c(, are also
95       Hilhert   Nullstellensatz   and dimension theory                      35

    algebraically independent over k, so that they form a transcendence         basis
    for &cc) over k. Now set S = k[X, , . , X,] - { 0); S is a multiplicative   set in
    k[X]    with PnS=@        and QnS=@.          Setting R=k[X,,...,X,]          and
    K=k(X,,...,     X,), we have R,  = K[X,+   r,. . . , X,], and
             Rs/PRsNk(cr,,...,x,)Ccr,+,,...,a,l,
    ~0 that R,/PR,     is algebraic over K = k(X,,. ..,X,) 2: k(a,,.. .,r,), and
    therefore by Theorem 1, PR, is a maximal ideal of R,; but this contradicts
.   the assumptions P c Q with P # Q and Q n S = fzr.
       Now let us prove that r < dim A by induction on r. If r = 0 then, by
    Theorem 1, A is a field, so dim A = 0 and the assertion holds. Now let r > 0,
    and suppose that A = k[a,, . , r,] with g1 transcendental over k; setting
    S = k[X,] - (0) and R = k[X,, . . .,X,] we get
             R, = k(X,)[X,,    . , X,]  and R,/PR, N k(x,)[a,, . . , a,].
    Hence R,/PR,      has transcendence degree r - 1 over k(X,), so that by
    induction dim R,/PR, > r - 1. Thus there exists a strictly increasing chain
    PRs=QocQl c~~~cQ,~,ofprimeidealsofR,.IfwesetPi=QinRthen
    Pi is a prime ideal of R disjoint from S; in particular, the residue class of
    X, in R/P,-,    is not algebraic over k, and so tr.deg, R/P,-, > 0. Then
    P,-, is not a maximal ideal of R by Theorem 3, and therefore R has a
    maximal ideal P, strictly bigger than P,- r. Hence dim A = coht P 3 r. n
\end{verbatim}
\end{proof}


\begin{theorem}[Printed Theorem 5.7]\label{mcrt-auto-theorem-5-7}\dcref{5.7}\source{matsumura-commutative-ring-theory:page-0050}\uses{}
\begin{verbatim}
Let A be a Noetherian         ring and M a finite A-module.     Set
           b(M) = sup {~(p, M) + coht pJp~Supp M};
    then M can be generated by at most b(M) elements.
       This theorem is a very important link between the number of local and
    global generators. However, there was room for improvement in the bound
    for the number of generators, and in no time R. Swan obtained a better
    bound. We will prove Swan?s bound. For this we need the concept ofj-Spec
    A introduced by Swan. This is a space having the same irreducible closed
    subsets as m-Spec A, but has the advantage of having a ?generic point?, not
    Present in m-Spec A, for every irreducible closed subset.
36        Prime ideals

    A prime ideal which can be expressed as an intersection of maximal
 ideals is called a j-prime ideal, and we write j-SpecA for the set of all
j-prime ideals. We consider j-Spec A also with its topology as a subspace
 of Spec A. Set M = m-SpecA and J = j-Spec A. If F is a closed subset of
J then there is an ideal I of A such that F = V(I)n J. One sees easily that
 a prime ideal P belongs to F if and only if P can be expressed as an
intersection of elements of F n M = V(Z)n M. Hence F is determined by
FnM, so that there is a natural one-to-one correspondence between
 closed subsets of J and of M. It follows that if M is Noetherian so is J,
and they both have the same combinatorial dimension. Now let B be an
irreducible closed subset of J, and let P be the intersection of all the
elements of B. If B = V(I)n J then I c P and we can also write
B = V(P)n J. If P is not a prime ideal then there exist f, gEA such that
f $P, gq!P and fgEP; but then
          B=(V(P+fA)nJ)u(V(P+gA)nJ),
and by definition of P there is a QEB not containing f and a Q?EB not
containing g, which implies that B is reducible, a contradiction. Therefore
P is a prime ideal. Hence PEB and B = V(P)n J. This P is called the
generic point of B. Conversely if P is any element of J then V(P)n J is
an irreducible closed subset of J, and is the closure in J of (P}. We will
write j-dim P for the combinatorial dimension of V(P)n J.
   For a finite A-module M and ~EJ we set
                          0 ifM,=O
           b(p,M) =
                         j-dim p + ~(p, M)   if M, # 0.
\end{verbatim}
\end{theorem}
\begin{proof}
\notready
\end{proof}


\begin{theorem}[Printed Theorem 5.8]\label{mcrt-auto-theorem-5-8}\dcref{5.8}\source{matsumura-commutative-ring-theory:page-0051}\uses{}
\begin{verbatim}
Swan [l]). Let A be a ring, and suppose that m-Spec A is
a Noetherian space. Let M be a finite A-module. If
         sup {b(p, M)lp~j-Spec A} = Y < cc
then M is generated by at most Y elements.
\end{verbatim}
\end{theorem}
\begin{proof}
\begin{verbatim}
Step 1. For p&pecA     and XEM, we will say that x is basic at p if x
has non-zero image in M 0 rc(p). It is easy to see that this condition is
equivalent to ~(p, M/Ax) = ~(p, M) - 1.
\end{verbatim}
\end{proof}


\begin{theorem}[Printed Theorem 6.1]\label{mcrt-auto-theorem-6-1}\dcref{6.1}\source{matsumura-commutative-ring-theory:page-0053}\uses{}
\begin{verbatim}
Let A be a Noetherian         ring and M a non-zero A-module.
   (i) Every maximal element of the family of ideals F = (ann (x) 10 # XE M)
is an associated prime of M, and in particular Ass(M) # a.
   (ii) The set of zero-divisors for M is the union of all the associated primes
of M.
\end{verbatim}
\end{theorem}
\begin{proof}
\begin{verbatim}
(i) We have to prove that if ann (x) is a maximal element of F then it
is prime: if a, be.4 are such that abx = 0 but bx # 0 then by maximality
ann (bx) = ann (x); hence, ax = 0.
   (ii) If ax = 0 for some x # 0 then asann (x)EF, and by (i) there is an
associated prime of M containing arm(x). n
\end{verbatim}
\end{proof}


\begin{theorem}[Printed Theorem 6.2]\label{mcrt-auto-theorem-6-2}\dcref{6.2}\source{matsumura-commutative-ring-theory:page-0053}\uses{}
\begin{verbatim}
Let S c A be a multiplicative    set, and N an As-module.
Viewing Spec (A,) as a subset of Spec A, we have Ass,(N) = Ass,~(N). If A
is Noetherian    then for an A-module         M we have Ass(M,) ==
Ass(M)nSpec(A,).
\end{verbatim}
\end{theorem}
\begin{proof}
\begin{verbatim}
For xeN we have annA(x)=ann,,(x)nA,          so that if PeAss,dN)
then PnAeAss,(N).       Conversely if p~Ass*(N) and we choose XEN
such that p = arm,(x) then x # 0, and hence, p n S = @ and pA, is a prime
ideal of A, with pA, = ann,,(x). For the second part, if p~Ass(M)n
Spec(A,) then pnS = a, and p = arm,(x) for some XEM; if (a/s)x = 0 in
M, then there is a YES such that tax =0 in M, and t$p, taE:p gives
aep, so that ann,,(x)=pA,         and pA,gAss(M,).    Conversely, if PE
Ass(M,) then without loss of generality we have P = annA.     with XEM.
Setting p = Pn A we have P = pA,. Now p is finitely generated since A is
Noetherian, and it follows that there exists some tES such that p =
ann,(tx). Therefore pass..,.        n
\end{verbatim}
\end{proof}


\begin{theorem}[Printed Theorem 6.3]\label{mcrt-auto-theorem-6-3}\dcref{6.3}\source{matsumura-commutative-ring-theory:page-0053}\uses{}
\begin{verbatim}
Let A be a ring and O+ M? -        M -M?   +O an exact
sequence of A-modules; then
         Ass(M) c Ass (M?) u Ass (M?).
\end{verbatim}
\end{theorem}
\begin{proof}
\begin{verbatim}
If PEASS (M) then M contains a submodule N isomorphic to A/P.
Since P is prime, for any non-zero element x of N we have arm(x) = P.
Associated primes and primary decomposition                        39
?6

Therefore if N n M? # 0 we have PEASS (M?). If N n M? = 0 then the image
of N in M? is also isomorphic to A/P, SO that PEAss(M?).      H
\end{verbatim}
\end{proof}


\begin{theorem}[Printed Theorem 6.4]\label{mcrt-auto-theorem-6-4}\dcref{6.4}\source{matsumura-commutative-ring-theory:page-0054}\uses{}
\begin{verbatim}
Let A be a Noetherian ring and M # 0 a finite A-module.
Then there exists a chain 0 = M, c M 1 c ... c M, = M of submodules of M
such that for each i we have M,/M,- 1 N A/Pi with PiESpeC A.
\end{verbatim}
\end{theorem}
\begin{proof}
\begin{verbatim}
Choose any P,EAss(M);      then there exists a submodule M, of M
with M1 N A/P,. If M, # M and we choose any P,gAss (M/M,) then there
exists Mz c M such that M,/M, N A/P,. Continuing in the same way and
using the ascending chain condition, we eventually arrive at M, = M. n
\end{verbatim}
\end{proof}


\begin{theorem}[Printed Theorem 6.5]\label{mcrt-auto-theorem-6-5}\dcref{6.5}\source{matsumura-commutative-ring-theory:page-0054}\uses{}
\begin{verbatim}
Let A be a Noetherian ring and M a finite A-module.
   (i) Ass(M) is a finite set.
   (ii) Ass (M) c Supp (M).
   (iii) The set of minimal elements of Ass (M) and of Supp (M) coincide.
\end{verbatim}
\end{theorem}
\begin{proof}
\begin{verbatim}
(i) follows from the previous two theorems; we need only note
that Ass (A/P) = {P}. For (ii), if 0 -+ A/P -M     is exact then SO is 0 +
A~/PA, + Mp, and therefore M, # 0. For (iii) it is enough to show that
if P is a minimal element of Supp (M) then PEASS (M). We have M, # 0
so that by Theorem 2 and (ii),
          121#Ass(M,)=Ass(M)nSpec(A,)cSupp(M)nSpec(A,)
                         = (PI.
Therefore we must have PEAss(M).          n
    Let A be a Noetherian ring and M a finite A-module. Let Pi,. . . , P, be the
minimal elements of Supp (M); then Supp (M) = V(P,) LJ?. u V(P,), and the
 v(PJ are the irreducible components of the closed set Supp(M) (see
Ex. 4.11). The prime ideals P,, . . . , P, are called the isolated associated
primes of M, and the remaining associated primes of M are called embedded
Primes. If I is an ideal of A then Supp,(A/I)       is the set of prime ideals
containing I, and the minimal prime divisors of I (that is the minimal
associated primes of the A-module A/Z) are precisely the minimal prime
ideals containing I. We have seen in Ex. 4.12 that there are only a finite
number of such primes, and Theorem 5 now gives a new proof of this.
(For examples of embedded primes see Ex. 6.6 and Ex. 8.9.)
\end{verbatim}
\end{proof}


\begin{theorem}[Printed Theorem 6.6]\label{mcrt-auto-theorem-6-6}\dcref{6.6}\source{matsumura-commutative-ring-theory:page-0055}\uses{}
\begin{verbatim}
Let A be a Noetherian ring and M a finite A-module. Then a
submodule N c M is primary if and only if Ass (M/N) consists of one
element only. In this case, if Ass (M/N) = {P] and ann (M/N) = I then I is
primary and JI = P.
\end{verbatim}
\end{theorem}
\begin{proof}
\begin{verbatim}
If Ass(M/N) = {P> then by the previous theorem Supp(M/N) =
 V(P), so that P = ,/(ann(M/N)).       N ow if aEA is a zero-divisor for M/N it
follows from Theorem 1 that aEP, so that aEJ(ann(M/N));             hence, N is a
primary submodule of M. Conversely, if N is a primary submodule and
PgAss(M/N)        then every aeP is a zero-divisor for M/N, so that by
assumption aEJI, where I = ann (M/N). Hence P c JI, but from the
definition of associated prime we obviously have I c P, and hence JI c P,
so that P = JI. Thus Ass (M/N) has just one element JI. We prove that in
this case I is a primary ideal: let a, SEA with b$Z; if abel then ab(M/N) = 0,
but b(M/N) # 0, so that a is a zero-divisor for M/N, and therefore
aEP=JI.         w
\end{verbatim}
\end{proof}


\begin{theorem}[Printed Theorem 6.7]\label{mcrt-auto-theorem-6-7}\dcref{6.7}\source{matsumura-commutative-ring-theory:page-0055}\uses{}
\begin{verbatim}
If N and N? are P-primary submodules of M then so is N n N?.
\end{verbatim}
\end{theorem}
\begin{proof}
\begin{verbatim}
We can embed M/(N n N?) as a submodule of (M/N) 0 (M/N?), so
that
           Ass(M/(NnN?))    cAss(M/N)uAss(M/N?)=          {P].    n

    If N c M is a submodule, we say that N is reducible if it can be written as
an intersection N = N, n N, of two submodules N,, N, with Ni # N, and
otherwise that N is irreducible; note that this has nothing to do with the
notion of irreducible modules in representation theory (= no submodules
other than 0 and M), which is a condition on M only.
    If M is a Noetherian module then any submodule N of M can be written
as a finite intersection of irreducible submodules. Proof: let 9 be the set
of submodules N c M having no such expression. If F # @ then it has a
maximal element N,. Then N, is reducible, so that N, = N, n N,, and
Ni#9.      Now each of the Ni is an intersection of a finite number of
irreducible submodules, and hence so is N,. This is a contradiction.
\end{verbatim}
\end{proof}


\begin{theorem}[Printed Theorem 6.8]\label{mcrt-auto-theorem-6-8}\dcref{6.8}\source{matsumura-commutative-ring-theory:page-0056}\uses{}
\begin{verbatim}
Let A be a Noetherian ring and M a finite A-module.
  (i) An irreducible submodule of M is a primary submodule.
  (ii) If
             N=N,n..,nN,      with Ass (M/N,) = (Pi}
is an irredundant primary decomposition of a proper submodule N c M
then Ass(M/N) = {P,, . . . , Pl}.
   (iii) Every proper submodule N of M has a primary decomposition. If N is
a proper submodule of M and P is a minimal associated prime of M/N then
the P-primary component of N is (pp l(Np), where (pp: M -M,         is the
canonical map, and therefore it is uniquely determined by M, N and P.
\end{verbatim}
\end{theorem}
\begin{proof}
\begin{verbatim}
(i) It is enough to prove that a submodule N c M which is not
primary is reducible: replacing M by M/N we can assume that N = 0. By
Theorem 6, Ass (M) has at least two elements P, and P,. Then M contains
submodules Ki isomorphic to A/Pi for i = 1,2. Now since ann (x) = Pi for
any non-zero x~K, we must have K, n K, = 0, and hence 0 is reducible.
   (ii) We can again assume that N = 0. If 0 = N, n...n N, then M is
isomorphic to a submodule of M/N, @... 0 M/N,, so that

            ASS(M)   c ASS                  ASS (M/NJ = {PI 2.. .) Pl}?

On the other hand N,n...nN,#O,          and taking O#xEN,n...nNN,        we
have arm(x) = 0:x = N, :x. But N, : M is a primary ideal belonging to PI, SO
that Pr M c N, for some v > 0. Therefore P; x = 0; hence there exists i 3 0
such that Pix # 0 but Py l x = 0, and choosing 0 # y~P;x we have
PIY = 0. However, since YEN, n..? n N, it follows that y#N,, and by
the definition of primary submodule ann (y) c P, , so that P, = ann (y) and
42       Prime ideals

P,EAss(M).      The same works for the other Pi, and this proves that
{P1,...,P,}    cAss(M).
   (iii) We have already seen thal a proper submodule has an irreducible
decomposition, so that by(i) it has a primary decomposition. Suppose that
N = N, n.-.n N, is a shortest primary decomposition, and that N, is the
P-primary component with P = P,. By Ex. 4.8 we know that N, =
V%n-        n(N,&,    and for i > 1 a power of Pi is contained in ann (M/N,);
then since Pi # P, we have (M/Ni)p = 0, and therefore (Ni), = M,. Thus
N, = (N1)P, and hence qP1(Np) = (PPI((N~)~); it is easy to check that
the right-hand side is N,. n
\end{verbatim}
\end{proof}


\begin{theorem}[Printed Theorem 6.10]\label{mcrt-auto-theorem-6-10}\dcref{6.10}\source{matsumura-commutative-ring-theory:page-0058}\uses{}
\begin{verbatim}
If    O+M? -M         -M?    -+O is an exact sequence of
A-modules, then Att (M?) c Att (M) c Att (M?)u Att (M?).
\end{verbatim}
\end{theorem}
\begin{proof}
\begin{verbatim}
The first inclusion is trivial from the definition. For the second, let
PEAtt (M) and let N be a submodule such that M/N is P-secondary. If
M? + N = M then M/N is a non-trivial quotient of M?, hence PEAtt (M?).
If M? + N #M then M/(M? + N) is a non-trivial quotient of M? as well
as Of M/N, hence M? has a P-secondary quotient and PEAtt(M?).             H
44       Prime ideals

  An A-module M is said to be sum-irreducible           if it is neither zero nor the
sum of two proper submodules.
\end{verbatim}
\end{proof}


\begin{theorem}[Printed Theorem 6.11]\label{mcrt-auto-theorem-6-11}\dcref{6.11}\source{matsumura-commutative-ring-theory:page-0059}\uses{}
\begin{verbatim}
If M is Artinian, then it has a secondary representation.
\end{verbatim}
\end{theorem}
\begin{proof}
\begin{verbatim}
Similar to the proof of Theorem 6.8, (iii). n
   The class of modules which have secondary representations is larger
than that of Artinian modules. Sharp [8] proved that an injective module
over a Noetherian ring has a secondary representation.
\end{verbatim}
\end{proof}


\begin{theorem}[Printed Theorem 7.1]\label{mcrt-auto-theorem-7-1}\dcref{7.1}\source{matsumura-commutative-ring-theory:page-0061}\uses{}
\begin{verbatim}
Let A -B       be a homomorphism      of rings and M a B-
module. A necessary and sufficient condition for M to be flat over A is that
for every prime ideal P of B, the localisation M, is flat over A, where
p = PnA (or the same condition for every maximal ideal P of B).
\end{verbatim}
\end{theorem}
\begin{proof}
\begin{verbatim}
First of all we make the following observation: if S c A is a
multiplicative  set and M, N are A,-modules, then M @*,N = M aA N.
This follows from the fact that in N OA M we have



for XEM, YEN, aEA and SEX (In general, if B is an A-algebra and
M and N are B-modules, it can be seen from the construction of the tensor
product that M QN is the quotient of M @,4N by the submodule generated
by {bx@y-x@byIx~M,        YEN and beB}.)
   Assume now that M is A-flat. The map A -B          induces A, -B,,
and M, is a BP-module, therefore an A,,-module. Let Y be an exact
sequence of A,-modules; then, by the above observation,
         YOAyMP=~POAMP=(~~*M)OBBP,
and the right-hand side is an exact sequence, so that M, is A,-flat.
   Next, suppose that M, is A,-flat for every maximal ideal P of B.
Let O+ N? -N       be an exact sequence of A-modules, and write K
for the kernel of the B-linear map N?@, M -N             GA M, so that
0 + K -N?     @ M ----f N @ M is an exact sequence of B-modules. For
any PEm-Spec B the localisation
Flatness                                                       47
?7
        O-K,-N?Q,M,-N&M,
is exact, and since N?@,M,=       N?ga(ApBApMP)=     NkBapMp,      and
similarly N0,J4P=NpO~DMP,     we have K,=O by hypothesis. There-
fore by Theorem 4.6 we have K = 0, and this is what we have to prove.
\end{verbatim}
\end{proof}


\begin{theorem}[Printed Theorem 7.2]\label{mcrt-auto-theorem-7-2}\dcref{7.2}\source{matsumura-commutative-ring-theory:page-0062}\uses{}
\begin{verbatim}
Let A be a ring and M an A-module. Then the following
conditions are equivalent:
   (1) M is faithfully flat over A;
   (2) M is A-flat, and N OAM # 0 for any non-zero A-module N;
   (3) M is A-flat, and mM # M for every maximal ideal m of A.
\end{verbatim}
\end{theorem}
\begin{proof}
\begin{verbatim}
(l)=(2). Let .Y be the sequence 0 -+ N -+O. If N @ M = 0 then
y@ M is exact, so .Y is exact, and therefore N = 0.
   (2)+(3). This is clear from M/mM = (A/m)@,M.
   (3)+(2). If N # 0 and 0 # XEN then Ax rr A/ann(x), so that taking a
maximal ideal m containing arm(x), we have M # mM 1 ann(x).M;
hence, Ax@ M # 0. By the flatness assumption, Ax @ M -N         @ M is
injective, so that N @ M # 0.
   (2)=>(l). Consider a sequence of A-modules
        Y:N?&V&N?.
If


is exact then gMof, = (gOf)M = 0, so that by flatness, Im(gof)@ M =
Im(g,of,)    = 0. By assumption we then have Im(gof) = 0, that is gof = 0;
hence Ker g I> Im f. If we set H = Ker g/Im f then by flatness,
          H 0 M = Ker (gM)/Im (fM) = 0,
so that the assumption gives H = 0. Therefore 9? is exact. n
   A ring homomorphism    f: A - B induces a map ?f : Spec B - Spec A,
under which a point pESpecA has an inverse image ?f-?(p) =
{PESpecBJPnA=p}        which is homeomorphic         to Spec(B OAic(n)).
Indeed, setting C = BOAT       and S = A - p, and defining g:B -C    by
g(b)= b@ 1, then since ~c(p)= (A/p)@ A,, we have
        C=@&WO~A~               =(BhB),   =(B/PB)~~~,.
Thus ?g: Spec c -Spec      B has the image

        (PESpecBJP      1 pB and Pnf(S)      = a}
           = {PESpecBIPnA=p),
which is ?f-?(p),    and ug induces a homomorphism      of SpecC with
?f-?(P). For this reason we call SpecC = Spec(B@K(p)) the fibre ouer p.
The inverse map a~ - l(p) - SpecC takes P@f - ?(p) into PC = PB,/pB,.
48           Properties     qf extension rings

For P*ESpecC           we set P = P*n B; then by Theorems 4.2 and 4.3, we have
             P* = PC and C,. = (B&B,),,           = B,/pB,   = B,@,K(~)
\end{verbatim}
\end{proof}


\begin{theorem}[Printed Theorem 7.3]\label{mcrt-auto-theorem-7-3}\dcref{7.3}\source{matsumura-commutative-ring-theory:page-0063}\uses{}
\begin{verbatim}
Let f:A -B       be a ring homomorphism      and M a B-
module. Then
   (i) M is faithfully flat over A*?f(Supp(M))     = Spec A.
   (ii) If M is a finite B-module then
M is A-flat and ?f(Supp(M))       3 m-Spec AoM is faithfully flat over A.
\end{verbatim}
\end{theorem}
\begin{proof}
\begin{verbatim}
(i) For p ESpec A, by faithful flatness we have M OAx(p) # 0. Hence,
if we set C=BQadp)            and M?= M@,k.(p)=      M&C,    the C-module
M? #O, so that there is a P*ESpecC such that M?,* #O. Now set
P= P*nB; then
             M~t=MOgCP*=MOg(BPOBpCP*)=MPOBpCP,
so that M, # 0, that is P~supp(M).         But P*~spec(B@ ti(p)), so that as
we have seen Pn A = p. Therefore p&?f(Supp(M)).
(ii) It is enough to show that M/mM # 0 for any maximal ideal m of A. By
assumption there is a prime ideal P of B such that Pn A = m and
M, # 0. By NAK, since M, is finite over B, we have MJPM,              # 0, and
a fortiori M,/mM,       = (M/mM),    # 0, so that M/mM # 0. W
    Let (A, m) and (B,n) be local rings, and f: A -B         a ring homomor-
phism; f is said to be a local homomorphism if f(m) c n. If this happens then
by Theorem 2, or by Theorem 3, (ii), we see that it is equivalent to say that f
is flat or faithfully flat.
    Let S be a multiplicative   set of A. Then it is easy to see that Spec(A,)
 -+ SpecA is surjective only if S consists of units, that is A = A,. Thus
from the above theorem, if A # A, then A, is flat but not faithfully flat over
A.
\end{verbatim}
\end{proof}


\begin{theorem}[Printed Theorem 7.4]\label{mcrt-auto-theorem-7-4}\dcref{7.4}\source{matsumura-commutative-ring-theory:page-0063}\uses{}
\begin{verbatim}
(i) Let A be a ring, M a flat A-module,and N,, N, two submodules of an
A-module       N. Then as submodules of N @ M we have
            (N,nN,)@M=(N,OM)n(N,@M).
   (ii) Let A -B            be a flat ring homomorphism,      and let I, and I, be
ideals of A. Then
            (I,nZ,)B=I,Bnl,B.
     (iii) If in addition    I, is finitely generated then
            (I,:I,)B      = Z,B:Z,B.
\end{verbatim}
\end{theorem}
\begin{proof}
\begin{verbatim}
(i) Define cp:N -N/N, @ N/N, by q(x) = (.Y+ N,, x + N,); then
O-+ N, n N, - N - NJN, @N/N, is exact, and hence so is
            O-+(N,nN,)@M-NOM-
                  (N 0 MW,         0 M) 0 (N 0 W/(Nz 0 M).
?7       Flatness                                                           49

This is the assertion in (i).
    (ii) This is a particular case of(i) with N = A, M = B, in view of the fact
that for an ideal I of A the subset I OAB of A OAB = B coincides with IB.
    (iii) If I, = Aa, + .. . + Aa, then since (I,:I,) = ni(ll:ai), we can use
(ii) to reduce to the case that I, is principal. For aeA we have the exact
sequence



and tensoring       this with B gives the assertion.   H
\end{verbatim}
\end{proof}


\begin{theorem}[Printed Theorem 7.5]\label{mcrt-auto-theorem-7-5}\dcref{7.5}\source{matsumura-commutative-ring-theory:page-0064}\uses{}
\begin{verbatim}
Let ,f:A -        B be a faithfully flat ring homomorphism.
    (i) For any A-module M, the map M-M               BAB defined by mint @ 1
is injective; in particular .f?: A -B   is itself injective.
   (ii) If I is an ideal of A then IBn A = I.
\end{verbatim}
\end{theorem}
\begin{proof}
\begin{verbatim}
(i) Let 0 # mEM. Then (Am) 0 B is a B-submodule of M 0 B which
can be identified with (ma l)B. But by Theorem 2, (Am)@B # 0, so that
m@l#O.
   (ii) follows by applying (i) to M = A/Z, using (A/I)@B = B/IB.
\end{verbatim}
\end{proof}


\begin{theorem}[Printed Theorem 7.6]\label{mcrt-auto-theorem-7-6}\dcref{7.6}\source{matsumura-commutative-ring-theory:page-0064}\uses{}
\begin{verbatim}
Let A be a ring and M a flat A-module. If aijEA and xj~M
(for 1 < i 6 I and 1 d j < n) satisfy
        TaijXj=O          for all   i,

then there exists an integer s and bjkEA, yk~M (for 1 d j 6 n and 1 6 k d s)
such that
           c aijbjk = 0 for all i, k, and xj = 1 bj,y, for all j.
            j                                     j
   Thus the solutions in a flat module M of a system of simultaneous linear
equations with coefficients in A can be expressedasa linear combination of
solutions in A. Conversely, if the above conclusion holds for the caseof a
single equation (that is for Y= l), then M is flat.
Procf. Set cp:A? -A?      for the linear map defined by the matrix (aij), and
let cp,,,,:M? - M? be the same thing for M; then ?pM= cp@1, where 1 is
the identity map of M. Setting K = Ker cpand tensoring the exact sequence

KA    A? A         A? with M, we get the exact sequence
50         Properties    of extension rings

By assumption        cpM(xl,.     , x,) = 0, so that we can write

           (XI,...,x.)=(iOl)(~~~B*~y*)              with &EK        and        y,cM.

If we write out Bk as an element of A? in the form Pk = (b,,, . . . , bnk) with
~,EA then the conclusion follows. The converse will be proved after the
next theorem.    n
\end{verbatim}
\end{theorem}
\begin{proof}
\notready
\end{proof}


\begin{theorem}[Printed Theorem 7.7]\label{mcrt-auto-theorem-7-7}\dcref{7.7}\source{matsumura-commutative-ring-theory:page-0065}\uses{}
\begin{verbatim}
Let A be a ring and M an A-module. Then M is flat over A if
and only if for every finitely generated ideal I of A the canonical map
I @,,M - A BAM is injective, and therefore I @ M N ZM.
Proqf. The ?only if? is obvious, and we prove the ?if?. Firstly, every ideal
of A is the direct limit of the finitely generated ideals contained in it, so
that by Theorems A 1 and A2 of Appendix A, I @ M -M           is injective for
every ideal 1. Moreover, if N is an A-module and N? c N a submodule,
then since N is the direct limit of modules of the form N? + F, with F
finitely generated, to prove that N? @ M -       N @ M is injective we can
assume that N = N? + Aw, + ... + AU,,. Then setting Ni = N? + AU,
 +. .. + Aoi (for 1 < i < n), we need only show that each step in the
chain
           N?@M-N1@M-N2@M-...-N@M
is injective, and finally that if N = N? + Ao then N? @M -N       @M
is injective. Now we set Z = {a~ A(u~EN?}, and get the exact sequence
           O+N?--+N-+A~I+O.
This induces a long exact sequence(see Appendix B, p. 279)
           . ..-Tor.4(M,A/Z)--+N?@M--*N@M-(A/Z)OM+O;
hence it is enough to prove that
     (*)   Torf(M,      A/Z) = 0.
For this consider the short exact sequence
        O+Z--+A-A/Z+0
and the induced long exact sequence
        Tor:(M, A) = 0 --+ Tort(M, A/Z) -               Z0 M -       M-         *. ?;
since Z0 M -M        is injective, (*) must hold. W
   From this theorem we can prove the converse of Theorem 6. Indeed, if
Z = Au, + ... + Au, is a finitely generated ideal of A then an element 5 of
Z@ M can be written as 5 = c; a, @mi with miEM. Supposethat < is 0 in M,
that is that Cairni = 0. Now if the conclusion of Theorem 6 holds for M,
there exist bijEA and y,eM such that
           Cuibij = 0 for all j,         and   mi = Cbijyj   for all      i.
            i
Flatness                                                              51
?7
Then 5 = xai@ m, = ~i~juihij@           yj = 0, so that I @ M -M        iS injec-
tive, and therefore M is flat.
\end{verbatim}
\end{theorem}
\begin{proof}
\notready
\end{proof}


\begin{theorem}[Printed Theorem 7.8]\label{mcrt-auto-theorem-7-8}\dcref{7.8}\source{matsumura-commutative-ring-theory:page-0066}\uses{}
\begin{verbatim}
Let A be a ring and M an A-module. The following
conditions are equivalent:
  (1) M is flat;
  (2) for every A-module N we have Tor:(M, N) = 0;
  (3) Torf(M, A/Z) = 0 for every finitely generated ideal I.
\end{verbatim}
\end{theorem}
\begin{proof}
\begin{verbatim}
(l)+-(2)   If we let ??-Li-Li-,-?.-Lo--?NjO
be a projective resolution of N then
         . ..---tLiOM-Li_.OM-...-L,OM
is exact, so that Tor:(M, N) = 0 for all i > 0.
   (z)+(3) is obvious.
   (3)=+(l) The short exact sequence 0 --f I -        A -A/I     + 0 induces a
long exact sequence
           Torf(M,A/I)=   O+I@M+M-+M@A/I+O,
and hence I @ M -M         is injective; therefore by the previous theorem M
is flat. W
\end{verbatim}
\end{proof}


\begin{theorem}[Printed Theorem 7.9]\label{mcrt-auto-theorem-7-9}\dcref{7.9}\source{matsumura-commutative-ring-theory:page-0066}\uses{}
\begin{verbatim}
Let 0 --t M? + M -M?            +O be an exact sequence of A-
modules; then if M? and M? are both flat, so is M.
\end{verbatim}
\end{theorem}
\begin{proof}
\begin{verbatim}
For any A-module N the sequence Tor,(M?, N) -Tor,(M,                     N)
 -Torl(M?,      N) is exact, and since the first and third groups are zero,
also Tor,(M, N) = 0. Therefore by the previous theorem M is flat. n
    A free module is obvious faithfully flat (if F is free and Y is a sequence of
A-modules then Y @F is just a sum of copies of Y in number equal to the
cardinality of a basis of F). Conversely, over a local ring the following
theorem holds, so that for finite modules flat, faithfully flat and free are
equivalent conditions.
\end{verbatim}
\end{proof}


\begin{theorem}[Printed Theorem 7.10]\label{mcrt-auto-theorem-7-10}\dcref{7.10}\source{matsumura-commutative-ring-theory:page-0066}\uses{}
\begin{verbatim}
Let (A,m) be a local ring and M a flat A-module. If
Xl,... ,&EM are such that their images Xi,. . .,X, in M = M/mM are
linearly independent over the field A/m then xi,. . . ,x, are. linearly
independent over A. Hence if M is finite, or if m is nilpotent, then any
minimal basis of M (see$2) is a basis of M, and M is a free module.
 proo! By induction on n. If n = 1, and EA is such that ax, = 0 then
 hY Theorem 6 there are b,, . . . ,b,EA such that abi = 0 and XE~ biM; by
 aSsumPtionx1 $mM, so that among the bi there must be one not contained
 in me This bi is then a unit, so that we must have a = 0.
    For n > 1, let caixi = 0 ; then there are bijgA and ~,EM (for 1 <<j < s)
 Such that C Uibij = 0 and xi = C b,y,. NOWx,$mM, sp that among the !J,,~at
least one is a unit. Hence a, is a linear combination of a,, . . . . , a,- 1, that
52           Properties of extension rings

is a, = c:Zf       aici for some c,EA. Therefore we have
        al(xl +clx,)+~~~+a,-,(x,-,         +c,-1x,)=O;
however, the (n - 1) elements X, + Fix,,, . . . ,X,- 1 + C, - I X, of R are linearly
independent over A/m, so that by induction, a, = . . . = a, _ 1 = 0. Hence also
a,=O.        n
\end{verbatim}
\end{theorem}
\begin{proof}
\notready
\end{proof}


\begin{theorem}[Printed Theorem 7.11]\label{mcrt-auto-theorem-7-11}\dcref{7.11}\source{matsumura-commutative-ring-theory:page-0067}\uses{}
\begin{verbatim}
Let A be a ring, M and N two A-modules,              and B a flat A-
algebra. If M is of finite presentation then we have
             Hom,(M,       N)O,B   = Hom,(M@,B,       N OAB).
\end{verbatim}
\end{theorem}
\begin{proof}
\begin{verbatim}
Fixing N and B, we define contravariant functors F and G of an
A-module M by
        F(M) = Hom,(M, N) aA B
and
        G(M) = Hom,(M @AB, N OA B);
then we can define a morphism of functors J:F --+ G by
        Il(f @b) = b.(f @ lB) for f EHom,(M,         N) and &B.
Both F and G are left-exact functors.
   Now if M is of finite presentation there is an exact sequence of the form
AP - A4 --+ M + 0, and from this we get a commutative diagram
        O+ F(M) - F(A4) - F(Ap)
                  11           ?I       ?1
             0 -+ G(M) -      G(Aq) -   G(AP)
having two exact rows. Now F(AP) = NP 0 B and G(AP) = (N @ B)j?, so that
the right-hand il is an isomorphism, and similarly the middle ;I is an
isomorphism.   Thus, as one sees easily, the left-hand i is also an
isomorphism.    n
\end{verbatim}
\end{proof}


\begin{theorem}[Printed Theorem 7.12]\label{mcrt-auto-theorem-7-12}\dcref{7.12}\source{matsumura-commutative-ring-theory:page-0067}\uses{}
\begin{verbatim}
Let A be a ring and M an A-module of finite presentation.
Then M is a projective A-module if and only if M, is a free A,,,-module for
every maximal ideal m of A.
\end{verbatim}
\end{theorem}
\begin{proof}
\begin{verbatim}
If M is projective it is a direct summand of a free module
and this property is preserved by localisation, so that M, is projective over
A,,,, and is therefore free by Theorem 2.5.
Proof of ?if?. Let N, --+ N, -+O be an exact sequence of A-modules.
Write C for the cokernel of
             HomAW,      N,) -     HomAW     N,);
Appendix to $7                                                                               53

then for any maximal          ideal m of A we have

          Cm = CokeriHomAJM,,,                 WA,)     -      Hom,+,(M,,,        (M,,)}       = 0.
Hence C = 0 by Theorem 4.6, and this is what we had to prove .                                    m
\end{verbatim}
\end{proof}


\begin{theorem}[Printed Theorem 7.13]\label{mcrt-auto-theorem-7-13}\dcref{7.13}\source{matsumura-commutative-ring-theory:page-0069}\uses{}
\begin{verbatim}
A submodule N of M is pure if and only if the following
condition holds: if xi = cj=, aijmj (for 1 < i < r), with mjEM, xi~N and
aijEA, then there exist yj~N (for 1 <j <s) such that xi = c;= la,iyj
(for 1 d i < r).
\end{verbatim}
\end{theorem}
\begin{proof}
\begin{verbatim}
Suppose N is pure in M. Consider the free module A? with basis
e,, . . . ,e, and let D be the submodule of A?generated by ~,aijei, 1 <j < s. Set
E = A*/D, and let ei denote the image of ei in E. Then in M @E we have




hence C xi 0 2, = 0 in N 0 E by purity. But this means that, in N 0 A?, the
element Cixi @ e, is of the form cjyjo Ciaijei for some YjE N.
    Conversely, suppose the condition is satisfied. Let E be an A-module of
finite presentation. Then we can write E = A?/D with D generated by a finite
number of elements of A?, say xi= iaijei, 1 <j d s. Then reversing the
preceding argument we can see that N @ E -+ M @ E is injective.         H
\end{verbatim}
\end{proof}


\begin{theorem}[Printed Theorem 7.14]\label{mcrt-auto-theorem-7-14}\dcref{7.14}\source{matsumura-commutative-ring-theory:page-0069}\uses{}
\begin{verbatim}
If N is a pure submodule and M,fN is of finite presentation,
 then N is a direct summand of M.
\end{verbatim}
\end{theorem}
\begin{proof}
\begin{verbatim}
We will prove that O+ N 2 M --% M,lN +O splits, where i
 and p are the natural maps. For this we need only construct a linear map
f :M/N -M         such that pf is the identity map of M/N. Let {tl,. . , tl>
 be a set of generators of M/N, so that M/N N A?IR, where R is the
 submodule of relations among the tj; let {(ail,. . . ,ail)I 1 < id sj be a
 set of generators of R. Choose a pre-image tj of tj in M for each j. Then
 set rji = xaij5jeN (for 1 <ids).    By the preceding theorem there exist
 ~JEN such that uli = xaij<J (for 1 d i <s). Then xaij(tj - <>) = 0 (for
1 < i < s), and setting ,f(tj) = tj - S;, we obtain a linear mapf:M/N     -M
which satisfies the requirement.       n
?8        Completion and the Artin-Rees       lemma                               55

      8   Completion    and the Artin-Rees     lemma
            Let A be a ring and M an A-module; for a directed set A, suppose
  that 9 = {M,} le,, is a family of submodules of M indexed by A and such
 that A< p 3 M, 1 M,. Then taking 5 as a system of neighbourhoods                of 0
 makes M into a topological group under addition. In this topology, for any
 XE M a system of neighbourhoods            of x is given by {x + M,),,,.      In M
 addition and subtraction are continuous, as is scalar multiplication xt-+ax
 for any aeA. When M = A each M, is an ideal, so that multiplication is also
 continuous:
           (a + M,)(b + M,) c ab + M,.
 This type of topology is called a linear topology on M; it is separated
 (that is, Hausdorff)     if and only if n,M, = 0. Each M, c M is an open set,
 each coset x + M, is again open, and the complement M -M,                of M, is a
 union of cosets, so is also open. Hence MA is an open and closed subset;
 the quotient module M/MA is then discrete in the quotient topology.
     M/nnM,     is called the separated module associated with M. Moreover,
 since for i. <p there is a natural linear map (P~~:M/M,, -M/M,,
 we can construct        the inverse system {M/M,;         (P,+} of A-modules; its
 inverse limit @M/M,           is called the completion of M, and is written M.
 We give each M/M, the discrete topology, the direct product n2 M/M, the
 product topology, and M the subspace topology in RAM/M,.                        Let
 $: M -+ M be the natural A-linear map; then $ is continuous, and $(M) is
 dense in M. Write p,:,@ -M/M,              for the projection, and set Ker pI =
 Mf; it is easy to see that the topology of M coincides with the linear
 topology defined by 9 = { Mz),,,.            The map pn is surjective (in fact
p,($(M)) = M/M,), so that &/MT = M/M,, and the completion of M
coincides with M itself. If $:M +M             is an isomorphism,    we say that M
is complete. (Caution: in Bourbaki           terminology     this is ?complete and
separated?; we shorten this to ?complete? throughout.)
    If 9? = {MI}+-      is another family of submodules of M indexed by a
directed set r, then 9 and 9? give the same topology on M if and only if for
each M, there is a YEI- such that Mt c M,, and for every MI, there is a PEA
such that M, c MI. It is then easy to see that there is an isomorphism of
topological modules $r M/M, 2 @r M/MI.. Thus M depends only on
the topology of M, as does the question of whether M is complete.
    When M = A, (M/M,;           qns} becomes an inverse system of rings, M = A
is a ring, and $:A -A              a ring homomorphism.         MX c A is not just
an A-submodule,         but an ideal of A^; this is clear from the fact that
pn: A^ -A/M,       is a ring homomorphism.
    If N c M is a submodule, then the closure N of N in M is given by the
56        Properties   of extension rings

following formula:
         N= n(N+M,).
Indeed,
          .xfzNo(x+M,)nN#IZ(          forallA.
               -xEN + M, for all 1.
   If we write M; for the image of M, in the quotient module M/N, the
quotient topology of M/N is just the linear topology defined by {MA),,,,.
In fact, let G c M be the inverse image of G? c M/N; then
   G? is open in the quotient topology of M/N
   OG is open in M
  0 for every XEG there is an M, such that x + M, c G
  0 for every x?EG? there is an M?, such that x? + M? c G?.
Hence the condition for M/N to be separated is that n,M; = 0, that is
n(N + MJ = N, or in other words, that N is closed in M. Moreover, the
subspace topology of N is clearly the same. thing as the linear topology
defined by {N n MA},,*. Set M/N = M?; then
          O-+ NJ(N n Ml) -       MJM, -     M?JM;   = MJ(N + MA)+0
is an exact sequence, so that taking the inverse limit, we,see that
          O-+fi-A?         -(M/N)-
 is exact. If we view N as a submodule of M, the condition that 5 = (<l)n,,~M
 belongs to N is that each tn can be represented by an element of N, or in
 other words that ~E$(N) + MR for each A. Hence N is the same thing as the
 closure of $(N) in M. In general it is not clear whether M -(M/Nj?is
 surjective, but this holds in the case A = { 1,2,. . . }. In fact then
             (M/N)-= lim M/(N + M,);
 given an elemen?-t? = (t;, t;, . .)E(MJN)~ with ~LEM/(N + M,), let
 x1 EM be an inverse image of 4;) and ~,EM an inverse image of 5;;
 then y, -x1 EN + M,, so that we can write
            y2-x1=t+m,         with tEN and rn,~M,.
If we set x2 = y, - t then X,E M is also an inverse image of 5;) and satisfies
x2 - xi EM~. Similarly we can successively choose inverse images X,EM of
ther:,insuchawaythatforn=1,2,...,wehavex,+,-X,EM,.Ifweset
(?GM/M, for the image of x,, then by construction r = (ti, c2,. . .) is an
element of lim M/M, = M which maps to 5? in (M/N)-. This proves the
following thzrem.
\end{verbatim}
\end{proof}


\begin{theorem}[Printed Theorem 8.1]\label{mcrt-auto-theorem-8-1}\dcref{8.1}\source{matsumura-commutative-ring-theory:page-0071}\uses{}
\begin{verbatim}
Let A be a ring, M an A-module with a linear topologY, and
N c M a submodule. We give N the subspace topology, and M/N the
quotient topology. Then these are both linear topologies, and we have:
    (i) 0 -+ N -ii?    + (M/N)- is an exact sequence, and N is the closure
of $(N) in I&, where $: M -M          is the natural map.
?8       Completion and the Artin-Rees          lemnra                         57

  (ii) If moreover the topology of M is defined by a decreasing chain of
submodules M i 3 M2 1. . ?, then
         O-&-w@-(M/N)-+0
is exact. In other words,       (M/N)^~fi/fl.      n


   Now suppose that M and N are two A-modules with linear topologies,
and let f: M - N be a continuous linear map. If the topologies of M and
N are given by {M,},,,            and (Ny}YEr, then for any YET there
exists JEA such that ML c f -?(NY).         Define cpv:fi+      NJN, as the
composite i6i -Q/M:        -N/N,,      where the first arrow is the natural
map, and the second is induced by f; one sees at once that (py does not
depend on the choice of A for which M, c f -?(NJ.        Also, for y < y? if we
let yQ,,* denote the natural map N/N,, -N/N,,          it is easy to see that
vpy= $yv?%?;    hence there is a continuous        linear map ?:I$ --+fl
defined by the ((~,,)~~r, and the following diagram is commutative          (the
vertical arrows are the natural maps):
                I
         M-N



 Moreover, f is determined uniquely by this diagram and by continuity.
Similarly, if A and B are rings with linear topologies, and f: A -B             is
a continuous ring homomorphism,               then f induces a continuous ring
homomorphism       J?: A -+ B.
   Among the linear topologies, those defined by ideals are of particular
importance. Let I be an ideal of A and M an A-module; the topology on M
defined by {I?M),= 1,2,... is called the I-adic topology. If we also give A the Z-
adic topology, the completions A and ti of A and M are called I-adic
completions; it is easy to see that fi is an A-module: for c1= (aI, a2,. . .)
EAIwitha,EA/l?and5=(x,,x           2,. . .)~fi with X,E M/PM (for all n), we can
just set
         at = (alxl,   a2x2,.    .)Efi.
   AS one can easily check, to say that M is complete for the I-adic topology is
equivalent to saying that for every sequence xi, x2,. . of elements of M
satisfying xi - xi + 1EZ?M for all i, there exists a unique XEM such that
x - xiEl?M for all i. We can define a Cauchy sequence in M in the usual
way ({xi} is Cauchy if and only if for every positive integer r there is an
00 such that x,+ 1 - x,EI?M      for n > no), and completeness can then be
expressed as saying that a Cauchy sequence has a unique limit.
\end{verbatim}
\end{theorem}
\begin{proof}
\notready
\end{proof}


\begin{theorem}[Printed Theorem 8.2]\label{mcrt-auto-theorem-8-2}\dcref{8.2}\source{matsumura-commutative-ring-theory:page-0072}\uses{}
\begin{verbatim}
Let A be a ring, I an ideal, and M an A-module.
  (i) If A is I-adically complete then I c rad(A);
58       Properties   of extension rings

   (ii) If M is I-adically complete and awl, then multiplication  by 1 + a is an
automorphism of M.
\end{verbatim}
\end{theorem}
\begin{proof}
\begin{verbatim}
(i) For u~l, 1 - a + a* - a3 + ... converges in A, and provides an
inverse of 1 + a; hence 1 + a is a unit of A. This means (see $1) that
I c rad (A).
   (ii) M is also an A-module, and 1 + a (or rather, its image in A) is a unit in
A^, so that this is clear. n
   The following two results show the usefulness of completeness.
\end{verbatim}
\end{proof}


\begin{theorem}[Printed Theorem 8.3]\label{mcrt-auto-theorem-8-3}\dcref{8.3}\source{matsumura-commutative-ring-theory:page-0073}\uses{}
\begin{verbatim}
Hensel?s lemma). Let (A, m, k) be a local ring, and suppose
that A is m-adically complete. Let F(X)E A[X] be a manic polynomial, and
let FEk[X]     be the polynomial obtained by reducing the coefficients of F
modulo m. If there are manic polynomials g, hek[X]               with (g, h) = 1
and such that P = gh, then there exist manic polynomials G, H with coeffi-
cients in A such that F = GH, G = g and E7= h.
\end{verbatim}
\end{theorem}
\begin{proof}
\begin{verbatim}
If we take polynomials G,, H, EA[X] such that g = Gi , h = 8, then
F 3 G,H, mod m[X]. Suppose by induction that manic polynomials G,,
H, have been constructed such that F E G,H, mocim?[X],              and G,, = y,
I?,, = h; then we can write
           F - G,H, = cw,U,(X),        with ODES? and deg Ui < deg F.
Since (g, h) = 1 we can find ui, w,ek[X]    such that Di = gui + hwi. Replacing
Di by its remainder modulo h, and making the corresponding correction to
wi we can assume deg vi < deg h. Then
           deg hWi = deg ( Ui - gUi) < deg F, hence deg wi < deg g.
Choosing Vi, W,EA[X]            such that Vi = ui, deg vi= degvi, lVi = Wi,
deg Wi = deg wi, and setting G,, i = G, + cwiWi, H,, 1 = H, + 10~ Vi,
we get

We construct in this way sequences of polynomials G,, H, for n = 1,2,. . ;
then lim G, = G and lim H, = H clearly exist and satisfy F = GH. Obvi-
ously, C = G, = g, R = R, = h. n
\end{verbatim}
\end{proof}


\begin{theorem}[Printed Theorem 8.4]\label{mcrt-auto-theorem-8-4}\dcref{8.4}\source{matsumura-commutative-ring-theory:page-0073}\uses{}
\begin{verbatim}
Let A be a ring, I an ideal, and M and A-module. Suppose
that A is I-adically complete, and M is separated for the I-adic topology. If
MIIM is generated over A/I by W,,.. ., o,, and o,gM is an arbitrary
inverse image of Oi in M, then M is generated over A by o1 , . . . , 0,.
\end{verbatim}
\end{theorem}
\begin{proof}
\begin{verbatim}
By assumption M = 1 Aoi + ZM, SO that M = 1 Aoi + I(1 Ami f
ZM) = 1 Aoi + Z?M, and similarly, M = c Aq + I?M for all v > 0. For
any (EM, write 5 = 1 uiwi + t1 with 5, EIM, then lI = c ui,i Oi + l2 with
u,,,EI and (2~12M, and choose successively u~.~EI? and ~;,EI?M to satisfy
         ~v=~ui,vui+~,+l        for v=1,2 ,....
Completion and the Artin-Rees       lemma                            59
?8

Then ai + ai,l + ai,z + ... converges in A. If we set bi for this sum then


   This theorem is extremely handy for proving the finiteness of M. For a
Noetherian ring A, the I-adic topology has several more important
properties, which are based on the following theorem, proved independ-
ently by E. Artin and D. Rees.
\end{verbatim}
\end{proof}


\begin{theorem}[Printed Theorem 8.5]\label{mcrt-auto-theorem-8-5}\dcref{8.5}\source{matsumura-commutative-ring-theory:page-0074}\uses{}
\begin{verbatim}
the Artin-Rees lemma). Let A be a Noetherian ring, M a
finite A-module, N c M a submodule, and I an ideal of A. Then there exists
a positive integer c such that for every n > c, we have
        I?Mn    N = I?-?(I?MnN).
\end{verbatim}
\end{theorem}
\begin{proof}
\begin{verbatim}
The inclusion 3 is obvious, so that we only have to prove c.
Suppose that I is generated by r elements a,, . . . , a,, and M by s elements
al,-.., w,. An element of PM can be written as ylfi(a)oi,                where
fi(X)=fi(xl,..~,xr)    is a homogeneous         polynomial   of degree  n  with
coefficients in A. Now set A[X i , . . . , X,] = B, and for each n > 0 set

                                   fi are homogeneous   of degree n
                                    and 1; ,fi(a)oiEN

let C c B? be the B-submodule generated by Un, ,J,. Since B is Noetherian,
C is a finite B-module, so that C = xi= i BUj, where each Uj is a linear
combination of elements of u J,; therefore C is generated by finitely
many elements of UJn. Suppose
         C = Bu, + ... + Bu,, where uj = (Ujl,. . , Ujs)EJd, for 1 <j d t.
Setc=max(d,,...          , d,}. Now if q~l?M n N, we can write q = 1 fi(a)wi with
vi,... ,f&J,,,       and hence
         (fi,...,f,)=Cpj(X)~j,            with PjEB=A[X1,...,X,].
The left-hand side is a vector made up of homogeneous polynomials of
degree n only, so that the terms of degree other than n on the right-hand side
must cancel out to give 0. Hence we can suppose that the Pj(X) are
homogeneous of degree n - dj. Then ye= 1 fi(a)wi = Cjpj(a)Ciuji(a)oi,
and CiUji(a)oiEZdJM n N, so that if n > c, pj(a)EZ?-CZc-dJ, giving
         ~EI?-?(Z?M~          N) for any n > c. w
\end{verbatim}
\end{proof}


\begin{theorem}[Printed Theorem 8.6]\label{mcrt-auto-theorem-8-6}\dcref{8.6}\source{matsumura-commutative-ring-theory:page-0074}\uses{}
\begin{verbatim}
In the notation of the above theorem, the I-adic topology of
N coincides with the topology induced by the I-adic topology of M on the
subspace N c M.
\end{verbatim}
\end{theorem}
\begin{proof}
\begin{verbatim}
By the previous theorem, for n > c, we have I?N c PM n N
C 1?-?N. The topology of N as a subspace of M is the linear topology
60       Properties of extension rings

defined by {Z?Mn N},= 1,2,..., and the above formula says that this defines
the same topology as {Z?Nj,= 1,2,,,,. m
\end{verbatim}
\end{proof}


\begin{theorem}[Printed Theorem 8.7]\label{mcrt-auto-theorem-8-7}\dcref{8.7}\source{matsumura-commutative-ring-theory:page-0075}\uses{}
\begin{verbatim}
Let A be a Noetherian ring, I and ideal, and M a finite A-
module. Writing A, A for the I-adic completions of M and A we have

Hence if A is I-adically complete, so is M.
\end{verbatim}
\end{theorem}
\begin{proof}
\begin{verbatim}
By Theorems 1 and 6, the I-adic completion of an exact sequence of
finite A-modules is again exact. Now given M, let A? --+ A4 -M    +O be
an exact sequence; the commutative diagram




has exact rows. Here the vertical arrows are the natural maps; since
completion commutes with direct sums, the two left-hand arrows are
obviously isomorphisms,   and hence the right-hand      arrow is an
isomorphism, as required. n
\end{verbatim}
\end{proof}


\begin{theorem}[Printed Theorem 8.8]\label{mcrt-auto-theorem-8-8}\dcref{8.8}\source{matsumura-commutative-ring-theory:page-0075}\uses{}
\begin{verbatim}
Let A be a Noetherian ring, Z an ideal, and A the I-adic
completion of A; then A is flat over A.
\end{verbatim}
\end{theorem}
\begin{proof}
\begin{verbatim}
By Theorem 7.7 it is enough to show that a@ A --tA^ is injective
for every ideal a c A; but a @ A = 6, and by Theorems 1 and 6,a -+ A is
injective.  n
\end{verbatim}
\end{proof}


\begin{theorem}[Printed Theorem 8.9]\label{mcrt-auto-theorem-8-9}\dcref{8.9}\source{matsumura-commutative-ring-theory:page-0075}\uses{}
\begin{verbatim}
Krull). Let A be a Noetherian ring, Z an ideal, and M
 a finite A-module; set (7n,0 Z?M = N. Then there exists SEA such that
a= 1modZ and aN=O.
\end{verbatim}
\end{theorem}
\begin{proof}
\begin{verbatim}
By NAK, it is enough to show that N = IN. By the Artin-Rees
lemma, Z?M n N c IN for sufficiently large n; now by definition of N, the
left-hand side coincides with N.
\end{verbatim}
\end{proof}


\begin{theorem}[Printed Theorem 8.10]\label{mcrt-auto-theorem-8-10}\dcref{8.10}\source{matsumura-commutative-ring-theory:page-0075}\uses{}
\begin{verbatim}
the Krull intersection theorem).
  (i) Let A be a Noetherian ring and Z an ideal of A with Z c rad A; then for
any finite A-module the I-adic topology is separated, and any submodule is
a closed set.
  (ii) If A is a Noetherian integral domain and Z c A a proper ideal, then
           .Q 1? = (0).
\end{verbatim}
\end{theorem}
\begin{proof}
\begin{verbatim}
(i) In this case the a of the previous theorem is a unit of A, so that
N = 0, and M is separated. If M? c M is a submodule then M/M? is also
I-adically separated, which is the same as saying that M? is closed in M.
   (ii) Setting M = A in the previous theorem, from 1$I we get that a # 0,
and since a is not a zero-divisor, N = 0. n
Completion and the Artin-Rees      lemma                           61
?f3
\end{verbatim}
\end{proof}


\begin{theorem}[Printed Theorem 8.11]\label{mcrt-auto-theorem-8-11}\dcref{8.11}\source{matsumura-commutative-ring-theory:page-0076}\uses{}
\begin{verbatim}
Let A be a Noetherian ring, I and J ideals of A, and M a
hnite A-module; write ^ for the completion of an A-module in the I-adic
topology, and ti :M ---+M for the natural map. Then
        (JM)^= J@ = the closure of $(JM) in M,
and
        (M/JM)-=     M/J@.
\end{verbatim}
\end{theorem}
\begin{proof}
\begin{verbatim}
By Theorems 1 and 6, (JMj is the kernel of &i -(M/JM)?        and
this is equal to the closure of $(JM) in M by Theorem 1. Now supposeJ
 =C?,a,A    and define cp:M?--+M       by (tl,...,<,.) ct cai&. Then the
sequence
        M?LM        2   MfJM+O,
where ~1is the natural map, is exact. The I-adic completion,



is again exact. On the other hand (j is given by the same formula
(t l,...,~,)~~ai~i       as cp,hence (JM)-= Ker(,r?)= Im(@) =ca,M = J&i. n
    As is easily seen,the (X 1, . , X,)-adic completion of the polynomial ring
ACX r,. . . , X,] over A can be identified with the formal power seriesring
A[[X,, . . . , X,,l] . Using this we get the following theorem.
\end{verbatim}
\end{proof}


\begin{theorem}[Printed Theorem 8.12]\label{mcrt-auto-theorem-8-12}\dcref{8.12}\source{matsumura-commutative-ring-theory:page-0076}\uses{}
\begin{verbatim}
Let A be a Noetherian ring, and Z = (a,,    , a,) an ideal of
A. Then the I-adic completion A of A is isomorphic to A[X I,. .,X,]/
(Xl -a,,..., X, - a,,). Hence A is a Noetherian ring.
\end{verbatim}
\end{theorem}
\begin{proof}
\begin{verbatim}
Let B = A[X,,...,    X,], and set I? =xX,B,    J =c(Xi - a,)B;
then BfJ N A, and the I?-adic topology on A considered as the B-module
B/J coincides with the I-adic topology of A. Now writing ^ for the I?-adic
completion of B-modules, we have
        AI=B/.F=&JB=A[x,           ,..., x,]/(x,-a,   ,..., Xn-a,).    n
\end{verbatim}
\end{proof}


\begin{theorem}[Printed Theorem 8.13]\label{mcrt-auto-theorem-8-13}\dcref{8.13}\source{matsumura-commutative-ring-theory:page-0076}\uses{}
\begin{verbatim}
Let A be a Noetherian ring, I an ideal, M a finite A-module,
and fi the I-adic completion of M; then the topology of M is the I-adic
topology of M as an A-module, and is the IA-adic topology of M as
an A-module.
\end{verbatim}
\end{theorem}
\begin{proof}
\begin{verbatim}
If we let M,* be th e k ernel of the map from A = @ (M/Z?M) to
MIZ?M, the topology of M is that defined by {Mzj. Thus it is enough to
Prove that M,* = Z?M. Since M/Z?M is discrete in the Z-adic topology, we
have (M/Z?M)^= M/Z?M and the kernel of ti --+(M/Z?M)^             is ZnM by
Theorem 11. Therefore M,* = I?M. Moreover, Z?fi can also be written
(I?&@, and Z?A = (lay, so that the topology of M is also the IA-adic
topology.
62       Properties of extension rings
\end{verbatim}
\end{proof}


\begin{theorem}[Printed Theorem 8.14]\label{mcrt-auto-theorem-8-14}\dcref{8.14}\source{matsumura-commutative-ring-theory:page-0077}\uses{}
\begin{verbatim}
Let A be a Noetherian ring and I an ideal. If we consider A
with the I-adic topology, the following conditions are equivalent:
   (I) I c rad (A);
   (2) every ideal of A is a closed set;
   (3) the I-adic completion A^ of A is faithfully flat over A.
\end{verbatim}
\end{theorem}
\begin{proof}
\begin{verbatim}
We have already seen (l)+(2).
   (2) a(3) Since A^is flat over A, we need only prove that VIA # A^for every
maximal ideal nr of A. By assumption, (0) is closed in A, so that we can
assume that A c A, and by Theorem 11, rn;i is the closure of m in A^.
However, m is closed in A, so that rnff n A = m, and so rn2 # A^.
   (3) *(l) By Theorem 7.5, ndn A = m for every maximal ideal m of A.
Now mA^ c A^is a closed set by Theorems 2, (i) and 10, (i), and since the
natural map A - A^ is continuous, m = mA^ n A is closed in A. If I $ m.
then I? + rn = A for every n > 0, so that m is not closed. Thus I c m. l
   If the conditions of the above theorem are satisfied, the topological ring A
is said to be a Zariski ring, and I an ideal of definition of A. An ideal of
definition is not uniquely determined; any ideal defining the same topology
will do. The most important example of a Zariski ring is a Noetherian local
ring (A, m) with the m-adic topology. When discussing the completion of a
local ring, we will mean the m-adic completion unless otherwise specified.
\end{verbatim}
\end{proof}


\begin{theorem}[Printed Theorem 8.15]\label{mcrt-auto-theorem-8-15}\dcref{8.15}\source{matsumura-commutative-ring-theory:page-0077}\uses{}
\begin{verbatim}
Let A be a semilocal ring with maximal ideals m,, . . , m,,
and set I = rad(A) = m1m2.. . m,. Then the I-adic completion A^ of A
decomposes as a direct product.
          ?24,     x-.x?&,
where Ai = A,,,, and Ai is the completion of the local ring Ai.
\end{verbatim}
\end{theorem}
\begin{proof}
\begin{verbatim}
Since for i #j and any n > 0 we have my + my = A, Theorem 1.4
gives
          A/Z?= A/m; x*.. x A/m: for n>O.
Hence taking the limit we get
          A^= IF A/I? = ( 9 A/m;) x ... x ( 1E A/m:).
If we set Ai for the localisation of A at mi, then, since A/m; is already local,
          A/ml = (A/m:),, = Ai/?(miAJ?,
and SO lim A,hr~ can be identified with Ai.           n
   We no;summarise          the main points proved in this section for a local
Noetherian ring. Let (A,nr) be a local Noetherian ring; then we have:
   (1) f-LO m? = (0).
   (2) For M a finite A-module and N t M a submodule,
                  + nt?M) = N.
          G (N
Completion and the Artin-Rees              lemma                                     63
?8

   (3) The completion A^ of A is faithfully flat over A; hence A c A^, and
IA~\A = I for any ideal I of A.
   (4) A^is again a Noetherian local ring, with maximal ideal m2, and it has
the same residue class field as A; moreover, A^/m?A^ = A/N? for all IZ > 0.
   (5) If A is a complete local ring, then for any ideal I # A, A/I is again a
complete     local   ring.

Remark 1. Even if A is complete,             the localisation        A, of A at a prime p may
not be.
Remark 2. An Artinian local ring (A, m) is complete; in fact, it is clear from
the proof of Theorem 3.2 that there exists a v such that my = 0, so that
A= 12 A/m?= A.
\end{verbatim}
\end{proof}


\begin{theorem}[Printed Theorem 9.1]\label{mcrt-auto-theorem-9-1}\dcref{9.1}\source{matsumura-commutative-ring-theory:page-0079}\uses{}
\begin{verbatim}
Let A be a ring and B an extension of A.
   (i) An element bEB is integral over A if and only if there exists a ring C
with A c C c B and bgC such that C is finitely generated as an A-module.
   (ii) Let A? c B be the set of elements of B integral over A; then A? is
a subring of B.
\end{verbatim}
\end{theorem}
\begin{proof}
\begin{verbatim}
(i) If b is a root of f(X) = X? + a,X?-l   + ... + a,, for any P(X)E
A[XJ let r(X) be the remainder of P on dividing by f; then P(b) = r(b)
and deg r < n. Hence
           A[b] = A + Ab + ... + Ah?-?,
SO that  we can take C to be A[b]. Conversely if an extension ring C of A is a
finite A-module then every element of C is integral over A: for if C = Ao,
 + ... + Aw, and bEC then
          bw,= CQijWj     with u~~EA,

so that by Theorem 2.1 we get a relation b? + uIbn-l + ... + a, = 0. (The left-
hand side is obtained by expanding out det (b6,, - aij).)
   (ii) If b, b?EA? then we see easily that A[b, b?] is finitely generated as an A-
module, so that its elements bb? and b ) b? are integral over A. w
   The A? appearing in (ii) above is called the integral closure of A in B; if
A = A?we say that A is integrally closed in B. In particular, if A is an integral
domain, and is integrally closed in its field of fractions, we say that A is an
integrally closed domain. If for every prime ideal p of A the localisation A, is
an integrally closed domain we say that A is a normal ring.
\end{verbatim}
\end{proof}


\begin{theorem}[Printed Theorem 9.2]\label{mcrt-auto-theorem-9-2}\dcref{9.2}\source{matsumura-commutative-ring-theory:page-0080}\uses{}
\begin{verbatim}
Let A be an integrally closed domain, K the field of fractions
of A, and L an algebraic extension of K. Then an element acL is integral
over A if and only if its minimal polynomial over K has all its coefficients
in A.
\end{verbatim}
\end{theorem}
\begin{proof}
\begin{verbatim}
Letf(X)=X?+a,X?-I+...              + a, be the minimal polynomial of tl
over K. We have f(a) = 0, so that if all the a, are in A then a is integral over
A. Conversely, if a is integral over A, then letting L be an algebraic closure of
L we have a splitting f(X) = (X - tli). . . (X - a,) of f(X) in E[X] into linear
factors; each of the a, is conjugate to a over K, so that there is an
isomorphism K[cl] 2: K[cq] fixing the elements of K and taking a into ai,
and therefore the a, are also integral over A. Then a,, . . . ,a,~A[cl~, . . . ,a,,],
and hence they are integral over A; but U,EK and A is integrally closed,
so that finally u,EA.
\end{verbatim}
\end{proof}


\begin{theorem}[Printed Theorem 9.3]\label{mcrt-auto-theorem-9-3}\dcref{9.3}\source{matsumura-commutative-ring-theory:page-0081}\uses{}
\begin{verbatim}
Let A be a ring, B an extension ring which is integral over
A and p a prime ideal of A.
   (i) There exists a prime ideal of B lying over p.
   (ii) There are no inclusions between prime ideals of B lying over p.
   (iii) Let A be an integrally closed domain, K its field of fractions, and L a
normal field extension of K in the sense of Galois theory (that is K c L is
algebraic, and for any tl~ L, all the conjugates of c(over K are in L); if B is the
integral closure of A in L then all the prime ideals of B lying over P
are conjugate over K.
\end{verbatim}
\end{theorem}
\begin{proof}
\begin{verbatim}
Localising the exact sequence O-+ A + B at p gives an exact
sequence O-+ A, --+ B, = BOaA, in which B, is an extension ring
89      Integral extensions                                                 67

    integral over A,. From the commutative      diagram
               A,-B
                t         TP
               A-B
    we see that the prime ideals of B lying over p correspond bijectively with the
   maximal ideals of BP lying over the maximal ideal pA, of A,. Hence, to prove
   (i) and (ii) it is enough to consider the case that p is a maximal ideal, which
   has already been done in Lemma 2.
       Now for (iii). Let P, and P, be prime ideals of B lying over p. First of all
   we consider the case [L:K] < cc; let G = (el,. . . ,(T,} be the group of K-
   automorphisms         of L. If P, # a,:?(P1) for any j then by (ii) we have
   p2 + o,:?(PJ,       so that there is an element XEP, not contained in any
   0; ?(PI) for 1 <j < I (see Ex. 1.6). Set y = (njaj(X))?,     where 4 = 1 if char
    K = 0, and 4 = p? for a sufficiently large integer v if char K = p > 0. Then
   y&, and is integral over A, so that YEA. However, the identity map of L is
   contained among the aj, so that YEP,, and hence YEP, n A = p c PI. This
   contradicts a,(x)$P1 for all j. Therefore P, = aj?(P1) for some j.
       If [L:K] = cc we need Galois theory for infinite extensions. Let K? c L be
   the fixed subfield of G = Aut (L/K); then L is Galois over K? and K c K? is a
   purely inseparable extension. If K? # K we must have char K = p > 0, and
   setting A? for the integral closure of A in K? we see easily that
               p? = (x~A?(x~~p for some q = p?>
   is the unique prime ideal of A? lying over p. Thus replacing K by K? we can
   assume that L is a Galois extension of K. For any finite Galois extension
   K c L? contained in L we now set

     then by the case of finite extensions we have just proved, F(E) # 0.
     Moreover, F(Z) c G is closed in the Krull topology. (Recall that the Krull
     topology of G is the topology induced by the inclusion of G into the direct
    product of finite groups n,.Aut (C/K); with respect to this topology, G is
    compact. For details see textbooks on field theory.) If Li for 1 < i < n are
    finite Galois extensions of K contained in L then their composite L? is
    also a finite Galois extension of K, and C)F(Li) 3 F(L?) # 0, so that the
    family (F(L?))L? c L is a finite Galois extension of K) of closed subsets
    of G has the finite intersection property; since G is compact,
     nF(L?) # 0. T a k?mg oe&F(L?)     we obviously have a(Pl) = P,. m
        For a ring A and an A-algebra B, the following statement is called the
 j going-up theorem: given two prime ideals p c p? of A and a prime ideal P of B
    lying over p, there exists P?ESpec B such that P c P? and P?nA = u?.
1:; Similarly, the going-down theorem is the following statement: given p c p?
68       Properties         of extension rings

and P?ESpec B lying over p?, there exists PESpec B such that P c P? and
PnA=p.
\end{verbatim}
\end{proof}


\begin{theorem}[Printed Theorem 9.4]\label{mcrt-auto-theorem-9-4}\dcref{9.4}\source{matsumura-commutative-ring-theory:page-0083}\uses{}
\begin{verbatim}
(i) If B 2 A is an extension ring which is integral over A then
the going-up theorem holds.
   (ii) If in addition B is an integral domain and A is integrally closed, the
going-down      theorem also holds.
\end{verbatim}
\end{theorem}
\begin{proof}
\begin{verbatim}
(i) Suppose p c p? and P are given as above. Since P n A = p, we can
think of B/P as an extension ring of A/p, and it is integral over A/p because
the condition that an element of B is integral over A is preserved by the
homomorphism B -+ B/P. By (i) of the previous theorem there is a prime
ideal of B/P lying over p?/p, and writing P? for its inverse image in B we have
P?ESpec B and P? n A = p?.
   (ii) Let K be the field of fractions of A, and let L be a normal extension
field of K containing B; set C for the integral closure of A in L. Suppose
given prime ideals p c p? of A and P? ESpec B such that P? n A = p?, and
choose Q?ESpec C such that Q?n B = P?. Choose also a prime ideal Q of C
over p, so that using the going-up theorem we can find a prime ideal Q1 of C
containing Q and lying over p?. Both Q1 and Q? lie over p?, so that by (iii) of
the previous theorem there is an automorphism aeAut(L/K) such that
c(Qi) = Q?. Setting o(Q) = Qz we have Qz c Q?, and Q2n A = Qn A = p,
so that setting P=Q,nB             we get PnA=p,         PcQ?nB=P?.        (For
a different proof of (ii) which does not useTheorem 3, (iii), see[AM], (5.16),
or [Kunz].)        q
   We now treat another important casein which the going-down theorem
holds.
\end{verbatim}
\end{proof}


\begin{theorem}[Printed Theorem 9.5]\label{mcrt-auto-theorem-9-5}\dcref{9.5}\source{matsumura-commutative-ring-theory:page-0083}\uses{}
\begin{verbatim}
Let A be a ring and B a flat A-algebra; then the going-down
theorem holds between A and B.
\end{verbatim}
\end{theorem}
\begin{proof}
\begin{verbatim}
Let p c p? be prime ideals of A, and let P? be a prime ideal of B lying
over p?; then B,. is faithfully flat over A,., so that by Theorem 7.3
Spec(B,.) --+ Spec(A,.) is surjective. Hence there is a prime ideal $1,
of B,, lying over PA,, and setting $@nB = P we obviously have P c P?
and PnA=p.       n
\end{verbatim}
\end{proof}


\begin{theorem}[Printed Theorem 9.6]\label{mcrt-auto-theorem-9-6}\dcref{9.6}\source{matsumura-commutative-ring-theory:page-0083}\uses{}
\begin{verbatim}
Let A c B be integral domains such that A is integrally closed
and B is integral over A; then the canonical map f:SpecB-Spec           A is
open. More precisely, for teB, let X? + a,X?- i + . . . + a, be amonic poly-
nomial with coefficients in A having t as a root and of minimal degree;then

        f(W)     = ij %4)>
                      i=l

where the notation D( ) is as in $4.
\end{verbatim}
\end{theorem}
\begin{proof}
\begin{verbatim}
(H. Seydi [4]). By Theorem 2, F(X) = X? + a,X?-?           + ... + a, is
?9        Integral extensions                                                      69

irreducible   over the field of fractions of A; if we set C = A[t] then
c s A cm(w) )>SO C is a free A-module with basis 1, t, t2,. . , tnpl, and
is hence faithfully flat over A. Suppose that PeD(t),        so that PESpecB
with t$P,    and set p = Pn A; then pEUiD(ai), since otherwise aigp for
all i, and SO t?EP, hence t EP, which is a contradiction.
   Conversely, given pEIJiD(aJ, supposethat t EJ(~C); then for sufficiently
large m we have t?? = EYE I bitnpi     with b+p. We can take m > n. Then X??
 - c?j biXnvi is divisible by F(X) in A[X], which implies that X? is divisible
by F(X) = x? + CiiiXn-i      in (A/p)[X];   since at least one of the 2, is non-
zero, this is a contradiction. Thus t$J(pC), so that there exists QESpec C
with t$Q and pC c Q. Setting Qn A = q we have p c q, so that by the
previous theorem there exists P, ESpec C satisfying P, n A = p and P, c Q.
Since B is integral over C there exists PESpec B lying over P,. We have
PeD(t), since otherwise tePnC = P, c Q, which contradicts t$Q. This
proves that


 Any open set of SpecB is a union       of open sets of the form D(r), and hence
f: Spec B -   SpecA is open.
\end{verbatim}
\end{proof}


\begin{theorem}[Printed Theorem 10.1]\label{mcrt-auto-theorem-10-1}\dcref{10.1}\source{matsumura-commutative-ring-theory:page-0087}\uses{}
\begin{verbatim}
Let R c R? c K be as above, let rtt be the maximal ideal ofR
and p that of R?, and suppose that R #R?. Then
      (i)pcmcRcR?andp#m.
    (ii) p is a prime ideal of R and R? = R,.
   (iii) R/p is a valuation ring of the field RI/p.
   (iv) Given any valuation ring Sof the field R?/p, let S be its inverse image
in R?. Then S is a valuation ring having the same field of fractions K as R?.
\end{verbatim}
\end{theorem}
\begin{proof}
\begin{verbatim}
(i)Ifx~pthenx-?4R?,sox-?~Randhencex~R;xisnotaunitofR,
so that xEm. Also, since R # R? we have p # m.
   (ii) We know that p c R, so that p = p n R, and this is a prime ideal of R.
Since R - p c R? - p = {units of R?} we have R, c R?, and moreover by
construction, the maximal ideal of R, is contained in the maximal ideal p of
R?. Thus by (i), R, = R?.
   (iii) Write cp:R? --* R?/p for the natural map; then for XER? - p, if XER
we have cp(x)~R/p, and if x$R we have q(x)-? = cp(x-?)cR/p,                 and
therefore R/p is a valuation ring of RI/p.
   (iv) Note that p c S and S/p = S, so that if XER? and x$S then x is
a unit of R?, and cp(x)$S Thus cp(x-?) = q(x)-?ES, and hence x-i~S.
If on the other hand XEK - R? then x-rep c S, so that we have proved
that SvS-? = K. n
   The valuation ring S in (iv) is called the composite of R? and S. According
to (iii), every valuation ring of K contained in R? is obtained as the
composite of R? and a valuation ring of RI/p.
   Quite generally, we write mR for the maximal ideal of a local ring R. If R
and S are local rings with R 1 S and mR n S = m, we say that R dominates Sj
and write R 3 S. If R >, S and R # S, we write R > S.
\end{verbatim}
\end{proof}


\begin{theorem}[Printed Theorem 10.2]\label{mcrt-auto-theorem-10-2}\dcref{10.2}\source{matsumura-commutative-ring-theory:page-0087}\uses{}
\begin{verbatim}
Let K be a field, A c K a subring, and p a prime ideal of A.
Then there exists a valuation ring R of K satisfying
         RxA      and m,nA=p.
\end{verbatim}
\end{theorem}
\begin{proof}
\begin{verbatim}
Replacing A by A, we can assume that A is a local ring with P = mA.
Now write 9 for the set of all subrings B of K containing A and such that
l$pB. Now AEF, and if $6 c F is a subset totally ordered by inclusion
then the union of all the elements of 9 is again an element of 9, so that, by
General valuations                                                  73
?lO
Zorn?s lemma, 9 has an element R which is maximal for inclusion. Since
pR # R there is a maximal ideal m of R containing pR. Then R c R,,,E~, so
that R = R,, and R is local. Also p c m and p is a maximal ideal of A, so that
m n A = p. Thus it only remains to prove that R is a valuation ring of K. If
 xEK and x#R then since R[x]$F               we have l~-pR[x],      and there
exists a relation of the form
          1 = a, + a,x + ... + a,x? with aiEpR.
Since 1 - a, is unit of R we can modify this to get a relation
   (*) 1 = b,x + ... + b,x? with &em.
Among all such relations, choose one for which n is as small as possible.
Similarly, if x-l $ R we can find a relation
   (**) 1 = ci .x-l + ... + c,,,x-~ with ciEm,
and choose one for which m is as small as possible. If n 2 m we multiply
(**) by b,x? and subtract from (*), and obtain a relation of the form (*)
but with a strictly smaller degree n, which is a contradiction; if 11< m then
we get the same contradiction on interchanging the roles of x and ,Y- l.
Thus if x&R we must have x~~ER.          n
\end{verbatim}
\end{proof}


\begin{theorem}[Printed Theorem 10.3]\label{mcrt-auto-theorem-10-3}\dcref{10.3}\source{matsumura-commutative-ring-theory:page-0088}\uses{}
\begin{verbatim}
A valuation ring is integrally closed.
\end{verbatim}
\end{theorem}
\begin{proof}
\begin{verbatim}
Let R be a valuation ring of a field K, and let XEK be integral over R,
sothatx?+a,x?-?     +...+a,=Owitha,~R.Ifx$Rthenx-?~m,,butthen
         1 +a,x-?   +...+a,x-?=O,
and we get 1 EmR, which is a contradiction. Hence XE R. n
\end{verbatim}
\end{proof}


\begin{theorem}[Printed Theorem 10.4]\label{mcrt-auto-theorem-10-4}\dcref{10.4}\source{matsumura-commutative-ring-theory:page-0088}\uses{}
\begin{verbatim}
Let K be a field, A c K a subring, and let B be the integral
closure of A in K. Then B is the intersection of all the valuation rings of K
containing A.
\end{verbatim}
\end{theorem}
\begin{proof}
\begin{verbatim}
Write 8 for the intersection of all valuation rings of K containing
A, SO that by the previous theorem we have B? 3 B. To prove the opposite
inclusion it is enough to show that for any element XEK which is not inte-
gral over A there is a valuation ring of K containing A but not x. Set
x-l = y. The ideal yA[y] of A[y] does not contain 1: for if 1 = a,y +
a,y2+ .-. + any? with a,eA then x would be integral over A, contradicting
the assumption. Therefore there is a maximal ideal p of A[y] containing
Y~[Y], and by Theorem 2 there exists a valuation ring R of K such that
R=A[y]      and m,nA[y]      = p. Now y = x-l~rn~, so that x$R. w
   Let K be a field and A c K a subring. If a valuation ring R of K contains A
we SaYthat R has a centre in A, and the prime ideal mR n A of A is called the
Centre of R in A. The set of valuation rings of K having a centre in A is called
the Zariski space or the Zariski Riemann surface of K over A, and written
74        Valuarion rings

Zar (K, A). We will treat Zar (K, A) as a topological space, introducing a
topology as follows.
   For x r ,..., x,EK, set U(x, ,..., x,)= Zar(K,A[x, ,..., x,]). Then since
             U(X I,... ,x,)nU(Y,,...,y,)=Uix,,...,x,,y,,       . . . . y,),
 the collection 9 = { U(x,, . . ,x,)ln > 0 and X,EK 3 is the basis for the
 open sets of a topology on Zar (K, A). That is, we take as open sets the
 subsets of Zar (K, A) which can be written as a union of elements of 9,
 As in the case of Spec, this topology is called the Zariski topology.
\end{verbatim}
\end{proof}


\begin{theorem}[Printed Theorem 10.5]\label{mcrt-auto-theorem-10-5}\dcref{10.5}\source{matsumura-commutative-ring-theory:page-0089}\uses{}
\begin{verbatim}
Zar (K, A) is quasi-compact.
\end{verbatim}
\end{theorem}
\begin{proof}
\begin{verbatim}
It is enough to prove that if .d is a family of closed sets of Zar (K, A)
 having the finite intersection property (that is, the intersection of any finite
 number of elements of .d is non-empty) then the intersection of all the
 elements of & is non-empty. By Zorn?s lemma there exists a maximal family
 of closed sets &? having the finite intersection property and containing d.
 Since it is then enough to show that the intersection ofall the elements of&
 is non-empty, we can take LX!?= &?. Then it is easy to see that & has the
following properties:
    (4 F ,,...,F,Ed=>F,n...nF,E~,
    (p) Zr,...,Z,, are closed sets and Z,u~~~uZ,,~.d+Z~~d                for some i;
    (y) if a closed set F contains an element of d then FE&.
 For a subset F c Zar (K, A) we write F? to denote the complement of F.
 If FEd and F? = UnUA then F = nJJ;, and moreover if u(x,, . . , x,)~ =
(Jr=, U(Xi)CE~ then by (a) ab ove one of the U(x,)? must belong to d.
Hence the intersection of all the elements of d is the same thing as the
intersection of the sets of the form U(x)c belonging to d. Set
            1- = (ydc u(y-?)?Ed.).
Now since the condition for REZar(K, A) to belong to U(y-?)?                   is that
yEmR, the intersection        of all the elements of d is equal to
            {REZar(K,A)lm,         =I r}.
Write I for the ideal of A[lJ generated by r. If 1 $1 then by Theorem 2, the
above set is non-empty, as required to prove. But if 1 EZ then there is a
finite subset (y,, . . , y,} c r such that 1 EC y,A[y,, . . . , y,]; but then
 U(y; ?)?n ... n U(y; l)c = 0, which contradicts the finite intersection
property of &.         n
    When K is an algebraic function field over an algebraically closed field k
of characteristic 0 (that is, K is a finitely generated extension of k), Zariski
gave a classification of valuation rings of K containing k, and using this and
the compactness result above, he succeeded in giving an algebraic proof in
characteristic 0 of the resolution of singularities of algebraic varieties of
dimension 2 and 3. However, Hironaka?s general proof of resolution ef
General valuations                                                    75

  singularities in characteristic 0 in all dimensions was obtained by other
  methods, without the use of valuation theory.
     As we saw at the beginning of this section, the ideals of R form a totally
  ordered set under inclusion. This holds not just for ideals, but for all R-
  modules contained in K. In particular, if we set
             G=(xR[xEK       and x#O],
  then G is a totally ordered set under inclusion; we will, however, give G the
  opposite order to that given by inclusion. That is, we define 6 by
           xR ,( yRoxR       =I yR.
  Mvreover, G is an Abelian group with product (xR)-(yR) = xyR. In general,
  an Abelian group H written additively, together with a total order relation
  2 is called an ordered group if the axiom
           x2y,z>t*x+z>y+t
  holds. This axiom implies
           (1) x>O,     y>O*x+y>O,          and (2) x>y+-y>-x.
  We make an ordered set Hu {co) by adding to H an element 00 bigger
  than all the elements of H, and fix the conventions co + LX= 00 for CCEH
  and a + a = cc. A mapv:K-Hu{co}               from a field K to HU{EJ}
  is called an additive valuation or just a valuation of K if it satisfies the
  conditions
      (1) VbY) = v(x) + v(y);
   :, (2) vb + y) 3 min { 44, V(Y));
    ? (3) v(x)= 000x=0.
  If we write K* for the multiplicative       group of K then u defines a
  bmomorphism       K * -H;    the image is a subgroup of H, called the value
  BtOUps of 0. We also set
   1.
               R?= {x~Klv(x) 3 0) and m, = (XEK~V(X) > 0},
     obtaining a valuation ring R, of K with m, as its maximal ideal, and call R,
     tie valuation ring of v, and m, the valuation ideal of u. Conversely, if R is a
   ?-balGation ring of K, then the group G = {xRlx~K*} described above is an
 ??tiered group, and we obtain an additive valuation of K with value group
  .l.pby defining v:K ---+ G u {CO} by u(0) = co and v(x) = xR for XEK* (there
   &Qn0 real significance in whether or not we rewrite the multiplication         in
::&
.*_I additively); the valuation ring of v is R. The additive valuation corres-
,-f:pMng        to a valuation ring R is not quite unique, but if u and v? are
*Y?$?oadditive valuations of K with value groups H and H? and both having
-$fie
 ;.y,r, valuation ring R then there exists an order-isomorphism        cp:H -H?
 .2@eh that v? = cpv(prove this!). Thus we can think of valuation rings and
         itive valuations as being two aspects of the same thing.
76       Valuation rings

   We now give some examples of ordered groups:
   (1) the additive group of real numbers R (this is isomorphic to the
multiplicative       group of positive reals), or any subgroup of this;
   (2) the group Z of rational integers;
   (3) the direct product Z? of n copies of Z, with lexicographical order,
that is
                                           the first non-zero element of
            (a,,...,a?)<(b,,...,b,)o       b -al,...,b     -a is positive
                                             1            n n
   An ordered group G is said to be Archimedean if it is order-isomorphic to
a suitable subgroup of R. The name is explained by the following theorem:
the condition in it is known as the Archimedean axiom. (Note that our
usage is completely unrelated to that in number theory, where non-
Archimedean fields are p-adic fields, as opposed to subfields of (wand @with
the usual ?Archimedean? metrics.)
\end{verbatim}
\end{proof}


\begin{theorem}[Printed Theorem 10.6]\label{mcrt-auto-theorem-10-6}\dcref{10.6}\source{matsumura-commutative-ring-theory:page-0091}\uses{}
\begin{verbatim}
Let G be an ordered group; then G is Archimedean if and
only if the following condition holds:
    if a, LEG with a > 0, there exists a natural number n such that na > b.
\end{verbatim}
\end{theorem}
\begin{proof}
\begin{verbatim}
The condition is obviously necessary, and we prove sufficiency. If
G = (0) then we can certainly embed G in R. Suppose that G # {O}.
Fix some 0 < XEG. For any LEG, there is a well-defined largest integer
n such that nx d y (if y 2 0 this is clear by assumption; if y < 0, let m be
the smallest integer such that - y d mx, and set n = - m). Let this be no.
Now set Y, = y - nOx and let n, be the largest integer n such that nx < 10 y,;
we have 0 d n, < 10. Set y2 = lOy, - nix and let n2 be the largest integer
n such that nx < IOy,. Continuing in the same way, we find integers no,
n1,n2,..., and set q(y) = CI,where tl is the real number given by the decimal
expression n, + O.n,n,n,. . . . Then it can easily be checked that cp:G -+ R
is order-preserving, in the sense that y < y? implies q(y) < q(y?).
   We also see that cp is injective. For this, we only need to observe that
if y < y? then there exists a natural number r such that x < lo*(y? -Y);
the details are left to the reader.
   Finally we show that cpis a group homomorphism.         For LEG, we write
n/10? with neZ to denote the number obtained by taking the first r
decimal places of q(y); the numerator n is determined by the property that
nx < 10?~ < (n + 1)x. For ~?EG, if n?x < 10?~? < (n? + 1)x then we have
           (n + n?)x d lol(y + y?) < (n + n? + 2)x,
and hence
           q(y+y?)-(n+n?).lO-?<2.10-?,
so that
           l~(Y+Y?)-~(Y)-~cp(Y?)l~4.~o-?~
and since r is arbitrary, q(Y + y?) = q(Y) + &?).     B
General valuations                                                                      77
     ?lO
       A non-zero group G order-isomorphic to a subgroup of Iwis said to have
     rank 1. The rational rank of an ordered group G of rank 1 is the maximum
     number of elements of G (viewed as a subgroup of [w) which are linearly
     independent over Z. For example, the additive group G = Z + ZJ2 c iw is
     an ordered group of rank 1 and rational rank 2.
\end{verbatim}
\end{proof}


\begin{theorem}[Printed Theorem 10.7]\label{mcrt-auto-theorem-10-7}\dcref{10.7}\source{matsumura-commutative-ring-theory:page-0092}\uses{}
\begin{verbatim}
Let R be a valuation ring having value group G. Then G has
     rank IoR has Krull dimension 1.
\end{verbatim}
\end{theorem}
\begin{proof}
\begin{verbatim}
(a) Since G # 0, R is not a field. Supposethat p is a prime ideal ofR
     distinct from mR. Let 5EmR be such that t$p, and set v(t) = x, where u is the
     additive valuation corresponding to R. Suppose that 0 # nip, and set JI =
     u(q);then yeG and x > 0, so that nx > y for some sufficiently large natural
     number n. This meansthat <?/PIER, so that YEI] R c p; then sincep is prime
     we have REP, which is a contradiction. Therefore p = (Cl).The only prime
     ideals of R are mR and (0), which means dim R = 1.
        (CL) If 0 # VEmR then mR is the unique prime ideal containing qR, and
     hence ,/(ryR) = mR. Thus for any 5~rrt~ there exists a natural number n such
     that l?~qR. From this one seeseasily that G satisfies the Archimedean
     axiom. n
\end{verbatim}
\end{proof}


\begin{theorem}[Printed Theorem 11.3]\label{mcrt-auto-theorem-11-3}\dcref{11.3}\source{matsumura-commutative-ring-theory:page-0095}\uses{}
\begin{verbatim}
Let R be an integral domain and I a fractional ideal of R.
Then the following conditions are equivalent:
   (1) I is invertible;
   (2) I is a projective R-module;
   (3) I is finitely generated, and for every maximal ideal P of R, the
fractional ideal I, = IR, of R, is principal.
\end{verbatim}
\end{theorem}
\begin{proof}
\begin{verbatim}
(l)-(2) If I-?I = R then there exist a@ and b&?             such that
x;aibi=     l.Thena,,..., a, generate I, sincefor any x~l we have C(xbi)ai =
D VRs and Dedekind rings                                                 81
 ?ll
  x, and xb,eR. Let F = Re, + ... + Re, be the free R-module with basis
  el,. . , e,,; we define the R-linear map cp:F -I       by cp(ei) = ai, so that cp is
  surjective. Then we defined $:I -F          by writing $i:I -R         for the map
  tii(x) = b,x, and setting I)(X) = C $;(x)ei. We then have q@(x) = x, SO that cp
  splits, and I is isomorphic to a direct summand of the free module F, and
  therefore projective.
       (2)+-(l) Every R-linear map from I to R is given by multiplication         by
  some element of K (prove this!). If we let cp:F - I be a surjective map from
  a free module F = @ Rei, by assumption there exists a splitting $:I + F
  such that cp$ = 1. Write $(x) = 1 A;( x )e,.f or x~l; then by what we have said,
  each Ai determines a b,EK such that &(x) = bix, and since for each x that
  are only finitely many i such that /$(x) # 0, we have hi = 0 for all but finitely
  many i. Letting b,, . , b, be the non-zero ones, we have c aibix = x for all
  xel, where ai = q(ei). Thus c; a,b, = 1. Moreover, since hiI = Ai c R we
  have biGI- ?, and therefore I ~? I = R.
      (l)=>(3) As we have already seen, I is finitely generated. Now if c aibi = 1
 and P is any prime ideal then at least one of aibi must be a unit of R,,
 and I, = a,R,. Hence I, is a principal fractional ideal.
      (3)*(l)    If I is finitely generated then (IP?),=(Z,))?.         Indeed, the
inclusion c holds for any ideal; for 3, if I = a, R + ... + a,R and x@Z,)) ?
 then xa,ER,, so there exist c+R - P such that xa,c,ER, so that setting
 C=C 1...c, we have (cx)a,ER for all i, which gives cxeZ-?           and x@~?),.
 From the fact that I, is principal, we get ZP.(lP)- ? = R,. Now if II- ? # R
 then we can take a maximal ideal P such that II-? c P, and then
 Zp.(Zp)- l = z,jz - ?)p c PR,, which is a contradiction.       Thus we must have
II-?=R.           n

Theorem 1 I .4. Let R be a Noetherian integral domain, and P a non-zero
prime ideal of R. If P is invertible then ht P = 1 and R, is a DVR.
Proof. If P is invertible the maximal ideal PR, of R, is principal, and
condition (3) of Theorem 2 is satisfied; thus R, is a DVR, and so dim R, = 1.

 Theorem I1 S. Let R be a normal Noetherian domain. Then we have
   (i) all the prime divisors of a non-zero principal ideal have height 1;
   (ii) R = nhtp=, R,.
p?!       (i) Suppose 0 # aER and that P is one of the prime divisors of aR;
then there exists an element bE R such that aR: b = P. We set PR, = nr, and
then aR,:b=m,           so that baPIEmP?   and ba-?q!R,.  If ba-?mcm    then
bY the determinant trick ba-? is integral over R,, which contradicts the
fact that R, is integrally closed. Thus ba- ?m = R,, so that m- ?m = R,, and
then by the previous theorem we get ht m = ht P = 1.
   (ii) It is sufficient to prove that for a, bER with a ~0, bEaR, for every
82       Valuation rings

height 1 prime PESpec R implies bEaR. Let Pr,..., P, be the prime
divisors of aR, and let aR = q, n . . . nq, be a primary decomposition of
aR, where qi is a Pi-primary ideal for each i. Then since ht Pi = 1, we have
bEaR,,nR=q,      for i= l,..., n, and therefore bE fi qi = uR. H
\end{verbatim}
\end{proof}


\begin{theorem}[Printed Theorem 11.6]\label{mcrt-auto-theorem-11-6}\dcref{11.6}\source{matsumura-commutative-ring-theory:page-0097}\uses{}
\begin{verbatim}
For an integral domain R the following conditions are
equivalent:
   (1) R is a Dedekind ring;
   (2) R is either a field or a one-dimensional Noetherian normal domain;
   (3) every non-zero ideal of R can be written as a product of a finite
number of prime ideals.
    Moreover, the factorisation into primes in (3) is unique.
\end{verbatim}
\end{theorem}
\begin{proof}
\begin{verbatim}
(l)-(2)     Every non-zero ideal is invertible, and therefore finitely
generated, so that R is Noetherian. Let P be a non-zero prime ideal of R;
then by Theorem 4, the local ring R, is a DVR and ht P = 1, and therefore
either R is a field or dim R = 1. Also, by Theorem 4.7 we know that R is the
intersection of the R, as P runs through all the maximal ideals of R, but
since each R, is a DVR it follows that R is normal.
   (2) *(l) If R is a field there is no problem. If R is not a field then for every
maximal ideal P of R the local ring R, is a one-dimensional Noetherian
local ring and is normal, so that by Theorem 2 it is a principal ideal ring.
Thus by Theorem 3, R is a Dedekind ring.
   (l)*(3) Let I be a non-zero ideal. If I = R then we can view it as the
product of zero ideals; if I is itself maximal then it is the product of just one
prime ideal. We have already seen that R is Noetherian, so that we can use
descending induction on 1, that is assume that I #R and that every ideal
strictly bigger than J is a product of prime ideals. If I # R then there is a
maximal ideal P containing I, and I c IPpl c R. If IP-l = I then using
P-'P = R we would have I = IP, and by NAK this would lead to a
contradiction. Hence ZP-l #I, so that by induction we can write IP-? =
Qi.. . Q,, with Q,&pec R. Multiplying both sides by P gives I = Qi.. . Q*P.
   The proof of (3) +(l) is a little harder, and we break it up into four steps.
D VRs and Dedekind. rings                                           83
?ll

   Step 1. In general, any non-zero principal ideal aR of an integral domain
R is obviously invertible. Moreover, suppose that I and J are non-zero
fractional ideals and B = IJ; then obviously I and J invertible implies B
invertible, but the converse also holds. TO see this, from I-?J-?B    c R we
get I-?J-?   c B-? 9 and also from B-?ZJ c R we get B-?Z c J-? and
B-?J C I-?; now if B is invertible then multiplying the last two inclusions
together we get B -? =Bm?B-lIJcl-?J-l,            and hence B-? =I-?Jpl.
Therefore
         R=BB-?=IJI-?J-?=(II-?)(JJ-?),
and we must have II-?           = JJ-?   = R.

   SreP 2. Let P be a non-zero prime ideal. Let us prove that if I is an ideal
strictly bigger than P then IP = P. For this it is sufficient to show that if
I = p + aR with a$P then P c IP. Consider expressions of 1? and a2R + P
as product of prime ideals, 1? = P,. . . P, and a2R + P = Q1 . . . Q,. Then Pi
and Qj are prime ideals containing I, and so are prime ideals strictly bigger
than P, We now set R = R/P, and write - to denote the image in R of
elements or ideals of R. Then we have
          (*) P,... Pr = a2R = (il.. Q,,
and applying Step 1 to the domain R we find that pi and gj are all invertible,
and are prime ideals of 8. We can suppose that P, is a minimal element of
the set {Pl,..., P?>. Moreover, at least one of Q, ,. . . ,Q, is contained
in P,, so that we can assume that Q, c P,, and, on the other hand, since
Q1 is also a prime ideal and P r...prccQ1 we must have Q,3Bi                for
some i. Then pi c Qi c B,, and by the minimality        of Pi we have pi =
Pl = Q,. M u hi p 1 .
                   ymg through both sides of (*) by P; ? gives
        P2~..P~=Q2...Qs.
proceeding in the same way, we see that Y = s, and that after reordering the
Qi we can assume that Pi = Qi for i = 1,..., r.FromthiswegetPi=Qi,and
a2R+P=(P+aR)2        =P2 +aP+ a2R. Thus any element XGP can be
Written
         x=y+az+a?t        with YEP?, ZEP and tER.
Smce a#P we must have td,              and then as required we have
PcP2+aP=(P+aR)P.
  Step 3. Let bER      be a non-zero element; then in the factorisation
bR=p 1 . . . P,, every Pi is a maximal ideal of R. Indeed, if I is any ideal
stfictlY greater than Pi then ZPi = Pi, and by Step 1 Pi is invertible, so that
I-R.
   SteP 4. Let P be a prime ideal of R, and 0 # aGP. If aR = P,. . . P, with
PieSPec R then P must contain one of the Pi, but from Step 3 we know that
Pi is maximal, so that P = Pi. We deduce that P is a maximal ideal and is
84           Valuation rings

invertible. If every non-zero prime ideal is invertible then any non-zero ideal
of R is invertible, since it can be written as a product of primes. This
completes the proof of (3)*(l).
   Finally, if(l), (2) and (3) hold, then as we have seen in Step 2 above, the
uniqueness of factorisation into primes is a consequence of the fact that
the prime ideals of R are invertible.      n
\end{verbatim}
\end{proof}


\begin{theorem}[Printed Theorem 11.7]\label{mcrt-auto-theorem-11-7}\dcref{11.7}\source{matsumura-commutative-ring-theory:page-0099}\uses{}
\begin{verbatim}
the Krull-Akizuki    theorem). Let A be a one-dimensional
Noetherian integral domain with field of fractions K, let L be a finite
algebraic extension held of K, and B a ring with A c B c L; then B is a
Noetherian ring of dimension at most 1, and ifJ is a non-zero ideal of B then
B/J is an A-module of finite length.
\end{verbatim}
\end{theorem}
\begin{proof}
\begin{verbatim}
We follow the method of proof of Akizuki [l] in the linear algebra
formulation of [BS]. First of all we prove the following lemma.
\end{verbatim}
\end{proof}


\begin{theorem}[Printed Theorem 12.1]\label{mcrt-auto-theorem-12-1}\dcref{12.1}\source{matsumura-commutative-ring-theory:page-0101}\uses{}
\begin{verbatim}
If A is a Krull ring and S c A a multiplicative set, then A, is
again Krull. If F = (R,},,, is a defining family of A then the subfamily
Pnln,r, where I- = (1~A1 R, 1 A,) is a defining family of A,.
\end{verbatim}
\end{theorem}
\begin{proof}
\begin{verbatim}
Setting m, for the maximal ideal of R, we have
        kreSn      m, = 0.
Let 0 # x~r)~,,R,;   there are at most finitely many REA such that Us < 0;
let A= (Al,..., 1n} be the set of these. If SEA then ,%$I-,hence we can find
t,Em, n S. Replacing tn by a suitable power, we can assumethat u,(t,x) b 0.
We then set t = nn,, t A, so that for every IeA we have u,(tx) 3 0, and
therefore txeA; but on the other hand tES so that XEA, and we have
proved that A, I nnEr R I. The opposite inclusion is obvious. The finiteness
condition (2) holds for A, so also for the subset r. n
   Krull rings defined by a finite number of DVRs have a simple structure.
Lemma 1 (Nagata). Let K be a field and R, , . . . , R, valuation rings of K;
Krull rings
 ?I*
 set A = 0 Ri. Then for any given a~ K there exists a natural number s 3 2
 such that
          (l+a+-~+a?-l)-?       and a.(1 +a+...+~?-~)-?
 both belong to A.
 pr&       We consider separately each Ri. Note first that (1 - a) (1 + a + ...
  +~-l)=l-u?.Ifu~Rithenu             -?urns, and any s >, 2 will do. If UER~ then
 Provided that there does not exist t such that 1 - u?~rn~, any s 2 2 will do.
 of I- aEmi then any s which is not a multiple of the characteristic of Ri/mi
 will do. If on the other hand 1 - u$mi but 1 - u?ern, for some t > 2, letting
 to be the smallest value of t for which this happens, we see that 1 - aSEmi
 only for s multiples of to, so that we only have to avoid these. Thus for each i
 the bad values of s (if any) are multiples of some number di > 1, so that
 choosing s not divisible by any of these di we get the result. H

 Theorem 12.2, Let K be a field and R,, . . . , R, valuation rings of K such
that Ri qkRj for i fj; set m, = rad (Ri). Then the intersection A = nl= 1Ri is
a semilocal ring, having pi = m, n A for i = 1,. . . , II as its only maximal
ideals; moreover A,, = R,. If each Ri is a DVR then A is a principal ideal
ring.
Proof. The inclusion A, c Ri is obvious. For the opposite inclusion, let
UCRi; choosing s > 2 as in the lemma, and setting u = (1 + a + ... + us- ?)- 1
we get UEA and UEA. Obviously u is a unit of Ri, so that UEA -pi
and a = (uu)/u~A,~. This proves that A,< =Ri. It follows from this
that there are no inclusions among pl, . . . , p,. If I is an ideal of A not
contained in any pi then (by Ex. 1.6) there exists xgl not contained in
u~=Ipi; then x is a unit in each Ri, and hence in A, SO that I = A. Thus
Pl,..., p, are all the maximal ideals of A.
   If each Ri is a DVR then we have mi # rn;, and hence pi # pi*?, (where p(*)
denotes p?A,n A). Thus there exists xi~pi such that x&1?, and xi$pj for
i #j; then pi = x,A. If1 is any ideal of A and ZRi = xyi Ri for i = 1,. . . , n then it
is easy to see that I = XII.. . x,yA. n
   If a Krull ring A is defined by an infinite number of DVRs then the
defining family of DVRs is not necessarily unique, but the following
theorem tells us that among them there is a minimal family.
\end{verbatim}
\end{proof}


\begin{theorem}[Printed Theorem 12.3]\label{mcrt-auto-theorem-12-3}\dcref{12.3}\source{matsumura-commutative-ring-theory:page-0102}\uses{}
\begin{verbatim}
Let A be a Krull ring, K its field of fractions, and p a height 1
lPdmeidealofA;thenif,9=={R       1>IsA is a family of DVRs of K defining A, we
 must have A,EB. If we set 9, = {A,Ip~Spec A and htp = l} then .9-o is a
 defining family of A. Thus 9-O is the minimal defining family of DVRs of A.
 %of. By Theorem 1, A, is a KrulI ring defined by the subfamily Y1 =
 (R,IA, c R,} c 9; if A, c R, then the elements of A - p are units of R,
88       Valuation rings

SO that p 3 m,nA.    If nt,n A = (0) then RI 3 K which is a contradiction,
hence m,n A # (0); since ht p = 1, we must have p = mln A. Thus if we fix
some 0 # x~p, then ul(x) > 0 for all RIeFI,    and hence F-, is a finite set,
Now by the previous theorem and Ex. 10.5, the elements of 9, correspond
bijectively with maximal ideals of A,, and 9, has just one element A,. Thus
A,EF, in other words 8, c 9.
   To prove that 9-O is a defining family of DVRs of A it is enough to show
that A 1 n,,,,=, A,. That is, it is enough to prove the implication
         for a, kA    with a # 0, beaA, for all A,ES,*~EUA.
As one sees easily, this is equivalent to saying that aA can be written as
the intersection of height 1 primary ideals. The set of REB such that
aR # R is finite, so we write R,, . . , R, for this. If we set
           aRinA = qi and rad(R,)nA = pi
then qi is a primary ideal belonging to pi for each i, and aA = ql.. . n q,.
Eliminating redundant terms from this expression, we get an irredundant
expression, say aA = q1 n. . . nq,. It is enough to show that then ht pi = 1
for 1 < i < r. By contradiction, suppose that ht p1 > 1. By Theorem 1, A,,1is
a Krull ring with defining family F? = {REBIA,,, c R), but is not itself
a DVR, and hence by Theorem 2, F? is infinite. Thus there exists R?6P-I
such that aR? = R?; we set p? = rad (R?)n A. We have a$p?, and ApI c R?
implies that p? c pl. Now by assumption aA #q, n...nq,.,         and RI is a
DVR, so that (rad(R,))? caR, for some v >O, and hence pi cql.
Therefore there exists an i 3 0 such that
           aA+p?,nq,n...nq,       and aAIpl+?nq,n...nq,
 Hence there exist by A such that b$aA but bp, c aA. In particular bp? c aA,
 but since a is a unit of R? we have
           (b/a)p? c An rad (R?) = p?.
 Taking 0 # cep? then for every n > 0 we have (b/a)?cEp? c A, and since A is
 completely integrally closed, b/aEA. This is a contradiction, and it proves
 that ht pi = 1 for 1 6 i < r. H
\end{verbatim}
\end{theorem}
\begin{proof}
\notready
\end{proof}


\begin{theorem}[Printed Theorem 12.4]\label{mcrt-auto-theorem-12-4}\dcref{12.4}\source{matsumura-commutative-ring-theory:page-0103}\uses{}
\begin{verbatim}
(i) A Noetherian normal domain is a Krull ring.
    (ii) Let A be an integral domain, K its field of fractions, and L an extension
field of K. If (Ai)i,, is a family of Krull rings contained in L satisfying
?12      Krull rings                                                           89

   the two conditions (1) A = nAi and (2) given any 0 # aeA we have
   a,-& = Ai for all but finitely many i, then A is a Krull ring.
       (iii) If A is a Krull ring then so is A[X] and A[Xj.
   ,J+oof. (i) 7%?IS f o 11ows from Theorem 11.5 and the fact that for any non-
   zero UEA there are only finitely many height 1 prime ideals containing aA
   (because these are the prime divisors of aA).
       (ii) is easy, and we leave it to the reader.
       (iii) K [X] .1s a p rincipal ideal ring and therefore a Krull ring. Moreover,
    if we let 9 be the set of height 1 prime ideals of A then for YES
   the ideal p[X] is prime in A[X], and by Theorem 11.2, (3), the local ring
   AL-Xl,,x, is a DVR of K(X). (If we write u for the additive valuation of K
   corresponding       to the valuation ring A,, we can extend u to an additive
   valuation of K(X) by setting u(F(X)) = min (u(q)} for a polynomial
            F(X) = a, f u,X + ... + u,X? (with u,EK),
   and v(F/G) = u(F) - v(G) for a rational function F(X)/G(X); then the
   ,valuation ring of u in K(X) is A [XI,,,,.)    Now we have K[X] A
   ACW,,x, = A,[X] (prove this!), and so


   by (ii) this is a Krull ring.
      Now for A[Xl], let CR,},,,, be a family of DVRs of K defining A; then
   inside Ic[yXj we have A[X]l=      nnRn[X$      also by Ex. 9.5, RA[Xj is an
   integrally closed Noetherian ring, and is therefore a Krull ring by (i).
   However, we cannot use (ii) as it stands, since X is a non-unit of all the rings
   Z&[Xj, so we set R,[XI][X-I]         =Bb and note that A[Xj =K[X]n
   (n$J; now the hypothesis in (ii) is easily verified. Indeed,
             ~~(X)=U,X?+U,+,X?~+~..EA~XI]            with a,#0
   is a non-unit of B, if and only if a, is a non-unit of R,, and there are only
   finitely many such j?. Therefore A[Xj is a Krull ring. n

   Remark 2. Note that the field of fractions of A[Xl        is in general smaller
   than the field of fractions of K[TXJ.

   Remark 2. The B, occurring above are Euclidean rings ([B7], 9 1, Ex. 9).
\end{verbatim}
\end{theorem}
\begin{proof}
\notready
\end{proof}


\begin{theorem}[Printed Theorem 12.5]\label{mcrt-auto-theorem-12-5}\dcref{12.5}\source{matsumura-commutative-ring-theory:page-0104}\uses{}
\begin{verbatim}
The notions of Dedekind ring and one-dimensional Krull
    ring coincide.
\end{verbatim}
\end{theorem}
\begin{proof}
\begin{verbatim}
A Dedekind ring is a normal Noetherian domain, and therefore a
    KruII ring. Conversely, if A is a one-dimensional Krull ring, let US prove that
.r -4 is Noetherian. Let I be a non-zero ideal of A, and let 0 # u~l. If we can
?- Prove that AIaA is Noetherian then IluA is finitely generated, and thus so is
;:.
a
90       Valuation rings

1. By the corollary of Theorem 3 we can write aA = q, n.*.nq,,    where
qi are symbolic powers of prime ideals pi and pi # pj if i #j; now since
dim A = 1 each pi is maximal and we have
         A/aA = A/q, x .? x A/q,
by Theorem 1.3 and Theorem 1.4. But A/q, is a local ring with maximal
ideal pi/qi, and hence A/qi % A,,/qi A,,; now since each A, is a DVR, A/aA is
Noetherian (in fact even Artinian). Hence A is one-dimensional Noetherian
integral domain, and is normal, and is therefore a Dedekind ring. w
\end{verbatim}
\end{proof}


\begin{theorem}[Printed Theorem 12.7]\label{mcrt-auto-theorem-12-7}\dcref{12.7}\source{matsumura-commutative-ring-theory:page-0105}\uses{}
\begin{verbatim}
Y. Mori and J. Nishimura).       Let A and 9? be as in the
previous theorem. If A/p is Noetherian        for every pe.9 then A is
Noetherian.
\end{verbatim}
\end{theorem}
\begin{proof}
\begin{verbatim}
(J. Nishimura). As in the proof of Theorem 5, it is enough to show
that A/p??? is Noetherian for p ~9 and any n > 0. If n = 1 this holds by
hypothesis. For n > 1 we proceed as follows. Applying the previous
theorem with r = 1 and e = - 1 we can find an element x in the field Of
fractions of A such that u,(x) = 1 and u,(x) < 0 for every other qE9. Set
B= .4[x]. If yip then y/x~A so that y~xxB, and conversely Bc A, and
xBcpA,, so that p=xBnA.       Moreover, B=A+xB        and so
         BfxB = A/p.
Krull rings                                                                   91
ii12
Now x?B/x?+ ? B ?v B/xB for each i, so that by induction on i we see that
B/x?B is a Noetherian B-module for each i, and hence a Noetherian ring.
NOW we have


and B/x?B is a finite A/(x?B n A)-module, being generated by the imagesof
l,X,...,P,     so that by the Eakin-Nagata theorem (Theorem 3.7),
A/(~?B~A)    is Noetherian ring; therefore its quotient A/p(?) is also
Noetherian. W
\end{verbatim}
\end{proof}


\begin{theorem}[Printed Theorem 13.2]\label{mcrt-auto-theorem-13-2}\dcref{13.2}\source{matsumura-commutative-ring-theory:page-0109}\uses{}
\begin{verbatim}
Let R = BnaoR, be a Noetherian graded ring with R,
Artinian, and let M be a finitely generated graded R-module. Suppose that
R= RO[xl,..., x,] with xi of degree di, and that P(M, t) is as above. Then
P(M, t) is a rational function of t, and can be written

        P(M, t) = f(t)/ fi (1 - tdi),
                        i=l

where f(t) is a polynomial with coefficients in Z.
\end{verbatim}
\end{theorem}
\begin{proof}
\begin{verbatim}
By induction on r?, the number of generators of R. When r = Q we
have R = R,, so that for n sufficiently large, M, = 0, and the power series
P(M, t) is a polynomial. When Y> 0, multiplication       by x, defines an Ro-
linear map M, - M, + d,; writing K, and L, +d,   for the kernel and cokernely
we get an exact sequence



Set K = OK,, and L = @ L,. Then K is a submodule of M, and L = MIxrM9
                                                                   tK
so that K and L are finite R-modules; moreover x,K = x,L = 0 SOtha

                                                                                 i
Graded rings, Hilbert function       and Samuel function                     95
913
and L can be viewed as R/x,R-modules,      and hence we can apply the
induction hypothesis to P(K, t) and P(L, t). Now from the above exact
sequence we get
        W,) - 4M?) + w, +d,) - 4L +d,) = 0.
If we multiply this by tnfdp and sum over n this gives
          tdrP(K, t) - tdrP(M, t) + P(A4, t) - P(L, t) = g(t),
where g(t)EZ[t].    The theorem follows at once from this. n
   A lot of information on the values of 1(M,) can be obtained from the
above theorem. Especially simple is the case d, = ... = d, = 1, so that R is
generated over R, by elements of degree 1. In this case P(M, t) = f(t)
(I - t)-?; if f(t) has (1 - t) as a factor we can cancel to get P in the form
          P(M,t)=f(t)(l        - t)-d with     ,f~Z[t],         d>O,
          and if d>O           then ,f(l)#O.
If this holds, we will write d = d(M). Since (1 - t)- 1 = 1 + t + t2 + . . . , we
can repeatedly differentiate both sides to get



(This can of course be proved in other ways, for example by induction on d.)
Iff(t)=a,+a,t+...+a,ts        then


(*I      W,) = a, (d~rr?)+a~(d:~~2)+...+a~(d+~~~-1);


                          = 0 for m < d - 1. The right-hand            side of (*) can be

formally rearranged as a polynomial        in n with rational coefficients, say cp(n);
then
                    f(l)
          ?(X)=(d- pXd-? I)!        + (terms of lower degree).

Since              coincides     with   the polynomial          m(m - 1). . . (m - d + 2)/
(d-   I)! for m >, 0, this implies the following      result.
\end{verbatim}
\end{proof}


\begin{theorem}[Printed Theorem 13.3]\label{mcrt-auto-theorem-13-3}\dcref{13.3}\source{matsumura-commutative-ring-theory:page-0113}\uses{}
\begin{verbatim}
Let A be a semilocal Noetherian ring, and 0 + M? -
M -M?   + 0 an exact sequence of finite A-modules; then
         d(M) = max(d(M?),       d(M?))
If I is any ideal of definition of A, then XL - xfu? and xf\l, have the same
 leading coefficient.
Pro@. We can assume M? = MJM?. Then since M?/l*Mn = MJ(M? + I?M)
we have
         l(M/I?M)   = l(M/M? + I?M) + l(M? + I?M/I?M)
                    = I(M?/I?M?) + l(M?/M? n I?M).
Thus setting q(n) = I(M?/M?n        I?+ ? M), we have & = XL,, + 9. Since
moreover both &, and cptake on only positive values, d(M) coincides with
whichever is the greater of d(M?) and deg cp. However, by the Artin-Rees
lemma, there is a c > 0 such that
         n~~~l?~lM?cM?nl?~lMcI?~c~lM?,
and hence
         xL,(n) 2 q(n)>   x',,(n - c);
therefore cpand XL, have the same leading coefficient. n
    We now define a further measure 6(M) of the size of M: let 6(M) be the
smallestvalue of n such that there exist x1,. .. ,X,EIII for which l(M/x, M +
... + x,M) < a. When l(M) < rx, we interpret this as 6(M) = 0. If I is any
ideal of definition of A then l(M/IM) < co, so that 6(M) < number of
generators of I. Conversely, in the case that A is a local ring and M = A,
then 1(A/Z) < co implies that I is an m-primary ideal. Therefore in this case
6(A) is the minimum of the number of generators of m-primary ideals.
    We have now arrived at the fundamental theorem of dimension theory.
\end{verbatim}
\end{theorem}
\begin{proof}
\notready
\end{proof}


\begin{theorem}[Printed Theorem 13.4]\label{mcrt-auto-theorem-13-4}\dcref{13.4}\source{matsumura-commutative-ring-theory:page-0113}\uses{}
\begin{verbatim}
Let A be a semilocal Noetherian ring and M a finite
A-module; then we have
         dim M = d(M) = 6(M).
\end{verbatim}
\end{theorem}
\begin{proof}
\begin{verbatim}
Step 1. Each of d(M) and 6(M) are finite, but the finiteness of dimM
has not yet been established. First of all, let us prove that d(A) 3 dim A
for the caseM = A, by induction on d(A). Set m = rad (A). If d(A) = 0 then
l(A/m?) is constant for n >>0, so that for some n we have m? = m?+? and
by NAK, mn = 0. Hence any prime ideal of A is maximal and dim A = 0.
O+A/po    %4/p,       -B+O,

we have d(B) < d(A), and so by induction
         dimB<d(B)<d(A)-              1.
(The values of d(B) and dim B are independent of whether we consider B as
an A-module or as a B-module, as is clear from the definitions.) In B, the
image of p1 c ... c p, provides a chain of prime ideals of length e - 1, so that
         e- 1 ddimB<d(A)-          1;
hence e < d(A). Since this holds for any chain of prime ideals of A, this
proves dim A < d(A). For general M, by Theorem 6.4, there are submodules
MisuchthatO=M,cM,c...cM,=Mwith
         M,/M,-    1 E A/p,    and         p,eSpecA.
Since for an exact         sequence        O-+ M? -M      +      M?+O      of finite   A-
modules we have
         Supp(M)       = Supp(M?)uSupp(M?)

and
         dim M = max (dim M?, dim M?),
it is easy to see that
         d(M) = max {d(A/p,)}         3 max {dim(A/p,)}         = dim M.
 Step2. We show that 6(M) 3 d(M). IfG(M) = Othen l(M) < CCso that x,(n)
is bounded, hence d(M) = 0. Next suppose that 6(M) = s > 0, choose
Xl ,...,x,EITI such that l(M/x,M+.~~+x,M)<~,      and set Mi=M/X,M+
..* + XiM; then clearly 6(Mi) = 6(M) - i. On the other hand,
         ~(M,/v~?M,)      = I(M/x,M        + m?M)
                          = I(M/m?M)         - I(x,M/x,Mnm?M)
                          = I(M/m?M)        - I(M/(m?M:x,))
                          3 1(M/m?M)        - /(M/m?- ?M),
so that d(M,) > d(M) - 1. Repeating this, we get d(M,) 2 d(M) - s, but since
6(M,) = 0 we have d(M,) = 0, so that s 3 d(M).
   Step 3. We show that dim M 3 6(M), by induction on dim M. If dim M = 0
then Supp(M) c m-Spec A = b?(m) so that for large enough n we have
m? C ann M, and l(M) -C co, therefore 6(M) = 0. Next suppose that
dim M > 0, and let pi for 1 < i ,< t be the minimal prime divisors of ann (M)
with coht pi = dim M; then the pi are not maximal ideals, so do not contain
tn. Hence we can choose X~E~ not contained in any pi. Setting
M, = M/x,M we get dim M, < dim M. Therefore by the inductive hypo-
thesis 6(M,) 6 dim M,; but obviously 6(M) d @MI) + 1, so that 6(M) 6
dimM, + 1 <dimM.        n
\end{verbatim}
\end{proof}


\begin{theorem}[Printed Theorem 13.5]\label{mcrt-auto-theorem-13-5}\dcref{13.5}\source{matsumura-commutative-ring-theory:page-0115}\uses{}
\begin{verbatim}
Let A be a Noetherian ring, and I = (a,, . . ,a,) an ideal
generated by Y elements; then if p is a minimal prime divisor of I we have
ht p < Y. Hence the height of a proper ideal of A is always finite.
\end{verbatim}
\end{theorem}
\begin{proof}
\begin{verbatim}
The ideal IA, c A, is a primary ideal belonging to the maximal
ideal, so that ht p = dim A, = 6(A,) < Y. w
\end{verbatim}
\end{proof}


\begin{theorem}[Printed Theorem 13.6]\label{mcrt-auto-theorem-13-6}\dcref{13.6}\source{matsumura-commutative-ring-theory:page-0115}\uses{}
\begin{verbatim}
Let P be a prime ideal of height r in a Noetherian ring A.
Then
   (i) P is a minimal prime divisor of some ideal (a,, . . , a,) generated by
r elements;
   (ii) if b 1,. . . , b,EP we have ht P/(b,, . . . , b,) 3 r - s;
   (iii) if a,, . . . , a, are as in (i) we have
            htP/(a,,...,u,)=r-i            for ldibr.
\end{verbatim}
\end{theorem}
\begin{proof}
\begin{verbatim}
(i) A, is an r-dimensional local ring, so that by Theorem 4 we can
choose r elements a,, . . . , u,EPA~ such that (a,, . , u,)A, is PAP-primary.
Each ui is of the form an element of P times a unit of A,, so that without
loss of generality we can assume that a+P. Then P is a minimal Prime
divisor of (a,, . . . , u,)A.
   (ii) Set A = A/(b,, . . . , b,), B = P/(bI , . . , b,) and ht P = t. Then by (i),
                 Graded rings, Hilbertfunction     and Samuel function                   101
     813
     there   exist   Cl,..., C,EP    such that P is a minimal prime divisor of
     (El,..., t ;)A. Then    P  is a  minimal prime divisor of (b,, . . . , b,, cl,. . . ,cJ,
     and hence Y d s + t by Theorem 5.
        (iii) The ideal P/(al , . . , a i) is a minimal prime divisor of (a,, 1,. . , &)
     kA/(al,...,     ai), hence ht P/(al, . , ai) d r - i. The opposite inequality was
     proved in (ii). W
\end{verbatim}
\end{proof}


\begin{theorem}[Printed Theorem 13.7]\label{mcrt-auto-theorem-13-7}\dcref{13.7}\source{matsumura-commutative-ring-theory:page-0116}\uses{}
\begin{verbatim}
Let A = @nBOAn be a Noetherian graded ring.
         (i) If I is a homogeneous ideal and P is a prime divisor of I then P is also
     homogeneous.
         (ii) If P is a homogeneous prime ideal of height Y then there exists a
     sequence P = P, 1 P, 3 ... 3 P, of length r consisting of homogeneous
     prime      ideals.
     J+YK$ (i) P can be expressed in the form P = ann (x) for a suitable element
     x of the graded A-module A/Z. Let UEP, and let x = x0 +x1 + ... + x,
     anda=a,+a,+,            + ... + a, be decompositions into homogeneous terms.
     Then since ax = 0,
                apxO = 0,   apxl +a,+,x,=O,          a,x,+a,+,x,       +ap+2x0=0,...,
     from which we get six, = 0, six, = 0,. . . , and finally a;+?~ = 0. It follows
      that aif EP, but since P is prime, apEP. Thus ap+l + ... + a,EP, so
    ?that in turn ap+ l~P. Proceeding in the same way, we see that all the homo-
     geneous     terms of a are in P, so that P is a homogeneous ideal.
        (ii) First of all note that we can assume that A is an integral domain.
     To see this, if we take a chain P = p0 2 ... 1 p, of prime ideals of length
     r then p, is a minimal prime divisor of (0), and so by (i) is a homogeneous
     ideal; so we can replace A by A/p,. Now choose a homogeneous element
     0 # b, EP; then by Theorem 6, ht (P/b, A) = r - 1, and so there is a minimal
    .prime divisor Q of b, A such that ht (P/Q) = r - 1; since Q # (0), it is a
     height 1 homogeneous prime ideal. By the inductive hypothesis on r
     applied to P/Q there exists a chain P = P, I P, I ... I> P,- 1 = Q of homo-
  ? geneousprime ideals of length r - 1, and adding on (0) we get a chain
    ,Of length r. w
,./
 :.     Let us investigate more closely the relation between local rings and
;- @aded rings.
-!l Theorem 13.8. Let k be a field, and R = k[<,,. . ., ~$1 a graded ring
1%. generated by elementstl,. . ,<, of degree 1; set M = c 5iR, A = R, and
*.**.
pgj mEMA.
?%     (i) Let x be the Samuel function of the local ring A, and cpthe Hilbert
   .function of the graded ring R; then q(n) = x(n) - x(n - 1);
  j (ii)dimR=htM=dimA=degcp+l;
      (iii) gr,(A) N R as graded rings.
\end{verbatim}
\end{theorem}
\begin{proof}
\begin{verbatim}
M is a maximal        ideal of R so that
            m?/m?+?     N Mn/Mni     1 2: R,;
hence x(n) - x(n - 1) = l(mn/m?+l) = 1(R,) = q(n), and so dim A = deg x = 1 +
deg cp. Then since A = R,, we have dim A = ht M. After this it is enough
to prove that dim R = ht M. First of all, assume that R is an integral
domain, so that by Example 3 in the section on Hilbert functions and by
Theorem 5.6, we have
          l+degcp>tr.deg,R=dimR>,htM;
putting this together with ht M = dim A = 1 + deg cp, we get dim R = ht M.
Next for general R, let P,, . . . , P, be the minimal prime ideals of R; then by
Theorem 7, these are all homogeneous ideals, and each R/Pi is a graded
ring. Choosing P, such that dim R = dim R/P, and using the above result,
we get
         dim R = dim R/P, = ht M/P, < ht M ,< dim R,
so that dim R = htM as required. We have R, c M? c m? with m?/m?+ ? 1
R,, and so taking an element x of R, into its image in m?/m?+ ? we obtain
a canonical one-to-one map R 1gr,,,A,    and it is clear from the definition
that this is a ring isomorphism.  n
\end{verbatim}
\end{proof}


\begin{theorem}[Printed Theorem 13.9]\label{mcrt-auto-theorem-13-9}\dcref{13.9}\source{matsumura-commutative-ring-theory:page-0117}\uses{}
\begin{verbatim}
Let (A, m, k) be a Noetherian local ring, and set G = gr,,A;
then dim A = dim G.
\end{verbatim}
\end{theorem}
\begin{proof}
\begin{verbatim}
Letting q be the Hilbert polynomial         of G, we have dim A =
 1 + deg cp (by Theorem 4), and by the previous theorem this is equal to
dim G. n
    In fact the following more general theorem holds: for I a proper ideal in a
Noetherian local ring A, set G = gr,(A); then dim A = dim G. This will be
proved a little later (Theorem 15.7).
\end{verbatim}
\end{proof}


\begin{theorem}[Printed Theorem 13.10]\label{mcrt-auto-theorem-13-10}\dcref{13.10}\source{matsumura-commutative-ring-theory:page-0119}\uses{}
\begin{verbatim}
Eagon). Let A be a Noetherian ring and M be an r x s
matrix (I d s) of elements of A. Let Ir be the ideal of A generated by the t x t
minors of M. If P is a minimal prime divisor of I, then we have
           ht P d (r - t + l)(s - t + 1).
\end{verbatim}
\end{theorem}
\begin{proof}
\begin{verbatim}
Induction on r. When Y = 1 we have t = 1, and so (r - t + l)(s - t +
 1) = s. The ideal I, is generated by s elements, so that the assertion is just
the principal ideal theorem (Theorem 5) in this case. Next assume that r > 1.
Localising at P we may assume that A is a local ring with maximal ideal P,
and that I, is P-primary.
    If t = 1, then I, is generated by rs elements and (r - t + 1) (s - t + 1) = rs,
so our assertion holds also for this case. Therefore we assume t > 1. If at
least one of the elements of M is a unit of A, then by what we said above, I, is
generated by (t - 1) x (t - 1) minors of a (r - 1) x (s - 1) matrix, and again
we are done. Therefore we assume that all the elements of M are in P. Now
comes the brilliant idea. Let M? be the matrix with elements in B = A[X]
obtained from M by replacing a, 1 by a,, + X, and let I? be the ideal of B
generated by the t x t minors of M?. Since t > 1 and aijgP for all i and j we
have I? c PB. We also have I? + XB = Z,B + XB since both sides have the
same image in B/XB = A. Therefore PB is a minimal prime divisor of I? by
the lemma. Since the element a, 1 + X of M? is not in PB, we have
ht PB < (r - t + l)(s - t + 1) by our previous argument. Since ht PB = ht P,
as we can see by Theorems 4 and 5, we are done. n

      14   Systems of parameters    and multiplicity
            Let (A, m) be an r-dimensional Noetherian local ring; by Theorem
13.4, there exists an m-primary ideal generated by r elements, but none
generated by fewer. If a,,. . . , a,~nt generate an m-primary             ideal,
Iu  I,?.., a,}  is  said to  be   a system    of parameters  of  A  (sometimes
abbreviated to s.0.p.). If M is a finite A-module with dim M = s, there exist
Yl,..., y,Em such that l(M/(y, ,..., y,)M) < cc, and then (y, ,..., y,} is
said to be a system of parameters of M.
   If we set A/m = k, the smallest number of elements needed to generate
m itself is equal to rank,m/m?;       (here rank, is the rank of a free module
over k, that is the dimension of m/m2 as k-vector space). This number is
called the embedding dimension of A, and is written emb dim A. In general
           dim A < emb dim A,
\end{verbatim}
\end{proof}


\begin{theorem}[Printed Theorem 14.1]\label{mcrt-auto-theorem-14-1}\dcref{14.1}\source{matsumura-commutative-ring-theory:page-0120}\uses{}
\begin{verbatim}
Let (A, m) be a Noetherian local ring, and x1,. . . , x, a system
  of parameters. Then
      (i) dim A/(x,, . . . , xi) = r - i for 1 < i < r.
      (ii) although it is not true that ht (x1,. . . , Xi) = i for all i for an arbitrary
  system of parameters, there exists a choice of x1,. . . , x, such that every
  subset Fc(xr,...,          x,] generates an ideal of A of height equal to the
 number of elements of F.
\end{verbatim}
\end{theorem}
\begin{proof}
\begin{verbatim}
(i) is contained in Theorem 13.6. We now prove the second half of
 (iii. If r< 1 the assertion is obvious; suppose that r > 1. Let poj (for 1 <
l<ee) be the prime ideals of A of height 0. Choosing xrom not contained
 in any poj, we have ht (xi) = 1. Next letting plj (for 1 <j 6 e,) be the minimal
 prime divisors of (xi), SO that ht Plj = 1, and choosing xZEm not contained
 in any poj or any plj, we have ht(xJ = 1, ht(x,,x,) = 2; if r = 2 we?re
 done. If r > 2 we choose x3Em not contained in any minimal prime
 divisor of (0), (x,), (x,), (x1, x,), and proceed in the same way to obtain
 the result.
      We now give an example where ht (x1,. . , xi) < i. Let k be a field and
 set R = k[X, Y, Zj; let I = (X)n( Y, Z), and write A = R/Z, and x, y, z for
 the images in A of X, Y, Z. The minimal prime ideals of A are (x) and
 TV,z); now A/(X) = R/(X) N k[lx Z] is two-dimensional                    and A/(y, z) N
 R/( y, Z) N k[X]      .is one-dimensional,    so that dim A = 2. {y, x + Z} is a
 system of parameters of A; in fact xy = xz = 0. so that x2 =x(x + a)~
(.Y,x + Z) and z* = z(x + z)~(y,x + z). However, y is contained in the
.minimal prime ideal (y, z) of A, and hence ht (y) = 0. n
\end{verbatim}
\end{proof}


\begin{theorem}[Printed Theorem 14.2]\label{mcrt-auto-theorem-14-2}\dcref{14.2}\source{matsumura-commutative-ring-theory:page-0120}\uses{}
\begin{verbatim}
Let (R, m) be an n-dimensional                    regular local ring, and
 Xl , . . . , Xi elements of m. Then the following conditions are equivalent:
     (1) Xl,... ,Xi is a subset of a regular system of parameters of R;
     (2) the images in m/m* of x r , . . . , xi are linearly independent over R/m;
     (3) R/(x,,... ,xJ is an (n - i)-dimensional regular local ring.
 ~00f.(1)~(2)1fx,,...,Xi,Xi+1,...           , x, is a regular system of parameters then
 their images generate m/m* over k = R/m, and since rank,m/m* = n they
 must be linearly independent over k.
     (I)*(3)              We know that dim R/(x l,...,xi)=n-ii,      and the images of
x1+1, . . . , x, generate the maximal ideal of R/(x,, . . . , xi).
     (3)=>(l) If the maximal ideal m/(x1,. . . , xi) of R/(x,, . . . , xi) is generated
bY the images of y I,..., y,-iEm                then m is generated by x1,. ..,xi,
:Y1,...,yn-i.
\end{verbatim}
\end{theorem}
\begin{proof}
\notready
\end{proof}


\begin{theorem}[Printed Theorem 14.3]\label{mcrt-auto-theorem-14-3}\dcref{14.3}\source{matsumura-commutative-ring-theory:page-0121}\uses{}
\begin{verbatim}
A regular local ring is an integral domain.
\end{verbatim}
\end{theorem}
\begin{proof}
\begin{verbatim}
Let (R,m) be an n-dimensional regular local ring; we proceed by
induction on n. If n = 0 then m is an ideal generated by 0 elements, so
that m = (0). This in turn means that R is a field. Thus a zero-dimensional
regular local ring is just a field by another name.
   When n = 1, the maximal ideal m = xR is principal and ht m = 1, so that
there exists a prime ideal p # m with m 3 p. If yen we can write y = xa with
aER, and since x$p we have aEp; hence p = xp, and by NAK, p = (0). This
proves that R is an integral domain. (There is a slightly different proof in the
course of the proof of Theorem 11.2; as proved there, a one-dimensional
regular local ring is just a DVR by another name.)
   When n> 1, let pi,... ,pI be the minimal prime ideals of R; then since
m $ m2 and m + pi for all i, there exists an element xEm not contained in
anyofm2,p,,.     . . , p,. (see Ex. 1.6). Then the image of x in m/m? is non-zero,
so that by the previous theorem R/xR is an (n - 1)-dimensional regular
local ring. By the induction hypothesis, R/xR is an integral domain, in
other words, xR is a prime ideal of R. If pi is one of the minimal prime ideals
contained in xR then since x$pi, the same argument as in the n = 1 case
shows that p1 = xp,, and hence p1 = (0). n
\end{verbatim}
\end{proof}


\begin{theorem}[Printed Theorem 14.4]\label{mcrt-auto-theorem-14-4}\dcref{14.4}\source{matsumura-commutative-ring-theory:page-0121}\uses{}
\begin{verbatim}
Let (A, m, k) be a d-dimensional       regular local ring; then
          gr,,V) 2: kCX,, . . .,XJ,
and if x(n) is the Samuel function of A then
                   n+d
          x(n)=      d          for all n 3 0.
                     (     >
\end{verbatim}
\end{theorem}
\begin{proof}
\begin{verbatim}
Since m is generated by d elements, gr,(A) is of the form
4X i , . . . , X,1/1, where I is a homogeneous ideal. Now if I # (0) let SEI be a
non-zero homogeneous element of degree r; then for n > r the homogeneous
piece of k[X]/Z       of degree n has length         at most

                     which is a polynomial       of degree d - 2 in n. This
implies that the Samuel function of A is of degree at most d - 1, and
contradicts dim A = d. Hence I = (0); the second assertion follows from
the first. n
   Let (A, m) be a Noetherian local ring. Elements yl,. . . ,y,Em are said to be
analytically independent if they have the following property; for every
\end{verbatim}
\end{proof}


\begin{theorem}[Printed Theorem 14.5]\label{mcrt-auto-theorem-14-5}\dcref{14.5}\source{matsumura-commutative-ring-theory:page-0122}\uses{}
\begin{verbatim}
Let (A,m) be a d-dimensional Noetherian local ring and
  xr,.i.,x,, a system of parameters of A; then x1,. . .,xd are analytically
  independent.
\end{verbatim}
\end{theorem}
\begin{proof}
\begin{verbatim}
Set q = xxiA. Since q is an ideal of definition of A, by
   Theorem 13.4, xi(n) = l(A/q?) is a polynomial of degree d in n for n >>0. Set
  ,@t = k; we say that a homogeneous form f (X)Ek[X1,.      . . ,X,] of degree n
  is a null-form    of q if F(x,....x,)~q?m     for any homogeneous form
  F(X)EACXl, . . ., X,] which reduces to f(X) modulo m. Write n for the ideal
  of&X,,...,   X,] generated by the null-forms of q. Then
              k[xlln = 0 q?/q?m,
  and writing cp for the Hilbert polynomial            of k[X]/n, we have q(n)
  = I(q?/q?m) for n >>0. The right-hand side is just the number ofelements in a
  minimal basis of q?, so that cp(n).l(A/q) >, l(q?/q?+?).   Now
              4q?lq??)         = x34   - x?dn - 1)
  is a polynomial in n of degree d - 1, so that deg q > d - 1, but if n # (0) this
  is impossible. Thus n = (0), and the statement in the theorem follows at
  once. n

  lirurtiplicity
   &et (A, m) be a d-dimensional Noetherian local ring, M a finite A-module,
   and q an ideal of definition of A (that is, an m-primary ideal). As we saw
   ti $13, the Samuel function ,!(M/q?+? M) = x$(n) can be expressed for
:. 5~0 as a polynomial in n with rational coefficients, and degree equal to
   dim M, and therefore at most d. In addition, this polynomial can only take
   tnteger values for n >>0, so it is easy to see by induction on d (using the fact
   that x(n + 1) - x(n) has the same property) that

              Xg,(n) = grid + (terms of lower order),

b.%th eeZ. This integer will be written e(q, M). By definition       we have the
'3dlOwing          property.
 :>A.=
C?
.::.Fom~la 14.1. e(q, M) = n+m
                             lim cl(M/q?M),   and in particular, if d = 0 then
                                  n*
$jyk M) = Z(M).
       rom this we see easily the following:
      rmula 14.2. e(q, M) > 0 if dim M = d, and e(q, M) = 0 if dim M < d;
     rmula 14.3. e(q?, M) = e(q, M)rd;
Formula 14.4. If q and q? are both m-primary                   ideals and q 3 q? then
eh M) d 44, Ml.
   We set e(q,A) = e(q),and define this to be the multiplicity of q. In addition,
we will refer to the multiplicity e(m) of the maximal ideal as the multiplicity
of the local ring A, and sometimeswrite e(A) for it. For example, if A is a
regular local ring then by Theorem 4, we can see that e(A) = 1.
\end{verbatim}
\end{proof}


\begin{theorem}[Printed Theorem 14.6]\label{mcrt-auto-theorem-14-6}\dcref{14.6}\source{matsumura-commutative-ring-theory:page-0123}\uses{}
\begin{verbatim}
Let O+ M? -M                -    M?+O       be an exact sequence of
finite A-modules. Then
        e(q, Ml = e(q,M?) + e(cr,M?).
\end{verbatim}
\end{theorem}
\begin{proof}
\begin{verbatim}
We view M? as a submodule of M. Then
         /(M/q? M) = I(M?/q? M?) + I(M?/M? n q? M),
and obviously q?M? c M?nq?M. On the other hand by ArtinRees, there
exists c > 0 such that
         M?nq?Mcq       n-CM?        for all   y1> C.
Hence
         /(Ml/q?-?M?)   < l(M?/M?n       q?M) < l(M?/q?M?).
From this and Formula 14.1 it follows easily that

         e(q, M) - e(q,M?) = lim f$(M?/M?n              q?M)   = e(q, M?).
                             n-sn
\end{verbatim}
\end{proof}


\begin{theorem}[Printed Theorem 14.7]\label{mcrt-auto-theorem-14-7}\dcref{14.7}\source{matsumura-commutative-ring-theory:page-0123}\uses{}
\begin{verbatim}
Let (pl,. . . ,p,} be all the minimal prime ideals of A such
that dim A/p = d; then



where si denotesthe image of q in A/p, and 1(M,) standsfor the length of M, as
A,-module.
Proqf (taken from Nagata [Nl]). We write CJ= &l(M,) and proceed by
induction on o. If CJ= 0 then dim M < d, so that the left-hand sideis 0, and
the right-hand side is obviously 0; now supposec > 0. Now there is some
p~{pi,. . ,p,} for which M, #O; then p is a minimal element of
Supp(M). Hence p~Ass(M), that is M contains a submodule N isomorphic
to A/p. Then
        e(q, W = 4% N) + e(s, MIN.
On the other hand, N, N Ap/pA, and N,: = 0 for pi # p, so that 1(N,) = 1,
and the value of g for M/N has decreasedby one, so that the theorem holds
for M/N. However, from the definition
         4% N) = e(q, A/p) = e(9,A/p), where 9 = (q + P)/P.
Putting this together, we seethat the theorem also holds for M. l
  Theorem 7 allows usto reduce the study of e(q, M) to the casethat A isan
T&orem 14.9. Let (A, m) be a Noetherian local ring, q an ideal of definition
   ofA,andx,,...,    xd a system of parameters of A contained in q. Suppose that
   xiEqv* for 1 < i 6 d. Then for a finite A-module M and s = 1,. . . ,d we have
             e(q/h ,..., x,),M/(xl ,..., x,)M)3v,v,...v,e(q,M).
   In particular if s = d, we have
               I(M/(x,,      . ,qJM) 3 v1v2.. . vde(q,M).
\end{verbatim}
\end{theorem}
\begin{proof}
\begin{verbatim}
It is enough to prove the case s = 1. We set A? = A/x, A,q? = q/x, A,
   &f=M/x,M            and v=vl.      By Theorem 1, we have dimA?=d-              1. On the
   other hand,
               l(M?/q??M?)= l(M/x, M + q?M)
                          = l(M/q?M) - 1(x, M + q?M/q?M).
     In addition, in view of (x1 M + q?M)/q?M 2: x1 M/x, Mnq?M                        N Ml
     (q?M:x,) and q?-?M c qnM:xl, we have
  .?           - 1(x, M + q?M/q?M) 3 - l(M/q?-?M),
 ?I and   therefore
     I
               &Vf'/q'"M')     > l(M/q?M)    - l(M/q?-?M).
 i;, When n >>0 the right-hand side is of the form
-7-t
   ?/       eh Ml
 d<
;p.?        T[nd        - (n - v)~] + (polynomial of degree d - 2 in n)

                   = pv.nd-l           + (polynomial     of degree d - 2 in n),

        that the assertion is clear. n
       A case of the above theorem which is particularly          simple, but important,


         eorem 14.10. Let (A, m) be a d-dimensional Noetherian local ring, let
             xd be a system of parameters of A, and set q = (x1,. . . ,x,); then

       d if in addition       xiEm? for all i then l(A/q) 2 v?e(m).

           orem 14.11. Let A, m, xi and q be as above. Let M be a finite A-module,
           set A? = A/x, A, M? = M/x, M and q? = q/x1 A = c$xiA?. Then ifx, is a
non-zero-divisor      of M, we have the following                   equality
          4-r, M) = 44, W.
Proof.   Since l(M?/q??+ 1M?) = l(M/x, M + q?+ 1 M) we have
          1(M/q? + r M) - l(M?/q?n+ r M?)=I(x,M+q?+?M/q?+?M)
            = l(x,M/x,Mnq?+?M)         = l(M/(q?+?M:x,))
            = l(M/q?M) - l((q?+?M:x,)/q?M).
On the other hand, setting a = x$xiA we have q = x,A + a and q?+? =
xlqn+a?+l,      and therefore
           qntl M:x, = q?M + (a?+?M:x,).
Moreover, by Artin-Rees, there is a c > 0 such that for n > c we have
a?+lMnx,M=a?-c(ac+?        Mnx,M),     and therefore a?+rM:xl c an-CM.
Thus
         (q?+?M:x,)/q?M   = (q?M +(a?+?M:x,))/q?M
                          c (q?M + a?-?M)/q?M
                          ~a?-?M/a?-?Mnq?M.
NOW an-?M/a?-?M n q?M is a module over A/q?, and since a is generated
                                            n-c+d-2
by d - I elements, an-c is generated by                   elements. Thus
                                                d-2     >
for n > c we have



where m is the number of generators of M. The right-hand side is a
polynomial of degree d - 2 in II, so that
          eh?, M ?) = (d - l)! lim l(M?/qln+ ? M?)/n?- ?
                               n-cc
                    = (d - l)! lim [l(M/q?+?M)                 - l(M/q?M)]/nd-?
                               n-30
                    = e(q, M). H
\end{verbatim}
\end{proof}


\begin{theorem}[Printed Theorem 14.12]\label{mcrt-auto-theorem-14-12}\dcref{14.12}\source{matsumura-commutative-ring-theory:page-0125}\uses{}
\begin{verbatim}
Lech?s lemma). Let A be a d-dimensional Noetherian
local ring, and x r , . . , ,xd a system of parameters of A; set q = (x1,. . .,xd), and
suppose that M is a finite A-module. Then
                                    l(M/(x;?,         . , x,y?)M)
          e(q, M) =           lim
                      min(vi)+m             VI...Vd
\end{verbatim}
\end{theorem}
\begin{proof}
\begin{verbatim}
If d = 0 then both sidesare equal to l(M). Ifd = 1 then the right-hand
side is exactly Formula 14.1 which defines e(q, M). For d > 1 we use
induction on d.
   Setting Nj= {meMlxjm=O)          we have N, c N, c . .. so that there is a
c > 0 such that N, = N,, 1 = .*.. If we set M? = xt M then x, is a non-zero-
divisor for M?, and there is an exact sequence0 + N, -M           -M?     40.
Since N, is a module over A/x?, A we have dim N, < d, and therefore
e,q, &I) = e(q, M?). On the other hand,
            l(M/(x;?, . . . ,xdYd)M) - l(M?/(x;?,.    . . ,xF)M?)
                = l(N, + (xy, . . ,xp)M/(x;?,           . ,xi?)M)
                = l(N,/N, n (x;? , . . , x;?)M)
                d l(N,/(x;?, . . . ,x;?)N,).
  If v1 7 c then x;l N, = 0, and N, is a module over the (d - 1)-dimensional
  Local ring A/x;A, so that by induction there is a constant C such that as
  min (vi) + CC we have
             l(N,l(x;?, . >x2) NJ = l(N,/(x?,Z, . . . , x;?) N,) < CT,. . . vd.
  Therefore,
            lim [l(M/(xT?, . . . , xi?)M) - l(M?/(xr?, . . . , x~?)h?f?)]/v,. . . vd = 0.
  This means that we can replace M by M? in the theorem, and so we can
  assume that x1 is a non-zero-divisor in M. Then by the previous theorem we
  have e(q, M) = e(q, J@),with q = q/x, A and ~ = M/x, M. If we furthermore
  set
           E=(xT,...,   xF)M      and F = M/E
  then by Theorem 9, we have
           e(q,   M).v, . ..vd < l(M/(x?,?,       . ,xi?)M)     = l(F/x;?F)

                               = i$ 1(x;- l F/x?,F)           ~v,~(F/x,F)=v,~(M/x,M+E)
                               = VI l(~/(X~)       . . ,xJy)Al).
   Then by induction on d we have
          lim l(M/(x;l,. . . , x;?)M)/v,.
                                       . . vd = lim l(R/(x;?, . . . , xdyd)fl)/vz.. . vd
                                              = e(q, M).     n

     Although we will not use it in this book, we state here without proof a
  remarkable result of Serre which shows that multiplicity can be expressed
  as the Euler characteristic of the homology groups of the Koszul complex
  (discussed in 5 16).
\end{verbatim}
\end{proof}


\begin{theorem}[Printed Theorem 14.13]\label{mcrt-auto-theorem-14-13}\dcref{14.13}\source{matsumura-commutative-ring-theory:page-0127}\uses{}
\begin{verbatim}
Let (A, m) be a Noetherian local ring, q an m-primary ideal
and b a reduction of q; then b is also m-primary, and for any finite A-module
M we have
         e(q, W = 46 Ml.
\end{verbatim}
\end{theorem}
\begin{proof}
\begin{verbatim}
If q?+l = bq? then q?+? c b c q, hence b is also m-primary.
Moreover,
         l(M/b?+?M)         3 I(M/q?+?M)   = l(M/b?q?)   3 f(M/b?M),
so that e(q, M) = e(b, M) follows easily.
\end{verbatim}
\end{proof}


\begin{theorem}[Printed Theorem 14.14]\label{mcrt-auto-theorem-14-14}\dcref{14.14}\source{matsumura-commutative-ring-theory:page-0127}\uses{}
\begin{verbatim}
Let (A, m) be a d-dimensional Noetherian local ring, and
suppose that A/m is an infinite field; let q = (u,, . . . , u,) be an nt-primary
ideal. Then if yi = xaijuj for 1 < i < d are d ?sufficiently general? linear
combinations of u Ir.. . , us, the ideal b = (y 1,. . . , yd) is a reduction of q and
 {yl,. . , yd} is a system of parameters of A.
\end{verbatim}
\end{theorem}
\begin{proof}
\begin{verbatim}
If d = 0 then q? = (0) for some Y > 0, hence (0) is a reduction of q
so that the result holds. We suppose below that d > 0.
    Step 1. Set A/m = k and consider the polynomial ring k[X,, . . . , X,] (or
k[X] for short). For a homogeneous form q(X) = I&X,, . . . ,X,)EA[X]               of
degree n, we write @(X)Ek[X]         for the polynomial obtained by reducing
the coefficients of q modulo m. As in the proof of Theorem 5 we say that
@(X)Ek[X]       is a null-form of q if cp(u1,. . . ,u,)Eq?m; this notion depends
not just on q, but also on ul,. .., u,. However, for fixed @ it does not
depend on the choice of cp.We write Q for the ideal of k[X] generated by
all the null-forms of q, and call Q the ideal of null-firms of q. One sees
easily that all the homogeneous elements of Q are null-forms of q, and
that the graded ring k[X]/Q has graded component of degree n isomorphic               )
to qn/qnm, so that we have
         KXIIQ = @q"/q"m = gr,(40A,,k.
                      n,O

Write q(n) for the Hilbert function of k[X]/Q; then
          44 = 4q?lq?m) f h?/q?+ ?1 G cpW Wd
(see the proof of Theorem 5). We know that for n >>0, the function l(q?/q??)
is a polynomial in n of degree d - 1 (where d = dim&            Thus from the
above inequality, cp is also a polynomial          of degree d - 1, so that by
Theorem 13.8, (ii), we have dimk[X]/Q         = d.
   Now set I/ = ct kX,, and let P,, . . . , p, be the minimal prime divisors Of
i=l

Hence we can take a linear form 1,(X)~l/ not belonging to any Pi. If
d 71 then similarly we can take ~,(X)EV such that 12(X) is not contained
m any minimal prime divisor of (Q, II(X)), and, proceeding in the same
way, we get II(X), . . .,&(X)E V such that (Q, 1,, . . . ,1,) is a primary ideal
belonging to (XI, . . . ,X,).
   Step 2. We let b be the ideal of A generated by linear combinations
                                                                   t


L,(U) = C       (for 1 < i < of u I , . . , u, with coefficients in A. Then if we
              Uijuj                     t)


f@t li(X)=ti(Xj    =CaijXj3   a necessary and sufficient condition for b to
be a- reduction of qis -that the ideal (Q, 1,). . . ,1,) of k[X] is (X,, . . . ,X,)-
primary.
Proof of necessity. Suppose that bqr = q?+l. Then if M = M(X ) is a
monomial of degree r + 1 in X,, . . . ,X,, we can write



where the F,(X) are homogeneous                    forms of degrees r with coefficients in
A. Thus
        ii?(X) - c li(X)Fi(X)gQ.
Hence

         (X l,...,Xs)rfl~(Q,ll,...,lt).
Proof of sufficiency.                We go through     the same argument    in reverse: if
~~--~I,F,EQ    then
         M(U)-              CL~(U)F~(U)E~?+~ITI,
so that qr+l c bqr + q?+?nt; thus by NAK, q*+? = bq?.
    Step 3. Putting together Steps 1 and 2 we see that q has a reduction
b=Cyl,... ,y,J generated by d elements. Both q and its reduction b are
m-primary ideals, so that y,, . . . ,y, is a system of parameters of A. We
are going to prove that there exists a finite number of polynomials D,(Zij)
for 1< c(6 v in sd indeterminates Z, (for 1 < i d d and 1 <j < s) such that
d linear combinations yi = C aijuj (for 1 < i < d) generate a reduction ideal
of 9 if and only if at least one of D,(aij) # 0. (The expression ?d sufficiently
general linear combinations? in the statement of the theorem is quite vague,
hut in the present case it has a precise interpretation as above,)
   Let G,(X) , . . . ,G,(X) be a set of generators of Q, with Gj homogeneous
of degree e> For any sd elements aij of k (for 1 < i f d and 1 <j < s), set
hw) = c aijxi. w e write I, for the homogeneous component of degree n
of a homogeneous ideal I c k[X,, . . . ,X,1, and in particular               we write
(X 1,. . . ,X,), = V,, so that
            w 1I..., XJc(Q,4   ,..., Id)0Vn=(Q,11,...,Id)n.
Set c, = dim, V,. We have

          (Q, 1l,...,Z~,={CliFi+        CGjHjlFiEV,-,           and    HjEV/,-,j}.

Let K1,..., K, be the elements obtained as xliFi + c GjHj as the Fi run
through a basis of V,-, and the Hj run independently through a basis
of Vnpe,; it is clear that they span (Q, I,, . . . ,I& Each of K,, . . . ,K, is a
linear combination of the c, monomials of degree n in the Xi, with linear
functions in the rxij as coefficients; we write out these coefficients in a
c, x w matrix. If qp,,,(aij) for 1 d v < p, are the c, x c, minors of this matrix
then the necessary and sufficient condition for (X,, . . . , X,)n c (Q, I,, . . . , lJ to
hold is that at least one of the 4pn(aij) is non-zero. Therefore the ideal
(Q,h,..., Id) will fail to be (X,, . . . , X,)-primary if and only if the
quantities uij satisfy (~?,,(a~~)= 0 for all n and all v. However, the ring k[Zij]
is Noetherian, so that the ideal of k[Zij] generated by all of the cp,,(Zij) is
generated by finitely many elements D,(Zij) for 1 f CId u. These D, clearly
meet our requirements.         n
\end{verbatim}
\end{proof}


\begin{theorem}[Printed Theorem 15.1]\label{mcrt-auto-theorem-15-1}\dcref{15.1}\source{matsumura-commutative-ring-theory:page-0131}\uses{}
\begin{verbatim}
Let cp:A -B             be a homomorphism     of Noetherian rings,
 and P a prime ideal of B; then setting p = Pn A, we have
   (i) ht P d ht p + dim B&B,;
   (ii) if q is flat, or more generally if the going-down theorem holds between
A and B, then equality holds in (i).
\end{verbatim}
\end{theorem}
\begin{proof}
\begin{verbatim}
We can replace A and B by A, and B,, and assume that (A, m) and
(B, n) are local rings, with mB c n. Rewriting (i) in the form
            dim B < dim A + dim B/mB
makes clear the geometrical content. To prove this, take a system of
parameters x1,. . . , x, of A, and choose yr,. . . , y,eB such that their images
in B/mB form a system of parameters of B/mB. Then for v, p large enough
 we have ny c mB + z y,B and rn@c xxjA, giving ny? c x y,B + C XjB.
Hence dim B 6 r + s.
   (ii) Let dim B/mB = s, and let n = P, 3 P, I ... 3 P, be a strictly decreas-
ing chain of prime ideals of B between n and mB. Obviously we have
P,nA=mforO~i~s.NowsetdimA=randletm=p,~p,~~~~~p,
be a strictly decreasing chain of prime ideals of A; by the going-down
theorem, we can construct a strictly decreasing chain of prime ideals of B
            Ps3Ps+lx?-,3Ps+r            such that P,, i n A = pi.
Thus dim B 2 r + s, and putting this together with (i) gives equality.         n
\end{verbatim}
\end{proof}


\begin{theorem}[Printed Theorem 15.2]\label{mcrt-auto-theorem-15-2}\dcref{15.2}\source{matsumura-commutative-ring-theory:page-0131}\uses{}
\begin{verbatim}
Let q:A -+ B be a homomorphism        of Noetherian rings?
and suppose that the going-up theorem holds between A and B. Then if
p and q are prime ideals of A such that p 3 q, we have
        dim B @ K(P) 2 dim B 0 K(q).
\end{verbatim}
\end{theorem}
\begin{proof}
\begin{verbatim}
Set r = dimB@Ic(q)     and s = ht(p/q). We choose a strictly
\end{verbatim}
\end{proof}


\begin{theorem}[Printed Theorem 15.3]\label{mcrt-auto-theorem-15-3}\dcref{15.3}\source{matsumura-commutative-ring-theory:page-0132}\uses{}
\begin{verbatim}
Let cp:A+B be a homomorphism    of Noetherian rings, and
      suppose that the going-down theorem holds between A and B. If p and q
      are prime ideals of A with p 1 q then
              dim B@c(p) <dim B@K(q).
\end{verbatim}
\end{theorem}
\begin{proof}
\begin{verbatim}
We may assume that ht(p/q)= 1, and it is enough to prove that,
      given a chain P, cP, c . . . c P, of prime ideals of B lying over p such that
      ht(P,/P, _ r) = 1 we can construct a chain of prime ideals Q. c Q 1 c . . . c Q,
      of B lying over ~7such that
              QicPi    (Odidr)   and ht(Qi/Qi-l)=l       (O<i<r).
      We can find Q. by going down. If r > 1 then take x~p -q and let ??, , . . . ,T,
      be the minimal prime divisors of Q,+xB.        Then ht(TJQ,)= 1, while
      ht(PJQ,) > 2, hence we can choose


      Let Qr be a minimal prime divisor of Qo+ yB contained in P,. Then
      ?ht(Qr/Q,J = 1, and Q1 # q for all i, hence c$x)#Q,.
         Therefore Q,nA#p,    and since ht(p/q)= 1 we must have QlnA=q.    By
      the same method we can successively construct Q1 , Qz, . . . , Q,. q




      2. polynomial and formal power seriesrings
\end{verbatim}
\end{proof}


\begin{theorem}[Printed Theorem 15.4]\label{mcrt-auto-theorem-15-4}\dcref{15.4}\source{matsumura-commutative-ring-theory:page-0132}\uses{}
\begin{verbatim}
Let A be a Noetherian        ring, and X,, . . . , X, indeterminates
     over A. Then
              dimA[X,,...      ,X,]=dimA[X,         ,..., X,j=dimA+n.
  + %of. It is enough to consider the case n = 1. For any p&pecA,                      the
:I?. hjt A[x]    &K(P)    = K(p)[x]    is a principal ideal ring, and therefore one-
2?:
dimensional; also A[X] is a free A-module, hence faithfully flat, so by
Theorem 1, (ii), dim A[X] = dim A + 1.
   For ,4[XJ it is not true in general that A[Xa@,rc(p)    and tc(p)[X]
coincide; however, if m is a maximal ideal of A we have
         AUxljOlc(m)=ABx40(A/m)=(A/m)[Cxn,
and this fibre ring is one-dimensional.    Also, as we saw on p. 4, every
maximal ideal ?%I of A[XJ is of the form W = (m,X), where m = sJJln~
is a maximal ideal of A. Thus for a maximal ideal 9.R of ,4%X] we have
         ht?9JI = ht(mnnA) + 1;
conversely, if m is a maximal ideal of A then ht(m, X) = ht m + 1, and
putting these together gives dim A [Xl = dim A + 1. m
Remark 2. It is not necessarily true that a maximal     ideal of A[X] lies
over a maximal ideal of A. For example, if A is a DVR and t a
uniformising element then A[t-?1 = K is the field of fractions of A, so
that A[X]/(tX - 1) 2: K, and (tX - 1) is a maximal ideal of A[X]; however,
(tX- l)nA=(O).
Remark 2. It is quite common         for fibre rings of A -+A[X,,.    . .,X,1
to have dimension strictly greater than n. For example, let k be a field and
set A = k[Y,Z]. It is well-known that the field of fractions of k[Xj has
infinite transcendence degree over k(X) (see[ZS], vol. II, p. 220). Let u(X),
u(X)~k[Xj      be two elements algebraically independent over k(X), and
define a k-homomorphism      (continuous for the X-adic topology)


bycp(X!=X,cp(Y)=u(X),cp(Z)=u(X).IfwesetKercp=PthenPnA=(O),
and A[XJ/P z k[XJ is one-dimensional.         Now every maximal ideal of
A[Xj has height 3, and, as we will see later, A[Xj is catenary, so that
ht P=2. Thus we see that the generic libre of A-+A[Xj            is two-
dimensional.

3. The dimension inequality
We say that a ring A is universally catenary if A is Noetherian and every
finitely generated A-algebra is catenary. Since any A-algebra generated
by n elements is a quotient of A[X,, . , . , X.1, and since a quotient of a
catenary ring is again catenary, a necessary and sufficient condition for a
Noetherian ring A to be universally catenary is that A[X,, . . . ,X,] is
catenary for every n 2 0. (In fact it is known that it is sufficient for A[XI]
to be catenary, compare Theorem 31.7.)
\end{verbatim}
\end{theorem}
\begin{proof}
\notready
\end{proof}


\begin{theorem}[Printed Theorem 15.5]\label{mcrt-auto-theorem-15-5}\dcref{15.5}\source{matsumura-commutative-ring-theory:page-0133}\uses{}
\begin{verbatim}
I. S. Cohen [3]). Let A be a Noetherian   integral domain, and
(*) htP + tr.deg,(,, K(P) d htp + tr.deg, B,
   where tr.deg,B is the transcendence degree of the field of quotients of B
   over that of A.
\end{verbatim}
\end{theorem}
\begin{proof}
\begin{verbatim}
We may assume that B is finitely generated over A. For if the right
   hand side is finite and m and t are non-negative integers such that m < htP
   and t < tr.deg,(,+(P),       then there is a prime ideal chain
   p=po3       P, X3??X P, in B. Take ai~Pi-Pi+l,          06i<m,      and let
   C1,...,~,~B    be such that their images modulo P are algebraically
   independent over A/p. Set C = A[{a,}, (cjj]. If the theorem holds for C,
   then we have m + t < htp + tr.deg, C < htp + tr.deg, B. Letting m and t vary
   we seethe validity of (*).
      We may furthermore assume,by induction, that B is generated over A by
   a single element: B = A[x]. We can replace A by A, and B by B, = AJx],
   and hence assumethat A is local and p its maximal ideal. Set k = A/p and
   w&e B = A[X]/Q. If Q = (0) then B = A[X] and by Theorem 1 we have
   htP = htp + ht(P/pB), and since B/pB = k[X] we have either P = pB or
   ht(P/pB) = 1. In both casesthe equality holds in (*).
      If Q # (0) then tr.deg, B = 0. Since A is a subring of B we have
   &A=(O),        so that writing K for the field of fractions of A we have
   htQ = htQK[X] = I. Let P* be the inverse image of P in A[X]. Then
   P=P*/Q,K(P)=K(P*),           and     htP<htP*-htQ=htP*-l=htp+l-
   tr.deg,&P*) - 1= htp -tr.deg,(,,lc(P).     n
\end{verbatim}
\end{proof}


\begin{theorem}[Printed Theorem 15.7]\label{mcrt-auto-theorem-15-7}\dcref{15.7}\source{matsumura-commutative-ring-theory:page-0137}\uses{}
\begin{verbatim}
Let A be a Noetherian ring and I a proper ideal; then
setting R = R(A, I) and G = gr,(A) we have
         dimR=dimA+         1, dimG<dimA.
If in addition A is local, then
         dim G = dim A.
\end{verbatim}
\end{theorem}
\begin{proof}
\notready
\end{proof}


\begin{theorem}[Printed Theorem 18.5]\label{mcrt-auto-theorem-18-5}\dcref{18.5}\source{matsumura-commutative-ring-theory:page-0161}\uses{}
\begin{verbatim}
We consider modules over a Noetherian ring A.
  (i) A direct sum of any number of injective modules is injective.
  (ii) Every injective module is a direct sum of indecomposable injective
modules.
  (iii) The direct sum decomposition in (ii) is unique, in the sense that if
         M = @ Mi (with indecomposable          Mi)
then for any PESpec A, the sum M(P) of all the Mi isomorphic to E(A/P)
depends only on M and P, and not on the decomposition M = @ Mi.
Moreover, the number of Mi isomorphic to E(A/P) is equal to
          dim,cp,Hom,,(rc(P), MJ, (where K(P) = AdPA,),
so that this also is independent of the decomposition.
\end{verbatim}
\end{theorem}
\begin{proof}
\begin{verbatim}
(i) Let M1 for 1~11 be injective modules. It is enough to prove that
for an ideal I of A, any linear map q : I -+ @MI can be extended to a linear
e I is finitely generated, q(Z) is contained
                                             roftheM,.Ifcp(l)cM,@...@M,and
                                            t of p(a) in Mi, then Cpi:I -+ Mi extends
                                            -@J%             by t41)=11/,(1)+.~.+$,(1)

b. (ii) Say that a family .F = {E,} of indecomposable injective submodules
& of M is fvee if the sum in M of the E, is direct, that is if, for any finite
$ number EA,,. . , E,,, of them,
$;
./I           E,, n(E,,    + ... + E,?) = 0.
&tit        %JIbe the set of all free families 8, ordered by inclusion. Then by
2: Zorn?s lemma !JJI has a maximal element, say FO. Write N = CEEFOE; then
g by (i), N is injective, hence a direct summand of M, and M = N 0 N?. If
g, NVf 0 then since it is a direct summand of M it must be injective, and
wa.for PEAss(N?), the proof of Theorem 4, (ii), shows that N? contains a
     direct summand E? isomorphic to E(A/P). Thus F0 u {E?} is a free family,
     contradicting the maximality of FO. Hence N? = 0 and M = N.
        (iii) If we can show that M(P) has the property that every submodule
     & of M isomorphic to E(A/P) is contained in M(P), then M(P) is the
     submodule of M generated by all such E, and therefore is determined by
   ?M and P only. To prove this, take any <EE; we can write t = 4, + ... +
     4, with ~iEM(Pi), where PI,. . . , P, are distinct prime ideals and P = P,.
     Setting~,-5=vl,and5i=rifor2~idrwehaver],+...+vl,=O,with
     )1IcM(P,) + E and qieM(Pi) for i >, 2. We need only prove that in this case
     each Y/i= 0. Suppose that P, is minimal among P,, . . , P,; then for any m
     we have (PI . . .Prml)??@Pl,        so that taking ag(P,...P,-,)?-P,      and m
     large enough, we get ar], = ... = ar],- 1 = 0. Then also q, = 0, but
     multiplication     by a is an automorphism         of M(P,), so that ql = 0. By
L ?mdnction on r we get vi = 0 for all i.
f       We now prove that if M(P) = M, @... @ M, with Mi N E(A/P) then
$                s = dim,(,,) Hom,,(rc(P), Mp).
I;
$ (we are writing this as ifs were finite, but, as one can see from the proof
k; below, the same works for any cardinal number.) By Theorem 4, (vi), both
     aides of M(P) = M, @...@ M, are &modules,                       and Mi = E(rc(P)).
     Moreover, by Theorem 4, (v), E(LI/Q)~ = 0 if Q $ P, so that

              MP = @ M(Q)p = @ M(Q).
                     QcP             QcP

      Hence we can replace A by A p, and assume that A is a local ring with P
      its maximal ideal; set k = K(P). If Q # P then any XEP - Q gives an
      automorphism of M(Q), but x.k = 0, so that Hom,(k, M(Q)) = 0. Hence
      Hom,(k, M) = Hom,(k, M(P)), so that there is no loss of generality in
assuming that M = M(P). For any A-module N, we can identify,
Hom,(k,N)     with the submodule {<ENJP~ =O}, but since E(k) is an
essential extension of k we must have dim, Hom,(k, E(k)) = 1, so that if
M = M, @...@ M, with Mi N E(k) then s = dim, Hom,(k, M). H
\end{verbatim}
\end{proof}


\begin{theorem}[Printed Theorem 18.6]\label{mcrt-auto-theorem-18-6}\dcref{18.6}\source{matsumura-commutative-ring-theory:page-0163}\uses{}
\begin{verbatim}
Let (A, m, k) be a Noetherian local ring, and E = E,(k) the
injective hull of k. For each A-module M set M? = Hom,(M, E).
    (i) If M is an A-module and 0 # XEM, then there exists cp~M? such that
q(x) # 0. In other words the canonical map 8: M --+ M? defined by Q)(,)
 = q(x) for XEM and cp~M? is injective.
    (ii) If M is an A-module of finite length, then l(M) = 1(M?) and the
canonical map M -M?            is an isomorphism.
    (iii) Let A^ be the completion of A; then E is also an A-module, and
is an injective hull of k as A-module.
    (iv) Hom,(E, E) = Hom,(E, E) = A^. In other words, each endomer-
phism of the A-module E is multiplication         by a unique element of 2.
    (v) E is Artinian as an A-module and also as an A-module. Assume
now that A is complete, and write N(resp. ~2) for the category of
Noetherian (respectively Artinian) A-modules. Then if ME&? we have
M?E~       and MN M?; if ME& we have M?EJ+?? and MN M?.
\end{verbatim}
\end{theorem}
\begin{proof}
\begin{verbatim}
(i) Let f:Ax -E           be the composite of the canonical maps
Ax N A/ann(x), A/ann(x) - A/m = k and k - E. Then f(x) # 0. Since Eis
injective we can extend f to cp:M -E.
   (ii) If l(M) = n < cc then M has a submodule M, of length n - 1, and
O-+M,-M-k-,Oisexact,sothatO-+k?--+M?--rM~+Oisexact.
However
            k? = Horn (k, E) = Horn (k, k) ru k,
so that by induction on II we get l(M) = n = [(M?). The canonical map
M -M?        is injective by (i), and I(M) = QM?) = l(M?), hence it must be an
isomorphism.
   (iii) Each element of E is annihilated by some power of m, so that the
canonical map E -E         @AA is surjective. However, since A^ is faithfully
flat over A it is also injective, so that E N E aAA, and we can view E as
an A-module. Let F be the injective hull of E as an A-module. Then F
is also the A-injective hull of k, so that every element of F is annihilated
by some power of rn2. As an A-module F splits into a direct sum Of E
and an A-module C. If xgC, and if mrAx = 0, then for each a*Ea we
can find aeA such that a* -a mod m?A and hence a*x =ax EC.
Therefore C is an A-module. But F is indecomposable as an A-module?
Hence C = 0 and E = F.
    (iv) For v > 0 set E, = (xEE)m?x = O}. Then we have (A/m?yS
Horn, (A/m?, E) rr E,, and Horn, (E,, E,) = Hom,(E,, E) = E:. = (Ah?)? =
A/m. Now El = E, = . . . and E = UYEY by Theorem 4, (v), hence E =
limE,. Therefore Hom,(E, E) = Horn,@     E,, E) = @ HomA(E,, E) =
GA/m?     = A^.
7?) Ifan A-module M is Artinian and XEM, then Ax in A/ann(x) is also
 Artinian and consequently my c arm(x) for some v. Therefore M can be
viewed as an A-module, and its A-submodules are precisely its A-
submodules. It is also clear that if an a-module M is Artinian then we have
the same conclusion. Therefore to prove (v) we may assume that A is
complete.
   JfM is a submodule of E set ML = (UEA jaM = O}. If Z is an ideal of A set
Zl= {x~EjZx =O}. Then clearly ML? 3 M. If XEE - M there exist
&E/M?)        satisfying cp(xmod M) # 0 by (i), and if we identify E? =
Hom,(E, E) with A then (E/M)? is identified with Ml. Thus cp(x mod M) =
ax for some aeM?,         and .x$M?.      Therefore Ml? = M. Similarly, if
acgA -I     then there exists cpe(A/I)? such that cp(umodZ) #O, and
(A/Z)? is identified with the submodule I1 of E = A?. Thus, setting x =
cp(1modZ) we have xel? and ux = cp(umod I) # 0. This proves a$Z?, so
that I = Z1l. Thus M w ML is an order-reversing bijection from the set of
submodules of E onto the set of ideals of A. Since A is Noetherian, it follows
that E is Artinian. By Theorem 3.1 finite direct sums E? of E are also
Artinian for all n > 0.
   If MEN then there is a surjection A? + M for some n, and so there
is an injection M? -(A?)?      = E?. Hence M? is Artinian. On the other
 hand, if ME& there is an injection M -En           for some n. This can be
seen as follows: consider all linear maps M + E?, where n is not fixed,
and take one cp:M -En         whose kernel is minimal among the kernels of
 those maps. Then, using (i) we can easily see that Ker(cp) = 0. Now, from
0 + M - E? we have (E?)? = A? - M? + 0 exact, hence M?eJlr. Now the
 assertion M z M? for M EJY or d can easily be checked using (iv) if M = A
 or E, and the general case follows from this and from (i). n
Zhnma 6. Let A be a Noetherian ring, S c A a multiplicative      set, M an A-
module and N c M a submodule. Assume that M is an essential extension
of N; then MS is an essential extension of N,.
proof. For (EM we write ts = l/l EMU; then any element of MS can be
written u*cs (with u a unit of A, and REM), so that it is enough to show that
for any non-zero & we have N,n A,.<, # 0. Suppose that ann (tot) is a
maximal element of the set of ideals {ann(    tES}; then ifwe set V= to 4, we
have ts = to 1qs, and hence v ~0. Now let b = {~EAI~?EN};                  by
assumption,
         bv=AqnN#O.
Suppose that b = (b ?,...,b,);   if b,q,=...   = b,qs =0 then there is a tcS
150      Regular sequences

such that tb,r = 0 for all i. Then tbv] = 0, but by choice of q we have
ann (q) = ann (tq), so that bq = 0, which is a contradiction. Thus biqs # 0
for some i, and
           biYls~As~V,nNs = As~s~Ns,
 as required.     n
    By Lemmas 5 and 6, if M is an injective hull of N then the As-module
M, is an injective hull of N,. Hence if O+ M - I0 -+I?          - ... is
a minimal injective resolution of an A-module M, then 0 + M, --+ I,0 -+
G -...      is a minimal injective resolution of the As-module M,. The
I? are determined uniquely up to isomorphism by M. We can therefore
define pi(P, M) to be the number of summands isomorphic to E(A/P)
appearing in a decomposition of I? as a direct sum of indecomposable
modules. We can write symbolically
         I?=          @       pi(P, M)E(A/P).
                  Pdpec   A

From what we have just proved, for a multiplicative          set S c A,
      Ili(P, M) = ~i(PA,, MJ   if P n S = a.
\end{verbatim}
\end{proof}


\begin{theorem}[Printed Theorem 18.7]\label{mcrt-auto-theorem-18-7}\dcref{18.7}\source{matsumura-commutative-ring-theory:page-0165}\uses{}
\begin{verbatim}
Let A be a Noetherian                 ring, M an A-module, and
PESpec A. Then
           Pi(P, W = dimKcp, Ext&(lc(P), MP) = dim,&Ext>(A/P,          M))p.
In particular, if M is a finite A-module then ~i(P, M) < co.
\end{verbatim}
\end{theorem}
\begin{proof}
\begin{verbatim}
Replacing A and M by A, and M, we can assume that (A, P, k) is
a local ring. Let O+M -+ZO-f+ll--%...                     be a minimal    injective
resolution of M, so that Exti(k, M) is obtained as the homology of the
complex
          . ..-+Hom.(k,I?-?)-Hom,(k,Z?)-Hom,(k,I?+?)
               +...
We can identify Hom,(k,I?)      with the submodule T?= {x~Z?/Px =0} c
I?. By construction of the minimal injective resolution, I? is an essential
extension of d(l?- ?), so that for XGT? the submodule AX 1: k intersects
d(l?-?), and xEd(Z?-?). Therefore, T? c d(I?-?), and dT?-? = dT?= 0, so
that Ext$(k, A) = T?. Also,
         dim, T? = dim, Hom,(k, I?),
and by Theorem 4, (iii), this is equal to pi(P, M). w
\end{verbatim}
\end{proof}


\begin{theorem}[Printed Theorem 18.8]\label{mcrt-auto-theorem-18-8}\dcref{18.8}\source{matsumura-commutative-ring-theory:page-0165}\uses{}
\begin{verbatim}
A necessary and sufficient condition for a ring A to be
Gorenstein is that a minimal     injective resolution O+ A -+ 1? -+
I? --+ .. . of A satisfies
         I? = @           E(A/P),
                 htP=i
or, in other words, /li(P, A) = Si,htP (the Kronecker 6) for every PESpec A.
\end{verbatim}
\end{theorem}
\begin{proof}
\begin{verbatim}
By Theorem 7 and condition (2) of Theorem 1 we have
                A, is Gorenstein e pi(P, A) = 8i,htP.
\end{verbatim}
\end{proof}


\begin{theorem}[Printed Theorem 18.9]\label{mcrt-auto-theorem-18-9}\dcref{18.9}\source{matsumura-commutative-ring-theory:page-0166}\uses{}
\begin{verbatim}
Let (A,m)       be a Noetherian     local ring, and M a finite
      A-module. Then
              inj dim M < r3 +inj dim M = depth A.
\end{verbatim}
\end{theorem}
\begin{proof}
\begin{verbatim}
Suppose that inj dim M = r < cc. If P is a prime ideal distinct from
      m, choose xEm - P. Then

ii            0 + A/P 2      A/P,
      together with the right-exactness      of Ext>( - , M) gives an exact sequence
              Ext?,(A/P,M)AExt:,(A/P,M)+O,
)..
(i.   so that by NAK Ext>(A/P, M) = 0. Putting this together with Lemma 1,
:.-
 :    we get Ext>(k,M)#O.             Set t=depthA,   and let xr,,..,x,Em        be a
$     maximal A-sequence; then setting A/(x,, . ,x,) = N we have mEAss (N).
$     Hence there exists an exact sequence O-+ k - N, and we must have
!-    Ext>(N, M) # 0. The Koszul complex K(x r, . . . ,x,) is a projective resolution
?1
2:
~:-   of N = A/(x,, . . . , x,), so computing Ext by means of it we see that
              Ext;(N, M) N M/(x,, . . ,x,)M,
;.,,.
$, and by NAK this is non-zero. Thus proj dim N = t, and from Ext\(N, M) #
?? 0 we get t d r, whereas from Ext>(N, M) # 0 we get t 2 r. Hence t = r. H
8 Remark (the Bass conjecture and the intersection theorem). Let (A, m, k) be
i,
; a Noetherian local ring of dimension d. H. Bass [l] conjectured the
1. following:
[     (B)     if there exists a finite A-module         M (#Q     of finite   injective
8:?           dimension, then A is a CM ring.
F According to Theorem 9 this is equivalent to asking that inj dim M = d.
i The converse of the Bass conjecture is true. Indeed, if A is CM, taking a
d maximal A-sequence x1,. . , x,, and setting B = A/(x,, . . . , xd) and E = E(k)
$?
!+ we have l,(B) < co. By Theorem 6, M = Horn,@, E) is also of finite length,
in hence is finitely generated. We prove inj dim,M < d; the Koszul complex
? O--bA-+Ad -+~~~-+Ad--tA+BL+O with respect to x1,. . . , x, provides an A-
! free resolution of B. Now applying the exact functor Hom,(-,E) to this
    gives    the exact sequence O--+M-+E+Ed-+..~+Ed--+E+O.            This proves
    injdim,M d d.
i.
L        (B) is a special case of the following theorem.
1 (C)         If I?: 0-10 -+?..-+Zd+O is a complex of injective modules such
c             that #(I?) is finitely generated for all i and I? is not exact, then
              Id # 0.
c4
Using the theory of dualizing complexes (see [Rob])            one can prove that (C)
is equivalent to the following
Intersection Theorem. If F.: O+F,-+~~~-+F,-+O         is a complex of finitely
generated free modules such that H,(F) has finite length for all i and F. is
not exact, then F, # 0.
   (B) was proved by Peskine and Szpiro [l] in some important cases,and
by Hochster [H] in the equal characteristic case (i.e. when A contains a
field) as a corollary of his existence theorem for the ?big CM module?, i.e. a
(not necessarily finite) A-module with depth = dim A, see [H] p.10 and
p.70. The intersection theorem was conjectured by Peskine-Szpiro [3] and
by P. Roberts independently. They pointed out that it was also a
consequence of Hochster?s theorem. Finally, P. Roberts [3] settled the
remaining unequal characteristic case of the intersection theorem by using
the advanced technique of algebraic geometry developed by W. Fulton
([Full). Therefore (B), which was known as Bass?sconjecture for 24 years,
is now a theorem. Some other conjectures listed in [H] are still open.
\end{verbatim}
\end{proof}


\begin{theorem}[Printed Theorem 20.1]\label{mcrt-auto-theorem-20-1}\dcref{20.1}\source{matsumura-commutative-ring-theory:page-0176}\uses{}
\begin{verbatim}
A Noetherian integral domain A is a UFD if and only if
every height 1 prime ideal is principal.
\end{verbatim}
\end{theorem}
\begin{proof}
\begin{verbatim}
Suppose that A is a UFD and that P is a height 1
prime ideal. Take any non-zero aeP, and express a as a product of prime
elements, a = nq. Then at least one of the ni belongs to P; if 7liEP then
(xi) c P, but (rq) is a non-zero prime ideal and ht P = 1, hence P = (xi).
Proof of ?iLf?. Since A is Noetherian, every element aEA which is neither
0 nor a unit can be written as a product of finitely many irreducibles.
Hence it will be enough to prove that an irreducible element a is a prime
element. Let P be a minimal prime divisor of (a); then by the principal
ideal theorem (Theorem 13.9, ht P= 1, so that by assumption we can
write P =(b). Thus a = bc, and since a is irreducible, c is a unit, so that
(a) = (b) = P, and a is a prime element.    H
\end{verbatim}
\end{proof}


\begin{theorem}[Printed Theorem 20.2]\label{mcrt-auto-theorem-20-2}\dcref{20.2}\source{matsumura-commutative-ring-theory:page-0177}\uses{}
\begin{verbatim}
Let A be a Noetherian integral domain, r a set of prime
elements of A, and let S be the multiplicative set generated by f. If A, is
a UFD then so is A.
\end{verbatim}
\end{theorem}
\begin{proof}
\begin{verbatim}
Let P be a height 1 prime ideal of A. If PnS # 0 then P contains
an element nil-, and since nA is a non-zero prime ideal we have P = nA.
If PnS = 0 then PA, is a height 1 prime ideal of A,, so that PA, = aA,
for some aEP. Among all such a choose one such that aA is maximal;
then a is not divisible by any nil-. Now if XEP we have S.X=ay for
some SES and YEA. Let s = or... x I with n,gT; then a$z,A, so that
y~qA, and an induction on r shows that y~sA, so that xEaA. Hence
P=aA.     n

Lemma 1. Let A be an integral domain, and a an ideal of A such that
 a@A?2:Anf1;          then a is principal.
Proof. Fix the basis e,, . , e, of A??,          and viewing a @ A? c A @A?, fix
fO,. . . , f, such that ,fO is a basis of A and fi,. . , f,, a basis of A?. Then
the isomorphism          q:A?+ ? -a@       A? can be given in the form q(ei) =
C;=o~iJj        Write di for the (i,O)th cofactor of the matrix (aij), and d
for the determinant, so that, since cp is injective, rl # 0, and 1 Uiodi = d,
C~ijdi =O if j # 0. H ence if we set eb = C;d,e, we have cp(eb) = dfo.
Moreover,       since the image of cp includes Jr,.. .,fn, there exist e;,...?
e;EA ?+? such that cp(e>) =fj. Now define a matrix (c,) by e; = x;=ocjkek
for j = 0,. . . , n (so c,,~ = C&J.Then we have
                         /d    0   . ..      O\

        (Cjk)(aij)   =    (f   ?   ..        O

                         i 0   0        .i   1?
SO that by comparing the determinants of both sides we get det tcjk)       =   I?
Therefore eb,. . , e; is another basis of A?+l, and afo = cp(Aeb) = dAf,,
  so that a = dA. n
     Let K be the field of fractions of the integral domain A; for a finite
  A-module M, the dimension of M aA K as a vector space over K is called
  the rank of M. A torsion-free finite A-module of rank 1 is isomorphic to
  an ideal of A. Lemma 1 can be formulated as saying that for an integral
  domain A, a stably free rank 1 module is free (see Ex. 19.2). The elementary
  proof given above is taken from a lecture by M. Narita in 1971.
\end{verbatim}
\end{proof}


\begin{theorem}[Printed Theorem 20.3]\label{mcrt-auto-theorem-20-3}\dcref{20.3}\source{matsumura-commutative-ring-theory:page-0178}\uses{}
\begin{verbatim}
Auslander and Buchsbaum [3]). A regular local ring is a
  UFD.
\end{verbatim}
\end{theorem}
\begin{proof}
\begin{verbatim}
Let (A,m) be a regular local ring: the proof works by induction
  on dim A. If dim A = 0 then A is a field and therefore (trivially) a UFD.
  If dim A = 1 then A is a DVR, and therefore a UFD. We suppose that
  dim A > 1 and choose xEm - m2; then since xA is a prime ideal, applying
  Theorem 2 to I- = {x}, we need only show that A, is a UFD (where
  A, = A[x- ?1 is as on p. 22). Let P be a height 1 prime ideal of A, and
  set p = Pn A; we have P = pA,. Since A is a regular local ring, the
  A-module p has an FFR, so that the AX-module P has an FFR. For any
  prime ideal Q of A,, the ring (A,), = A,,, is a regular local ring of
  dimension less than that of A, so by induction is a UFD. Thus P, is free
  as an (A&-module,     so that by Theorem 7.12, the AX-module P is projective;
, hence by Ex. 19.2, P is stably free, and therefore by the previous lemma,
  P is a principal ideal of A,.    n
     The above proof is due to Kaplansky. Instead of our Lemma 1, he used
  the following more general proposition, which he had previously proved:
  if A is an integral domain, and Zi,Ji are ideals of A for 1 d i d r such that
  a=, Ii N @= 1 Ji, then I,. . I, N J,. . . J,. This is an interesting property
  of ideals, and we have given a proof in Appendix C.
\end{verbatim}
\end{proof}


\begin{theorem}[Printed Theorem 20.4]\label{mcrt-auto-theorem-20-4}\dcref{20.4}\source{matsumura-commutative-ring-theory:page-0178}\uses{}
\begin{verbatim}
Let A be a Noetherian integral domain. Then if any finite
  A-module has an FFR, A is a UFD.
\end{verbatim}
\end{theorem}
\begin{proof}
\begin{verbatim}
By Ex. 19.4, A is a regular ring. Let P be a height 1 prime ideal
  of A. Then A,,, is a regular local ring for any mESpec A, so by the previous
  theorem, the ideal P,,, is principal, and is therefore a free A,,,-module. Hence
  by Theorem 7.12, P is projective. Therefore by Ex. 19.2, P is stably free,
  and so by Lemma 1 is principal.         n
     Let A be an integral domain; for any two non-zero elements a, bEA,
  the notion of greatest common divisor (g.c.d.) and least common multiple
  (1.c.m.) are defined as in the ring of integers. That is, d is a g.c.d. of a and
  b if d divides both a and b, and any element x dividing both a and b
164      Regular rings

divides d; and e is an 1.c.m. of a and b if e is divisible by both a and b,
and any y divisible by a and b is divisible by e; this condition is equivalent
to (e) = (a) n(b).
Lemma 2. If an 1.c.m. of a and b exists then so does a g.c.d.
Proof. If (a)n(b) = (e) then there exists d such that ab = ed. From ee(a)
we get bQd) and similarly a@d), so that (a, b) c (d). Now if x is a common
divisor of a and b then a = xt and b = xs, so that xst is a common multiple
of a, b, and is hence divisible by e. Then from ed = ab = x.xst we get that
d is divisible by x. Therefore, d is a g.c.d. of a and b. n

Remark 1. If A is a Noetherian    integral domain which is not a UFD then
A has an irreducible element a which is not prime. If xyE(a) but x$(a),
 y+(a) then the only common divisors of a and x are units, so that 1 is a
g.c.d. of a and x. However, xye(a)n(x),          but xy$(ax), so that (a)n
(x) # (ax), and there does not exist any 1.c.m. of a and x. Thus the converse
of Lemma2 does not hold in general.
Remark 2. If A is a UFD then an intersection        of an arbitrary collection
of principal ideals is again principal (we include (0)). Indeed, if n ill aiA # 0,
then factorise each a, as a product of primes:



with ui units, and p, prime elements such that pJ # p,A for x # 8.
Then r) a,A = dA, where d = ~~~~~~~~~~~~~~~~~~
                                             (We could even allow the a, to
be elements of the field of fractions of A.)
\end{verbatim}
\end{proof}


\begin{theorem}[Printed Theorem 20.5]\label{mcrt-auto-theorem-20-5}\dcref{20.5}\source{matsumura-commutative-ring-theory:page-0179}\uses{}
\begin{verbatim}
An integral domain A is a UFD if and only if the ascending
chain condition holds for principal ideals, and any two elements of A have
an 1.c.m.
ProoJ: The ?only if? is already known, and we prove the ?if?. From the first
condition it follows that every element which is neither 0 nor a unit can
be written as a product of a finite number of irreducible elements, so that
we need only prove that an irreducible element is prime. Let a be an
irreducible element, and let xyE(a) and x$(a). By assumption we can
write (a)n(x) = (2); now 1 is a g.c.d. of a and x, so that one sees from the
\end{verbatim}
\end{theorem}
\begin{proof}
\begin{verbatim}
that YE(a). Therefore (a) is prime.    n
\end{verbatim}
\end{proof}


\begin{theorem}[Printed Theorem 20.6]\label{mcrt-auto-theorem-20-6}\dcref{20.6}\source{matsumura-commutative-ring-theory:page-0179}\uses{}
\begin{verbatim}
Let A be a regular ring and u, UEA. Then uAnuA is a
projective ideal.
\end{verbatim}
\end{theorem}
\begin{proof}
\begin{verbatim}
A,,, is a UFD for every maximal ideal m, so that (uAnuA)A,, =
uA,nvA,     is a principal ideal, and hence a free module. n
?Remark.     There are examples where A is a UFD but A[X] is not.
     i It is easy to see that a UFD is a Krull ring. For any Krull ring A we
     : can define the divisor class group of A, which should be thought of as a
     ?Ineasure of the extent to which A fails to be a UFD. We can give the
     ,fdeIinition in simple terms as follows: let 9 be the set of height 1 prime
       ideals of the Krull ring A, and D(A) the free Abelian group on 9. That
     ?is, D(A) consists of formal sums 1 ptBnp?p (with n,EZ and all but finitely
     ?many n, = 0), with addition defined by
                (Cn;p)+(Cnb.P)=C(n,+n6)p.
          K be the field of fractions of A, and K* the multiplicative   group of
               o elements of K, and for EK* set div(a) = CPE,~uP(u).p, where up
                ormalised additive valuation of K corresponding to p. Then
               = div(a) + div (b), so that div is a homomorphism from K* to D(A).
          write F(A) for the image of K*; this is a subgroup of D(A), so that
          can define C(A) = D(A)/F(A) to be the divisor class group of A.
             US~Y, if A is a UFD then each PEP is principal, and if p = aA then
             element of D(A) we have p = div(a), so that C(A) = 0. Conversely,
            ) = 0 then each PEP is a principal ideal, and putting this together
           he corollary of Theorem 12.3, one sees easily that A is a UFD. Hence
                A is a UFD~C(A)       = 0.
Now let A be any ring, and M a finite projective A-module. For each
PtsSpecA, the localisation M, is a free module over A,, and we write
n(P) for its rank. Then n is a function on Spec A, and is constant on every
connected component (since n(P) = n(Q) if P 2 Q). This function n is called
the rank of M. If the rank is a constant r over the whole of SpecA then
we say that M is a projective module of rank r. We write Pit(A) for the
set of isomorphism    classes of finite projective A-modules of rank 1; cl(M)
denotes the isomorphism class of M. If M and N are finite projective rank
1 module then so is MOAN; this is clear on taking localisations. Thus
we can define a sum in Pit(A) by setting
       cl(M)+cl(N)=cl(M@N).
We set M? = Hom,(M,      A), and define cp:M 0 M* -A         by
         44x   mi   0   fi)   =   C   h(F);

then cp is an isomorphism (taking localisations and using the corollary to
Theorem 7.11 reduces to the case M = A, which is clear). Hence cl(M*) =
 -cl(M),    and Pit(A) becomes an Abelian group, called the Picard group
of A. If A is local then Pit(A) = 0.
    If A is an integral domain with field of fractions K, then MC,, = M 0 K,
so that the rank we have just defined coincides with the earlier definition
(after Lemma 1). If M is a finite projective rank 1 module, then since M
is torsion-free    we have M c MC,, N K, so that M is isomorphic as an
A-module to a fractional ideal; for fractional ideals, by Theorem 11.3,
projective and invertible are equivalent conditions, so that for an integral
domain A, we can consider Pit(A) as a quotient of the group of invertible
fractional ideals under multiplication.   A fractional ideal I is isomorphic
to A as an A-module precisely when I is principal, so that




   Suppose in addition that A is a Krull ring. Then we can view Pit(A)
as a subgroup of C(A). To prove this, for p~;/p and I a fractional ideal.
set
         v,(Z) = min {v,(x)Ix~l};
this is zero for all but finitely many PEP (check this!), so that we can set
         div (I) = 1 v,,(l).p~D(A).
                   PEP
For a principal ideal I = zA we have div(I) = div(cx). One sees easily that
div(Z1?) = div(1) + div(I?), and that div(A) = 0, so that if I is invertible,
div (I) = - div (I - ?).
For invertible I we have (I-?))?          = I: indeed, I c(Z-?)-I       from the
  definition, and I = Z.A 2 Z(Z-?(I-?)-?)          1 (I-?)-?. If Z is invertible and
  ~v(Z)=Othendiv(Z-?)=O,sothatZcA,Z-?cA;henceAc(Z-?)-1=Z,
  and Z = A. It follows that if I, I? are invertible, div(Z) = div(Z?) implies Z = I?.
  Thus we can view the group of invertible fractional ideals as a subgroup of
  D(A), and Pit(A) as a subgroup of C(A).
     If A is a regular ring then as we have seen, p~9? is a locally free module,
  and so is invertible. Clearly from the definition, div(p) = p. Hence, in the
~ case of a regular ring, D(A) is identified with the group of invertible
1,fractional ideals, and C(A) coincides,with Pit(A).
     The notions of D(A) and Pit(A) originally arise in algebraic geometry.
,: Let V be an algebraic variety, supposed to be irreducible and normal. We
1,write 9 for the set of irreducible codimension 1 subvarieties of V, and
Sdefine the group of divisors D(V) of I/ to be the free Abelian group on 9?;
; a divisor (or Weil divisor) is an element of D(V). Corresponding to a
       onal function f on I/ and an element WEP, let z+(f) denote the
       er of zero off along W, or minus the order of the pole if f has a pole
       ng W. Write div (f) = ~,,9uw(f). Wfor the divisor of f on I? (or just
         For WEP, the local ring CO, of W on I/ is a DVR of the function
         of V, and ow is the corresponding valuation. We say that two divisors
            D(V) are linearly equivalent if their difference M - N is the divisor
              ction, and write M - N. The quotient group of D(V) by - , that
          quotient by the subgroup of divisors of functions, is the divisor class
            of I/ (up to linear equivalence), and we write C(V) for this. (In
         ion to linear equivalence one also considers other equivalence
      ations with certain geometric significance (algebraic equivalence,
      merical equivalence,. . . ), and d ivisor class groups, quotients of D(V) by
           rresponding subgroups.)
           ivisor M on I/ is said to be a Cartier divisor if it is the divisor of
         ction in a neighbourhood of every point of I/. From a Cartier divisor
          onstructs a line bundle over V, and two Cartier divisors give rise to
      morphic line bundles if and only if they are linearly equivalent. Cartier
         ors form a subgroup of D(V), and their class group up to linear
          alence is written Pit(V); this can also be considered as the group of
      morphism classes of line bundles over V (with group law defined by
       sor product). If T/ is smooth then (by Theorem 3) there is no distinction
     tween Cartier and Weil divisors, and C(V) = Pit (V).
         e reader familiar with algebraic geometry will know that the divisor
      ss group and Picard group of a Krull ring are an exact translation of
        corresponding notions in algebraic geometry. If I/ is an affine variety,
     th coordinate ring k[ I?] = A then C( I?) = C(A) and Pit (I?) = Pit (A). In
this case, to say that A is a UFD expresses the fact that every codimension
  1 subvariety of I/ can be defined as the intersection of I/ with a hypersurface.
 If I/ c P? is a projective algebraic variety, defined by a prime ideal
 Zck[X,,...,      X,], and we set A = k[X]/I = k[to,. . . , t,,] (with ti the
 class of Xi) then A is the so-called homogeneous coordinate ring of I/. If
 A is integrally closed we say that V is projectively normal (also
 arithmetically normal). This condition is stronger than saying that V is
 normal (the local ring of any point of I/ is normal). If A is a UFD then
 every codimension 1 subvariety of Vcan be given as the intersection in
 P? of V with a hypersurface. Let m = (Co,. . . ,&,) be the homogeneous
 maximal ideal of A, and write R = A,,, for the localisation. The above
 statement holds if we just assume that R is a UFD; see Ex. 20.6. All the
 information about V is contained in the local ring R.
     Thus C(A), Pit(A) and the UFD condition are notions with important
 geometrical meaning, and methods of algebraic geometry can also be used
 in their study. For example, in this way Grothendieck [G5] was able to
 prove the following theorem conjectured by Samuel: let R be a regular
 local ring, P a prime ideal generated by an R-sequence, and set A = R/P;
 if A, is a UFD for every p~Spec A with ht p < 3 then A is a UFD.
     We do not have the space to discuss C(A) and Pit(A) in detail, and we
just mention the following two theorems as examples:
     (1) If A is a Krull ring then C(A) 2: C(A[X]).
     This generalises the well-known theorem (see Ex. 20.2) that if A is a UFD
then so is A[X].
     (2) If A is a regular ring then C(A) N C(A[XJ).
     This generalises Theorem 8.
     Finally we give an example. Let k be a field of characteristic 0, and set
 A = k[X, Y, Z]/(Zn - X Y) = k[x, y, z] for some n > 1. Then A/(z, x) N
 k[X, Y, 21/(X, Z) N k[ Y], so that p = (x, z) is a height 1 prime ideal of A.
 In D(A) we have np = div(x), and it can be proved that C(A) z Z/d (see
 [S2], p. 58). The relation xy = z? shows that A is not a UFD.
     For those wishing to know more about UFDs, consult [K], [S21 and
IFI.
\end{verbatim}
\end{proof}


\begin{theorem}[Printed Theorem 21.1]\label{mcrt-auto-theorem-21-1}\dcref{21.1}\source{matsumura-commutative-ring-theory:page-0185}\uses{}
\begin{verbatim}
Let (A,m, k) be a Noetherian            local ring, and A its
completion.
  (i) E&A) = ~~(2) for all p 3 0.
  (ii) el(A) 3 em bdimA-dimA.
  (iii) If R is a regular local ring, a an ideal of R and A N R/a, then
          p(a) = dim R - emb dim A + cl(A).
\end{verbatim}
\end{theorem}
\begin{proof}
\begin{verbatim}
(i) is clear from the fact that a minimal basis of m is a minimal
basis of mA^, so that applying OaA^ to the complex E, made from A
gives that made from A. Then since A is A-flat, H,(E.) @ A^ = H,(E. 0 A^),
and mH,(E.) = 0 gives H,(E.) @ A^= H,(E.).
   (ii) If A is a quotient of a regular local ring, then as we have seen above,
there exists a regular ring (R, n) such that A = R/a with a c n2, so that
el(A) = p(a) > ht a = dim R - dim A = emb dim A - dim A,           where    the
equality for ht a comes from Theorem 17.4, (i). Now A itself is not necessarily
a quotient of a regular local ring, but we will prove later (see 429) that A
always is, and we admit this in the section. Having said this, the two sides
of (ii) are unaltered on replacing A by A, and the inequality holds for A.
   (iii) Set n = rad(R). If a c n2 then, as we have seen above, p(a) = &i(A),
and dim R = embdim A, so that we are done. If a$ n2, take xEa - nZ;
we pass to R/xR and a/xR, each of dim R and p(a) decreases by 1, SO
    nduction completes the proof. n
    ition. A Noetherian local ring A is a complete intersection ring
    eviated to c.i. ring) if cl(A) = emb dim A - dim A.

heorem 21.2. Let A be a Noetherian        local ring.

 (ii) Let A be a c.i. ring and R a regular local ring such that A = R/a;
hen a is generated by an R-sequence. Conversely, if a is an ideal generated
   an R-sequence then R/a is a c.i. ring.
  iii) A necessary and sufficient condition for A to be a c.i. ring is that
     completion A^ should be a quotient of a complete regular local ring
        n ideal generated by an R-sequence.

      By Theorem 1, (iii), p(a) = dim R -embdim          A+ cl(A), and by
    rem 17.4, (i), ht a = dim R - dim A, so that A is a c.i. ring is equivalent
                         Theorem 17.4, (iii), this is equivalent to a being

        e sufficiency is clear from (i) and (ii). Necessity follows from the
    hat A^is a quotient of a complete regular local ring (see ?29), together


  eorem21.3. A c.i. ring is Gorenstein.
  oaf. If A is c.i. then so is A, and if A^is Gorenstein then so is A, so that we
    assume that A is complete. Then we can write A = R/a, where R is
   gular local ring and a is an ideal generated by a regular sequence. Since
  1s Gorenstein, A is also by Ex. 18.1. w
                                  g chain of implications for Noetherian local

       regular * c.i. + Gorenstein s CM.
      A be a c.i. ring, and p a prime ideal of A. If A is of the form
                xI), where R is regular and .x1,. . . , x, is an R-sequence, then
         can be written A, = R,/(x,, . . . , x,), where R, is regular and
         is an R,-sequence, it follows that A, is again a c.i. ring. The
          f deciding whether A, is still a c.i, ring even if A is not a quotient
        ar local ring remained unsolved for some time, but was answered
         ely by Avramov [l], making use of AndrC?s homology theory
        1. This theory defines homology and cohomology                     groups
           ) and H?(A, B, M)for n 3 0 associated with a ring A, an A-algebra
           module M. The definition is complicated, but in any case these
    B-modules having various nice functorial properties. If A is a
Noetherian     local ring with residue lield k then
           A is regular-H&          k, k) = 0,
and
           A is c.i.-H,(A,     k, k) = 0;
for n 2 3 the statements H,(A, k, k) = 0 and H,(A, k, k) = 0 are equivalent.
Thus Andrk homology is particularly relevant to the study of regular and
c.i. rings.
\end{verbatim}
\end{proof}


\begin{theorem}[Printed Theorem 22.2]\label{mcrt-auto-theorem-22-2}\dcref{22.2}\source{matsumura-commutative-ring-theory:page-0188}\uses{}
\begin{verbatim}
Let A be a ring, B a Noetherian              A-algebra,    and M
                                                                              171
a finite B-module; suppose that b is an M-regular element of rad (B).
Then if MIbM is flat over A, so is M.
\end{verbatim}
\end{theorem}
\begin{proof}
\begin{verbatim}
For each i > 0 the sequence 0 -+ M/b?M A M/b?+ 1M -+
M/bM + 0 is exact, so that by Theorem 7.9 and an induction on i, every
M/b?M is flat over A. Thus we can just apply the previous theorem.   n
\end{verbatim}
\end{proof}


\begin{theorem}[Printed Theorem 22.3]\label{mcrt-auto-theorem-22-3}\dcref{22.3}\source{matsumura-commutative-ring-theory:page-0189}\uses{}
\begin{verbatim}
In the above notation, suppose that one of the following
two conditions is satisfied:
    (a) I is a nilpotent ideal;
or (p) A is a Noetherian ring and M is I-adically ideal-separated. Then
the following conditions are equivalent.
   (1) M is flat over A;
   (2) Tort(N, M) = 0 for every A,-module N;
   (3) M, is flat over A, and ZO,M = ZM;
   (3?) M, is flat over A, and Torf(A,, M) = 0;
   (4) M, is flat over A,, and Y,, is an isomorphism for every n 3 0;
   (4?) M, is flat over A, and y is an isomorphism;
   (5) M, is flat over A,, for every n > 0.
   In fact, the implications (1) 3 (2)0(3)0(3?)   e(4)+(5) hold without any
assumption on M.
\end{verbatim}
\end{theorem}
\begin{proof}
\begin{verbatim}
First of all, let M be arbitrary.
   (l)*(2) is trivial.
   (2)=>(3) If N is an A,-module then we have
        N&M       = (NO,,A,)O,M         = NO,,M,,

and hence for an exact sequence O+ N, --+ N, -+N,                 -+O of Ao-
modules we get an exact sequence
        0 = Torf(N,,    M) -     N, OaoMo --+Nz OA,,MO -+
            N, O,&-fo -+ 0;
*herefore M, is flat over A,. Also, from the exact sequence 0-t I -+
 A - A0 --* 0 we get an exact sequence
          O=Tor:(A,,M)--tZ@M-+M+M,-+O,
?so that 10 M = IM.

                      an A,-module, we can choose an exact sequence of
  AO-modules 04 R +FO       -N     +O with F, a free A,-module. From this
i we get the exact sequence
           Tor:(F,,    M) = 0 -            Torf(N,    M) -       R @,,M,       -     F0 OAoMO,

                      is flat over         A,, the final arrow          is injective,      so that

                      we have Tor<(I/I?,  M) = 0, so that from O+ 1? +
                                   O~12~M-I~M-(1/1*)~M-t0
i is exact. From IBM = IM we get 1* @ M =1*&f              and (Z/1?)@ M ?v
\ IM/Z2M. Proceeding similarly, from 0 + I?+ ? + I? - Y/l?+ 1 + 0 we get
   by induction 1?? ? @ M = I?+ ?M and (,?/I?+ ?) @ M N I?M/I?+?M.     (4?) is
  just a restatement of (4).
     (4)+(5) We fix an n > 0 and prove that M, is flat over A,. For i d n we
:I have a commutative diagram
             (I?+?/~?+?)@M             +         (I?+?/I?)OM-(I?/I?+?)OM~O
                                                          ai I                       Yij
 ~-,~?+?JM,~~?+?M~~?+?M-~?~M,~~?~~?+?M                                  -          l?MII?+?M+O
 with exact rows. By assumption               yi is an isomorphism,          and since a, + 1 is an
                                             downwards      induction        on i we see that E,,,
 a,,-,, . . . ,CI~ are isomorphisms.         In particular,
           ~l:(l/I?+?)@,M        = IA,&M,=IM,,
 So that the conditions in (3) are satisfied by A,,, M, and I/I?+l. Therefore
 by (z)-(3),  we have To$(N,        M,) = 0 for every A,-module N. Now if
 N is an A,-module then IN and N/IN are both A,_,-modules,                   and
                          -+ 0 is exact, so that by induction on i we get finally
 that Tor:?(N, M,) = 0 for all An-modules N. Therefore M, is a flat A,,-

    Next, assuming either (a) or (p) we prove (5)+(l).    In case (a) we have
 A = A, and M = M, for large enough IZ, so that this is clear. In case (p),
                                 rove that the standard map J: a @ M -      M
                                   y hypothesis we have n,,Z?(a 0 M) = 0, so
 that we need only prove that Ker (i) c Y(a @ M) for all n > 0. For a fixed n,
 by the Artin-Rees   lemma, 1k n a c Z?a for suffkiently large k > n. We now
 consider the natural map
n~M~(njlknn)~M4-,(a/I?a)~M=(a~M)/l?(a~M).
Since M,- 1 is flat over A, _ I = AIlk, the map
         (a/lkna)O,M=(a/lkna)O,,_,Mk-,               -+M,-,
is injective, so that from the commutative      diagram
         a@ M L(a/lkna)@             M
          I
           I                  I
          M-                  M,-,
we get Ker (j) c Ker(f) c Ker(gf) = Z?(a@ M). This is what we needed to
prove. H
   This theorem is particularly effective when A is a Noetherian local ring
and I is the maximal ideal, since if A, is a field, M, is automatically flat
over A, in (3)-(4?). Also, in this case, requiring M, to be flat over A,, in
(5) is the same as requiring it to be a free AR-module, by Theorem 7.10.
   We now discuss some applications of the above theorem.
\end{verbatim}
\end{proof}


\begin{theorem}[Printed Theorem 22.4]\label{mcrt-auto-theorem-22-4}\dcref{22.4}\source{matsumura-commutative-ring-theory:page-0191}\uses{}
\begin{verbatim}
Let (A, m) and (B, n) be Noetherian local rings, A and fi
their respective completions, and A -+ B a local homomorphism.
   (i) For M a finite B-module, set I%?= M ORB; then
           M is flat over A-&?   is flat over AoQ    is a flat over A.
   (ii) Writing M* for the (mB)-adic completion of M we have
           M is flat over AoM*     is flat over A-M*    is flat over A.
\end{verbatim}
\end{theorem}
\begin{proof}
\begin{verbatim}
(i) The first equivalence comes from the transitivity law for flatness,
together with the fact that fi is faithfully flat over B; the second, from
the fact that both sides are equivalent to fi/m?fi      being flat over A/m?
for all n > 0.
   (ii) All three conditions are equivalent to M/m?M being flat over A/m?
for all n.
\end{verbatim}
\end{proof}


\begin{theorem}[Printed Theorem 22.5]\label{mcrt-auto-theorem-22-5}\dcref{22.5}\source{matsumura-commutative-ring-theory:page-0191}\uses{}
\begin{verbatim}
Let (A,m, k) and (B,n, k?) be Noetherian local rings,
A --+ B a local homomorphism,  and U: M -+ N a morphism of finite
B-modules. Then if N is flat over A, the following two conditions are
equivalent:
   (1) u is injective and N/u(M) is flat over A;
   (2) U: M @Ak - N @Ak is injective.
\end{verbatim}
\end{theorem}
\begin{proof}
\begin{verbatim}
(1) * (2) is easy, so we only give the proof of (2) + (1). Suppose that
EM is such that u(x) = 0; then U(X) = 0, so that X = 0, in other words,
xEmM.          Now assuming xEm?M,         we will deduce xEm?+rM.         Let
{a 1,. . . ,a,} be  a minimal basis of the A-module  m?, and  write  x = 1 &Yi
with y+M; then 0 = C aiu(y,). Since N is flat over A, by Theorem 7.6 there
exist cij~A and zj~N such that
             1 aicij = 0 for all j, and u(y,) = c cijzj for all i.
                                                  i
By choice of a,,...,a,,    all the cijEm, and hence u(y,)EmN     and
U(yi) =O, SO that yi =O, and y,EmM. Therefore xEnt?+?M.     We have
proved that x~n,m?M       = 0, and hence u is injective. Now from
O+ M --+ N - N/u(M)+0         we get Torf(k, N/u(M)) = 0, so that by
Theorem 3, N/u(M) is flat over A. n
\end{verbatim}
\end{proof}


\begin{theorem}[Printed Theorem 22.6]\label{mcrt-auto-theorem-22-6}\dcref{22.6}\source{matsumura-commutative-ring-theory:page-0192}\uses{}
\begin{verbatim}
Let A be a Noetherian ring, B a Noetherian A-algebra, M
a finite B-module, and DEB a given element. Suppose that M is flat over
A and that b is M/(P n A)-regular for every maximal ideal P of B; then b is
M-regular and M/bM is flat over A.
\end{verbatim}
\end{theorem}
\begin{proof}
\begin{verbatim}
Write K for the kernel of M 2 M; then K = 00 K, = 0 for all P.
Hence b is M-regular if and only if b is M,-regular  for all P. Moreover,
according to Theorem 7.1, A-flatness is also a local property in both A
and B, so that we can replace B by B, (for a maximal ideal P of B), A by
A (Pna) and M by M,, and this case reduces to Theorem 5. n
\end{verbatim}
\end{proof}


\begin{theorem}[Printed Theorem 23.3]\label{mcrt-auto-theorem-23-3}\dcref{23.3}\source{matsumura-commutative-ring-theory:page-0195}\uses{}
\begin{verbatim}
Let (A,m, k) and (B,n, k?) be Noetherian local rings, and
cp:A -B   a local homomorphism.   Let M be a finite A-module, N a finite
B-module, and assume that N is flat over A. Then
        depth,(M    OA N) = depth, M + depth,(N/mN).
\end{verbatim}
\end{theorem}
\begin{proof}
\begin{verbatim}
Let x1,. . , x,Em be a maximal M-sequence, and yr, . . . , ysEn a
maximal N/mN-sequence.         Writing xi for the images of xi in B, let US
prove that xi,. . . , XL, y,,. . . ,y, is a maximal M @N-sequence. Now
Xi,..., XL is an M Q N-sequence, and if we set M, = M/zXiM       then

        mEAssA(M,),          and   (M $9 N)/ i    xxM @ N) = M, 0 N.
                                            i=l

Moreover,   by the corollary        of Theorem      22.5, y, is N-regular,   and
N1 = N/ylN   is flat over A, so that from the exact sequence O+
     NLN-+N,+O            we get the exact sequence O-+M,@N-M,@N
     -M,    @ N, 40. Proceeding in the same way we see that y,, . , y, is an
     M,@ N-sequence. After this we need only prove that the B-module.

             (MON)/?(C-$(MO        N) + CYj(MO     N)) = M,O N,

F has depth 0, that is tt~Ass~(M,@ NJ; however, m~Ass,,(M,)                   and
: n~Ass,(N,/mN,),        so that this follows at once from the previous theorem.
i:
 :, Corollary. Let A --+B be a local homomorphism         of Noetherian rings as
     in the theorem, and set F = B/mB. Assume that B is flat over A. Then
a?      (i) depth B = depth A + depth F;
        (ii) B is CMoA    and F are both CM.
I, Proof. (i) is the case M = A, N = B of the theorem. From (i) and (*) we have
 ,?,
                dim B - depth B = (dim A - depth A) + (dim F - depth F)
,, and in view of dim A > depth A and dim F 3 depth F, (ii) is clear. n
,+

r, Theorem 23.4. Let A +B         be a local homomorphism         of Noetherian
i? local rings, set m = rad(A) and F = B/mB. We assume that B is flat over A;
.? then
             B is GorensteinoA    and F are both Gorenstein.
   Proof (K. Watanabe Cl]). By the corollary just proved, we can assume
j that A, B and F are CM. Set dim A = r and dim F = s, and let {xi,. . .,x,}
   be a system of parameters of A, and {yl,. . . , y,} a subset of B which
   reduces to a system of parameters of F modulo mB. Then as we have seen
   in the proof of Theorem 3, (x1,. . . ,x,, y,,. . ., y,} is a B-sequence, and
: therefore a system of parameters of B, and B= B/(x, y)B is flat over
   A= A/(x)A. Thus replacing A and B by A and B, we-can reduce to
: the case-dim A = dim B = 0. Now in general, a zero-dimensional local ring
 _ (R, M) is Gorenstein if and only if Hom,(R/M,      R) = (0: M)R is isomorphic
   to R/M. Now set
             rad(B) = n, rad(F) = n/mB = ii and (O:m), = I.
  Then I is of the form I N_(A/m)? for some t, and (O:mB), = ZB N (A/m)* 0
  B = F?. Furthermore, we have (O:n), = (O:n),, N ((O:fi),)?, and hence if we
: set(O:n), N (F/ii)? = (B/n)? then (O:n), N (B/n)??. Therefore
             B is Gorensteino     tu = 1o t = u = 1o A are F are Gorenstein.    n
\end{verbatim}
\end{proof}


\begin{theorem}[Printed Theorem 23.5]\label{mcrt-auto-theorem-23-5}\dcref{23.5}\source{matsumura-commutative-ring-theory:page-0196}\uses{}
\begin{verbatim}
If A is Gorenstein then so are A[X] and A[Xl.
\end{verbatim}
\end{theorem}
\begin{proof}
\begin{verbatim}
We write B for either of A[X] or A[X& so that B is flat over A.
For any maximal ideal M of B we set MnA = p and AJp,4, = qP). In
case B = A[X], the local ring B, is a localisation of BaAA, = 4,[~],
and the libre ring of A, --+BM is a localisation of !c(p)[X], hence
regular. In case B = A[XJ then XEM, and p a maximal ideal of 4, so
that K(P) = A/p and
         B @A4~) = (Ah)[Xl = +?)BXI/.
This is a regular local ring, and is the libre ring of A, --+ B,. Thus in
either case B, is Gorenstein by the previous theorem.      H
\end{verbatim}
\end{proof}


\begin{theorem}[Printed Theorem 23.6]\label{mcrt-auto-theorem-23-6}\dcref{23.6}\source{matsumura-commutative-ring-theory:page-0197}\uses{}
\begin{verbatim}
Let A be a Gorenstein ring containing a field k; then for
any finitely generated field extension K of k, the ring A &K is Gorenstein.
\end{verbatim}
\end{theorem}
\begin{proof}
\begin{verbatim}
We need only consider the case that K is generated over k by one
element x. If x is transcendental over k then A QK is isomorphic to a
localisation of A &k[X]      = A[X], and since A[X] is Gorenstein, so is
A@ K. If x is algebraic over k then since K CTk[X]/(f(X))   with f(x)Ek[X]
a manic polynomial, we have
          A 0 K = A[xl/U-(9);
now A[X] is Gorenstein and f(X) is a non-zero-divisor       of A[X], so that
we see that A @ K is also Gorenstein. w
\end{verbatim}
\end{proof}


\begin{theorem}[Printed Theorem 23.7]\label{mcrt-auto-theorem-23-7}\dcref{23.7}\source{matsumura-commutative-ring-theory:page-0197}\uses{}
\begin{verbatim}
Let (A, m, k) and (B, n, k?) be Noetherian local rings, and
A -B        a local homomorphism;    set F = B/mB. We assume that B is flat
over A.
   (i) If B is regular then so is A.
   (ii) If A and F are regular then so is B.
\end{verbatim}
\end{theorem}
\begin{proof}
\begin{verbatim}
(i) We have Torf(k, k) @AB = Torf(B@ k, B@ k), and the right-
hand side is zero for i > dim B. Since B is faithfully flat over A, we have
TorA(k, k) = 0 for i >>0, so that by 5 19, Lemma 1, (i), proj dim, k < Q and
since proj dim k = gl dim A, by Theorem 19.2, A is regular.
   (ii) Set r = dim A and s = dim F. Let {x1,. . . , xl} be a regular system
of parameters of A, and (yr,. . . , y,} a subset of n which maps to a
regular system of parameters of F. Since A --+ B is injective, We can ?lew
                                                                          but
A as a subring A c B. Then {x1,. . . , xv, y1 , . . . , ys) generates nj
dim B = r + s, so that B is regular. n
Flatness andfibres                                               183
\end{verbatim}
\end{proof}


\begin{theorem}[Printed Theorem 23.8]\label{mcrt-auto-theorem-23-8}\dcref{23.8}\source{matsumura-commutative-ring-theory:page-0198}\uses{}
\begin{verbatim}
Serre). (R,) + (S,) are necessary and sufficient conditions
    for a Noetherian   ring A to be normal.
\end{verbatim}
\end{theorem}
\begin{proof}
\begin{verbatim}
We defined a normal        ring (see 99), by the condition that the
    localisation at every prime is an integrally closed domain. The conditions
    (Ri) and (Si) are also conditions on localisations, so that we can assume
    that A is local.
    Necessity. This follows from Theorems 11.2 and 11.5.
    Suficiency. Since A satisfies (R,) and (S,) it is reduced, and the shortest
    primary decomposition of (0) is (0) = P, n ... n P,, where Pi are the minimal
    primes of A. Thus if we set K for the total ring of fractions of A, we have
               K=K,   x . . x K,, with Ki the field of fractions of A/P,.
    First of all we show that A is integrally closed in K. Suppose that we
    have a relation in K of the form
             (a/b)?+c,(a/b)?-?+...+c,=O,
    with a, b, cl ,..., C,EA and b an A-regular element. This is equivalent    to
    a relation
            a?+fc,a?-?b?=O
                   1

    in A. Let PESpecA be such that ht P = 1; then by (R,), A, is regular,
    and therefore normal, so that a,gb,A,, where we write ap, b, for the
images in A, of a, b. Now b is A-regular, so that by (S,), all the prime
divisors of the principal ideal bA have height 1; thus if bA = q1 n...nq,
is a shortest primary decomposition and we set pi for the prime divisor
of qi, then aebApcnA = qi for all i, and hence agbA, so that a/be-A.
Therefore A is integrally closed in K; in particular, the idempotents e, of
 K, which satisfy ef - ei = 0, must belong to A, so that from 1 = zei and
eiej = 0 for i fj we get
          A = Ae, x ... x Ae,.
Now since A is supposed to be local, we must have Y= 1, so that A is an
integrally closed domain.      n
\end{verbatim}
\end{proof}


\begin{theorem}[Printed Theorem 23.9]\label{mcrt-auto-theorem-23-9}\dcref{23.9}\source{matsumura-commutative-ring-theory:page-0199}\uses{}
\begin{verbatim}
Let (A, m) and (B, n) be Noetherian local rings and A -+B
a local homomorphism.         Suppose that B is flat over A, and that i 3 0 is a
given integer. Then
   (i) if B satisfies (Pi), so does A;
    (ii) if both A and the libre ring B@,k(p) over every prime ideal p of
A satisfy (Ri), so does B.
   (iii) The above two statements also hold with (Si) in place of (Ri).
\end{verbatim}
\end{theorem}
\begin{proof}
\begin{verbatim}
(i) For pESpec A, since B is faithfully flat over A, there is a prime
ideal of B lying over p; if we let P be a minimal element among these
then ht(P/pB) = 0, so that ht P = ht p. Hence ht p < i*B,          is regular, so
that by Theorem 7, A, is regular. Also, by the corollary to Theorem 3,
depth B, = depth A,, so that one seeseasily that (Si) for B implies (Si) for A.
   (ii) Let PESpec B and set P n A = p. If ht P d i then we have ht p < i
and ht (P/pB) ,< i, hence A, and B&B,         are both regular, so by Theorem
7, (ii), B, is regular. Hence B satisfies (Ri). Moreover, for (Si) we have
            depth B, = depth A, + depth B&B,
                      3 min(ht p, i) + min(ht P/pB, i)
                      3 min (ht p + ht P/pB, i) = min (ht P, i). w
\end{verbatim}
\end{proof}


\begin{theorem}[Printed Theorem 24.2]\label{mcrt-auto-theorem-24-2}\dcref{24.2}\source{matsumura-commutative-ring-theory:page-0201}\uses{}
\begin{verbatim}
Let A be a Noetherian          integral domain, R a finitely
generated A-algebra, and S a finitely generated R-algebra; we let E be a
finite S-module, M c E an R-submodule which is finite over R, and N C R
an A-submodule which is finite over A, and set D = E/(M + N). Then there
exists 0 # a6A such that D, is a free A,-module.
\end{verbatim}
\end{theorem}
\begin{proof}
\begin{verbatim}
Write A? for the image of A in R, and suppose that R =
A?[u,,...,   uJ; similarly, write R? for the image of R in S, and suppose
that S = R?[v,, . . , vJ. We work by induction on h + k; if h = k = 0 then
D is a finite A-module, and we have already dealt with this case,
   Write Rj = A?[u,, . . . , Uj] for O<j<h,      and Sj=R?[v,,...,vj]   for
O<j<k.
   Suppose first that k > 0; set M + N = M? c E. We have a filtration
           S,M? c S,M? c . ..c&M?=SM?cE.
the successive quotients of which are S,M?, S,M?/S,M?,          . . , S,- ,M?/
S,- zM?, S,M?&       r M?, EISM?. We can apply the induction hypothesis
to each of these except the last two. By virtue of the lemma, S,M?/S,-, M?
has a filtration with (up to isomorphism)     just a finite number of finite
S,- ,-modules appearing as quotients, and so we can apply the induction
hypothesis again. For the final term, write E? = E/SM?, and let e,, , , , e,
be a set of generators of E? over S; write E,-, = S,-,e, + ... + Sk-len.
Then SE,- r = E?, so that the lemma again gives a filtration of E? with
essentially finitely many finite s,- ,-modules appearing as quotients, and
we can apply the induction hypothesis to this term also.
   If k = 0 then E is a finite R-module, and replacing E by E/M we can
assume that M = 0. The preceding proof then applies almost verbatim to
this case, with Rj instead of Si. n
\end{verbatim}
\end{proof}


\begin{theorem}[Printed Theorem 24.3]\label{mcrt-auto-theorem-24-3}\dcref{24.3}\source{matsumura-commutative-ring-theory:page-0202}\uses{}
\begin{verbatim}
Let A be a Noetherian ring, B a finitely generated A-algebra,
and M a finite B-module. Set U = {PESpecBIM,               is flat over A);
 then U is open in Spec B.
\end{verbatim}
\end{theorem}
\begin{proof}
\begin{verbatim}
We verify the conditions (1) and (2) of Theorem 2.
    (1) If P XJQ are prime ideals of B then for an A-module N we have
 NgAMQ = (N OaMP)OBpBQ, so that if M, is flat over A then so is M,.
    (2) Let PEU and p = Pn A; set A = A/p. Now if QEV(P), we have
gB, c: rad(BJ, and hence by Theorem 22.3, M, is flat over A if and only
 if M&M,       is flat over A and Tor:(Mo, A) = 0. Now Tor:(M,, A) = 0,
;and the left-hand side is equal to Tor$(M, A)@Bp. By computing the
?Tar by means of a finite free resolution of A over A, we see that
,Torf(M, A) is a finite B-module, so that there is a neighbourhood         W
 pf P in Spec B such that Torf(M,,A)          = 0 for QE W. Moreover, by
~$%eorem 1, there exists aEA - p such that M,/pM, is a free &-module,
 so that if Q$V(aB), then M&M,          is flat over A. Thus the open set
{Wn V(P)) - V(aB) of V(P) is contained in U. n
 ?pemark. If A is Noetherian and B is a finitely generated A-algebra which
Ijs flat over A then it is also known that the map Spec B - Spec A is open;
,,pe [M], p. 48 or [G2], (2.4.6).
     Let A be a ring, and P a property of local rings; we define a subset
:kA) t SpecA by P(A) = {p&pec AIP holds for Ap}. For example, if
$? = regular, complete intersection, Gorenstein or CM we write Reg (A),
 %I (A), Gor (A) or CM (A) for these loci. The question as to whether P(A)
     open is an interesting and important question. For Reg(A) this is a
    assical question, but for the other properties the systematic study was
      tiated by Grothendieck.
     The following proposition is called the (ring-theoretic) Nagata criterion
    r the property P, and we abbreviate this to (NC).
     (NC): Let A be a Noetherian ring. If P(A/p) contains a non-empty open
   ubset of Spec (A/p) for every pe Spec A, then P(A) is open in Spec A.
  - The truth or otherwise of this proposition depends on P; in the
    mainder of this section we discuss some P for which (NC) holds. In
        24.2 and Ex. 24.3 we illustrate how (NC) can be applied to prove
         ness results.
   Orem 24.4 (Nagata). (NC) holds for P = regular.
  OOf. Let U = Reg (A). A localisation of a regular local ring is again
regular, so that U satisfies condition (1) of Theorem 2. We now check
condition (2). If PE U then A, is regular, so that we can take x1,. . . , X,EP
to form a regular system of parameters of A, (where n = ht P). Then there
exists a neighbourhood    W of P in Spec A such that
          PA, = (x,, . . , x,)Ao
for all QE W. (In fact, if aeA is an element not contained in P, but
contained in every other prime divisor of (x,, . . ,x,) then PA, =
(x 1,. , x,)A,.) Moreover,       by the hypothesis in (NC) there exists a
neighbourhood    W? of P in V(P) such that A,/PA,      is regular for QEW'.
Then A, is regular for QE W?n W, SO that W?n W c u. n
\end{verbatim}
\end{proof}


\begin{theorem}[Printed Theorem 24.5]\label{mcrt-auto-theorem-24-5}\dcref{24.5}\source{matsumura-commutative-ring-theory:page-0203}\uses{}
\begin{verbatim}
(NC) also holds for P = CM.
\end{verbatim}
\end{theorem}
\begin{proof}
\begin{verbatim}
As with the previous proof, we reduce to checking condition (2)
 of Theorem 2. Let PECM(A).             If we take UEA -P and replace A by A,
 then we are considering a neighbourhood           of P in Spec A, so that we will
 refer to this procedure as ?passing to a smaller neighbourhood of P. Since
 A, is CM, if ht P = n we can choose an A,-sequence y,, . . , y,,~p. One
 sees easily that after passing to a smaller neighbourhood           of P, we can
 assume that
    (a) y, , . . , y, is an A-sequence; and
    (b) I = (y, , . . , y,)A is a P-primary ideal.
Then for QEV(P), it is equivalent to say that A, is CM or that AJIA,
is CM. Thus replacing A by A/I we can assume that 0 is a P-primary ideal.
Then P? = 0 for some r > 0. Now consider the filtration 0 c Pr- ? c ... c
P c A of A. Each Pi/Piil IS a finite A/P-module, but A/P is an integral
domain, so that passing to a smaller neighbourhood           of P we can assume
that Pi/P i+l is a free A/P-module for 0 < i < r. It is then easy to see that
if x1,..., X,EA is an A/P-sequence, it is also an A-sequence. However,
according to the hypothesis in (NC), passing to a smaller neighbourhood
of P, we can assume that A/P is a CM ring. Then for QE V(P) the ring
A,/PA,      is CM, so that from what we have said above,
           depth A, 3 depth A,/PA, = dim A,/PA, = dim A,,
and A, is CM.      n
   Let A be a Noetherian ring and I an ideal of A; we set B = A/I and
write Y for the closed subset V(I) c Spec A. Let hil be a finite A-module.
We say that M is normally -flat along Y if the B-module gr,(M) =
 @ff& I?MJI?+?M      1s fl at over B. If B is a local ring, this is the same as
saying that each I?M/I?+ ?M is a free B-module. Normal flatness is an
important notion introduced by Hironaka, and it plays a leading role in
the problem of resolution of singularities; we have used it in the above
proof in the statement that if P is nilpotent and A is normally flat along
then an A/P-sequence is an A-sequence. However, in this book we
         ot have space to discuss the theory of normal flatness any further,
           e refer to Hironaka [l] and [G2], (6.10).
\end{verbatim}
\end{proof}


\begin{theorem}[Printed Theorem 24.6]\label{mcrt-auto-theorem-24-6}\dcref{24.6}\source{matsumura-commutative-ring-theory:page-0204}\uses{}
\begin{verbatim}
(NC) holds for P = Gorenstein.
              Once more we reduce to verifying condition (2) of Theorem 2.
             se that PEGor (A); if ht P = n then since A, is CM, we can take
       . ..,x,EP forming an A,-sequence. Passing to a smaller neighbour-
            of P, we can assume that x1,. . . , x, is an A-sequence. Moreover,     ?
            ing A by A/(x,, . . , x,J we can assume that ht P = 0. In addition, we
         assume that P is the unique minimal prime ideal of A. Since A, is a
         -dimensional Gorenstein ring, we have
               Ext;(A/P, A)BAAP = Ext:JK(P), AP) = 0

             Hom,(A/P, A)OA A, = Hom,,(ti(P), AP) = K(P).
    Thus passing to a smaller neighbourhood,   we can assume that Exti(A/P,
    A)=0    and Hom,(A/P,     A) N A/P. In addition, as in the proof of the
    previous theorem, we can assume that Pi/Pi+ ? is a free A/P-module for
              Y - 1, where P? = 0. Then using
             o+P?/P?+?+P/P?+?+P/P?+.o,
    we get by induction       that Exti(P, A) = 0; from this it follows       that
    Exti(A/P, A) = 0, and in turn by induction that Exti(P, A) = 0, so that
    Ext;(A/P, A) = 0. Proceeding in the same way we see that ExtL(A/P, A) =
   0 for every i > 0. If we take an injective resolution 0 -+ A -I?     of A as an
$ A-module, and consider the complex obtained by applying Hom,(A/P,             -)
2 to it, from what we have just said we obtain an exact sequence
1 O+ A/P + Hom,(A/P,           I?), and this is an injective resolution of A/P as
E; an A/P-module. The same thing holds on replacing A by A, for QE V(P),
* and then setting k = IC(Q), we get Extiaipa,(k, A,/PA,) = Extip(k, Aa).
; Thus it is equivalent to say that A, is Gorenstein or that A,/PA,                is
   Gorenstein.     Therefore from the hypothesis          in (NC) we have that
   Gor(A)n     V(P) contains a neighbourhood      of P in V(P).    w
      The above proof is due to Greco and Marinari [l]. Their paper also
   proves that (NC) also holds for P = complete intersection.
\end{verbatim}
\end{theorem}
\begin{proof}
\notready
\end{proof}


\begin{theorem}[Printed Theorem 25.1]\label{mcrt-auto-theorem-25-1}\dcref{25.1}\source{matsumura-commutative-ring-theory:page-0208}\uses{}
\begin{verbatim}
First fundamental              exact sequence). A composite k L
 A --%B    of ring homomorphisms               leads to an exact sequence of B-
;modules

   (1)     Q,&B-~-?~&RB,A+O.
 where the maps are given by cr(d,,,a 0 b) = bd,,,g(a) and B(d,,,b)= d,,,b
 for aeA and &B. If moreover B is O-smooth over A then the sequence
    (2) 0 -+ R,,, 0 B -         R,,, -   QBiA -+ 0,
 obtained from (1) by adding 0 +at the left, is a split exact sequence.
\end{verbatim}
\end{theorem}
\begin{proof}
\begin{verbatim}
In order for a sequence N? AN           -%N? of B-modules to be
?exact, it is sufficient that for every B-module T, the induced sequence
           Hom,(N?,      T) L      Hom,(N,    T) z    Hom,(N?,   T)
is exact. Indeed, taking T = N?, we get a*B*(l T) = 0, and therefore /XX= 0;
and taking T = N/Imcr, we see easily that Ker p = Ima. From this, to
prove that (1) above is exact, it is enough to show that for any B-module
 T,
    (3) Der,(A, T)-     Der,(B, T)+-Der,(B,    T)+O
is exact, but this is obvious.
   Now suppose that B is O-smooth          over A. Choose DEDer,(A,                 T) and
consider the commutative diagram



         gt T
         AABeT
                            with q(a) = (ga, Da).


Then by assumption, there exists k:B -B*T         which can be added to
the diagram, leaving it commutative.     If we write k(b) = (6, D?h) then
D?: B - T is a derivation of B such that D = D?oy, and D? corresponds
to a B-linear map a?:Q,,, -   T. Now take T to be fiZAik@ B, and define
D by D(a) = dAik(a)@ 1, so that D = D?og implies that CL?C(= 1T. Thus (2)
is split. n
   Now consider     the case k A A 5 B when y is surjective; set
Ker g = m, so B = A/m. Then in (1) of the previous theorem we of course
have RsiA = 0, and we want to determine Kercc.
\end{verbatim}
\end{proof}


\begin{theorem}[Printed Theorem 25.2]\label{mcrt-auto-theorem-25-2}\dcref{25.2}\source{matsumura-commutative-ring-theory:page-0209}\uses{}
\begin{verbatim}
Second fundamental            exact     sequence).         In   the   above
notation, we have an exact sequence

  (4)   m/m2 LQ       ,wChB     AQ,,d
where 6, is the B-linear map defined by 6(xmodm2)                = dAlkx@ 1. If B is
O-smooth over k then
   (5) O+m/m2-QZ,,,@B-Cl,,,-+0
is a split exact sequence.
\end{verbatim}
\end{theorem}
\begin{proof}
\begin{verbatim}
We once more take an arbitrary          B-module     T and consider
  (6) Hom,(m/nt2,     T),   ?    Der,(A,   T)4: il*    Der,(B,     T).

For DEDerdA,     T), to say that 6*(D) = 0 is just to say that D(m) = 0, SO
that D can be considered as a derivation from B = A/m; hence (6) is exact.
If B is O-smooth over k then the extension


of the k-algebra B by m/m2 splits, that is there exists a homomorphism
of k-algebras s:B -+ A/m2 such that gs = 1,. Now sg: A/m2 --P A/m2 is a
homomorphism      vanishing on m/m2, and g(1 - sg) = 0, SO that if we set
A --+ A/m2 5           m/m? 5      T

is an element of Der,(A, T) satisfying S*(D?) = $. Indeed, for xEm, if we
let x = x mod m2 then
         D?(x) = t,b(D(X)) = Ii/(X - sg(X)) = I+@).
Therefore 6* is surjective. If we set T = m/m2 then we see that (5) is a
split exact sequence. n
\end{verbatim}
\end{proof}


\begin{theorem}[Printed Theorem 25.4]\label{mcrt-auto-theorem-25-4}\dcref{25.4}\source{matsumura-commutative-ring-theory:page-0211}\uses{}
\begin{verbatim}
Let K be a field of characteristic p, and let 0 # DEDer (K).
    (i) 1,D,D2 ,..., Dpm ? are linearly independent over K;
   (ii) the only way in which c0 + c, D + ... + cp-, Dp- ? with c+K can be
a derivation is if co = cz = ... = cpm 1 = 0.
Pro~$ For UEK, write aL for the operation of multiplying         by a; then
the property D(nx) = D(u).x +a.Dx        of a derivation means that DC-a,=
D(u), + CID. We can write the Leibnitz formula as

         D?oa, = aD?+     i.D(a)D?-?   +

our proof exploits this formula.
    (i) For some i < p suppose that 1, D,. . . , D?- ? are linearly independent
over K, but that l,D ,..., D? are not. Then we can write D? = ci- I Dim ? +
... + co, with c,EK. If we choose some UEK such that D(a) # 0, then in
view of DiouL = ciPl Die?ou, + ..., we get
           aDi + i.D(a)D?-? + ... = ci- ,aD?-?    + ...,
where    . indicates a linear combination of 1, D, . . . , Dim 2. Subtracting     a
times our original relation from this gives a relation of the form
        i.D(a)D?-?   = ... ,
and this contradicts     the assumption      that 1, D, . . . , Die? are linearly
independent.
   (ii) Suppose that E = ciD? + ... + c1 D + c0 is a derivation of K, with i C/J
and ci # 0. Then E(1) = co, so that c0 = 0. Now if i > 1 then take UEK
such that D(u) # 0, and substitute both sides of Eoa, = ciLhclL + ..? in the
Leibnitz formula: we get
         aE + E(a)L = aciDi + [i,q.D(a)      + ac,_,]D?-?    + ..?,
but then in view of the linear independence of 1, D, . . . , Dp- ?, the coefficients
of D?-? on both sides must be equal; therefore i.ci.D(u) = 0, which is a
contradiction.   H

Remarks. (i) The theorem also holds if char K = 0.
   (ii) If K is not a field, this result does not necessarily hold. For example,
let k be a field of characteristic       p, and set A = k[X]/(Xp)   = k[xl, with
$ = 0; then every derivation of k[X] will take the ideal (Xp) into itself,
and therefore induces a derivation of A. In particular, the derivation
Xp-l.a/dX           of k[X]   induces DEDer,(A)       such that D(x)=x~-~,     but
    i) = i.xi-lxP-l       = 0 if i > 1, and therefore  for p > 2 we have 0? = 0.
\end{verbatim}
\end{theorem}
\begin{proof}
\notready
\end{proof}


\begin{theorem}[Printed Theorem 25.5]\label{mcrt-auto-theorem-25-5}\dcref{25.5}\source{matsumura-commutative-ring-theory:page-0212}\uses{}
\begin{verbatim}
the Hochschild formula). Let A be a ring of characteristic
p; then for aEA and DED~~(A) we have
         (aD)? = apDP + (aD)?-?(a).D.
\end{verbatim}
\end{theorem}
\begin{proof}
\begin{verbatim}
Set E = aD. Then E2 = E~a,~D = (aE + E(a))D = a2 D2 + E(u)D,
and proceeding by induction, we get a relation of the form
                                           k-l
         Ek=ukDk+                           c bk,iDi+Ek-l(u)D,
                                           i=2

where bk,i are elements of A given                               by a pUrdy       fOrma   computation,     so

         bk,i         = fk,i(a,           D(U),   D2(a), . . .) D?-?(a)),
where the fk,i are polynomials with coefficients in Z/(p) not depending on
A, on a or on D. Now to prove our theorem we need only show that ,j?,,i = 0
for 1 < i < p. Let k be a field of characteristic p, and let x1, x2,. . . be a
countable number of indeterminates over k; set K = k(x,, x2,. .). Define
a k-derivation D of K by Dxi = xi+ r. (Since C12,,,is the free K-module with
basis dx,, dx,,. .., given any ,fieK there exists a unique DEDer,(K)
such that Dxi = fi.) For this D, we set E = x,D; then since EP- xpDp =
bp,p-lDP-l + ... + b,,2D2 + EP-l(u).D is a derivation, by the previous
theorem we must have b,,i = 0 for 1 < i < p. Therefore
         bp,i=fg,i(X1,X2,...,Xp~i+1)=0,


and this proves that f,,i = 0. w
  This formula is known as the Hochschild formula, although it is also
reported to have been first proved by Serre. Be that as it may, it is an
important fact that (uD)? is a linear combination of Dp and D.
\end{verbatim}
\end{proof}


\begin{theorem}[Printed Theorem 26.3]\label{mcrt-auto-theorem-26-3}\dcref{26.3}\source{matsumura-commutative-ring-theory:page-0215}\uses{}
\begin{verbatim}
If k is a perfect field then every extension field K of k is
 separable over k, and a k-algebra A is separable if and only if it is reduced.
\end{verbatim}
\end{theorem}
\begin{proof}
\begin{verbatim}
Recall that a field k is perfect if every algebraic extension of k is
 separable. Ifk has characteristic 0 then every extension field K is separably
 generated, and therefore separable. In characteristic p, perfect implies
 k = kl?p, so that by the previous theorem, all sublields of K finitely
generated over k are separable, so that K itself is separable over k. (Caution:
K may fail to be separably generated over k; for a counter-example, let x
be an indeterminate over k, and K = k(x, xp-l, xp-?, . .).) Now we show
that if A is a reduced k-algebra then A is separable. We can assume that
A is finitely generated over k. Then A is a Noetherian ring, and the total
ring of fractions K of A is of the form K = K, x ... x K, by Ex. 6.5. Each
Ki is separable over k, so that K is also separable, and hence so is its
 subring A. n
   In general two subfields K, K? of a given field L are said to be linearly
disjoint over a common subfield k if the following three equivalent
conditions are satisfied:
   (a)ifa,,...,     CI,EK are linearly independent over k they are also linearly
independent over K?;
   (b) the same with K and K? interchanged;
   (c) if we write K [K?] for the subring of L generated by K and K?j the
natural map K OkK? -+ K[K?] is an isomorphism.
Proofofequivalence.        (a)*(c) Let 4 = Cyxi o yi be an element of the kernel
of K OkK? - K[K?].             Suppose that x1,. . , x, are linearly independent
over k, and that the remainder x,, l,. . . , x, arc linear combinations of them7
and rewrite 4 = C; xi myi. The image in K[K?] of 5 is CxiY:, but if this is ?
then by (a) we have y; = 0 for all i, so that 4 = 0. This proves (C).
   (c)+(a)      is also easy; finally, since (c) is symmetric in K and K 12 we Of
course also get (a+(b).
~,,EK are linearly independent over k. If
                               we set k, = k(cl,. . . , {,), so that k, is a finite
 xtension of k; for some sufficiently large n we have kp? c k, and if we set
   = Kgkk,,      then A is a reduced ring. However, A is finite as a K-module,
    is a zero-dimensional ring, but the p?th power of any element of A is in K,
     that we see that A has only one prime ideal. Hence A is a field, and
   N K[k,]. From this we get x&i 0 ci = 0, that is ti = 0 for all i.
   (ii) If K and kP?? are linearly disjoint, then kPm? c kPm?,and hence K and kP-?
are also linearly disjoint over k, so that K QkkPm? is a field. If K? is a subfield
of K which is finitely generated over k then by Theorem 2, K? is separable
over k. Hence K is also separable over k. w



Let K be an extension field of a field k; then QKjk is a vector space over
K, and is generated by {dx(xeK},              so that there exists a subset B c K
such that {dxl xEB} forms a basis of the vector space Q,,,. A subset
Bc K with this property is called a differential basis of K over k. The
following condition (*) is necessary and sufficient for a subset {x~}~~~ c K
to form a differential basis for K over k.
   (*) if ~,EK are specified for every LEA in an arbitrary way, then there
exists a unique DEDer,(K) such that D(x,) = y, for all i.
    For x~,...,x,EK,       let us study the condition for dxl,...,dx,ERKik      to be
linearly independent over k. If k has characteristic 0 then this is equivalent
            x, being algebraically independent over k. Indeed, if there exists
                  1,. . . ,X,] such that J-(x,, . . . ,x,) = 0, then choose such a
relation of smallest degree; suppose for instance that X, actually appears in
f, SO that f1 = af/a X, is non-zero, but of smaller degree than f, and hence
f 1(x) # 0. Then f(x) = 0 gives
           0 = df = 1 fi(x)dxi,
                    ,dx, are linearly dependent. Conversely, if x1,. . . ,x, are
algebraically independent over k then there exists a transcendence basis
B of K/k containing        these, so that there exists k-derivations Di of the Purely
transcendental        extension k(B) satisfying
              Di(Xi) = 1 and Of(v) = 0 for Xi #DEB,
(namely, a/ax,). Moreover, K is a separable algebraic extension of k(D),
and so by Theorem 25.3 is 0-etale, so that the derivations Di extend to
derivations from K to K. Then since Di(xj) = 6ij, the differentials
dx 1, . . . , dw%,k       are clearly linearly independent. Thus in this case a
differential basis is the same thing as a transcendence basis.
    Now consider a field k of characteristic p. We say that elements
X1,...,-   Y,EK of an extension field K are p-independent over k if
 CKW, xl,. , x,):KP(k)]          = p?, and a subset B c K is p-independent if any
finite subset of B is p-independent. This condition means precisely that the
set
                               x1,. . . ,x, are distinct
                               elements of B and 0 6 Cli < p
is linearly independent over F?(k); the elements of rB are called the
p-monomials of B. If B is not p-independent, we say it is p-dependent.
The condition of p-independence is not just a property of B and k, but
also depends on K. If B c K is p-independent over k and K = KP(k,B),
we say that B is a p-basis of K/k. If C c K is p-independent over k then
one can easily show by Zorn?s lemma that there exists a p-basis of K/k
 containing C.
    One sees easily that B is a p-basis of K/k is equivalent to IB being a
basis of K over K?(k) in the sense of linear algebra. If this holds then any
map D: B-K             has a unique extension to an element DEDer,(K). Indeed,
for a p-monomial of B we set
           D(x?,?. . . x~~)   = Cr=   1 OriX9?.   . . X~I-   ?. . . x~?D(x~),

and extend D to K as a KP(k)-linear map; then D is a k-derivation. Thus
a p-basis B is a differential basis of K/k. Conversely, if B? is a differential
basis of K/k then B? is p-independent over k; for if xi,. ., x,EB? are
p-dependent, we can assume that xl~KP(k, x2,. ., x,), so that we can
write x1 = S(x,, . . . , x,), where f is a polynomial with coeflicients in KP(k).
Then in QKjk we get dx, = CZ(af/axi)d x i, which contradicts the linear
independence of dx, , . . , dx,. Now if we take a p-basis B of KJk containing
B?, then since both B and B? are differential bases, we have B = B?. We
summarise the above as follows:
\end{verbatim}
\end{proof}


\begin{theorem}[Printed Theorem 26.5]\label{mcrt-auto-theorem-26-5}\dcref{26.5}\source{matsumura-commutative-ring-theory:page-0217}\uses{}
\begin{verbatim}
The notion ofdifferential basis coincides with transcendence
basis in characteristic 0, and with p-basis in characteristic p.
  Now we look at the relation between separability and differential bases.
Letting II c k denote the prime subfield of k, we write Qk for !&,.
\end{verbatim}
\end{theorem}
\begin{proof}
\notready
\end{proof}


\begin{theorem}[Printed Theorem 26.8]\label{mcrt-auto-theorem-26-8}\dcref{26.8}\source{matsumura-commutative-ring-theory:page-0219}\uses{}
\begin{verbatim}
Let K/k be a separable extension of fields of characteristic
p, and let B be a p-basis of K/k. Then B is algebraically independent over
k, and K is 0-etale over k(B).
\end{verbatim}
\end{theorem}
\begin{proof}
\begin{verbatim}
Suppose by contradiction       that b,, . . . , ~,EB are algebraically
dependent over k. Suppose that 0 # fEk[X,,       . . . ,X,1 is a polynomial of
minimal degree among all those with f(b) = 0, and set deg f = d. Then write
         f(xl~~~~~Xn)=o~i~i                    ipgil..,i,(XE,...rX~)X;I...X~

                                 .    ,r..rn


Then since b,, . . . , b, are p-independent over k and f(b) = 0, we have
gi ,,.,i,(bP) = 0 for all values of i,, . . , i,. However, since
         d 3 deggi,...i,(XP)         + il + .?. + i,,
by choice of f we must have f(X) = g,,,e(XJ?). Hence we can write f in
the form f(X) = h(X)P, with hEkl?p[Xl,.      . ,XJ. However, since K and
klip are linearly disjoint over k, the monomials of degree < d in bl , . . . , bn,
being linearly independent over k, must also be linearly independent over
kllp. Thus h(b) # 0, but this contradicts h(b)P = f(b) = 0. For the second
half, see the proof of the following theorem.      n
\end{verbatim}
\end{proof}


\begin{theorem}[Printed Theorem 26.9]\label{mcrt-auto-theorem-26-9}\dcref{26.9}\source{matsumura-commutative-ring-theory:page-0219}\uses{}
\begin{verbatim}
If K is a separable field extension of a field k, then K is
O-smooth, and conversely.
\end{verbatim}
\end{theorem}
\begin{proof}
\begin{verbatim}
Let B be a differential basis of K/k. If K/k is separable then by
Theorem 5 and Theorem 8, k(B) is purely transcendental Over k. Then?
as one sees at once from the definition, k(B) is O-smooth over k. Moreover7
KIlJR\ is O-etale. In characteristic 0 this follows from Theorem 25.3; in
j   ?26     Separability                                                                          205

  characteristic p, by Theorem 6 we have an exact sequence O-tR,O K
   -aK       -+fiZKlk+O,  so that putting together an absolute p-basis of k
; with B we have an absolute p-basis of K, and this is clearly also an absolute
  p-basis of k(B). Hence K/k(B) is 0-etale by Theorem 7. Therefore K/k is
; O-smooth.
     Conversely, if K/k is O-smooth then by Theorem 25.1, QR,@ K -0,
  is injective, so that by Theorem 6, K/k is separable. n
5

$ Imperfection    modules and the Cartier           equality
FQtg ui e enerally, if k --+ A + B are ring homomorphisms,      we write rsiAik
i for the kernel of nAikO,,,B -fiZejk,    and call it the imperfection module
1; of the A-algebra B over k. If k = Z or k = Z/p (the prime field of
1 characteristic p) we omit k, and write IBIA.
1
1: Lemma 1. Let k + K -L         --+ L? be field homomorphisms.      Then there
1, exists a natural exact sequence
.*
            0 + rLiKlk 0, c - rLsIKIk - hk
               - QLjrc &Ii -t&f,,      - QL?,L + 0.

    Proof We have a commutative            diagram with exact rows.
            04-L,K,kc3Lfi -qlkcw                    -~L,k~L~            -Q,,,cW+o
                     II 1                      II              12I             f3 1
            O-+      r   L?/K/k     -%,k   @ L?        -         %?,k      -     clL.,I;   -+o.

    We abbreviate this as
            0+X-A-B-P-,0
                   i   II             I    1
            O+Y-A-C-Q+O.
    and from it construct
            o?Ai?-ia+o

            O-A/Y        -C-Q-to;
    applying the snake lemma gives the exact sequence
              O-+Y/X+Kerf,+Kerf,+O,
    from which we easily get
              0+X--+    Y-Kerf,     +P+Q-+Cokerf,+O.
    This is just what we wanted to prove. n
\end{verbatim}
\end{proof}


\begin{theorem}[Printed Theorem 26.10]\label{mcrt-auto-theorem-26-10}\dcref{26.10}\source{matsumura-commutative-ring-theory:page-0220}\uses{}
\begin{verbatim}
the Cartier equality). Let k be a perfect field, K an
    extension of k and I. a finitely generated extension field of K; then
       (*I   rb.QL,K = tr.deg,L + rkLrLIKIk,
    (where rk, denotes the dimension of a vector space over L).
\end{verbatim}
\end{theorem}
\begin{proof}
\begin{verbatim}
Suppose that k c K c L c L?, with L finitely generated over K and
L? finitely generated over L. If the theorem holds for k c L c L? and for
kc K c L then both rEILlk and r,,K, 0 L are finite-dimensional    over L,
so that by the lemma, f-L,,K,L is also finite-dimensional, and
           rkL,QLs,K - rkL,,K,k
              = (rk,,&,,, - rk,G~,L,k) + (rk,%,, - rWLiKik)
              = tr.deg,L?+    tr.deg,L= tr.deg,L?;
thus the theorem also holds for k c K c L?. Now every finitely generated
field extension can be obtained by a succession of the following three
kinds of simple extensions:
    (1) L = K(a) where c1is transcendental over K;
    (2) L = K(E) where c1is separable algebraic over K;
    (3) L = K(cc) where char K = p, and G??= aeK, but cr$K.
Hence we need only prove (*) in each of these special cases. (1) and (2)
are easy. For (3), if we write L = K[X]/(Xr        - a) we see that
           fhlk = (QKLXIIk 0 WLda
                = (R&K     da) 63 L 0 Lda,
and dcc# 0. Furthermore, since k is a perfect field, we have a$ Kr = kKP,
so that in QZKlk(= 0,) we have da # 0, rk R,,, = 1 and rk rLIKIL = 1, so that
(*) also holds in this case. n
\end{verbatim}
\end{proof}


\begin{theorem}[Printed Theorem 27.1]\label{mcrt-auto-theorem-27-1}\dcref{27.1}\source{matsumura-commutative-ring-theory:page-0223}\uses{}
\begin{verbatim}
If the ring A is O-smooth over a ring k, then a higher
derivation of length m < cc over k from A to an A-algebra B can be
extended to a derivation of length GO.

  This theorem can be applied for example to the case of a field k and a
separable extension field A of k.
If A is a ring of characteristic p then an ordinary derivation of A is
  zero on the subring AP, but higher derivations need not vanish on AP. If
  Q = (Do, D,, D,, . .) is a higher derivation of length m > p then from
 E,(d) = E,(a)P = up + D1(a)P.tP + ..., we get
            Dp(aP) = D, MP, and in general D,,(&) = D, (u)p?.
 For example, it follows from this that if k is a field of characteristic p, and
 K= k(a) with aPEk but a$k, then although there exists DEDerk(K)
 such that D(a) = 1, this D cannot be extended to a higher k-derivation of
 length B p. Since K is separable over the prime sublield, D can be extended
 to a higher derivation of length GO (over the prime subfield), but this
 extension cannot be trivial on k.
    We say that a higher derivation g = (Do, D,, . . .)EHS,(A) is iterative if
 it satisfies the following conditions:

                           Di+ j   for all    i,j.

 This condition   is equivalent    to asking that the following     diagram    is
 commutative:




  where E,(u) = C t?D,( a) and E,(x tYuv) = C t?E,(u,), and the right-hand
: vertical arrow is the inclusion map. Indeed,
           U-W))     = E,(~~?D&))    = 1 t?x u?D,D,W>
                                        ? B
: and

            E,+,(u) = T(t + u)?D,(u) = 1 t?x u?                 D,+,(u).
                                               v fl
;. If A contains the rational field Q, then one sees by induction on n that
   an iterative higher derivation            satisfies D, = D:/n!,   and is hence
  determined by D, only. Conversely, for DEDerk(A), we see that
  (1, D, D2/2!, D3/3!, . . .) is an iterative higher derivation. If A has charac-
i teristic p then for an iterative D = (D,, D,, . . .) we have Di = Di/i! for
; i <P, and DI = 0. Thus one cannot hope to extend any derivation to an
1 iterative higher derivation, even if A is a field.

i Th eorem 27.2. Let k LA 2 B be ring homomorphisms,    and suppose
? that B is 0-etale over A. Then given D = (D,,D,, . . .)EH&(A, B),
 there exists a unique D? = (Db,D;,          . .)EH?~(B)   such that Di(g(a)) =
D,(a) for all i. Moreover,    if D* = (D;F,DT,. ..)eHS,(A)      is such that
Di = goD* for all i, and if D* is iterative, then D? is also iterative.
\end{verbatim}
\end{theorem}
\begin{proof}
\begin{verbatim}
There is no problem about De = 1,. Now assume that (Db,. ..,D&)E
HS,(B, m) has been constructed so that Diog = Di for i < m; then if we define
                                                                           m+l
         h:A   -B,+            1   = B[t]/(t??)               by h(a) = c        t?DJa),

andu:B--+B,byu(b)=                    C~t?D~,(b),             we obtain the left-hand commutative
diagram.




                                                          h
          A h--?&+1                               A   -          B,+,


Hence by the 0-etale assumption, there exists a unique o:B -+ B,, , which
makes the right-hand diagram commutative. Repeating this we see that
D exists and is unique. If r>* is iterative, and we consider the homo-
morphisms    E,: A -+ A[r] and E;:B - B[t] corresponding      respectively
to Q* and D?, then we know E,o E, = E,,,, and we need only prove that
E;oE; = Ej,,. By induction on m, assume that
         EL(E;(b)) = E:+,,(b)mod(t,u)?+?                          for all bEB;
then from the commutativity                  of
                            C.?
          B           -                           But, uj/(t, u)?+ ?

          I,                                        II
          A -f%           A[t,uj      --+B[t,u]/k4m+2,

from E,,, = E,oE, and from the assumption                                that B is 0-etale over A, we get
         EI(E;(b)) = E;+,(b) mod (t, u)~+~                              for all beB.
This proves that EuoE; = E;,,.    m
\end{verbatim}
\end{proof}


\begin{theorem}[Printed Theorem 27.3]\label{mcrt-auto-theorem-27-3}\dcref{27.3}\source{matsumura-commutative-ring-theory:page-0225}\uses{}
\begin{verbatim}
(i) Let A be a ring of characteristic            p, and suppose
that XEA, DEDer(A)           satisfy Dx= 1 and Dp=O;           set A, = {~EAI
Da = 0). Then A is a free module over A,, with basis 1, x,. . , xp- ?.
   (ii) Let k be a field of characteristic p, and K a separable extension of k;
let DEDer,(K)      be such that DfO,           DP=O. Set K,={a~KlDa=o).
Then there exists an XE K such that Dx = 1, and a subset B, c K, such
that B = {x} uB, is a p-basis for K over k.
\end{verbatim}
\end{theorem}
\begin{proof}
\begin{verbatim}
(i) Suppose that a0 + zlx + ... + xixi = 0 for some i < p and ccje,4,.
Then applying D? we get i!a, = 0, hence ri = 0. Hence by downward
induction on i we see that 1, x, . . . , xp- ? are linearly independent over A,.
Now in view of Dp = 0, for every aeA we have D?+ ?a = 0 for some
0 < i <p. If i = 0 we have LIEA,,. If i > 0 then D?(a - x?D?a/i!) = 0, so that we
see by induction that
         D??u   = O=aEA,      + A,x + ... + A,x?;
?Setting i = p - 1 in this gives A = A, + A,x + ... + A,xP-?.
    (ii) Since D # 0 we can find ZCSK such that Dz # 0. Now in view of
DPz = 0, there is an i such that D?z # 0 but D?+?z = 0. If we set y = D?z
 and x = (D?- ?z)/y then Dx = 1, so that by (i) we have K = K,(x) and
[K:K,J =p. Now if we had xPgKP,k, then x~K,k?~,                    and we could
write x=x;wiai        with q,...,       w,EK~ linearly independent over k, and
aEk?lP. Now kc K, and x#K,,                   so that x,ol ,..., o, are linearly
independent over k, and hence by the assumption that K is separable over
k, they are also linearly independent over kllp, which contradicts x = 1 qq.
 Thus xP$Kgk. Hence we can choose a p-basis C of K, over k such that
 x~EC; set B0 = C - {xp). Then if y,, . . . ,y, are distinct elements of B,,
 we have [KP,k(xp,y,,...,y,):K$k]          = p?+l, and together with K = K,(x)
 this gives [KP k(y , , . . . , y,):KPk] =p?. Thus B, as a subset of K is
p-independent      over k. Since also K, = KP,k(xP, B,), we have K =
K,(x) = KPk(x, II,), so that setting B = B, u {x} we get a p-basis of K
 over k. n
\end{verbatim}
\end{proof}


\begin{theorem}[Printed Theorem 27.4]\label{mcrt-auto-theorem-27-4}\dcref{27.4}\source{matsumura-commutative-ring-theory:page-0226}\uses{}
\begin{verbatim}
Let K be a field of characteristic p, and k c K a subfield
such that K is separable over k. Then a necessary and sufficient condition
for DEDer,(K)    to extend to an iterative element of H&(K) is that Dp = 0.
\end{verbatim}
\end{theorem}
\begin{proof}
\begin{verbatim}
We have already seen necessity, and we prove sufficiency. We can
assume that D #O; if DP = 0 then we can choose K,,x and B, as in
Theorem 3, (ii). We set K? = k(B,); then D is a K?-derivation,    K is 0-etale
over K?(x), and K?(x) is a purely transcendental extension of K?. We define
a homomorphism       @K?(X) -K?(x)[t]      by setting E,(a) = a for CCEK? and
J%(X) = x + t; then
          E?(E,(X)) = x + 24+ t = E,+,(x),
so that E,oEt= Et+u holds over the whole of K?(x). Thus E, defines a
iterative higher derivation D of K?(x) over K?. Since K is O-&ale over K?(x),
bY Theorem 2 there is an extension of D to an iterative higher derivation of
K over K?; the term of degree 1 in D is D, so that Q is an extension of D (or
more precisely, of (1, D)). n
\end{verbatim}
\end{proof}


\begin{theorem}[Printed Theorem 28.1]\label{mcrt-auto-theorem-28-1}\dcref{28.1}\source{matsumura-commutative-ring-theory:page-0229}\uses{}
\begin{verbatim}
Transitivity).    Let A 5 B LB?         be ring homomor-
phisms, and suppose that g? is continuous for the I-adic topology of B
and the I?-adic topology of B?; if B is I-smooth over A, and B? is I?-smooth
over B then B? is I? smooth over A. The same thing holds with I-unramified
in place of I-smooth.
\end{verbatim}
\end{theorem}
\begin{proof}
\begin{verbatim}
Suppose that u is given in the diagram;




then since ug?:B -C/N         is continuous, by the I-smoothness of B, there
exists a lifting w:B - C. Next by the It-smoothness of B? over B, we can
lift u to v:B? --+ C. Also if B is I-unramihed over A, and the map v in the
diagram exists, then w = ug? is unique, and if in addition B? is I?-unramified
over B then v is unique. n
\end{verbatim}
\end{proof}


\begin{theorem}[Printed Theorem 28.2]\label{mcrt-auto-theorem-28-2}\dcref{28.2}\source{matsumura-commutative-ring-theory:page-0229}\uses{}
\begin{verbatim}
Base-change). Let A be a ring, B and A? two A-algebras,
and set B? = B@, A?. If B is Z-smooth over A then B? is IB?-smooth over
A?. The same thing holds for I-unramified.
\end{verbatim}
\end{theorem}
\begin{proof}
\begin{verbatim}
We have the diagram
             B p       B? -L    C/N


             T        4J
             A -A?         -    C,

    where p, q are the natural homomorphisms.      Then if u satisfies u(Z?B?) = 0,
    there is a lifting v: B --+ C of up. Then if we define v?: B? = B @AA? - C,
    by U? = v @A, this is a lifting of u to C. For unramilied, this is clear from
    the fact that v? is uniquely determined by v. n
\end{verbatim}
\end{proof}


\begin{theorem}[Printed Theorem 28.3]\label{mcrt-auto-theorem-28-3}\dcref{28.3}\source{matsumura-commutative-ring-theory:page-0230}\uses{}
\begin{verbatim}
Let (A, m, K) be an equicharacteristic      local ring. Then
(i) A has a quasi-coefficient field;
   (ii) if A is complete, it has a coefficient field;
   (iii) if the residue field K of A is separable over a subfield k c A then
A has a quasi-coefficient field K? containing k;
   (iv) if K? is a quasi-coefficient field of A, then there exists a unique
coefficient field K? of the completion A containing K?.
\end{verbatim}
\end{theorem}
\begin{proof}
\begin{verbatim}
(iii) Suppose that B = (r1,t2,. . .} is a differential basis of K/k,
and for each ti, choose an inverse image x+A. Then by Theorem 26.8,
t135 2,. . . are algebraically independent over k, so that the subring
4X1,x2,... ] of A meets m in (O}, and hence A contains the field
K? = k(x,, x2,. . .). We identify K? with its image k(B) in K, so that K is
clearly 0-etale over K?, and K? is a quasi-coefficient field, as required.
    (i) By assumption A contains a field, so that it contains a perfect field
(for example, the prime subfield). We need only apply (iii) to this.
            K =       ;i/ti
           t        t,
           K?-      A
  (iv) In the diagram above, there exists a unique lifting of the identity
map K - A/lit to K - 2, and its image is the required coefficient field.
  (ii) follows from (i) and (iv). n
  The next lemma will be made more precise in Theorem 28.7.
Lemma I. Suppose that (A, m, K) is a Noetherian local ring containing a
field k. If A is m-smooth over k then A is regular. The converse holds if
the residue field K is separable over k.
Proof. Take a perfect subfield k, c k; then k is O-smooth over k,, so that
by transitivity, A is also m-smooth over k,, so that we can assume that
k is a perfect field. Also, replacing A by A, we can assume that A is
complete. Then A has a coefficient field containing k; for ease of notation
we write K for this, and identify it with the residue field. If {x1,. . . ,x,)
is a minimal basis of m then as K-algebras we have
           A/m2 N K[X,, . . . ,X,]/(X,, . . . ,X,J2.
The composite
        A-     A/m2 ZK[X,,       . . . ,X,]/(X)?   Z   K[X,,.   . ., X,j/(X)2
lifts to A-KIX1,...,      X,1, and by Theorem 8.4, this is surjective. Thus
dimA3dimK[X,,...,X,J=           n, and together with embdim A = n this
gives the regularity of A.
    Conversely, if A is regular and K is separable over k, then A has a
coeficient tield K containing k. Let {x1,. . . ,x,) be a regular system
I-smoothness                                                         217

    of parameters       of A, and define a homomorphism          of K-algebras
                         -+ A^ by It/(X,) = xi; then once more by Theorem 8.4,
    Ic/ is surjective, and comparing dimensions, we see that
              K[X,,...,X,]=AI.
    Therefore A^ is ti-smooth over K, and since K is O-smooth over k, we
    see that A^is ti-smooth over k, and therefore A is m-smooth over k. n
       Let k -A       -B       be ring homomorphisms,      and let I be an ideal of B;
    we consider B in the I-adic topology. We say that B is I-smooth over A
    relative to k if the following condition holds: for any A-algebra C, and an
                           ch that N2 = 0, given an A-algebra homomorphism
                           fying u(Z?) = 0 for sufficiently large v, if u has lifting
    v?:B - C as a k-algebra homomorphism,              it also has a lifting v?:B - C




    Theorem 28.4 Let k L       A5     B and I c B be as above; then the following

       (1) B is l-smooth over A relative to k;
                  is a B-module such that Z?N = 0 for sufficiently large v, then
    Der,(B, N) - Der,(A, N) is surjective;
       (3) for every v > 0, the mapR,/, OA(B/ZY) -!&,,@,(B/I?)             has a left
    inverse (that is, it maps injectively onto a direct summand).
    Proof. (1)~(2)IfI?N=0,setC=(B/Z?)*N,andletu:B-B/I?=C/Nbe
    the natural map. Given DEDer,(A, N), define ).:A -C                   by j?(a) =
    (us(a), D(a)). If we consider C as an A-algebra via & then be(u(h),O)eC
    is a k-algebra homomorphism        from B to C lifting U, so that by assumption
    there exists a lifting u?: B -+ C of u as an A-algebra homomorphism;        then
    writing v? in the form
               v?(b) = (u(b), D?(b)), with D?EDer,(B, N),
    we have v?g = ;C,so that D?g = D.
       (2)*(l)    Suppose given a commutative diagram

                       Bu       C/N
with j the natural map, and ~(1?) = 0; if u:B -+C               is a k-algebra
homomorphism        satisfying jo = u and ugf = 1J, then setting D = 1*- ug, we
can view D as an element of Der,(A, N). By assumption there exists
D?EDer,(B, N) such that D = D?g. Using this, we set u? = u + D?; then
           u?g = cg + D?g = 1- D + D = R, and jv? = u.
   (2)0(3) comes from observing the general fact that for a ring R, a
map cp:M +M?          of R-modules has a left inverse if and only if for every
R-module N the induced map
           Horn&W, N) --+ Hom,(M, N)
is surjective. n
Theorem 285. Let A be a ring, B an A-algebra, and I an ideal of B; set
B = B/I, and assume that B is I-smooth over A. Then n,,,@,B         is
projective as a B-module.
Proof. It is enough to show that for an exact sequence L LM        40 of
B-modules, the sequence
          Homg(Qr,, @B, L) -+ Horn&J,,,      @ 8, M) -+ 0
is exact, that is that
          Der,(B, L) - Der,(B, M) + 0
is exact. Set C = B* L and N = Ker cp.If we view both L and N as ideals of
C, we have L2 = N2 = 0 and C/N N B*M. Now for any DtzDer,(B, M)
we have an A-algebra homomorphism        B-C/N      given by
         b t-+(6, D(b))eB*      M,
and lifting this to B -+ C is equivalent      to lifting D to an element of
Der,(B, L). n
Lemma     2. Let B be a ring and I an ideal of B, and let u: L - M be a
map of B-modules; assume that M is projective. Suppose also that one
of the following two conditions hold:
      (a) I is nilpotent;
or (p) L is a finite B-module and I c rad(B).
Then
           u has a left inverseoti:L/IL-  M/IM    has a left inverse.
Proof. (-) is trivial. To prove (G), suppose that V:M/lM       -L/IL  is a
left inverse of U. Since M is projective, there is a map u: M -+ L such
that the diagram
                       ?
             M--+               L

                I -L            1
                             L/IL
         M/IM
commutes. Set w = uu. Then w induces the identity on L/IL, so that
L = w(L) + IL, so that by NAK, L = w(L). Hence if L is a finite B-module,
then by Theorem 2.4, w is also injective. Furthermore, if I? = 0 we do the
following: if xEKer w then 0 = w(x) = xmod IL, so that XEZL, and we
can write x = Caiyi with a,~Z and y,~L. Then
             0 = W(X)   =   C   Ui W(yi)   E   Jf aiyi   = X mod   Z?L,
so that XEZ~L, and proceeding in the same way we arrive at XEI?L = 0.
Hence also in this case w is an automorphism   of L, and w-lo is the
required left inverse of U. n
\end{verbatim}
\end{proof}


\begin{theorem}[Printed Theorem 28.6]\label{mcrt-auto-theorem-28-6}\dcref{28.6}\source{matsumura-commutative-ring-theory:page-0234}\uses{}
\begin{verbatim}
Let k -+ A -B         be ring homomorphisms,     I an ideal of B,
and suppose that B is Z-smooth over k. Set B, = B/Z. Then the following
conditions are equivalent:
   (1) B is Z-smooth over A;
   (2) fi.4,kOABl-Q B,k @Bl has a B,-linear left inverse.
\end{verbatim}
\end{theorem}
\begin{proof}
\begin{verbatim}
(l)*(2)    is contained in Theorem 4. Conversely, suppose that (2)
holds. For any v > 0, set B, = B/Z?; then since Z-smoothness and Iv-smooth-
ness are the same, by Theorem 5, Qsik @ B, is a projective B,-module. Now
set I, = Z/Z?; then BJZ, = B,, and (I,,)? = 0, so that applying Lemma 2, we
see that R,,,OAB,     -QRIk    &BY has a left inverse. By Theorem 4, B is
Z-smooth over A relative to k, but since it is also Z-smooth over k, it is also
Z-smooth over A. w
\end{verbatim}
\end{proof}


\begin{theorem}[Printed Theorem 28.7]\label{mcrt-auto-theorem-28-7}\dcref{28.7}\source{matsumura-commutative-ring-theory:page-0234}\uses{}
\begin{verbatim}
Let (A,m, K) be a Noetherian                           local ring, and kc A a
subfield; then
         A is m-smooth              over ko        A is geometrically     regular over k.
\end{verbatim}
\end{theorem}
\begin{proof}
\begin{verbatim}
(a) Let k? be a finite extension field of k. Then A Okk? = A? is
 mA?-smooth      over k? by base-change. Let n be any maximal ideal of A?;
 then A? is a finite A-module, hence integral over A, so that n 3 mA?. Thus
 if we set A? = AL and m? = nA? then A? --+ A? is continuous for the mA?-
 adic topology of A? and the m?-adic topology of A?, but as a localisation A?
is O-etale over A?, so that by Theorem 1, A? is m?-smooth over k?, and hence
by Lemma 1, A? = A; is a regular local ring. This is what was required to
prove.
(c?) According to Lemma 1, there is only a problem if k is of characteristic
p. The proof we now give was discovered in 1977 by G. Faltings Cl] while
he was still a student.
   By the corollary of the previous theorem, we need only prove
that i&&K      -R,     oaK is injective. For this, we let x1,. . , , x,Ek be
p-independent elements, and prove that dx, , . . . , dx,EQA @ K are linearly
independent over K. Write ai for pth roots of the xi, and set k? = k(cc,, . . . ,a,).
Then
          B = A&k?            = ACT,, . . . , T,]/(T;   -x1,.   . . , T; - x,)
is a Noetherian local ring; write n for its maximal ideal, and L for the residue
field L = B/n. By Theorem 25.2 the sequence


is exact. Similarly,
          O-m/m2 -Q,OaK                     -R,-+O

is exact. Now consider the commutative diagram:
              0 + n/n2 -     R,@,L    -   Q,+O
                     VI                   v2/            ml
                          I
          O-r(m/m2)OKL-RAOAL-SZKOKL-,0
where cpl, (p2 and (p3 are the natural maps. Then by the snake lemma, we
get a long exact sequence of L-modules
            O-+Kerqp --+Kerq2       -Kerq3
              - Coker cp1 - Coker ?p2 - Coker (p3 -+ 0.
By assumption A and B are regular local rings of the same dimension, so
that rank m/m? = dim A = rank n/n2, so rank Kerql and rank CokerqI
are finite and equal; moreover, since L is a finite extension field of K, both
rank Ker (p3 and rank Coker (p3 are finite and equal (the Cartier equality).
Therefore by the above exact sequence, we get
            rank Ker (p2 = rank Coker cp2.
However, Coker (p2 = QBiA C&L, and by Theorem 25.2,
            R ,,,=BdT,+...+BdT,-B?,
so that both of Ker ?p2 and Coker cp2have rank equal to r. Now if we set
J=(T:-x,,...,          T; - x,) we have an exact sequence

          J/J2 ~QA[T,,...,T,I OB=SZ,OBOCBdTi---*S1,-*0,
and since this remains exact after performing CC&L, we see that Ker CPZ
is generated by dx, , . . . , dx,. Therefore dx, , . . . , dx,EQA 0 L are linearly
independent over L, so that they must also be linearly independent over
K as elements of Q, @ K. This is what we had to prove. n
Let B be an A-algebra, I an ideal of B, and consider B with the I-adic
   topology. Let N be a B-module such that I?N = 0 for some v > 0; in what
   follows we will say that a B-module with this property is discrete. An
   A-bilinear map f:B x B -+ N will be called a continuous symmetric
   2-cocycle if it satisfies the three conditions.
      (4 xf(y, 4 - f(xy, 4 + ./Xx, YZ) - ./TX, Y)Z = 0 for all x, y, =B,
      (B) .!I-% Y) = f(Y, 42
      (y) there exists p > v such that f(x, y) = 0 if either x~l? or y~l?.
   If this holds, we set f(1, 1) = r; then substituting y = z = 1 in (a) gives
   xz = f(x, 1).
      Deline a product on the A-module C = (B/I?)@ N by
:?            6, MY, 4 = GE - .0x, Y) + xv + YO
1? for x, DEB; then C is a commutative ring with unit (1, z), and N is an
: ideal of C satisfying N2 = 0. If we define a map A -C             by a~(G,az)
,? then this is a ring homomorphism,       and the diagram
i            B LB/I?        = C/N
              T                 T
                 A-C
    is commutative; then a necessary and sufficient condition for u to have a
    lifting to B --+ C is that there should exist an A-linear map g:B -N
    such that
         (d) f(x, Y) = xgb) - gW + &)y            for all x9 YEB.
    For if g exists, then defining v:B + C by v(x) = &g(x)) we find that v
    is a lifting of U, and conversely, if there is a lifting v of u one checks easily
    that one can find a g as above.
         We say that the 2-cocycle f splits if there exists g satisfying (cr?). For
    any A-linear mapg:B ---+ N, we write 6g for the bilinear map B x B
     -N        given by the right-hand side of (N?); this satisfies (a) and (/I), and
    if g is continuous (that is, 3~ such that g(1*) = 0), then it also satisfies (y).
\end{verbatim}
\end{proof}


\begin{theorem}[Printed Theorem 28.8]\label{mcrt-auto-theorem-28-8}\dcref{28.8}\source{matsumura-commutative-ring-theory:page-0236}\uses{}
\begin{verbatim}
Let A be a ring and B an A-algebra with an I-adic topology.
       (i) If B is Z-smooth over A then every continuous symmetric 2-cocycle
    f:B x B -N        with values in a discrete B-module N splits.
       (ii) If B/I? is projective as an A-module for infinitely many n, and if
    every continuous symmetric 2-cocycle with values in a discrete B-module
    splits, then B is Z-smooth over A.
\end{verbatim}
\end{theorem}
\begin{proof}
\begin{verbatim}
(i) is what we have just said.
      (ii) Suppose that we are given a commutative diagram
             B:CIN
             T T
             A-C
with N2 = 0 and u(P) = 0; then in view of N2 = 0, the C-module N can
be viewed as a C/N -module, and by means of u as a B-module; but then
1?N = 0, SO that N is a discrete B-module. Take an integer n > v such that
B/Z? is projective as an A-module. Then u can be lifted as a map of
A-modules to i:B -+ C such that ,%(I?) = 0. For x, yeB we set
         fk Y) = ~XY) - WW,
and since ,I is a ring homomorphism modulo N, we have f(x, y)~ N. Now
for (EN and XEB, by definition we have n(x).< = x.< (both sides are
products evaluated in C), and using this one computes the left-hand side
of (a) to be zero. The symmetry (p) is obvious. Also n(P) = 0, so that we
also get (y) . Thus f is a continuous symmetric 2-cocycle, and hence by
assumption there is an A-linear map g:B - N satisfying
       S(X>Y) = XdY) - dXY) + S(X)Y.
Now if we set v = i + g, we have
        4XY) = GY) + &Y)
              = 4XMY) + Sk Y) + dXY)
              = &MY) + WdY) + &V(Y)
              = 4+(Y)~
so that u is an A-algebra homomorphism, and is a lifting of U.       n

Theorem28.9 ([Gl 1, 19.7.1). Let (A, m, k) and (B, n, k?) be Noetherian
local rings, and cp:A --+ B a local homomorphism; set B, = B @Ak = B/mB
and n, = n/mB. Then the following conditions are equivalent:
   (1) B is n-smooth over A,
   (2) B is flat over A and B, is n,-smooth over k.
   This is an extremely important theorem, but the proof is long and
difficult, and we refer to [Gl] for it. We content ourselves with proving
the following analogous theorem, which is all that we will need in what
follows.
\end{verbatim}
\end{proof}


\begin{theorem}[Printed Theorem 28.10]\label{mcrt-auto-theorem-28-10}\dcref{28.10}\source{matsumura-commutative-ring-theory:page-0237}\uses{}
\begin{verbatim}
Let (A, m, k) be a local ring, and B a flat A-algebra; suppose
that B, = B@,k is O-smooth over k. Then B is mB-smooth over A.
\end{verbatim}
\end{theorem}
\begin{proof}
\begin{verbatim}
As one seesfrom the definition of mB-smoothness, it is enough to
prove that B/m?B is O-smooth over A/m? for every v > 0. Since B/m?B is
flat over A/m?, we can assumethat m is nilpotent. Then a flat A-module
is free (by Theorem 7.10), so that B is a projective A-module, and hence
by Theorem 8, we need only show that every symmetric 2-cocycle
f:B x B --f N with values in a B-module N splits. First of all, in the case
that N satisfies mN = 0, then since f is A-bilinear f is essentially a 2-
cocycle over B,, that is there is a map fo:B, x B, -N                such that
f(x, Y) = m-c Y).
   Now B, is O-smooth over k, so that by Theorem 8, So splits, that is
there is a map qo:B, -N           such that ,f, = 6g,. Thus setting g(x) = g,,(.X)
     we get
              .t? = &l.
     In the general case, write cp for the natural map N ---+ N/mN, and consider
     cp0.L then this splits, that is there exists @:B - N/mN such that


     Now since B is projective over A, we can lift S to an A-linear map g:B -+ N,
     and then f - 6y is a 2-cocycle with values in mN. Doing the same thing
     once more, we find h:B - mN such that f - S(g + h) is a 2-cocycle with
     values in nt2 N. Proceeding in the same way, since m is nilpotent we finally
     see that ,f splits.  n
\end{verbatim}
\end{proof}


\begin{theorem}[Printed Theorem 29.1]\label{mcrt-auto-theorem-29-1}\dcref{29.1}\source{matsumura-commutative-ring-theory:page-0238}\uses{}
\begin{verbatim}
Let (A, tA, k) be a DVR and K an extension field of k; then
     there exists a DVR (B, tB, K) containing A.
\end{verbatim}
\end{theorem}
\begin{proof}
\begin{verbatim}
Let {x~}~~~ be a transcendence basis of K over k, and set k, =
     k({x,}). We take indeterminates {XL},,, over A in bijection with the
      {XL}, and setA[ {XA^}] = A? and A, = (A?)tAf. Now A? is a free A-module, and
     henceseparatedfor the t-adic topology; therefore, so is A,, and A, is a DVR
     with A,/tA, N k,. Hence, replacing A and k by A, and k,, we can assume
     that K is algebraic over k. Let L be an algebraic closure of the field of
     fractions of A, and let 9 be the set of all pairs (B, cp), where B is an
intermediate ring A c B c L and cp:B --+ K is an A-algebra homomorph-
ism, satisfying the conditions
    (*) B is a DVR and Ker cp= rad (B) is generated by t.
We introduce an order on 9 by defining (B, cp) < (C, $) if B c C and
$1, = cp. Let us show that F has a maximal element. Suppose that

is a totally ordered subset of 8, and set

then one sees easily that Be is a local ring with maximal ideal tB,. If
0 # XEB,, then XEB~ for some i, and since Bi is a DVR we can write
x = t?u with u an unit of Bi and some n. From this we get x#t?+lBo,        so
that B, is t-adically separated. Hence B, is a DVR, and if pPo:B, + K is
defined to be equal to 4oi on Bi then (B,, (p&9. Hence by Zorn?s lemma
P? has a maximal element; suppose that (B, q) is one. Then if p(B) # K, take
~EK - q(B), let f(X) be the minimal polynomial of a over q(B), and take
a manic f(X)sB[X]       which is an inverse image of f(X). Then f(X) is
irreducible in B[X], and hence (by Ex. 9.6) also irreducible over the field
of fractions of B. Let tl be a root of f(X) in L, and set B? = Bra]; then
B? = B[X]/(f).   Therefore
          B?lt B? = BCxl/(t, f) = cp(B)CXIATI = dW4
is a field; since B? is integral over B, every maximal ideal of B? contains
tB?, so that B? is a local integral domain with maximal ideal tB?, and is
Noetherian because it is finite over B, therefore a DVR. This contradicts
the maximality of B, so that we must have p(B) = K.
\end{verbatim}
\end{proof}


\begin{theorem}[Printed Theorem 29.2]\label{mcrt-auto-theorem-29-2}\dcref{29.2}\source{matsumura-commutative-ring-theory:page-0239}\uses{}
\begin{verbatim}
Let (A, m, K) be a complete local ring, (R,pR, k) a p-ring,
and qo:k -K     a field homomorphism;      then there exists a local homo-
morphism p:R --+A which induces q,, on the residue fields.
\end{verbatim}
\end{theorem}
\begin{proof}
\begin{verbatim}
Set S = ZPz, and let kO c k be the prime subfield. Since cp,(k,) c K.
the prime number p, viewed as an element of A, belongs to m. Hence the
standard homomorphism      Z -+A     extends to a local homomorphism
S --+ A. Now R &k, = R/pR = k is a separable extension of k,, and
hence O-smooth over k,; also R is a torsion-free S-module, hence flat over
at by Theorem 28.10, R is pR-smooth       over S.


                         j          fqi:??i
                                    s -LA@+         1
           re, as we discussed at the beginning of $28, we can lift
           /m = K successively to R --+ A/d, and using the fact that A =
           i, we get cp:R --+ A making the left-hand diagram commute.  n

                         te p-ring is uniquely determined up to isomorphism    by

            Suppose that R and R? are both complete p-rings with residue
           ; then by the theorem there exists a local homomorphism
   cp:R -R?     which induces the identity map on the residue field. We have
   R? = q(R) + pR?, and of course cp(p)= p, so that by the completeness of R we
   see that 50 is surjective, and is also injective, since p?R is not contained in
   Ker cp for any IZ. Therefore R N R?. w
      Let (A, m, k) be a complete local ring of unequal characteristic, and let
   p = char k. We say that a subring A, c A is a coefficient ring of A if A,
   is a complete Noetherian local ring with maximal ideal pA, and
             A = A, + m, that is, k = A/m N A,/pA,.
   By Theorem 1 applied to the residue field k of A, there exists a p-ring S
   such that S/pS = k; write R for the completion of S, so that R is a complete
   p-ring with residue field k. By Theorem 2, there exists a local homo-
   morphism cp:R -+ A inducing an isomorphism on the residue fields. If
   we set q(R) = A, then this is clearly a coefficient ring of A. If A has char-
   acteristic 0 then cpis injective and A, N_R. If A has characteristic p? then
   A, N R/p?R. We summarise the above discussion in the following theorem.
                    f (A, m, k) is a complete local ring and p = char k then A
   has a coefficient ring A,. If A has characteristic 0 then A, is a complete

I     In what follows, in order to include the case of equal characteristic in
i our discussion, we also consider coefficient fields as being coefficient rings.
[? From the previous theorem and Theorem 28.3, we get the following
i important result.
\end{verbatim}
\end{proof}


\begin{theorem}[Printed Theorem 29.4]\label{mcrt-auto-theorem-29-4}\dcref{29.4}\source{matsumura-commutative-ring-theory:page-0240}\uses{}
\begin{verbatim}
(i) If (A, m) is a complete local ring and m is finitely
   generated, then A is Noetherian.
      (ii) A Noetherian complete local ring is a quotient of a regular local
   ring; in particular it is universally catenary.
      (iii) If A is a Noetherian complete local ring (and in the case of unequal
characteristic, A is an integral domain), then there exists a subring A? c A
with the following properties: A? is a complete regular local ring with the
same residue field as A, and A is finitely generated as an A?-module.
\end{verbatim}
\end{theorem}
\begin{proof}
\begin{verbatim}
We choose a coefficient ring A, of A. If m = (x1,. . ,x,) then every
element of A can be expanded as a power series in x1,. . . , x, with coefficients
in A,, so that A is a quotient of A,I[X,,. . .,X,,J and hence Noetherian.
Now A, is a quotient of a p-ring R, so that A is a quotient of R[X, , . . . , XJ,
which is a regular local ring, hence a CM ring, and therefore according
to Theorem 17.9, A is universally catenary. To prove (iii), set dim A = II,
and in the case of equal characteristic let (y,, . . . , y,} be any system of
parameters of A; if A is an integral domain of characteristic 0 and
char k = p, we can choose a system of parameters {y, = p, y2,. . . , y,) of
A starting with p. In either case R = A,, so that we set A? = R[y]. Then
A? is the image of (p:R[Y] -A,        where R[Y] is the regular local ring
          R[Yj = R[IY, ,..., Y,], or R[Y2 ,..., Y,,] if y, =p,
and cp is the R-algebra homomorphism        defined by cp(Yi) = yi. Set m? =
c; y,A?. Since A/m = A?J m?, every A-module of finite length has the same
length when viewed as an A?-module. In particular A/m?A is a finite module
over A?/,? and A is m?-adically separated, so that by Theorem 8.4 A is a finite
A?-module. Therefore, we have
          dim A? = dim A = n.
R[Y] is an n-dimensional integral domain, and if Ker cp# 0 we would have
dim A? < n, which is a contradiction.        Therefore cp is injective and
A? _NR[Y].
\end{verbatim}
\end{proof}


\begin{theorem}[Printed Theorem 29.5]\label{mcrt-auto-theorem-29-5}\dcref{29.5}\source{matsumura-commutative-ring-theory:page-0242}\uses{}
\begin{verbatim}
Let (A, m, K) be a local ring, and suppose char K = p. Let
   Cc A be a subring, and assume that C is a Noetherian local ring with
; maximal ideal PC, and that K = A/m is separable over C/PC. Then there
i exists a quasi-coefficient ring S of A containing C; moreover, if A is flat
1 over C, then it is also flat over S.
i Proof. Let {Bn}ncl\ be a p-basis of K over C/PC, and choose an inverse
e image blgA for each PA. Setting C[{b,}] = c?, by Theorem 26.8 we see
i? that C?/(mnC?) = (C/pC)[{/?,>]         is a polynomial ring over C/PC. Hence
; iff(...X,...)?     is a non-zero polynomial with coefficients in C which satisfies
!; f(b,)Em, then setting p? for the highest common factor of the coefficients
p off, we have
;<
kt              f(X) = ~?.fdW      and Jb&) Z 0.
i:
& Thus f,(b,)$m, and we must have r > 0, so that we have shown that
t mnC?=pc?.       Setting S=(C),.       we have ScA, mnS=pS          and SIPS=
                    rice c? is p-adically separated, so is S, and hence all the
                      the form (0) or (p?). Thus S is Noetherian, and it satisfies
  all the conditions for a quasi-coefficient ring of A. If A is flat over C, then

              p?C BcA N p?A.

     and hence the composite
             p?C@A     = (p?C@,S)@,A          -p?S@A     -p?A

     is injective; but the first arrow is surjective, so that the second arrow
     p?SO A -p?A       is injective. By Theorem 7.7, this proves that A is flat

        All the assumptions of this theorem are satisfied by taking C to be the
     image of Z,,, +     A, so that this proves that every local ring has a
     Quasi-coefficient ring (including the quasi-coefficient field of a local ring
     in the equal characteristic case).
\end{verbatim}
\end{theorem}
\begin{proof}
\notready
\end{proof}


\begin{theorem}[Printed Theorem 29.6]\label{mcrt-auto-theorem-29-6}\dcref{29.6}\source{matsumura-commutative-ring-theory:page-0242}\uses{}
\begin{verbatim}
Let (A, m, K) be a local ring, and A^ its completion, and
     a?PPose char K = p. Let S be a quasi-coefficient ring of A, and write S
     for its image in A^; then there exists a unique coefficient ring A, of A^

             lnce s? is a quasi-coefficient    ring of A, we can assume that A is
a complete local ring. If A has characteristic 0 then S is a DVR, and by
Theorem 1 there exists a complete p-ring R containing S and with residue
field K. Now K is 0-etale over S/p& and R is flat over S, so that R is
pR-etale over S (by Ex. 28.1). Hence there exists a unique S-algebra homo-
morphism R -A           which induces the identity map on the residue
fields; we write A,, for the image of R. Then R 2: A, and A, is a coefficient
ring of A. If A had another coefficient ring B containing S then B would
also be a complete p-ring, and for the same reason as above, there exist
unique S-algebra homomorphisms        B --+ A, and B -A,      so that we must
have B = A,.
    If A has characteristic p? then S is an Artinian ring, and therefore
complete, so that applying Theorem 3 to S itself, we can write S = R&J?)
with R, a complete p-ring. Let R be a complete p-ring containing R, and
with residue field K; then by Ex. 28.1, R is pR-etale over R,, so that there
exists a unique R,-algebra homomorphism        R -A     inducing the identity
map on the residue field K. The image is a coefficient ring of A
containing S; the uniqueness is proved as in the case of characteristic 0. n
    Next we study the structure of complete regular local rings. Let (A, m, K)
be a local ring of unequal characteristic, and suppose that char K = p;
then A is said to be ram$ed if przm2 and unramified if p$m?. We will
also say that A is unramitied in the case of equal characteristic.
\end{verbatim}
\end{theorem}
\begin{proof}
\notready
\end{proof}


\begin{theorem}[Printed Theorem 29.7]\label{mcrt-auto-theorem-29-7}\dcref{29.7}\source{matsumura-commutative-ring-theory:page-0243}\uses{}
\begin{verbatim}
An unramified complete regular local ring is a formal power
series ring over a field or over a complete p-ring.
\end{verbatim}
\end{theorem}
\begin{proof}
\begin{verbatim}
Let R be a coefficient ring of A. In the case of equal characteristic,
R is a field, and if x1,. . . ,x, is a regular system of parameters of A then
A=R[.x,    ,..., x,j-RIXI      ,..., X,] (see the proof of Theorem 4). In the
case of unequal characteristic, R is a complete p-ring, and since pern - m2,
we can choose a regular system of parameters {p,x2,. . .,x,) of A
containing p. Then A = R[xz, . . . , x,l] ?v R[X2,. . . , X,].  n
   In the ramified case, it is not necessarily the case that A can be expressed
as a formal power series ring over a DVR. To give the structure theorem
in this case we need the notion of an Eisenstein extension.
Lemma 1 (Eisenstein?s     irreducibility criterion). Let A be a ring, and
f(X) = X? + a,X?-? + .*. + a,, with u,EA. If there exists a prime ideal P
of A such that a,,..., a,,~p but a,$p2 then f is irreducible in A[X]. If
in addition A is an integrally closed domain then the principal ideal (f)
is a prime ideal of A[X].
Proof. If f is reducible, we can write f =(X?+ b,X?-? + ... + 6,)
(X? -t clXs-? + ... + c,) with 0 < r < n, s = n - r and b,, cjeA. Reducing
the coefficients     on either side modulo     p we have
            x?=(x?+6,x~-?+~~~+&,)(x~+clx?-?+~~~+c,)
  in (A/p)[X],  so that we must have bi, cj~p for all i,j, but then a,, = b,c,q?,
  which contradicts the assumption. If A is an integrally closed domain and
  K is the field of fractions of A, then by Ex. 9.6, f remains irreducible in
  K[X]. Also, f is manic, so that we have f.A[X]        = f.K[X]   nA[X],      and
  this is a prime ideal of A[X].    n
     Let (A,m) be a normal local ring; then an extension ring
            B = 4x1/m         = ax1
  defined by an Eisenstein polynomial
           f=X?+a,X?-?+...+a,          with a+nt for all i, and an$m2
  is called an Eisenstein extension of A. By the lemma, B is an integral
  domain, and is integral over A. We have B/mB = (A/m)[X]/(X?),           so that
  B has just one maximal ideal n = mB + xB. Hence B is a local ring, and
  its residue field coincides with that of A.
\end{verbatim}
\end{proof}


\begin{theorem}[Printed Theorem 29.8]\label{mcrt-auto-theorem-29-8}\dcref{29.8}\source{matsumura-commutative-ring-theory:page-0244}\uses{}
\begin{verbatim}
(i) If (A, m) is a regular local ring, then an Eisenstein
  extension of A is again a regular local ring.
     (ii) If A is a ramified complete regular local ring and R is a coefficient
  ring of A then there is a subring A, c A with the following properties:
     (1) A, is an unramified complete regular local ring containing R, and
  hence can be expressed as a formal power series ring over R;
     (2) A is an Eisenstein extension of A,.
\end{verbatim}
\end{theorem}
\begin{proof}
\begin{verbatim}
(i) Let B = A[x], and x? + alxn-? + ... + a,, = 0, with aiEm and
: a&m?. Then there exists a regular system of parameters {y,, . . . , y, = a,}
, of A with a, as an element. As we have seen above, the maximal ideal of
  B is mB + xB, but a,ExB, so that {y,,...,y,-,,x}            is a regular system
  of parameters of B.
     (ii) Since htpA = 1, by a skilful choice of a regular system of parameters
  {x1 ,..., xd) of A, we can arrange        that (p,xz,. ..,x,}   is a system of
  Parameters of A. If we set A, = R[x,, . . . , XJ then A, is a complete
  unramified regular local ring, and A is a finite module over A, (see the
  proof of Theorem 4). We set m, for the maximal ideal of A,. Now
  A = m,A+ AJx,],          so that by Theorem 8.4 (or by NAK), A = A,,[x,].
  Let
           f(X)      = X? + a,X?-?    + ... + a,, witha,EA,
  be the minimal polynomial           of x1 over A,. Then a,ExlA
                                                             cm, so that
  a&m,,. Therefore by Hensel?s lemma (Theorem 8.3), all the aiEm,. We
I are left to prove that a,,$mi. Write p = ct bixi with bigA, and express
the bi in the form b, = qi(xl), with qi(X)~A,[X];                             then xi is a root of
          F(X)   =   (PICxlx        +   Cd2   ?Pilxlxi         -     P3

so that F(X) is divisible by f(X). Hence the constant term F(0) of F is
divisible by a,. However, F(0) = ci qi(0)xi - p, and p, x2,. . . , xd is a regular
system of parameters, so that F(O)@t$, hence also a,,$nti.           n
\end{verbatim}
\end{proof}


\begin{theorem}[Printed Theorem 30.1]\label{mcrt-auto-theorem-30-1}\dcref{30.1}\source{matsumura-commutative-ring-theory:page-0245}\uses{}
\begin{verbatim}
Nagata-Zariski-Lipman).          Let (A, m) be a complete Noether-
 ian local ring with QccA. Suppose that xi,...,x,Em               and D1,...,DIe
 Der(A) are elements satisfying det(Dixj)$m. Then
    (i) There is a subring C c A such that
            A = qx,, . . .) xr] N qx,, . .,X,]
Therefore x i , . . ,x, are analytically independent over C, and A is Z-smooth
over C, where I = 1; Axi, and therefore also m-smooth over C.
   (ii) If g = C; ADi .is a Lie algebra, (that is if [Oi, Dj]~g for all i,j) then we
can take C to be {a~AjD,a = ... = D,a = O}.
\end{verbatim}
\end{theorem}
\begin{proof}
\begin{verbatim}
Letting (cij) be the inverse matrix of (Dixj), and setting 0: = CCijDj,
we have D;xj = 6,,, so that we can assume that Dixj = 6ij. Quite generally,
for an element tEm and a derivation DEDer (A), we define a map
E(D, t): A + A by

          w4 t) = f.           p;

by our assumptions, E(D, t)(u) = ~(t?/n!)D?(u) is meaningful, and one sees
easily that E(D, t) is a ring homomorphism.   Now set
         E, = EP,, -xi)       and C, = Im(E,);
then C, is a subring of A, and by computation we see that

          4($(y)&)
            +w-lD?                                  +f(-Xl)nD"+l=O
              1 (n- l)!                         ?        0           n!   1      ?
so that C, c {UEA (D,u = 0). Conversely,                                  if D,u = 0 then E,(u)=a,
bollary.       Let (A, m) be a reduced n-dimensional local ring containing Q,
&nd suppose that the completion A^ of A is also reduced. If there exist
 Ielements D r ,..., D,EDer(A) and x1 ,..., x,Em such that det(Dixj)$m
 a?
ithen A is a regular local ring and x 1,. . . , x, is a regular system of parameters
?Of      A. Suppose in addition that g = x;ADi              is a Lie algebra; then
 ?k= (aeAID,a=...       = D,a = O> is a coefficient field of A^.
  Proof.    Consider 2, with each of the Di extended to A. If [Di, Dj] =
 Fxaij,D, holds in A then it also holds in 2, so that if g is a Lie algebra,
  so
   4 is CAD,,. By the theorem, A^= C[x,,. .., x& and C is isomorphic to
  ?/xxiAI,    SO is a zero-dimensional   local ring. Now by assumption C is
 also reduced, so that C must be a field. Therefore A^ is a formal power
,$?eriesring over a field, and hence is regular, so that A is also regular. The
rother assertion is also clear. w
\end{verbatim}
\end{proof}


\begin{theorem}[Printed Theorem 30.2]\label{mcrt-auto-theorem-30-2}\dcref{30.2}\source{matsumura-commutative-ring-theory:page-0247}\uses{}
\begin{verbatim}
Let k be a field, and A = k[x l,. . . , x,] a finitely generated
ring over k. If A, is O-smooth over k for every pgm-SpecA, then A is o-
smooth over k.
\end{verbatim}
\end{theorem}
\begin{proof}
\begin{verbatim}
Write k[X] for k[X,, . . , X,1, and let I= {f(-V~kCXIIf(x)=o},
so that A = k[X]/Z. Suppose that I = (fi , . . . .f,). Consider a commutative
diagram
         A 2            C/N
         T               TV
         k--+            C,
where C is a ring, and N c C is an ideal satisfying N2 = 0. TO lift $ to
 A -C,            we first of all choose U,EC such that I     = Cp(Ui). If f~1 then
f(u  i,.   . . , u,)EN.   Now     if we can choose ~,GN   for  1 d i d n such that
fj(u + y) = 0 for all j, th e h omomorphism        A --+ C defined by Xi++rli + Yi
 is a lifting of $. We have fj(u + y) = A(U) + x1= r (afj/dXi)(U).yi, SO that we
 are looking for solutions in N of the system of linear equations in y,, . . , yn:

         (*)     fj(u) + i$l
                         2      (u).y, = 0 for j = l,...,s.
                       (     .>
For each maximal ideal p, the local ring A, is O-smooth over k, so that if we
set S= $(A - p) and S = cp-?(S) then in the diagram
         A, A            WV,   = GIN,
          T                       T
          k     -                cs
there exists a YP:A, + Cs lifting $,. From this we see that (*), as a system
of equations in N,, has a solution in N,. If we view N as a C/N-module then
Ns = N,. Thus the theorem reduces to the following lemma.
Lemma I. Let A and B be rings, $ :A --+ B a ring homomorphism,    and N
a B-module; suppose that b,+B and &EN. If the system of linear equations

         fl     bij Yj = pi    (for i = 1,. . . , s)
                    .
has a solution m N,, + for every pEm-Spec A, then it has a solution in N.
Proof. The assumption that there is a solution in NIL(A-p) means
that there exist qjp~:N for 1 <j < n, and t,EA - p such that
         C bijrl,, - $(t,)B; = 0 for 1 < i < S.
          j
Now since ?&tPA = A, there is a finite set {pr , . . . , p,} c m-Spec A and a,EA
such that
          i    a&   = 1.
        V=l
Hence if we set



we get Cjbijqj = pi for 1 3 i 2 s.    n
\end{verbatim}
\end{proof}


\begin{theorem}[Printed Theorem 30.3]\label{mcrt-auto-theorem-30-3}\dcref{30.3}\source{matsumura-commutative-ring-theory:page-0248}\uses{}
\begin{verbatim}
Let k be a field, S = k[X,, . . . ,X,], and I, P ideals of S such
that I c PESpecS. Set
           S, = R, rad(R) = PR = M, and R/IR = A, rad(A) = m,
           R/M=A/m=          K,
and suppose that ht ZR = r and I = (jr(X), . . . , f,(X)). Then the following
conditions are equivalent.
   (1) A
   (2) rank
          is @(fl,.
             O-smooth. . , LYf(kxl ). . . yXJmod P) = r;
                                  2
   (3) A is m-smooth over k;
   (4) RAjk is a free A-module of rank II - r;
   (5) A is an integral domain, its field of fractions is separable over k, and
a,,, is a free A-module.
\end{verbatim}
\end{theorem}
\begin{proof}
\begin{verbatim}
Note that R is a regular local ring. (1) a(2) By assumption, for a
suitable choice of I elements D,, . . . ,D, from a/8X,, . ..,a/aX,, and of r
elements gl,. . . , gr from jr,. . . , f,, we have det (Digj)$ M. We observe
that taking REM into (Or fmod M, . . . , D,f mod M)EK? induces a linear
map M/M2 --+ K?, so that the images of gr, . . . , g1 in M/M2 are linearly
independent. Therefore x;giR is a height r prime ideal contained in IR,
and hence c; giR = IR. Given a commutative diagram




         k-C
with N2 =O, write X,EA for the image of Xi, and choose U,EC
such that cp(u,) = $(x~)EC/N. Then there is a homomorphism               R +C
defined by Xi I-+z+, and this induces a lifting of $ to C if and only if gi(u) = 0
 for 1 6 i < r. Therefore, as in the proof of the previous theorem, we need
 only solve the system of equations in unknowns y,, . . . , y,,gN:

         (*) gi(U)+ j$l ( g .I> (?)?Yj=o        for   l<i<r.

However, we view N as a C/N-module,           then via $ as an A-module, so
that we can replace the coefficients (ag,/aX,)(u)        of this system by
(&J~/c?x~)(x)EA. Now by assumption, one Y x Y minor of this r x n matrix
of coefficients is a unit of A, so that (*) can always be solved.
   (2)=>(3) is trivial.
   (3)*(l)    By $28, Lemma 1, A is regular, so that IR is generated by ,.
elements forming a subset of a regular system of parameters of R, and
the image of the natural map IR --f M - M/M2 is an r-dimensional
K-vector space. Let k, c k be the prime subfield; then by Theorem 26.9
K = R/M is O-smooth over k,, so that by Theorem 25.2, the sequence
         O-+M/M?--tC&@K-Cl,+0
is exact. We can write 0, = (a, @ S)@ F, where F is the free S-module with
 basis dX, , . . , dX, (for example, by Theorem 25. l), and localising we get


hence Q2, @ K = (Q, &K)      @ (F @ K ). However,    from


we get the exact sequence (I/I?)@,K            -R,&K         -Q,@,K-+O,
and if A is m-smooth over k then by the corollary to Theorem 28.6,
Qk@ K -R,@           K is injective, so that V = Im((l/l?)@K      -Q2,@K}
maps isomorphically      to its projection W c F @ K in the second factor
of the decomposition     Q, 0 K = (Q, 0 K) @ (F 0 K). Now factor I/1? 0 K
 -OR        0 K as the composite I/I20 K - M/M2 -Cl,@               K; as we
have seen above, the first arrow has rank r, and the second is injective, SO
that rankV=r.         Now F@K=KdX,+...+KdX,,                 and if we write
(afi/aXj)     modP =clij then W is spanned by ~~ai,dXj           for 1 6 i < t.
Therefore rank(aij) = T, and this proves (1).
    (2)*(j)    By Theorem 25.2,
            O+IR/12R     -+QRikOA     - fiAik-+O
is a split exact sequence, and since QRlk 0x A is a free A-module with basis
dX 1,. . . , dX,, the A-module RAjk is projective; but A is a local ring, SO that
QAjk is free. Also A is a regular local ring, therefore an integral domain,
and if L is its field of fractions then L is O-smooth over A, hence also
O-smooth over k, so that L is separable over k.
   (j)*(4)     By Theorem 26.2, the field of fractions L of A is separably
generated over k, so that a separating transcendence basis of L over k is
a differential basis, and ranka,,, = tr.deg,l = n - r; but C& = Qa,k @A~*
and hence rank.Q,,,        = rank,&,,  = n - r.
   (4)+(l)     In the exact sequence
         IR/12R -R,,,@        A -QAlk+O,
set E = Im{IR/l*R    -C&,k  0 A). Then R,,k @A N E @A?-?,               so that
E N A?, and therefore E 0 K N K?. This gives (1). H
Remark 1. In the above proof, the equivalence          of (l), (2), (4) and (5)
was comparatively easy. The proof of (3)+(l)           used the corollary to
Theorem 28.6, and so is not very elementary.
Remark 2. If k is a perfect field, or more generally if the residue field K
is separable over k, then m-smoothness is equivalent to A being a regular
local ring, so that Theorem 3 gives a criterion for regularity. In the case
of an imperfect field k, if A is O-smooth over k then so is the field of
fractions L of A, but the residue field K is not necessarily separable over
k; for example, A = k[X](,,_.,     with aEk - kp. For the case of an imperfect
field k, we give a regularity criterion for A in Theorem 5.
    Quite generally, let A be a ring and P a prime ideal of A, and let
D 1,..., D,EDer(A)andf,      ,... &A;thenwewriteJ(f,     ,..., ft;D, ,..., D,)(P)
 = (D&mod P). This is an s x t matrix with entries in the integral domain
A/P.
\end{verbatim}
\end{proof}


\begin{theorem}[Printed Theorem 30.4]\label{mcrt-auto-theorem-30-4}\dcref{30.4}\source{matsumura-commutative-ring-theory:page-0250}\uses{}
\begin{verbatim}
Let R be a regular ring, PESpec R, and let I c P be an
ideal of R; suppose that ht IR, = r.
   (i) for any D1,..., D,EDer(R) and fi,. . .~,EI we have
           rank J(fl,. . . , ft; D,, . . . , D,)(P) d r;
   (ii) if D 1,..., D,EDer(R) and fi ,... f$Z are such that det(Difj)$P
then IR, = (fl,. . . , f,)Rp and R,/ZR, is regular.
\end{verbatim}
\end{theorem}
\begin{proof}
\begin{verbatim}
(i) If Q is a prime ideal of R with I c Q c P and ht Q = r then
           rankJ(f,, . . . ; D,, . . .)(P) < rankJ(f,, . . .; D,, . . .)(Q),
and if we set QRe = nr then R, is an r-dimensional regular local ring, so
that m is generated by r elements, m = (gl,. . . , g?). Working in R,, we can
Write
            f j = 2 g,,a,j        with      cc,j~R,
                    1
SO   that
            Difi~        i    (Digy).Clyj    modQ,
                        v=1

and therefore

               <rankJ(g, ,..., gl;D, ,..., D,)(Q)<r.
   (ii) Set M = PR,; then if det(Difj)$M        one sees easily that the images
inM/M*off,,...,f        I are linearly independent over R,/M = K(M), SO tha?t
c;f iR, is a prime ideal of height r, and therefore coincides with IR,.
Also, RJIR, is regular.        n
\end{verbatim}
\end{proof}


\begin{theorem}[Printed Theorem 30.5]\label{mcrt-auto-theorem-30-5}\dcref{30.5}\source{matsumura-commutative-ring-theory:page-0251}\uses{}
\begin{verbatim}
Zariski). Let k be a field of characteristic p, S = k[X,, , . . ,X,1,
and I and P ideals of S with I c PESpec S. Set Sp = R and RJIR = A, and
suppose that ht IR = r and I = (fr,. . . ,f,). Then the following three
conditions are equivalent.
   (1) A is a regular local ring.
   (2) For any p-basis (u,},,,- of k, define D,EDer(S) by D,(u,.) = 6,,. (the
Kronecker 6) and II,                = 0; then there are a finite number of elements
a, P, * *. 7YET such that
            rank .J(fi, . . . , ft; D,, D,, . . . , D,, 8/8X,, . . . ,8/8X,)(P) = r.
   (3) There exists a subfield k? c k with the following properties: I@?C k,
[k: k?] < co, and QA,k, is a free A-module, with
            rank OA,k. = n - r + rank a,,,. .
\end{verbatim}
\end{theorem}
\begin{proof}
\begin{verbatim}
(2)+(l) comes from Theorem 4, (ii).
   (1) a(2), (3) If A is regular then IR is generated by r elements, and these
form an R-sequence, so that IR/Z?R is a free A-module of rank r. Set
M = PR, m = M/IR and K = R/M = A/m; then the image of IR --+ M/M2
is an r-dimensional          K-vector space, so that the natural map
(IR/I?R)@, K -M/M?             is injective. From the exact sequence
         IR/l?R-R,Q,A-Cl,+0
we get the exact sequence
  (*)    (IR/I?R)O,K      -S&&K          -C&,@AK+O.
Now by Theorem 25.2,
         O+MfM2-42,@,K-Cl,+0
is an exact sequence. The first arrow in (*) is the composite (ZR/I?R) --+
M/M2 -QR       OK, and so is injective. Thus
  (**)   O+(ZR/12R)@K         -f&&K           -R,&K-,O
is exact. Let {du,l yEI-} be a basis of f& over k; then 0, is the free
S-module with basis {du,JyEr} u {dX,, . . . , dX,}, and Sz, = C&&R,
C&@&K =0,&K.            Now reorder fI ,..., f, so that fI ,..., fi are gen-
erators of IR; then df,, . . . , df,Ef&&K           can be expressed using
dX,,. . . ,dX, together with finitely many elements du,, . . .,duy, and this
gives (2). Now if we let k? be the field obtained by adjoining to kP all the
elements u, for CTET other than the du,, . . . , du, just used, then from (**)
we see that
           0 + (IR/I?R) @ K - it,,,, @ K - R,,,, Q K --+0
is also exact. In the exact sequence
   (t)   IR/12R -i&k,@          A -%,,,-*O,
the middle term is the free A-module with basis du,, . . , du,, dXr,. . . ,dXn;
now the generators .f?i,. . . ,,f, of IR map to df,, . . . ,df,.&,,,. Q A, and
since these are linearly independent over K, they base a direct summand
of R R,k, 0 A. Thus RAIk, is also a free A-module, and rank SZAIkP+ r =
rank QZRlk,= rank !& + n.
   (3)*(l) In the exact sequence (T) the two terms SZRlk,0 A and R,,,. are
both free modules, and the difference between their ranks is r, so that
IR/I?R maps onto a direct summand of GRjk. @ A, which is a free module of
rank Y. Thus we can choose fi,. . . ,f,eIR such that df, 0 1,. . . , df,@ 1
base this direct summand, and then by NAK we see that Rdf, + . .. +
Rdf, is a direct summand of !Z&,. Hence there exist D,, . . , D,EDerk.(R)
such that det(Difj)$M.    Therefore, by Theorem 4, A is regular. n
\end{verbatim}
\end{proof}


\begin{theorem}[Printed Theorem 30.6]\label{mcrt-auto-theorem-30-6}\dcref{30.6}\source{matsumura-commutative-ring-theory:page-0252}\uses{}
\begin{verbatim}
M. Nomura). Let (R, m) be an equicharacteristic n-dimen-
sional regular local ring, and R* the completion of R; suppose that k is a
quasi-coefficient field of R and K a coefficient field of R* containing k.
Let x 1,. . .,x, be a regular system of parameters of R.
   (i) R*=K[xl,...,         x,J, and if we write a/ax, for the partial derivatives
in this representation,         then Der,(R*) = Der,(R*)     is the free R*-module
with basis i3/dxi for 1 d i d n.
   (ii) The following conditions are equivalent:
   (1) a/ax, maps R into R for 1 ,< i d n, so that they can be considered
as elements of Der,(R);
   (2) there exist D,,...,D,EDer,(R)                and a,,...,a,~R      such that
Diaj = S,,;
   (3) there exist D, ,..., D,ED~~,(R)              and a, ,..., a,,eR such that
det (Diaj)$m;
   (4) Der,(R) is a free R-module of rank n;
   (5) rank Der,(R) = n.
\end{verbatim}
\end{theorem}
\begin{proof}
\begin{verbatim}
(i) Since K is 0-etale over k, any derivation of R* which vanishes
on k also vanishes on K. If DEDer,(R*),                 set Dx, = yi; then for any
PER*=          K[x~,...,x,JJ       we have D(f) = EYE l(af/axi).yi,      and hence
D = xy,d/dx,. Conversely, for any given yi we can construct a derivation
by this formula, SO that Der,(R*)              is the free R*-module     with basis
alax l,...,a/ax,.
   (ii) (l)=>(2)*(3)     and (4)*(5) are trivial. If (3) holds then D,, . . , D, are
linearly independent over R and over R*, so that by (i), any DEDer,(R)
can be written as a combination D = CciDi of the Di with coefficients in
the field of fractions of R*, but since Daj = cciDi(aj)         for j = 1,. . . , n, we
have c,ER; therefore D,, . . . , D, form a basis of Der,(R), which proves (4).
   (5)*(l)     If D1,..., D, are linearly independent over R then there exist
a,,..., a,ER such that det(Diaj) # 0. Therefore D,, . . . , D, are also linearly
independent over R*. Thus writing L! for the field of fractions of R* we
can write a/ax, = ccijDj, with cijeL!. From this we get 6i, = CCijDjXh,
so that the matrix (cij) is the inverse of(Djx,), and cij~L, where L is the field
of fractions of R; therefore (a/ax,)(R) c Ln R* = R. n

Lemma 2. Let R be a regular ring and PESpec R with ht P = r; then the
following two conditions are equivalent:
    (1) there   exist D1,..., D,EDer(R)     and fl,. .,fpP       such that
 det(Di.fj)$P;
   (2) for all QESpec R with htQ = s d r such that Q c P and Rp/QRp is
regular, there exist D, ,..., D,eDer(R)       and g1 ,..., g,EQ such that
det(Digj)$P.
Proof. (2) contains (1) as the special case P = Q; conversely, suppose that
(1) holds. By Theorem 4 we see that fl,. . . ,f, is a regular system of
parameters of R,. If Rp/QRp is regular then we can take gl,. . .,gseQ
forming a minimal basis of QRp, and then take gs+ i, . . . , g,EP such that
gl,. . . , g1 is a regular system of parameters of R,. Then from det(Difj)#P
 we deduce that det(Digj)$P,               or in other words rank J(g,,. . . , gl;
 D i,. . . , D,)(P) = r. Therefore
              rankJ(g, ,..., gs;D, ,..., D,)(P)=s.         n

    If the above condition (1) holds, we say that the weak Jacobian condition
 (WJ) holds at P. If (WJ) holds at every PESpec R then we say that (WJ)
 holds in R. If this holds, then for any P, QESpec R with Q c P, setting
 htQ = s we have
                                        there exist D,, . . . , D,EDer(R) and
              Rp/QRp is regular-
                                        f 1,. . . , f,EQ such that det(Di fj)$P,
(the implication (=z=)is given by Theorem 4). This statement is the Jacobian
criterion for regularity. We can use Der,(R) in place of Der(R), and we
then write (WJ),.
\end{verbatim}
\end{proof}


\begin{theorem}[Printed Theorem 30.7]\label{mcrt-auto-theorem-30-7}\dcref{30.7}\source{matsumura-commutative-ring-theory:page-0254}\uses{}
\begin{verbatim}
Let (A, m) be an n-dimensional Noetherian local integral
 domain containing CC.&    and let k c A be a subfield such that tr.deg,(A/nt) =
 r < co. Then Der,(A) is isomorphic to a submodule of A?+?, and is therefore
 a finite A-module, with
            rank Der,(A) < dim A + tr. deg,(A/m).
\end{verbatim}
\end{theorem}
\begin{proof}
\begin{verbatim}
We write A* for the completion of A; choose a quasi-coefficient
field k? of A containing k, and let K be a coefficient field of A* containing
k?. Let ui,... , u, be a transcendence basis of k? over k, and xi,. . . , x, a
 system of parameters of A. We define q:Der,(A) - A?+? by q(D) =
(Du 1,..., Du,,Dxl ,..., Dx,); then cpis A-linear, so that we need only prove
that it is injective. Suppose then that DUi = Dxi = 0 for all i and j; then D
 has a continuous           extension     to A*, and this vanishes on
 B=K[x,,...       ,x,1. Now  we do  not know  whether A* is an integral domain,
 but it is finite as a B-module, so that any UEA is integral over B, and if
 f(X)eB[X]       has a as a root, and has minimal degree, then f(u) = 0,
f?(a) #O. Then 0 = D(f(u)) = f?(u)*Du,          and DaEA, so that, since a
 non-zero element of A cannot be a zero-divisor in A*, we must have
 Da=O. Hence D=O.            n
\end{verbatim}
\end{proof}


\begin{theorem}[Printed Theorem 30.8]\label{mcrt-auto-theorem-30-8}\dcref{30.8}\source{matsumura-commutative-ring-theory:page-0254}\uses{}
\begin{verbatim}
Let (R, m) be a regular local ring containing          U& and k a
quasi-coefficient field of R. Then the following three conditions are
equivalent:
    (1) (WJ), holds at m;
    (2) rank Der,(R) = dim R;
    (3) (WJ), holds at every PESpec R.
Furthermore, if these conditions hold, then for any PESpecR, every
element of Der,(R/P) is induced by an element of Der,(R) and
           rank Der,(R/P) = dim R/P.
\end{verbatim}
\end{theorem}
\begin{proof}
\begin{verbatim}
(l+(2)      is known from Theorem 6, and (1) is contained in (3).
    (l)+(3) Write R* for the completion of R, and let K be the coefficient
field of R* containing k; then by Theorem 6, if x1,. . . , x, is a regular system
of parameters of R then the derivations 8/8x, of R* = K [x1,. . . , xn] belong
to Der,(R) for 1 < i < n, and form a basis of it. Now let PESpec R, and write
qp:R -R/P       for the natural homomorphism;       to say that D?EDer,(R/P)
is induced by DEDer,(R) means that there is a commutative diagram




Suppose then that D? is given; then D?ocpEDer,(R, R/P), and this has
a unique extension to an element of Der,(R*, R*/PR*),                             so that
it is uniquely determined by its values on x1,. . . , x,. Therefore if we choose
b r,. . . , b,ER such that D?(cp(x,)) = cp(b,), and set D = xb,a/ax,, then D? is
induced by D.
    Now Der,(R, R/P) is a free R/P-module with basis cp08/axi for 1 < i < II,
and Der,(R/P) can be identified with the submodule
              N = {sEDer,(R,R/P)(S(f)           = 0 for all f~Pj.
Therefore, if P = (fI, . , . ,ft) and htP = Y then
              rank Der,(R/P) = n - rank J(fr,. . . , ,f,;a/ax,, . . . , dpx,)(P);
according to Theorem 4, the right-hand side is 3 n - Y, and by Theorem 7
the left-hand side is <dim R/P = n - r, so we see that
              rank J(f, , . . ,ft; a/&,, . . . , iT@x,)(P) = r,
              rank Der,(R/P) = dim R/P. w
   This theorem has many applications. For example, in the ring R of
convergent power series in n variables over k = Iw or @ (R is denoted by
k((X, ,..., X,> in [Nl],              and elsewhere by k{X, ,..., X,}) we have
aXi/aXj = 6,,, so that the Jacobian criterion for regularity holds in R. In
characteristic 0, the assumptions of the theorem are satisfied by the formal
power series ring over any field, or more generally by the formal power
series ring R [ Y, , . . . , Y, 1 over a regular local ring R which satisfies the
assumptions of the theorem. Even if R is not local, but is a regular ring
containing a field k of characteristic       0, and such that the residue field at
every maximal ideal is algebraic over k, then if (WJ), holds at every
maximal ideal of R, it in fact holds at every prime ideal of R, so that the
Jacobian criterion for regularity holds for example in rings such as
k[X,,...,X,ljlY,,...,         Y,]. A further extension of this theorem can be
found in Matsumura           [2].
    Now we are going to prove the results analogous to Theorem 5 for formal
power series rings over a field of characteristic p. This is a difficult theorem
obtained by Nagata [6]. First of all we have to do some preparatory work.
    Let k be a field and k? c k a subfield; we say that k? is cofinite in k if
[k:k?] < 00. We say that a family 9 = {k,},,r             of subfields of k is a
directed family if for any LX,/Jo I there exists y ~1 such that k, c k, n k,.
    Let K be a field of characteristic      p, and k c K a subfield. Then there
exists a directed family 9 = {ka}as, of intermediate fields k c k, c K
cotinite in K such that n k, = k(KP). To construct this, we let B be a fixed
p-basis of K over k, and let I be the set of all finite subsets of B; then we
only need take k, = k(KP, B - a) for CCEI.

Lemma 3. Let K be a field, (k,},Er a directed family of subfields of K, and
set k = f7 k,. Then if V is a vector space over K, and vr,. . . , V,E I/ are
linearly independent over k, there exists cr~l such that vl,. . . , v, are also
linearly independent over k,.
Proof. For each CZIZZ,write q(a) for the number of linearly independent
elements over k, among vr, . . , v,; let a be such that q(a) is maximal, and
set q = q(x). Now if q < n, and we assume that ur,. . , u, are linearly
independent over k,, we have v, = c; civi with c+k,. Since the vi are
linearly independent over k, at least one of the ci does not belong to k,
SO that we can assume c1 $k. Hence there exists /?EZ such that cr $k,, and
also 7~1 such that k, c k, n k,, so that vr,. . . , vq and IJ, are linearly inde-
pendent over k,; this contradicts the maximality of q, so that q = n. H

Lemma 4. Let k c K be fields of characteristic p, and let 9 = {k,},,r be
a directed family of intermediate fields k c k, c K; then the following
conditions are equivalent:
   (1) n,k,(K?)     = k(KP);
   (2) the natural map QZKlk+ @r flKjk, is injective;
   (3) if a finite subset {ur , . . . , un} c K is p-independent over k, then it is
also p-independent over k, for some a;
   (4) there exists a p-basis B of K over k such that every finite subset of
B is p-independent over k, for some CI.
Proqf:(1)=>(3)If{u,,...,        u,,) is p-independent over k then the p? monomials
u;l...u;? for 0 d vi < p are linearly independent over k(KP), so that by the
previous lemma they are also linearly independent over k,(KP) for some c(.
   (3)+(4) is trivial; any p-basis will do.
   (4)*(2)IfO       # WEQ,,k and B is a p-basis then there is a unique expression
o = xcidKlkbi, with {h,, . . , b,} c B a finite set and C+ K. If we take CIsuch
that b,, . . . , h, are p-independent over k, then dKikmbifor 1 d i d n are linearly
independent as elements of QKik,, so that o has non-zero image in G,!, .
   (2)*(l)    Let ~EK be such that a$k(KP); then d,,,a # 0, so that dKlk,u #b
for some k,, in other words a$k,(K?).              n

 Lemma 5. Let k c K and 9 = {k,},,, be as in the previous lemma, and
 assume that nn k,(KP) = k(Kp). Then if L is an extension field of K which
is either separable over K or finitely generated over K, we again have
 na k,(LP) = k(Lp).
Proqf. (i) If L/K is separable, choose a p-basis B of K/k and a p-basis C
 of L/K; then in view of the exact sequence (Theorem 25.1)


 BuC is a p-basis of L/k. If b,,...,b,EB               and ci,...,c,~C     are finitely
 many distinct elements, then by the previous lemma, b,, . . . , b, are
p-independent over some k,, and from this (replacing k by k, in the above
 exact sequence) we see that b,, . . . , b,, cl,. . . , c, are p-independent over k,.
    (ii) If L/K is finitely generated, then since any finitely generated extension
 can be obtained as a succession of elementary extensions of type
    (a) L = K(x) with x separable over K,
 or
     (b) L = K(x) with xp = ~EK,
we need only consider case (b). We further divide this into two subcases:
in the first, d,,,a = 0, and then from Theorem 25.2 and the fact that
L N K[X]/(Xp        - a) we get RLjk = (a,,, @ L) @ Ldx; from this one sees easily
 that QLjk -       lim RLjk ~is injective. In the second subcase, d,,ka # 0 and
 now we have -    fiL+ -w ((a,,, 0 L)/L.d,,,a)  @ Ldx. Hence if we take a p-basis
of K/k of the form {u} u B? with a$B?, then {x> uB? is a p-basis of L/k. NOW
L/k satisfies condition (4) of the previous lemma, since for b,, . . , b,EB?, if
we choose c( such that {a, b,, . . . , b,,,} c K is p-independent over k,, then
{x,bl,...,b,}      CL is p-independent over k,. w

Lemma 6. Let K be a field of characteristic p, and let {K,j be a directed
family of colinite subfields K, c K such that n,K, = KP. Then if L is a
finite field extension of K, there exists SI such that
           rdnk,Q,,,. = rank&,,,,
for all subfields K?c    K, with [K:K?]      < ~0.
Proof. Let K = K, c K, c ... c K, = L be a chain of intermediate fields
   such that Ki = Ki_,(x,), with xi either separable algebraic over KimI or
   xfEKi- 1? Then by the previous lemma we have &K,(Kf)               = Kf, so
   that we need only prove the lemma for t = 1; hence suppose that L = K(x).
   If L is separable over K then in view of !&,,. = R,,,. OK L, the assertion is
   clear (any cxwill do). Thus suppose that xp = aEK, but a$KP; then there
   exists LY such that a$K,. If K? c K, then computing by means of
   L = K[X]/(Xp - a) we see that
            R LIP? = (%j~s OKLO L dx)/L.dxiKTa,
  and if rank!&,.     < co then rank&,,.      = rank!&,,,.      w

,; Theorem 30.9 (Nagata). Let k be a field of characteristic p, S =
;: k[Y1,...,        Y,] and PESpec S; suppose that 9 = {k,},,, is a directed
;,, family of colinite subfields of k such that /-jak. = kP. Then there exists
! a~1 such that for every intermediate field kP c k? c k, with [k:k?] -=zco
i the following formula holds:
::
:,$<              rank Der,(S/P) = dim (S/P) + rank Der,.(k).
     Proof. Set A = S/P, let L be the field of fractions of A, and dim A = n.
!, Choose a system of parameters x1,. . . ,x, of A, set B = k[x,, . . . , x,I], let K
1: be the field of fractions of B and mB its maximal ideal. Then A is a finite
 ! B-module, and hence [L:K] < co. If k? is an intermediate field kP c k? c k
 9 with [k:k?] = p?< co, and if ur,. . . , U, is a p-basis of k over k?, then
 t {lJ l,...,Ur,Xl,...,        x,} is a p-basis of B over C?= k?[xf,...,xE],     in the
 ;- sense that B is the free C?-module with basis the set of p-monomials in
 i/ 241,. . . ,u,, xi,. . . , x,. A derivation from B to B is continuous in the m,-adic
 j, topology, and any element of Der,.(B) is 0 on C?. Therefore Der,.(B) =
 ,; Der,.(B) = Hom#&,                B), and Qe,c, is the free B-module of rank n + I
  i with basis du,, . . . , du,, dx,, . . . , dx,. Therefore Der,.(B) is also a free
     B-module of rank n + r, and the theorem holds in case A = B.
        We write F? for the field of fractions of C?, and for k,E9 we set
  : C.=k,[x;         ,..., x;]l and write K, for the field of fractions of C,. As
  ,: above we have
                  Der,.(A) = Der,(A) = Hom,(R,,c,,         A),
  ? and since A is a finite C?-module, Q,,, is a finite A-module. Hence
                  Der,.(A)@ L = Hom,(Q,,,.         0 L, L) = Hom,(!F&,,  L),
     so that rank, Del-,.(A) = rank,R,,,..          Moreover,
                  n + rank Der,(k) = rank R,,,. = rank RKIF,,
     so that the conclusion of the theorem can be rewritten
                  rank RLIF, = rank aZKIF,.
     Any element of K, can be written with its denominator in
kP[xf,...,x[jj=      BP.

 However, B is faithfully flat over C,, so that K, n B = C,. From this one
deduces easily that n .K, = KP. Thus the theorem follows from the previous
\end{verbatim}
\end{proof}


\begin{theorem}[Printed Theorem 30.10]\label{mcrt-auto-theorem-30-10}\dcref{30.10}\source{matsumura-commutative-ring-theory:page-0259}\uses{}
\begin{verbatim}
Nagata?s Jacobian criterion). Let k be a field of charac-
teristic p, S = k[X,,.          .,X,,j, and let 1, P be ideals of S such that
I c PESpecS; set S,= R, R/IR = A, and suppose that htfR = r and
I = (fi,. . . ,f,). Choose a p-basis (uy}yer of k, and define D,EDer(k) by
D&f) = 6,,,. We make any element DEDer(k) act on the coefficients of
power series, thus extending D to an element of Der(S).
   Then the following conditions are equivalent:
   (1) A is a regular local ring;
   (2) there exists a finite number of elements tx,p,. . . , y~r such that
          rank J(fI, . . . , f,; D,, . . . , D,, 8/8X,, . . . , B/8X,)(P) = Y.
\end{verbatim}
\end{theorem}
\begin{proof}
\begin{verbatim}
(2) a(l) follows from Theorem 4.
   (1) * (2) As we see from the proof of Lemma 2, we need only prove (2)
in the case I = P. We let B be the family of subfields of k obtained by
adjoining all but a finite number of {uy}ysr to kP; then the conditions
of Theorem 9 are satisfied, and there exists /z?eF such that
          rank Der,.(S/P) = n - ht P + rank Der,.(k).
If we set [k:k?] = ps then by construction there exist yr,. . .,ys~r such
that k = k?(u y,, . . . , Us,);now set C = k?[X;, . . . , Xi], so that Q,,, is the free S-
module with basis duy,, . . . , duy,, dX,, . . . , dX,, and arguing as in the proof
of Theorem 8, we see that the derivations of S/P over k? are all induced by
derivations of S over k?, and that if gl,. . . , gm are generators of P, we have
           rank J(g,, . . . ,gm; D.fI,. . . ,Dps, @3X,, . . . ,3/8X,)(P) = ht P. n
\end{verbatim}
\end{proof}


\begin{theorem}[Printed Theorem 31.1]\label{mcrt-auto-theorem-31-1}\dcref{31.1}\source{matsumura-commutative-ring-theory:page-0261}\uses{}
\begin{verbatim}
Let A be a Noetherian ring and P&pec(A). Then there are
at most finitely many prime ideals P? of A satisfying P c P?, ht(P?/P) = 1
and ht P? > ht P + 1. (Ratliff [3] in the semi-local case, and McAdam [3] in
the general case.)
ProoJ: Let htP = n and take a,, . . . , U,EP such that ht(a,, . . . , a,) = n. Set
I=@,,...,     a,) and let P, = P, P,, . . . , P,. be the minimal prime divisors of
1. In general, if {Qd}l is an infinite set of prime ideals such that QA 1 P and
ht(Q,/P) = 1, then Q QA = P. This is because 0 Q1 is equal to its own
radical, hence is a finite intersection   of prime ideals containing    P, hence
either


or

           nQi = QA,n.. . nQkt;
but the second case cannot occur since Q1 $ QA, n. . .nQA, forJ$[1,, . . . , A,}.
Therefore if there were an infinite number of prime ideals Qi such that
QA 1 P, ht(Q,/P) = 1 and ht Qn # n + 1, then 0 Q1 = P. Hence there would
exist a QI, say Qo, which does not contain P,, . . . , P,. Then P would be the
only prime ideal lying between Q. and I. Let ~EQ,, - P. Then Q. would be
a minimal prime divisor of I + bA = (a,, . . . , a,, b), so that ht Q,, 6 n + 1, a
contradiction.   H
\end{verbatim}
\end{theorem}
\begin{proof}
\notready
\end{proof}


\begin{theorem}[Printed Theorem 32.2]\label{mcrt-auto-theorem-32-2}\dcref{32.2}\source{matsumura-commutative-ring-theory:page-0262}\uses{}
\begin{verbatim}
Ratliff?s weak existence theorem). Let A be a Noetherian
ring, and p, P&pec A be such that p c P, ht p = h and ht (P/p) = d > 1;
then there exist infinitely many p?&pecA with the properties
          pcp?czP,       htp?=h+  1 and ht(P/p?)=d-1.
\end{verbatim}
\end{theorem}
\begin{proof}
\begin{verbatim}
We first observe that if PI p1 I p2 I .?. 2 pd = p is a strictly
decreasing chain of prime ideals, and if pd- z I p? I p with ht p? = h + 1
then ht (P/p?) = d - 1. Now there exist infinitely many p?ESpec A such
that pd-2 I p? 2 p and ht (p?/p) = 1: for if pi,. . . , pk are a finite set of
these, let aep,-, - u pi, and let pk+r be a minimal prime divisor of
p + aA contained in pde2; then ht(&+ r/p) = 1. By the previous theorem,
all but finitely many of these satisfy ht p? = ht p + 1. w
Lemma I. Let A be a Noetherian           ring, and PESpecA with ht P = h > 1;
 suppose that UEP is such that ht (uA) = 1. Then there exist infinitely many
prime ideals Q c P such that
             u#Q and htQ=h-           1.
Proof. Suppose that pr,. . . , pt are the minimal prime ideals of A, and let
P,, . . . , P, be finitely many given height 1 prime ideals not containing U.
Let Qr,..., Q, be the minimal prime divisors of uA, so that these are also
height 1 prime ideals of A. Since h > 1, there exists UEP not contained
in any pi, Pj or Qk. Then
             ht (u, u) = 2 and ht (0) = 1.
Now let P,+l ,..., P,,, be the minimal prime divisors of (0); continuing
in the same way we can find infinitely many height 1 prime ideals not
containing u, so that ,if h = 2 we are done. If h > 2 we set A = A/(u) and
P= P/(u), so that ht 8= h - 1, and since the image U of u in A satisfies
ht (ri) = 1, by induction on h we can find infinitely many prime ideals ijz
of A satisfying
             ii&ls,cP      and htp==h-2.
The inverse image P, of r?, in A, does not contain u, and from
P,J(u) = Fa we have ht P;= h - 1. n
\end{verbatim}
\end{proof}


\begin{theorem}[Printed Theorem 31.3]\label{mcrt-auto-theorem-31-3}\dcref{31.3}\source{matsumura-commutative-ring-theory:page-0262}\uses{}
\begin{verbatim}
RatlifT?s strong existence theorem). Let A be a Noetherian
integral domain, p, PESpec A, and suppose that ht p = h > 0 and
ht (P/p) = d. Then for each i with 0 < i < d the set
             {p?ESpecAIp?cP,ht(P/p?)  =d - iandhtp?=   h + i}
is infinite.
\end{verbatim}
\end{theorem}
\begin{proof}
\begin{verbatim}
For i > 0 this follows at once from the weak existence theorem, so
that we consider the case i = 0.
   Step I. Replacing A by A, we can assume that (A, P) is a local integral
domain. Choose a, ,..., a,q such that ht(a, ,..., aj) =j for j= l,..., h,
and set
          a=(a,,...,     a,,) and b=(a,,...,a,-,).
Then p is a minimal prime divisor of a. Let
          a = a, n?..na,       and b = 6, n...nb,s
be shortest primary decompositions of a and b, and let pi, pi be the prime
divisors of ai and bj, respectively; we can assume that p = pr. Suppose
that b t+ 1,. . , b, are all the bj not contained in p. Then since
          a2n...na,nb,+,n...nb,~p,
we can choose an element YEP contained in the left-hand side and not
contained in p. Now p) c p for 1 d j < t, so that y$pl, and hence
        a:yA = a,          and     b:yA = b, n...nb,.
Now set
          B = A[x, , . . . ,x,],      where   Xi = ai/y,
          Z=(x,,...,       x,,)B   and    Q =PB+        1.
   Step 2. We prove that
          B/I-A/a,,      QESpecB       and htQ=h+d.
A general element of B can be written in the form
          sly?, with a@a + yA)?, for some v 3 0.
Now B = A + I, so that B/I N A/(Z n A). Now if cr~l n A, then there exists
v>O such that y?aEa(a+yA)?-?.              Since a:yy = a, we have C(EQ~.
Conversely, ya, c a, giving a, c In A, and hence In A = a,. Therefore
          BJI ?y A/a,.
Under this isomorphism, the prime ideals p/a, and P/a, of A/a, correspond
to (pB + I)/1 and (PB + 1)/I respectively in B/I; hence, setting
          q=pB+Z=p+I             and Q=PB+Z=P+I,
we have q, QESpec B, with B/q = A/p and B/Q = A/P, and Q is a maximal
ideal of B. Also, B/I = A/a, and since I is generated by h elements and al
is a p-primary ideal, we get
          htQ=dimB,gh+ht(P/a,)=h+ht(P/p)=h+d.
On the other hand, from y$p we get q = pA[y-?]nB:                indeed, if
a~pA[y-?1 n B, then we can write a = c/y? with c~pn(a + yA)?; since
a c p and y$p, we have
          pn(a+yA)?=a(a+yA)?-?+y?p,
and therefore cr~l+ p = q. Also, A[y-?I] = B[y-?I,  so that
        htq = htqA[y-?1   = htpA[y-?1    = htp = h.
However, Q/q = P/p, so that
        WQ/q) = W-?/p) = 4
and therefore htQ 3 d + h. Putting this together      with the previous
inequality, we see that ht Q = d + h.
   Step 3. For v = 1,2,. . ., set
           Jv=(X1,...,Xh-l,X*-y?)B,
let Q? be a minimal prime divisor of J? satisfying
          Q, = Q and WQ/QJ = WQIJJ,
and set Y? = Q?nA. We will complete the proof by showing that
P;, P;, . . . are all distinct, and that
          htP? = h, ht(P/P?) = d for all v.
Now B = A + J,, so that B/Q, N AJP: and Q/Q? N P/P:, hence ht(P/P?) =
ht(Q/Q?) = ht(Q/J?); moreover, ht Q = d + h, and since J, is generated by
h elements,
        ht(P/P:)>d+h-h=d.
Therefore we have
          d d ht(P/P:) = ht (Q/Q?) < ht Q - ht Q? = d + h - ht Q?,
and so if we can prove that
   (*) ht Q? = ht P: > h
then this will show simultaneously that ht(P/P?) = d and htP? = h.
   We have already seen that q = pA[y-?1 n B, so that qn A = p, and
hence y$q. Also, a, is a p-primary ideal with B/I N Ala,, so that I is a
q-primary ideal, and from htq = h we get
          ht(Z+yB)=ht(x,,...,x,,y)B=h+           1.
In addition, we have I + yB = J, + yB, and since Q? is a minimal prime
divisor of the ideal J,, which is generated by h elements, ht Q? < h, and
hence y$Q?, and
          Q,=Q,A[y-?]nB             and P:=Q,A[y-?]nA,
so that, by p.2Q, Example 1,
         ht Q? = ht P:.
Furthermore
        (b:yA)B c (xl.. . .,x~-~)B,
        (b:yA)+(a,-y?+?)AcJ,nAcQ,nA=P:,
and since all prime divisors of (b:yA) are also prime divisors of 6, we
have ht(b:yA) = h - 1. Now a,q and y$p, so that ah- y?+l$p, and
since all minimal prime divisors of (b : yA) are contained in p they do not
contain ah - y? + I. Therefore htPV 2 h. This completes the proof of (*).
   Step 4. If v < /.Lthen (Q + Qp)Ba contains y? = (y? - y@)/(l - yyeLc),and
therefore contains (J, + y?B)B, = (I + y?B)B,, so that
          ht(Q, + Q,)& > h + 1, but htQY < h,
and therefore Q, # Q,. From Q,A[y-?1            =P:A[y-?1     we see that
p;,p;,...    must also all be distinct. n
\end{verbatim}
\end{proof}


\begin{theorem}[Printed Theorem 31.4]\label{mcrt-auto-theorem-31-4}\dcref{31.4}\source{matsumura-commutative-ring-theory:page-0265}\uses{}
\begin{verbatim}
A Noetherian           local integral domain (A, m) is catenary if and
only if
          ht p + coht p = dim A for all p&pec A.
\end{verbatim}
\end{theorem}
\begin{proof}
\begin{verbatim}
?Only if? is trivial, and we prove ?if?. Let dim A = n; if A is not
catenary, then there exist p, PESpec A such that
          p c P, ht (P/p) = 1 but ht P > htp + 1.
Set ht(m/P) = d. Applying the strong existence theorem to A/p we see that
there exist infinitely many P,&pec A such that
          p c P,, ht(P,/p) = 1 and ht(m/P,) = d.
However, by assumption, ht(m/PJ + htP, = n, so that
          htP,=n-d=htP>htp+l.
But according to Theorem 1, there are only finitely many such P,, and
we have a contradiction.        n
   If A is a ring of finite Krull dimension, we say that A is equidimensional
if dim A/p = dim A for every minimal prime p of A.

Lemma 2. If an equidimensional     local ring (A,m) is catenary then
         htp, = htp, + ht(p,/p,)     for all pI,p2 ESpec A with p1 c p2.
Proof. If we choose a minimal         prime ideal p c p1 then ht(p,/p)=
ht(m/p) - ht(m/p,) = dim A - ht(m/p,), and this is independent of the
choice of p, so that htp, = ht(p,/p). Similarly, htp, = ht(p,/p),    so that
ht pz = ht(p,/p,)     + W,.      n
\end{verbatim}
\end{proof}


\begin{theorem}[Printed Theorem 31.5]\label{mcrt-auto-theorem-31-5}\dcref{31.5}\source{matsumura-commutative-ring-theory:page-0265}\uses{}
\begin{verbatim}
Let A, B be Noetherian local rings, and A -B         a local
homomorphism.      If B is equidimensional and catenary and is flat over A
then A is also equidimensional and catenary, and B/pB is equidimensional
for every p E Spec A.
\end{verbatim}
\end{theorem}
\begin{proof}
\begin{verbatim}
Write m and llJz for the maximal ideals of A and B. For any minimal
prime ideal p. of A there exists a minimal prime ideal P, of B lying over
po; then dim B/P, = dim B, so that dimB/p,B         = dim B, and then by
Theorem 15.1 we have
         ht(m/p,) = ht(fm/p,B) - ht(%R/mB) = dim B - ht(YJVm@.
This is independent of the choice of p,,, so that A is equidimensional.   If
pESpecA and P is a minimal prime divisor of pB then from the
going-down theorem (Theorem 9.5) we see that PnA = p, so that by
Theorem 15.1, htP = ht p, and therefore ht(!J.R/P) = ht!lJ - htP = htYJI -
htp is determined by p only. That is, B/pB is equidimensional.      Also, if
p?&pec A is such that p? c p and ht(p/p?) = 1, we let P? be a minimal
prime divisor of p?B contained in P; then B/p?B is also equidimensional
and flat over A/p?, so that ht(P/P?) = ht(P/p?B) = ht(p/p?) = 1. However,
B is equidimensional     and catenary, so that ht(P/P?) = htP - ht P? =
htp - htp?, and therefore htp = htp? + 1, and so A is catenary. n
\end{verbatim}
\end{proof}


\begin{theorem}[Printed Theorem 31.6]\label{mcrt-auto-theorem-31-6}\dcref{31.6}\source{matsumura-commutative-ring-theory:page-0266}\uses{}
\begin{verbatim}
Let (A, m) be a formally equidimensional Noetherian local
ring.
   (i) A, is formally equidimensional        for every p&pec A.
   (ii) If I is an ideal of A, then
               A/Z is equidimensional o A/I is formally equidimensional.
   (iii) If B is a local ring which is essentially of finite type over A (see
p. 232) and if B is equidimensional then it is also formally equidimensional.
   (iv) A is universally catenary.
\end{verbatim}
\end{theorem}
\begin{proof}
\begin{verbatim}
(i) Let PESpec(A*) be such that PnA= p, and set B =(A*)p
Since B is flat over A,, by Theorem 22.4, B* is flat over (A,)*. Now B is a
quotient of a regular local ring, and is equidimensional, so that by the above
corollary, B* is also equidimensional.         Hence by Theorem 5, (A,)* is also
equidimensional.
   (ii) follows easily from Theorem 5.
   (iii) B is a localisation of a quotient of A[X,, . . . ,X,1 for some n, so
that by (ii) we need only show that a localisation B of A[X,, . . . ,X,1 is
formally equidimensional.           Now A*[X,, . . . ,X,1 is faithfully flat over
ACX     i, . .  ,X,1,  so that there is a local ring C which is a localisation of
A*[X,, . . ,X,], and a local homomorphism           B + C such that C is flat over
B, and hence C* is flat over B*. By the remark after Theorem 15.5, C is
equidimensional, and is a quotient of a regular local ring, so that by the
corollary of Theorem 5, C* is also equidimensional; hence by Theorem 5,
B* is also equidimensional.
   (iv) Any local integral domain essentially of finite type over A is formally
equidimensional, and hence catenary, so that any integral domain which
is finitely generated over A is catenary. H
   We say that a Noetherian local ring A is formaZly catenary ([G2],
(7.1.9)) if A/p is formally equidimensional    for every pcSpecA. One sees
easily from the above theorem that formally catenary implies universally
catenary. The converse of this was proved by Ratliff [2]. Universally
catenary is a property of finitely generated A-algebras, and we have to
deduce from this a property of the completion, so that the proof is difficult.
Before giving Ratliff?s proof we make the following observation.
   Let (R, m) be a Noetherian local integral domain, K the field of fractions
of R, and R? the integral closure of R in K; let S be an intermediate ring
R c S c R? such that S is a finite R-module. S is a semilocal ring, and its
completion S* (with respect to the m-adic topology, which coincides with
the rad(S)-adic topology) can be identified with R* &S. Now R* is flat
over R, so that R c S c R? c K gives R* c S* c R* C& R? c R* QR K. The
ring R* 6&K is the localisation of R* with respect to R - (O}, so that
writing T for the total ring of fractions of R*, we can consider R* 6&K
 c T, and hence R* c S* c T. This leads to the possibility that properties of
R* will be reflected in some S.
\end{verbatim}
\end{proof}


\begin{theorem}[Printed Theorem 31.7]\label{mcrt-auto-theorem-31-7}\dcref{31.7}\source{matsumura-commutative-ring-theory:page-0267}\uses{}
\begin{verbatim}
For a Noetherian local ring A, the following conditions are
equivalent:
   (1) A is formally catenary,
   (2) A is universally catenary.
   (3) A[X] is catenary.
\end{verbatim}
\end{theorem}
\begin{proof}
\begin{verbatim}
(l)*(2) was proved in Theorem 6 and (2)*(3) is trivial. We prove
(3) + (1). Suppose then that A[X] is catenary; we will prove that A* is
equidimensional     by assuming the contrary and deriving a contradiction.
The proof breaks up into several lemmas.

Lemma 3. Let (R, m) be a catenary Noetherian local integral domain, and
let R* be its completion. Let dim R = n, and suppose that there exists a
minimal prime Q of R* such that
          1 < dim (R*/Q) = d c n.
For i = 1,2 ,... ., d - 1, write Oi for the set of p~Spec R satisfying the
conditions
   (1) htp = i,
and (2) there exists a minimal prime divisor P of pR* such that Q c P and
dim (R*/P) = d - i.
    Then mi is non-empty for each i.
Proof. We work by induction on i. If 0 # u~m then any minimal prime
divisor P of aR* + Q satisfies ht(P/Q) = 1, P n R # 0 and contains a. Hence
if we set M = {P&pec(R*)IQ            cP, ht(P/Q)= 1 and PnR ZO}, then
m = UpE,(PnR). Now ht(mR*/Q) = d > 1, so that nrR*$M, and hence
m itself is not of the form PnR for PEM; therefore both M
and {P~RIPEM)              are infinite    sets. By Theorem 1, M?= {PE
MI ht P = 1) is also infinite; choose any PEM?, and set p = PnR. Then
0 < htp = ht pR* < htP = 1, so that htp = 1, and P is a minimal prime
divisor of pR*. Since R* is catenary, dim (R*/P) = dim (R*/Q) - ht(P/Q) =
d - 1; hence ~4~ and the assertion is true for i = 1.
    If i > 1, take p as above, and set R = Rfp and p = P/pR*; then since
R is catenary, dim R = n - 1, and P is a minimal          prime divisor of
R* = R*/pR*, with dim (R*/F) = dim (R*/P) = d - 1 < n - 1. Hence by
induction there exists a prime ideal pi = pi/p of R of height i - 1, and
a minimal prime divisor Pi of $,R* such that P c Pi and dim(R*/
pii, = (d - 1) -(i - 1) = d - i. If pi = PJpR* then Pi is a minimal prime
divisor of piR* containing P, and hence Q, and R*/Pi = ITS/pi is
(d - i)-dimensional,    so that from the fact that R is catenary we get
ht pi = ht iji + ht p = i, and pi~~i.    n


Lemma 4. Let (R,m)         be a Noetherian local integral domain, R* its
completion, and let 0 = q1 n . . . n q, be a shortest primary decomposition
of 0 in R*, with Pi = ,/qi for 1 6 i < r. Suppose that P, satisfies
          htP, =O, cohtP, = 1 <dimR.
Then there exist b, cam and 6E(qZn...nq,)           -PI with the following
properties:
     (1) b - 6Eq,,
     (2) (b, 6)R* = (b, c)R*,
and (3) c/b#R but is integral over R.
Proof
  Step 1. P, is a minimal    prime ideal and Y> 1, so that q2.. . q, + P, ,
and we can choose G?E(q,n..*nq,)      - P,. Since coht P, = 1, it follows
that q1 + 6?R* is (mR*)-primary,  and hence if we set
        a = (ql + 6?R*)n R,
then this is an m-primary ideal. Now if 0 # bEa is any element, then in
R* we can write
         b=i+/?6?     with [Eq,   and /IER*.
Since b is not a zero-divisor in R* we have pS?$P,, so that setting 6 = /?S?,
 we get (1).
    Step 2. We prove that &$bR*. By contradiction,      suppose that 6 = b5
 with PER*; then b-6 = b(l- &ql,              and since b$P, we have
 1 - 5Eq1 c mR*, so that < is a unit of R*, and b&R* c q2. This
contradicts the fact that b is not a zero-divisor of R*.
    Step 3. If b is a non-zero ideal of R such that m is not a prime divisor
 of b then &bR*. Indeed, b:m = b, so that bR*:mR* = (b:m)R* = bR*,
 and so mR* is not a prime divisor of bR*; if P is any prime divisor of
 bR*, then P # mR* and P # P, (in view of P, n R = 0), so that coht P, = 1
implies P + P,, hence P $ ql, and there exists aeq, -P. If we write
 Q for the P-primary component of bR* then cr6 = OEQ, so that SEQ.
 Thus finally kbR*.
    Step 4. By the previous two steps m is a prime divisor of bR. Hence
we can write
          bR = I n J with I an m-primary ideal and J:m = J,
and then bR* = IR*nJR*             with ~EJR* and 6q!IR*. Moreover,
(ZR* + JR*)/IR* N (I + J)/Z, so we can choose CEJ such that 6 - ceIR*.
Then
        6 - CEIR* n JR? = bR*,
so that (b,c)R* = (b, 6)R*, and (2) is proved. If cebR we would have
(b, 6)R* = bR*, contradicting  Step2, so that c/b$R. On the other hand,
we have b-6Eq,      so that 6(b -6)=0,    that is b6 =d2, and c-d~bR*.
Set c = 6 + by; then c2 = d2 + 2b6y + b2y2Eb(d,b)R* = b(c, b)R*, so that
        c?~(bc, b2)R* n R = (bc, b2)R.
From this, we get c2 = bcu + b2v with u, VER, which proves that c/b is
integral over R. Thus we have proved (3). n

Lemma 5. In the notation      and assumptions of Lemma 4, set S = R[c/b];
then S has a maximal ideal of height 1.
Proof. Write T for the total ring of fractions of R*; then we can view S*
as an intermediate ring R* c S* = R*[c/b] c T, and T is the total ring of
fractions of S*. In Lemma 4 we had Ass(R*) = {PI,. . .,P,}, so that
setting Qi = P,Tn S*, we get
          Ass(S*)={Q,,...,Q,}      with htQi=htPi.
Moreover, S* is integral over R*, so that S*/Q, is integral over R*/Pi,
and hence also coht Qi = coht Pi. Let P* be any maximal ideal of S*
containing Qr. Then from (b,c)R* = (b, 6)R* we get S* = R*[c/b] =
R*[G/b], and since deq,n...nq,      and 6 - bcq, we have
          d/bEQ,n...nQ,       and 6/b- ~EQ,,
so that Qi + Qi = S* for all i > 1. Therefore Q1 is the only minimal prime
ideal contained in P*. However, coht Q1 = 1 and P* n R* = mR*, so that
ht P* = 1. Setting P = P* n S, we have ht P = ht P* = 1, and P is a maximal
ideal of S, since P* is a maximal ideal of S*. n
   Now we return to the proof of Theorem 7. Let A be a Noetherian local
integral domain,     and suppose that A* is not equidimensional;
then by Lemma 3, there exists a prime ideal p of A such that (A/p)* =
A*/pA* has dimension > 1, and has a minimal prime ideal of coheight
1. Set R = A/p; then by Lemma 4 and Lemma 5 applied to R, there exists
a subring S of the integral closure R? of R generated by one element, and
having a maximal ideal P with ht P = 1 < dim R. Let f: R[X] -S          be a
surjective homomorphism       of R-algebras and let q be its kernel. Set
Q=f-?(P).     Then P=Q/q and QnR=m.           Since RcS we have qnR=
(0), hence ht Q = ht m + 1 and ht q = 1 by the remark after Theorem 15.5.
Thus htQ - htq = htm = dim A/p > 1 = ht(Q/q), so that R[X] (and hence
also A[X]) is not catenary. l
Corollary 1. A Noetherian ring A is universally catenary if and only if
A[X]   is catenary.
Proof. Suppose A[X] is catenary and set B = A[X,, . . . ,X,]. In order to
prove that B is catenary it suffices to prove that BP is catenary for every
PESpec B. Let p = P n A. Then B, is a localisation of A,[X1,. . . , X,]. Since
A,[X,]    is catenary, A,[X, , . . . , X,] is catenary by the theorem.  l

Corollary 2. A Noetherian ring of dimension d is catenary if d d 2 and is
universally catenary if d d 1.
Proof. The first assertion is obvious from the definitions and the second
assertion follows from the first because dim A[X] = d + 1. n


      32 The formal fibre
          Let (A, m) be a Noetherian local ring and A* its completion. The
libre ring of the natural homomorphism       A --+ A* over any p&pecA is
called a formal jibre of A (although strictly speaking we should distinguish
between the libre and the libre ring, we will not do so in what follows).
If I is an ideal of A then (A/I)* = A*/ZA*, so that a formal fibre of A/I is
also a formal libre of A.
    Let A be a Noetherian ring and k c A a subfield. We say that A is
geometrically regular over k if A gkk? is a regular ring for every finite
extension k? of k (see $28). This is equivalent to saying that A, is
geometrically regular over k for every maximal ideal p of A.
    We say that a homomorphism          cp:A -+B     of Noetherian    rings is
regular if cp is flat, and for every peSpec A, the libre B OArc(p) of q over p
is geometrically regular over the field K(P).
   A Noetherian        ring A is said to be a G-ring (here G stands for
Grothendieck)      if A, -(A,)*    is regular for every prime ideal p of A; this
means that all the formal libres of all the local rings of A are geometrically
regular.
\end{verbatim}
\end{proof}


\begin{theorem}[Printed Theorem 32.1]\label{mcrt-auto-theorem-32-1}\dcref{32.1}\source{matsumura-commutative-ring-theory:page-0271}\uses{}
\begin{verbatim}
Let A 5 B 2 C be homomorphisms                     of Noetherian
rings; then
   (i) if cp and $ are regular then so is (1/q;
   (ii) if $cp is regular and $ is faithfully flat then cp is also regular.
\end{verbatim}
\end{theorem}
\begin{proof}
\begin{verbatim}
(i) Clearly $cp is flat. For pESpec A, write K = K(P), and let L be a
finite extension field of K. Set BBaL = B,> and C oAL = C,; then the
homomorphism         tiL:B, --+ C, induced by $ is also regular. Indeed, if P is
a prime ideal of B, then C&B,               = COB(BOA L) = CO,,!, = C,, and
hence if F is a finite extension of ti(P) then C,OB,.F = C@,F; setting
PnB = Q, since B, is a finite B-module we have [rc(P):ic(Q)] < co,
and hence [P:K(Q)] < a, so that C&F               is a regular ring. Now cp is
regular, so that B, is a regular ring, and hence by Theorem 23.7, (ii), C, is a
regular ring.
   (ii) The flatness of cp is obvious. If we let p, K and L be as above, then
 C, is a regular ring and is flat over B,, so that by Theorem 23.7, (i), B,
is also regular.      n
\end{verbatim}
\end{proof}


\begin{theorem}[Printed Theorem 32.3]\label{mcrt-auto-theorem-32-3}\dcref{32.3}\source{matsumura-commutative-ring-theory:page-0272}\uses{}
\begin{verbatim}
A complete Noetherian local ring is a G-ring.
\end{verbatim}
\end{theorem}
\begin{proof}
\begin{verbatim}
Let A be a complete Noetherian local ring and pESpec A; set
B = A,,, and let B* be the completion of B. We prove that B --+B*               is
a regular homomorphism;       that is, for any prime ideal p? of B, we need to
show that B* @ ti(p?) is geometrically regular over I.       However, A/p? n A
is also a complete local ring, so that we can replace A by A/p/n A and
reduce to the case p? = (0). Thus assume that A is an integral domain, and
let L be the common field of fractions of A and B; we must show that
B*&L      is geometrically regular over L.
   The problem can be further reduced to the case when A is a regular
local ring. In fact, by Theorem 29.4, A contains a complete regular local
ring R and is a finite module over R. Set p nR = q, R, = S and
B? = A,, = A 0,s; then B? is a semilocal ring, and B is a localisation of B? at a
maximal ideal, so that B* is one direct factor of B?* = B?@,S*. Write K for
the common field of fractions of R and S.
                    B?* = B?@.S* ----tB*

         A-
                     T
                    B? = A&S        -B=A,-L
                                              T
         T           ?                                      T
         R-S             =R,                              +K
Now B* C&L can be written as B* @esL, and it is hence a direct factor of
B?* OS, L = S* QL = (S* &K)&L,              so that we need only show that
S* OS K is geometrically regular over K.
    Now R, S and S* are regular local rings, and S* OS K is a localisation
of S*, so is a regular ring. Hence if char K = 0, there is nothing to prove.
We assume that char K = p in what follows. Then R has a coefficient
field k, and can be written R = k/TX,, . .,X,1. Choose a directed family
{k,] of cofinite subfields k, c k such that n,ka = kP, and set R, =
k,[[X:, . . . , X,Pj; write K, for the field of fractions of R,. Then one sees
easily (compare the proof of Theorem 30.9) that n,K, = KP.
    We set qa = qnR,; then since R, 3 RP, we see that q is the unique
prime ideal of R lying over q,. Hence if we let S, = (R,),% then S =
R, = R &,Sclr and S is a finite module over S,. Hence S* = S @,Sz; let us
prove that S* is O-smooth over S relative to S, (see $28). Suppose we are
given a commutative diagram of the form

         I?-rTrN


         s,-s        -      c
where C is a ring and N is an ideal of C with N2 = 0. If there is a lifting
u?:S* -+ C of u as a homomorphism        of S,-algebras, set w = u?,~*, and
let v? = u@ w:S* = S@s12Sz-C;       then one sees easily that v? is a*lifting
of v over S. Hence S* is O-smooth over S relative to S,. Now for
QESpec(S*), let QnS = (0); then (S*), is a local ring of S* @K, and
conversely, every local ring of S* &K is of this form. From the diagram




         S -K                 ---+C


        I               I
        S, -            Km,
one sees that (S*), is O-smooth over K relative to K,. Set E = (S*). and
m = rad(E); then E is m-smooth over K relative to K,, so that according
to Theorem 28.4,
         QKIK, @dElm)         --+QEjK 1@EWm)
is injective for every CI. Moreover,     since n,K,   = KP, by $30, Lemma 4,
           QK -   lim
                   t RKIK,
is injective, and hence
          % 0 (Elm) -       l@ (QKIK, 0 (Elm))
is also injective. Therefore, from the commutative      diagram
          QK O&V4      - QE 0 (E/m)
              1                     1
          I@ (Q,,,x@(E/m)) -       l@ (Q,,,,@Wm)),
we finally see that &@(E/m)         -fin,@(E/m)       is injective. Hence it
follows from the corollary to Theorem 28.6 that E is m-smooth over K, and
thus is geometrically regular. Since E is an arbitrary local ring of S* &K,
we see that S*@K is geometrically regular over K. n
\end{verbatim}
\end{proof}


\begin{theorem}[Printed Theorem 32.4]\label{mcrt-auto-theorem-32-4}\dcref{32.4}\source{matsumura-commutative-ring-theory:page-0273}\uses{}
\begin{verbatim}
Let A be a Noetherian    ring; if A,,, --+(A,,,)* is regular for
every maximal ideal m of A, then A is a G-ring.
\end{verbatim}
\end{theorem}
\begin{proof}
\begin{verbatim}
Since (A,,)* is a G-ring, by Theorem 2, A,, is also a G-ring. For
any peSpec A, if we let m be a maximal ideal of A containing p then A,
is a localisation of the G-ring A,,,, and hence A,, -(A,,)* is regular. n
   Theorem 4 makes it much easier to distinguish G-rings. For example,
the next theorem is based on Theorem 4.
\end{verbatim}
\end{proof}


\begin{theorem}[Printed Theorem 32.5]\label{mcrt-auto-theorem-32-5}\dcref{32.5}\source{matsumura-commutative-ring-theory:page-0274}\uses{}
\begin{verbatim}
Let A be a Noetherian semilocal ring; then a sufficient
 condition for A to be a G-ring is that if C is a finite A-algebra which is
 an integral domain, m is a maximal ideal of C and we write B = C,, then
 (B*)o is a regular local ring for every QESpec(B*) such that Q nB = (0).
\end{verbatim}
\end{theorem}
\begin{proof}
\notready
\end{proof}


\begin{theorem}[Printed Theorem 32.6]\label{mcrt-auto-theorem-32-6}\dcref{32.6}\source{matsumura-commutative-ring-theory:page-0274}\uses{}
\begin{verbatim}
H. Mizutani).          Let R be a regular ring. If the weak Jacobian
 condition (WJ) of$30 holds for R[X,, . . . ,X,1 for every n 3 0, then R is a G-
 ring.
\end{verbatim}
\end{theorem}
\begin{proof}
\begin{verbatim}
Since (WJ) is inherited by any localisation we can assume that R is
local. We prove that the condition of Theorem 5 holds. Set R,=
R(X,,... ,X,1. Any integral domain C which is finite as an R-module can be
expressed as C= R,/Q with Q&pec(R,) for some n. Let m be a maximal
ideal of C, M the maximal ideal of R, corresponding to m, and S=(R,),;
then it is enough to show that (S*),/Q(S*), is regular for every P&pec(S*)
with PnS=QS.             If htQ =r then htQ(S*),= r, and by assumption
 there exist D,, . . . ,D,EDer(R,) and fi,. . . ,fi~Q such that det(D&#Q.
Now Di has a natural extension to S, and then to S*, and since PnR,=Q,
we have det(Difj)#P, so that by Theorem 30.4, (S*),/Q(S*), is regular. n
\end{verbatim}
\end{proof}


\begin{theorem}[Printed Theorem 11.2]\label{mcrt-auto-theorem-11-2}\dcref{11.2}\source{matsumura-commutative-ring-theory:page-0278}\uses{}
\begin{verbatim}
3), and their fields of fractions are L and K respectively. Let
K? and K be their residue fields; then [K?:K] 6 [L:K] by Ex.10.8. Since A?/Q
is contained in the integral closure of A/P = k[X,, . . . ,X,1 in K?, by the
induction hypothesis, A?/Q is finite over A/P. Since
        Qi/Qi+l   =yliA!/yli+l
                                 A? N A?/Q and Qq= PA?

 we see that Al/PA? is finite over A/P. Moreover, A? is separated in the P-adic
topology (which is the same as the Y,-adic topology), because
 A? c kllq[Yl,. . . , ml. Since A is P-adically complete, A? is finite over A by
?Theorem 8.4. n

   From this result it is easy to derive the following theorem: Zf A is a
Noetherian local ring whose completion A* is reduced, then the integral
closure A? of A in its total ring of j-actions is finite ouer A. See [M] p. 237.
  On the other hand, if A is not reduced and if the maximal ideal m contains
a regular element (that is, a non-zero-divisor), then A? is not finite over A
(Krull [2]). In fact, if x # 0 is nilpotent, take a regular element s such that
 x$sA; (it is possible to find such an s since 0 mi = (0)). Then the elements
x/sj (forj = 1,2,3,. . . ) belong to A?. If A? is finite over A then there must be
some integer n such that s?(x/sj)~- A for all j. But then XE n,., Os?A = (0), a
contradiction.      n
    Suppose (A, m) is a one-dimensional Noetherian local integral domain.
Then A* is reduced if and only if A? is finite over A (Krull[2]).     In fact, if A? is
finite over A then it is a semilocal ring, and for each maximal ideal P of A?
the local ring(A?), is a DVR. The rad(A?)-adic topology of A? coincides with
 the m-adic topology, and the completion A?* of A? with respect to this
topology is a direct product of complete DVRs by Theorem 8.15, hence is
reduced. On the other hand it coincides with A* BAA? by Theorem 8.7.
 Since O+A-+A?         is exact, O--+A* --+ A* @ A? is also exact. Therefore
 A* is reduced. The converse holds, as already mentioned, without the
restriction on dimension.         n
    Rees [9] proved that a reduced Noetherian local ring A has reduced
 completion A* if and only if for every finite subset T of the total ring of
fractions K of A, the integral closure of A[lJ in K is finite over A[r].
    Akizuki Cl] constructed the first example of a one-dimensional Noeth-
erian local integral domain with non-reduced completion (see also Larfeldt-
Lech [ 11). To avoid such pathology, Nagata [Nl] defined and studied the
class of pseudo-geometric rings, which were called ?anneaux universelle-
ment japonais? by Grothendieck ([Gl], [G2]). These are now known as
?Nagata rings? ([Ml, [B9]). A Noetherian ring A is called a Nagata ring if
for every prime ideal P of A and for every finite extension field L of the field
 of fractions rc(P) of A/P, the integral closure of A/P in L is finite over A/P.
 For the basic properties of Nagata rings, see [M], $31. An alternative
definition is the following: a Noetherian ring A is a Nagata ring if (1) for
every maximal ideal m the formal tibres of A, are geometrically reduced,
 and (2) for every finite A-algebra B which is an integral domain, the set
 Nor(B) = {PESpecBIB,         is normal) is open in SpecB. The equivalence of
 these two definitions can easily be proved from [G2], (7.6.4) and (7.7.2).
    Although the integral closure A? of a Noetherian integral domain A may
fail to be finite over A, it is a Krull ring by the theorem of MoriCNagata
mentioned in 5 12. Because of its importance we quote here the theorem in
full.
Mori-Nagata       integral closure theorem. Let A be a Noetherian integral
domain and let A? be its derived normal ring. Then (1) A? is a Krull ring, and
(2) for every prime ideal P of A there are only finitely many prime ideals P? of
A? lying over P, and for each such P? the field of fractions rc(P?) of A//P? is
finite over K(P).

For a proof, see [Nl], (33.10). This proof depends on I.S. Cohen?s structure
\end{verbatim}
\end{theorem}
\begin{proof}
\notready
\end{proof}


\begin{theorem}[Printed Theorem 4.7]\label{mcrt-auto-theorem-4-7}\dcref{4.7}\source{matsumura-commutative-ring-theory:page-0042}\uses{}
\begin{verbatim}
Let A be an integral domain with field of fractions K; set
      x = m-Spec A. We consider any ring of fractions of A as a subring of K.
?     Then in this sense we have

              A = n A,,,.
                   IIEX
\end{verbatim}
\end{theorem}
\begin{proof}
\begin{verbatim}
For XEK the set I= (aEAlaxEA}          is an ideal of A. Now
      XEA, is equivalent to I $ p, so that if XEA,,, for every maximal ideal
      m then 1~1, that is XEA.  w
\end{verbatim}
\end{proof}


\begin{theorem}[Printed Theorem 5.2]\label{mcrt-auto-theorem-5-2}\dcref{5.2}\source{matsumura-commutative-ring-theory:page-0047}\uses{}
\begin{verbatim}
Let k be a field and A = k[cc 1,. . . , a,] an integral domain, and
write r = tr. deg, A for the transcendence degree of A (that is, of its field of
fractions) over k. Then if r > 0, A is not a field.
\end{verbatim}
\end{theorem}
\begin{proof}
\begin{verbatim}
Suppose that czl,, . . , q is a transcendence basis for A over k, and set
K = k(cr,, . . . , a,). Then since c(,+ 1,. . . , a, are algebraic over K, there exist
polynomials       ,fi(X,+ I ,                            , X,)EK [X,,                1,. . . , Xi], manic of degree di in Xi,
such that
            ~C~,+I~...,hJ 2.KCX,+1,...,X,ll(f,+l,...,fn)
and
            di = CKCar+ 13. . . 1ai):K(xr+    1) ?. .2 + 111.
The coefficients of fi are in K, so that for suitable 0 # gsk[a,, . . . , r,] we
have gfi~k[cr, )...) cI,][X,+1,...         ,X,1. In other words, if we set B =
4-a 1,. . . , a,, g-?1 then the fi are polynomials with coefftcients in B. We
are now going to show that A[g- ?I= B[cl,+ 1,. . , n,] is g free module over
B with basis in;=;+ 1 $?I 0 < e, < di}. Every element of B[a,+ r , . . . , a,] can
be written as P(x,.+~,. . . ,a,,) for some PcB[X,+,           , . ,X,1; dividing P as a
?5       Hilbert   Nullstellensatz   and dimension theory                       33

    polynomial in X, by f,, and replacing P by the remainder, we can assume
   that P has degree at most d, - 1 in X,; then dividing P as a polynomial
   in X,-i by f,- i, and replacing by the remainder, we can assume that P
   has degree at most d, _ 1 - 1 in X,- r. Proceeding in the same way, we
   can assume that P has degree at most di - 1 in Xi for each i; in addition,
   the elements (a110 < e < di} are linearly independent over K(cr,+ i, . . . ,
   gi- i). Hence A[g - ?1 is a free B-module. However, B is not a field; for
   Ma 1,. . . , a,] is a polynomial ring in Yvariables over k, and hence it contains
   infinitely many irreducible polynomials (the proof of this is exactly the
   same as Euclid?s proof that there exist infinitely many primes). Hence,
   there is an irreducible polynomial hek[a,, . . . , a,] which does not divide
   g, and then obviously h-?$k[a,,...         ,a,,g-I].  Therefore B contains an
   ideal Z with Z # 0, B, and since A[g-?1 is a free module over B, ZA[g-?1
   is a proper ideal of A[g-?1. Thus A[g-?1 is not a field. But if A were a
   field then we would have Al-g-?] = A, and hence A is not a field. n
\end{verbatim}
\end{proof}


\begin{theorem}[Printed Theorem 12.6]\label{mcrt-auto-theorem-12-6}\dcref{12.6}\source{matsumura-commutative-ring-theory:page-0105}\uses{}
\begin{verbatim}
Let A be a Krull ring, K its field of fractions, and write 9
for the set of height I prime ideals of A. Suppose given any pl,. . . p,.e~
and e , , . . , ~,EZ. Then there exists XEK satisfying
         q(x) = e, for               1 <i < r
and
         u,(x)>0           for all     PEP-{pl,...,pr}.

Here ui and u,, stand for the normalised additive valuations of K
corresponding to pi and p.
\end{verbatim}
\end{theorem}
\begin{proof}
\begin{verbatim}
If y,eA is chosen so that y,ep, but y,$p?:?up,   u??upr,     then
o&i) = 6,, for 1 <if r. Similarly we choose y,,...,y,~A        such that
Ui(yj) = bij. Then we set

         y=    fi   yp';
              i=l

let pi,. . . , P: be all the primes PEP - (p,,...,pr}   for which o,(y) ~0.
Then choosing for each j = 1,. . . , s an element tjep; not belonging to
PI U?? u P,, and taking v to be sufficiently large, we see that
          x = y(t1 . . t,)?
satisfies the requirements of the theorem.     n
\end{verbatim}
\end{proof}
